\documentclass[11pt, twoside, openright]{book}

% Trim size: 6" x 9" (US Trade paperback)
\usepackage[
  paperwidth=6in, paperheight=9in, inner=0.875in, outer=0.75in, top=0.75in, bottom=0.75in
]{geometry}

% Typography
\usepackage[T1]{fontenc}
\usepackage[french]{babel}
\usepackage[utf8]{inputenc}
\usepackage{palatino}          % Elegant serif font
\usepackage{microtype}         % Better typography
\emergencystretch=1em          % Allow extra stretch to avoid overflow with hyphenated/apostrophe words
\usepackage{setspace}
\setstretch{1.15}              % Standard book leading

% Graphics
\usepackage{graphicx}
\graphicspath{{../figures/}{../figures/book/}}

% Links (hidden for print - no colored text)
\usepackage{xcolor}
\usepackage[hidelinks, bookmarks=false]{hyperref}
\hypersetup{
  pdftitle={La simulation que vous appelez « Je »},
  pdfauthor={Matthias Gruber},
  pdfsubject={Consciousness, Computation, and the Cosmos},
}

% Chapter styling
\usepackage{titlesec}
\titleformat{\chapter}[display]
  {\normalfont\huge\bfseries}
  {\chaptertitlename\ \thechapter}{20pt}{\Huge\raggedright}
\titleformat{name=\chapter,numberless}[display]
  {\normalfont\huge\bfseries}
  {}{0pt}{\Huge\raggedright}
\titlespacing*{\chapter}{0pt}{50pt}{40pt}

% Section styling
\titleformat{\section}
  {\normalfont\Large\bfseries}{}{0pt}{}
\titlespacing*{\section}{0pt}{20pt}{10pt}

% Headers and footers
\usepackage{fancyhdr}
\pagestyle{fancy}
\fancyhf{}
\fancyhead[LE]{\small\itshape La simulation que vous appelez « Je »}
\fancyhead[RO]{\small\itshape\leftmark}
\fancyfoot[C]{\thepage}
\renewcommand{\headrulewidth}{0.4pt}

% For plain pages (chapter starts, front matter)
\fancypagestyle{plain}{
  \fancyhf{}
  \fancyfoot[C]{\thepage}
  \renewcommand{\headrulewidth}{0pt}
}

% Quote formatting
\usepackage{csquotes}
\renewenvironment{quote}
  {\list{}{\leftmargin=1.5em\rightmargin=1.5em}\item\relax\itshape}
  {\endlist}

% Prevent widows and orphans
\widowpenalty=10000
\clubpenalty=10000

% Tables
\usepackage{booktabs}
\renewcommand{\lightrulewidth}{0.8pt}   % KDP minimum line weight is 0.75pt
\usepackage{longtable}
\usepackage{array}
\usepackage{tabularx}

% Figure placement
\usepackage{float}

% Landscape pages for wide tables
\usepackage{pdflscape}
\renewcommand{\textfraction}{0.1}
\renewcommand{\topfraction}{0.9}
\renewcommand{\bottomfraction}{0.9}

% Make blank pages (from \cleardoublepage) truly blank
\makeatletter
\def\cleardoublepage{\clearpage\if@twoside \ifodd\c@page\else
  \thispagestyle{empty}\hbox{}\newpage
  \if@twocolumn\hbox{}\newpage\fi\fi\fi}
\makeatother

\begin{document}

% ==== FRONT MATTER ====
\frontmatter
\pagestyle{empty}

% ---- Half-title page (recto) ----
\vspace*{3in}
\begin{center}
{\Huge\bfseries La simulation que vous\\[0.3cm] appelez « Je »\par}
\end{center}
\cleardoublepage

% ---- Full title page (recto) ----
\vspace*{2in}
\begin{center}
{\Huge\bfseries La simulation que vous\\[0.3cm] appelez « Je »\par}
\vspace{0.8cm}
{\Large L'architecture de la\\[0.2cm] conscience, du calcul et du cosmos\par}
\vspace{2cm}
{\large Matthias Gruber\par}
\end{center}
\clearpage

% ---- Copyright page (verso of title) ----
\thispagestyle{empty}
\vspace*{\fill}
{\small
\noindent \textcopyright\ 2026 Matthias Gruber. All rights reserved.\par
\vspace{0.5cm}
\noindent No part of this publication may be reproduced, distributed, or transmitted
in any form or by any means without the prior written permission of the author,
except for brief quotations in reviews and certain noncommercial uses
permitted by copyright law.\par
\vspace{0.5cm}

\vspace{0.5cm}
\noindent Deuxième édition, 2026\par
\vspace{0.5cm}
\noindent www.matthiasgruber.com\par
}
\cleardoublepage

% ---- Dedication page (recto) ----
\thispagestyle{empty}
\vspace*{3in}
\begin{center}
\textit{À toutes celles et ceux qui se sont un jour demandé pourquoi quoi que ce soit fait l'effet de quelque chose.}
\end{center}
\cleardoublepage

% ---- Table of contents ----
\pagestyle{plain}
\tableofcontents
\cleardoublepage





\textit{À toutes celles et ceux qui se sont un jour demandé pourquoi quoi que ce soit fait l'effet de quelque chose.}



\chapter*{Le point qui n'existe pas}
\addcontentsline{toc}{chapter}{Le point qui n'existe pas}
\markboth{Le point qui n'existe pas}{}

Avant de vous dire quoi que ce soit --- avant la théorie, avant qui je suis, avant pourquoi tout cela devrait vous importer --- je veux que vous vous prouviez quelque chose à vous-même. Cela prend dix secondes.

Regardez les deux marques ci-dessous : un \textbf{X} à gauche, un \textbf{point} plein à droite.


\begin{figure}[htbp]
  \centering
  \includegraphics[width=0.95\textwidth]{../figures/blind-spot-test.png}
\end{figure}


Tenez le livre à bout de bras. Fermez l'œil \textbf{gauche}. De l'œil droit, fixez le X --- rien que le X. Ne trichez pas, ne jetez pas un coup d'œil vers le point ; vous le sentez posé là, à la lisière de votre champ de vision, et cela suffit. Maintenant, sans quitter le X des yeux, approchez lentement le livre de votre visage.

À une trentaine de centimètres environ de votre visage, observez le point à la limite de votre champ de vision.

Il a disparu.

Pas flou. Pas estompé. \textit{Disparu} --- et voici le détail qui mérite qu'on s'y arrête : il n'y a pas de trou là où il était. Pas de lacune, pas de tache sombre, pas de vide. La page est simplement, uniformément, continûment blanche à l'endroit exact où un gros point noir y est imprimé, devant votre œil ouvert et parfaitement fonctionnel. Votre œil est braqué droit dessus. Et vous ne pouvez pas le voir, parce que votre cerveau l'a discrètement repeint en blanc.

Voici ce qui vient de se passer. Chaque œil possède une tache aveugle --- une zone de la rétine, exactement là où le nerf optique perce vers le cerveau, entièrement dépourvue de cellules photosensibles. Il y a dans votre vision un trou assez large pour avaler neuf pleines lunes, dans chaque œil, en permanence. Vous ne l'avez jamais remarqué, pas une seule fois. Non parce qu'il serait petit --- il est énorme ---, mais parce que votre cerveau ne vous montre jamais le trou. Il lit ce qui entoure la lacune et en \textit{fabrique} le milieu : il invente le blanc de la page, invente le papier peint, invente ce qui devrait s'y trouver, et vous tend l'image achevée comme s'il s'agissait de la vérité brute.

Éloignez de nouveau le livre. Le point revient. Recommencez plusieurs fois. Vous êtes en train de regarder votre propre cerveau basculer entre rapporter le monde et l'inventer, et vous ne pouvez pas sentir le basculement. Il ne s'annonce jamais. Il ne s'est jamais annoncé.

Voici maintenant ce qui semble insensé et qui pourtant est un fait établi des neurosciences : ce remplissage n'est pas un tour de passe-passe que votre cerveau réserve à la tache aveugle. C'est ce que votre cerveau fait sur l'ensemble de votre champ visuel, à chaque seconde d'éveil. La tache aveugle est simplement le seul endroit assez maladroit pour se trahir --- l'unique zone sans aucune lumière entrante, où l'invention doit partir de zéro et trébuche. Partout ailleurs, le signal brut est presque aussi mince et bruité, étoffé par la même machinerie de suppositions ; vous ne le remarquez pas, uniquement parce que les suppositions sont généralement justes. Ce point qui s'éclipsait, c'était votre cerveau pris en flagrant délit. Tout le reste de ce que vous voyez est le même délit, commis trop habilement pour être remarqué.

Suivez cette idée jusqu'au bout. Vous n'avez jamais, de toute votre vie, vu la réalité directement. Pas une seule fois. Ni la page entre vos mains, ni la pièce autour de vous, ni le visage d'aucune des personnes que vous avez aimées. Chacune de ces choses était une reconstruction --- un modèle, construit à l'intérieur de votre crâne à partir d'un mince filet d'impulsions nerveuses et d'une montagne de suppositions préalables, puis vous a été présenté avec une telle fluidité que la pensée \textit{ceci n'est peut-être pas la réalité} ne vous a, jusqu'à cet instant précis, probablement jamais traversé l'esprit.

Vous vivez à l'intérieur d'une simulation. Elle est générée par environ un kilo et demi de tissu humide dans l'obscurité, elle est le seul monde auquel vous aurez jamais accès, et au cours d'une vie entière d'activité, elle n'a pas une seule fois mentionné qu'elle était là.

Ce vertige que vous ressentez peut-être --- retenez-le. Ce sentiment est le sujet tout entier de ce livre.

Je le connais de l'intérieur. Je l'ai éprouvé sur un pont d'Innsbruck, en plein jour, à vingt-cinq ans exactement, des larmes coulant sur mon visage tandis que je riais sans pouvoir m'arrêter et que la pierre la plus lourde de toute ma vie tombait de mes épaules. Je ne sais pas si quelqu'un m'a vu. Cela m'aurait été bien égal. Ce qui venait de se mettre en place, tout d'un coup, c'était un cadre qui expliquait non seulement comment le cerveau construit ce monde sans la moindre couture, mais pourquoi le construire \textit{fait un effet, quel qu'il soit} --- pourquoi cela fait quelque chose d'être vous, en train de lire ces lignes. À partir de cet instant, en ce qui me concernait, ma liste de choses à faire pour le reste de ma vie était bouclée.

En 2015, je l'ai couché par écrit : un livre de 300 pages, en allemand, autoédité, bourré d'appareil technique. Il s'en est vendu zéro exemplaire. Pas un seul. Je ne dis pas cela pour qu'on me plaigne --- je le dis parce que c'est pertinent pour l'histoire. La théorie contenue dans ce livre jamais lu dissout ce qui est peut-être le problème ouvert le plus difficile de la science, formule des prédictions que les théories concurrentes ne parviennent pas à formuler, et livre un schéma directeur concret pour une machine authentiquement consciente --- non pas un chatbot qui imite un esprit, mais une machine qui en \textit{a} un. Puis il est resté sur un disque dur pendant une décennie, tandis que les preuves, discrètement, s'accumulaient en sa faveur. Ce livre est la seconde tentative : plus court, écrit pour un public plus large, pour quiconque s'est un jour demandé pourquoi quoi que ce soit fait un effet quelconque. Si j'ai raison, le Problème Difficile de la conscience n'est pas difficile. C'est une erreur de catégorie qui se dissout à l'instant où on la voit --- comme ce point s'est dissous une page plus tôt.


\chapter*{À propos de l'auteur}
\addcontentsline{toc}{chapter}{À propos de l'auteur}
\markboth{À propos de l'auteur}{}

Je devrais vous dire qui avance ces affirmations. Je ne suis affilié à aucune université. Je n'ai pas de doctorat --- seulement un master en bio-informatique. Je n'ai jamais obtenu de financement de recherche, jamais travaillé dans un laboratoire de neurosciences. Si vous vérifiez les titres avant d'accorder votre confiance à un argument --- et c'est un réflexe légitime ---, c'est ici que vous risquez de reposer ce livre. Utilisez-le au moins pour caler une table bancale : les arbres ne seront pas morts en vain.

Ce que j'ai à la place, c'est un parcours d'outsider autodidacte : des décennies de lectures qui ont bifurqué chaque fois qu'elles butaient sur un mur, et qui n'ont jamais appris à respecter les frontières entre les domaines. Cela compte plus qu'il n'y paraît. Cette théorie n'existe que parce que personne ne m'a jamais dit quelles disciplines n'avaient pas le droit de se toucher --- les mathématiques et les automates cellulaires, l'apprentissage automatique, les neurosciences, de la neurologie clinique à la psychopharmacologie, la biologie de l'évolution, la philosophie de l'esprit. La plupart des théories de la conscience vivent dans un ou deux de ces mondes. Celle-ci tente de les relier tous --- ce qui est, à bien y réfléchir, exactement ce que fait le cerveau : il prend des flux disparates issus de sources complètement différentes et les tisse en une expérience unique. Une théorie de la conscience incapable de traverser les disciplines comme le cerveau les traverse devrait éveiller votre méfiance. La mienne est le produit de la même boucle récursive que celle qu'elle décrit --- le savoir nourrit la performance, qui nourrit la motivation, qui alimente à son tour le savoir --- et je déplierai cette boucle plus loin dans ce livre. Quant à savoir si la théorie vaut quelque chose, vous en jugerez vous-même. Venons-en maintenant à la théorie.


\mainmatter
\pagestyle{fancy}
\chapter*{Chapitre 1 : Le problème le plus difficile de la science}
\addcontentsline{toc}{chapter}{Chapitre 1 : Le problème le plus difficile de la science}
\markboth{Chapitre 1 : Le problème le plus difficile de la science}{}

Vous lisez cette phrase. Vous vivez une expérience.

Cette expérience (l'impression visuelle des lettres sur la page, la voix intérieure qui lit les mots, le sentiment de compréhension ou de confusion) est la chose la plus familière de votre vie et la plus mystérieuse de l'univers. Nous en savons plus sur l'intérieur des trous noirs que sur la raison pour laquelle lire fait l'effet de quelque chose.

Ce n'est pas une exagération. (Même si, en toute honnêteté, je dois préciser que les mathématiques comportent des problèmes que je tiens pour plus difficiles encore --- mais ceux-là n'empêchent pas la plupart des gens de dormir.) Les physiciens ont le Modèle standard. Les biologistes ont l'évolution et la génétique. Les chimistes ont le tableau périodique. Mais la conscience (le fait que « cela fasse quelque chose » d'être vous, en ce moment même, en train de lire ces lignes) n'a ni théorie établie, ni cadre dominant, ni explication communément admise.

Ce n'est pas faute d'avoir essayé. Depuis les années 1990, lorsque la conscience est redevenue un sujet scientifique respectable après des décennies d'exil béhavioriste, des milliers d'articles ont été publiés, des dizaines de théories proposées et des centaines de millions de dollars dépensés. Le résultat ? Un domaine dans ce que le philosophe des sciences Thomas Kuhn appelait un « état pré-paradigmatique » --- beaucoup d'idées concurrentes, aucun consensus, et le sentiment croissant qu'il manque peut-être quelque chose de fondamental.

\section*{Ce que le Problème Difficile demande réellement}

En 1995, le philosophe David Chalmers a donné au mystère son nom canonique : le Problème Difficile de la conscience.

Voici ce qu'il demande. Considérez l'expérience de voir du rouge. Les neuroscientifiques peuvent vous en dire long sur ce qui se passe dans le cerveau quand vous voyez du rouge : une lumière d'une certaine longueur d'onde frappe les cônes de votre rétine, des signaux voyagent le long du nerf optique, ils sont traités dans le cortex visuel, et diverses régions cérébrales se coordonnent pour produire la perception. Tout cela est bien compris, du moins dans les grandes lignes.

Mais rien de tout cela n'explique \textit{pourquoi voir du rouge fait l'effet de quelque chose}.

Vous pourriez, en principe, construire un modèle neuronal complet de la réponse du cerveau à la lumière rouge --- chaque neurone, chaque synapse, chaque voie de signalisation. Vous auriez une description fonctionnelle parfaite. Et vous n'auriez pas expliqué la sensation du rouge. Le « quel effet cela fait ». Le \textit{quale}, comme disent les philosophes.

Chalmers a distingué cela des « problèmes faciles » de la conscience (qui ne sont pas faciles du tout, seulement abordables en principe) : comment le cerveau intègre-t-il l'information ? Comment dirige-t-il l'attention ? Comment rend-il compte de ses propres états ? Ce sont des problèmes de mécanisme. Ils sont difficiles, mais d'une difficulté que les neurosciences savent aborder. Le Problème Difficile est différent : il demande pourquoi ces mécanismes s'accompagnent, tout simplement, d'une expérience. Pourquoi le cerveau ne se contente-t-il pas de traiter l'information « dans le noir », comme un ordinateur ?

\section*{L'état des lieux}

Ce n'est pas faute de candidats. La \textbf{Théorie de l'Information Intégrée} (IIT) réduit la conscience à un seul nombre, $\Phi$ --- et paie cette élégance au prix de l'implication qu'une grille de portes logiques suffisamment vaste est un tout petit peu consciente ; en 2023, plus de 120 chercheurs ont signé une lettre ouverte la qualifiant d'infalsifiable. La \textbf{Théorie de l'Espace de Travail Neuronal Global} (GNW) affirme que l'information devient consciente lorsqu'elle est diffusée à l'échelle du cerveau entier --- une explication limpide du \textit{quand}, et un silence éloquent sur le \textit{pourquoi} : pourquoi cette diffusion ferait-elle l'effet de quoi que ce soit ? Le \textbf{Traitement Prédictif} fait du cerveau un moteur de prédiction et de la perception une « hallucination contrôlée » --- brillant sur le \textit{contenu} de l'expérience et, de l'aveu même de ses architectes, sans même viser la partie difficile. Théories d'ordre supérieur, schéma de l'attention, traitement récurrent, théories du champ électromagnétique --- chacune une intuition véritable enroulée autour du même centre manquant.

Puis, en 2025, l'arbitre est entré sur le terrain. La collaboration COGITATE a mené pendant plusieurs années un test en confrontation directe des deux favorites, l'IIT contre la GNW, et a publié le résultat dans \textit{Nature}. Le verdict : aucune n'a gagné. Les signaux les plus fortement liés à la conscience sont apparus à l'arrière du cerveau, là où aucun des deux camps n'avait planté son drapeau. Après trois décennies et des centaines de millions de dollars, le domaine est sans doute plus loin du consensus qu'à ses débuts. Ce n'est pas là l'image d'une science qui approche de la solution. C'est l'image d'une science à laquelle il manque une pièce.

\section*{Deux dogmes qui bloquent le progrès}

Avant de vous dire ce qui manque à mon sens, je dois nommer deux préjugés qui sapent sourdement le domaine depuis des décennies. Je leur ai donné des noms dans mon livre d'origine, parce que je crois que les biais sans nom sont plus difficiles à combattre.

Le premier est ce que j'appelle le \textbf{dogme nSAI} --- « no strong artificial intelligence », pas d'intelligence artificielle forte. C'est la conviction répandue que des machines véritablement intelligentes sont impossibles --- une conviction enracinée non dans une preuve, mais dans l'échec des débuts de la recherche en IA dans les années 1960 et dans le retour de bâton qui a suivi. Quiconque croit que l'IA forte est possible apprend à garder cela pour soi s'il veut être pris au sérieux dans la recherche dominante. Ce n'est pas du scepticisme rationnel. C'est une cicatrice de vieilles défaites, durcie en doctrine.

Le second est plus profond et plus pernicieux. Je l'appelle le \textbf{dogme nSU} --- « no self-understanding », pas de compréhension de soi. C'est la croyance que l'esprit humain, la conscience humaine, ne peut pas, par principe, être compris par ce même esprit. On invoque les théorèmes d'incomplétude de Gödel, ou de vagues analogies avec les limites de l'observation cosmologique depuis l'intérieur de l'univers, ou bien --- c'est le cas le plus honnête --- on trouve simplement la perspective d'être entièrement expliqué trop effrayante pour l'envisager. Si la conscience n'est qu'une machine, que devient l'âme ? Que devient le sens ? Que devient la singularité d'être humain ?

Ces dogmes se renforcent mutuellement. Si l'on ne peut pas comprendre la conscience (nSU), alors on ne peut certainement pas en construire une (nSAI). Et si l'on ne peut pas en construire une (nSAI), alors peut-être la conscience est-elle réellement hors de portée de la compréhension (nSU). C'est une boucle fermée de pessimisme institutionnel, et elle a dissuadé un nombre considérable de chercheurs intelligents de seulement s'atteler à la tâche.

Je ne dis pas que ces dogmes sont défendus de mauvaise foi. Beaucoup de chercheurs y croient sincèrement. Mais ni l'un ni l'autre n'a jamais été prouvé. Ce sont des articles de foi, et ils ont fait plus de mal à la recherche sur la conscience que n'importe quelle expérience ratée.

\section*{Il manque quelque chose}

Je pense que si aucune théorie n'a percé le Problème Difficile, c'est que la plupart des gens cherchent la conscience au mauvais endroit. Ils regardent la machinerie neuronale (les neurones, les synapses, les oscillations, la connectivité) et demandent : « Lequel de ces processus est conscient ? »

La bonne question, à mon sens, est différente : « À quel niveau du traitement de l'information, et dans quelle architecture, l'expérience se produit-elle ? »

C'est le point de départ de la Théorie des Quatre Modèles --- le fait que vous vous êtes prouvé à vous-même avant même que ce livre ne commence vraiment. Vous n'avez jamais, de toute votre vie, fait l'expérience directe de la réalité. Le point que vous avez effacé de la page n'était pas un tour joué par le papier ; c'était votre cerveau comblant un trou de la rétine avec du contenu fabriqué, comme il comble chaque trou, à chaque regard, si parfaitement que vous n'avez jamais eu le moindre soupçon. Ce n'est pas une thèse marginale : que la perception soit construite plutôt que directe relève des neurosciences les plus classiques, admises par pratiquement tous les chercheurs du domaine. Ce que la Théorie des Quatre Modèles ajoute, c'est le pas que presque personne ne franchit --- l'idée que ce simple fait, suivi jusqu'au bout, dissout le Problème Difficile.

\section*{Un seul rasoir}

J'ai construit la théorie sous une seule discipline visible : \textbf{le rasoir d'Occam.} Pas de physique nouvelle, pas de magie quantique dans les microtubules, pas de proto-conscience saupoudrée dans la matière --- seulement des ingrédients que nous possédons déjà (réseaux de neurones, apprentissage, simulation, auto-référence), agencés dans la bonne architecture. Deux autres vieux principes œuvrent en silence, en coulisses, et plutôt que de vous les présenter tous maintenant, je ferai entrer chacun exactement là où il mord : l'un quand nous arriverons aux animaux, l'autre quand nous arriverons aux zombies et au libre arbitre. Ici, une seule lame suffit.

Il est temps, maintenant, de regarder les quatre modèles.


\chapter*{Chapitre 2 : Les quatre modèles}
\addcontentsline{toc}{chapter}{Chapitre 2 : Les quatre modèles}
\markboth{Chapitre 2 : Les quatre modèles}{}

Imaginez que vous regardez une pomme.

La pomme est posée sur une table devant vous. Rouge, ronde, brillante, à une quinzaine de centimètres de votre main. Vous la voyez, vous savez ce que c'est, vous pourriez tendre le bras et la saisir. Cela paraît simple : vous voyez une pomme.

Mais ce qui se passe réellement est bien plus compliqué.

La lumière réfléchie par la surface de la pomme pénètre dans vos yeux, où elle frappe les cellules photoréceptrices de vos rétines. Ces cellules convertissent la lumière en signaux électriques. Les signaux voyagent le long de vos nerfs optiques jusqu'au cortex visuel, à l'arrière de votre cerveau, où ils sont traités à travers une hiérarchie de détecteurs de traits de plus en plus sophistiqués : contours, orientations, couleurs, textures, formes et, finalement, objets. Quelque part dans cette cascade, l'activité neuronale correspondant à « pomme » s'active. Simultanément, votre système moteur prépare des actions potentielles (tendre le bras, saisir), votre système mnésique active des associations (le goût, la texture, la dernière fois que vous avez mangé une pomme), et votre système spatial suit la position de la pomme par rapport à votre corps.

Tout cela se produit en moins d'une seconde. Et rien de tout cela n'est ce que vous \textit{vivez}. Vous ne vivez pas les photons frappant les cônes, ni les signaux parcourant les axones, ni les détecteurs de traits qui s'activent. Vous vivez \textit{une pomme}. Un objet unifié, stable, tridimensionnel, posé dans un environnement spatial cohérent, doté d'une apparence, d'une texture et d'une signification particulières. Ce que vous vivez, c'est un \textit{modèle} : une simulation en temps réel de la pomme, engendrée par votre cerveau à partir des données brutes et de tout ce qu'il a précédemment appris sur les pommes, les objets, les tables et la physique.

Comme je l'ai soutenu au chapitre 1, cela ne prête pas à controverse en neurosciences. Tout neuroscientifique et tout philosophe de la perception s'accordent à dire que ce que vous vivez est un modèle, et non la réalité elle-même. La pomme que vous voyez est la \textit{meilleure hypothèse} du cerveau sur ce qui se trouve là-dehors, nourrie par les données sensorielles mais non identique à elles. (Les illusions d'optique en sont la preuve vivante : quand une illusion s'effondre --- quand vous la voyez soudain des deux façons ---, vous prenez la simulation sur le fait. Vous n'avez jamais vu la réalité directement. Vous avez toujours vu le modèle. L'illusion l'a simplement rendu évident.)

Mais c'est ici que ma théorie commence : le cerveau ne modélise pas seulement la pomme. Il modélise \textit{vous en train de regarder la pomme}. Et c'est ce second modèle (le modèle du soi) qui transforme le traitement de l'information en conscience.

\section*{Les quatre représentations de votre cerveau}


\begin{figure}[htbp]
  \centering
  \includegraphics[width=0.95\textwidth]{../figures/figure2-real-virtual-split-bw-fr.png}
  \caption{La séparation réel/virtuel. Le substrat (versant réel) stocke le savoir dans les poids synaptiques : physique, structurel, inconscient. La simulation (versant virtuel) engendre l'expérience en temps réel : fugace, dynamique, consciente.}
\end{figure}


Même de simples réseaux de neurones à trois couches peuvent apprendre à modéliser ce qu'ils reçoivent : montrez-leur assez d'exemples et ils construisent des représentations internes des motifs qu'ils rencontrent. Votre cerveau fait exactement cela, mais à une échelle infiniment plus riche. Il ne construit pas un modèle ; il en construit de multiples, couvrant tout, du champ visuel à la position de vos membres, du son d'une voix à la pression de vos pieds sur le sol. Ces modèles englobent aussi bien le monde extérieur \textit{que} votre propre corps, intégrant toutes les entrées sensorielles disponibles en représentations cohérentes.

Les neurosciences connaissent ces modèles depuis plus d'un siècle. Dans le cortex moteur et le cortex somatosensoriel, le corps est littéralement déployé sous forme de carte déformée --- mains et lèvres grotesquement agrandies parce qu'elles comptent davantage de terminaisons nerveuses, tronc et jambes comprimés en minces fuseaux. Ces cartes corticales, appelées \textit{homoncules}, ont d'abord été dressées par Wilder Penfield dans les années 1930, par stimulation électrique directe au cours d'opérations du cerveau. Elles ne sont que les exemples les plus frappants ; le cerveau entretient des cartes et des modèles semblables dans toute son architecture. (Voir l'annexe A pour en savoir plus sur l'organisation corticale.)


\begin{figure}[htbp]
  \centering
  \includegraphics[width=0.95\textwidth]{../figures/book/homunculi.en.png}
  \caption{L'homoncule cortical de Penfield. Le cortex somatosensoriel consacre une surface bien plus grande aux mains, aux lèvres et à la langue qu'au tronc tout entier --- une carte du corps déformée qui reflète la densité des terminaisons nerveuses, non la taille des parties du corps.}
\end{figure}


Je les appelle les \textbf{modèles implicites} : le Modèle Implicite du Monde (IWM) et le Modèle Implicite du Soi (ISM). Ils sont stockés dans la structure du cerveau --- dans la force des connexions synaptiques, l'architecture des circuits neuronaux, l'apprentissage accumulé d'une vie entière. Ils sont le disque dur du cerveau. Vous ne les vivez jamais directement, pas plus que vous ne vivez le silicium de votre téléphone. Mais ils encodent tout ce que le cerveau sait du monde et de vous.

Voici maintenant l'intuition clé. Ces modèles implicites ne restent pas simplement là, inertes. Ils \textit{engendrent} quelque chose. En ingénierie, un \textbf{jumeau numérique} est une réplique virtuelle en temps réel d'un système physique (un réacteur d'avion, un réseau électrique, une chaîne de production), continuellement mise à jour par les données de capteurs afin que les ingénieurs puissent surveiller le système et interagir avec lui sans le toucher directement. C'est exactement ce que font vos modèles implicites. Ils produisent une simulation virtuelle en temps réel du monde, et une simulation virtuelle en temps réel de vous. Ce sont les \textbf{modèles explicites} : le Modèle Explicite du Monde (EWM) et le Modèle Explicite du Soi (ESM). Tout ce que vous voyez, entendez, ressentez et pensez se déroule à l'intérieur de ces simulations, et non dans le monde lui-même.

Deux groupes de modèles (implicites et explicites), chacun contenant à la fois un modèle du monde et un modèle du soi. Quatre modèles en tout, et avec eux un langage pour parler de ce que la conscience fait réellement. (Une remarque à l'intention des neuroscientifiques et des lecteurs à l'esprit technique : le nombre « quatre » est un minimum de principe, non un décompte littéral de ce que le cerveau entretient. Si cela vous préoccupe, lisez l'annexe E avant de poursuivre --- elle traite ce point directement.)

Passons maintenant aux quatre modèles.

\textbf{Le Modèle Implicite du Monde (IWM)} est tout ce que vous savez du monde. Non pas ce à quoi vous pensez en ce moment --- tout ce à quoi vous \textit{pourriez} penser. Les lois de la physique (vous savez que les objets lâchés tombent). L'agencement de votre appartement (vous pouvez y naviguer dans le noir). La grammaire de votre langue maternelle (vous pouvez juger si une phrase est grammaticale sans connaître les règles). Les visages de tous ceux que vous avez connus. Le goût du chocolat. Le bruit de la pluie.

Tout ce savoir est stocké dans les connexions synaptiques de votre cerveau --- la force des liens entre neurones. Il s'est constitué au fil de toute votre vie, par l'expérience et l'apprentissage. Et vous n'en avez jamais, au grand jamais, directement conscience. Vous ne pouvez pas introspecter vos connexions neuronales. Vous ne pouvez pas sentir vos synapses. Le Modèle Implicite du Monde est comme une immense bibliothèque où vous n'entrez jamais --- vous ne faites que lire les livres qu'elle vous apporte à votre bureau.

\textbf{Le Modèle Implicite du Soi (ISM)} est tout ce que vous savez de vous-même. Votre schéma corporel --- la représentation inconsciente de l'endroit où se trouvent vos membres, de leur taille, de la façon dont ils bougent. Vos aptitudes motrices (faire du vélo, taper au clavier, jouer d'un instrument). Vos traits de personnalité, vos compétences sociales, vos schémas émotionnels, vos habitudes. La structure de votre mémoire autobiographique (le cadre qui ordonne vos souvenirs en une histoire de vie).

Comme le modèle du monde, le modèle du soi est stocké dans des poids synaptiques et n'est jamais directement conscient. Vous ne vivez pas votre schéma corporel ; vous vivez le corps que votre schéma engendre. Vous ne vivez pas votre personnalité ; vous vivez les pensées et les sentiments que votre personnalité produit. Le Modèle Implicite du Soi est l'équipe en coulisses --- indispensable à la représentation, mais jamais vue du public.

\textbf{Le Modèle Explicite du Monde (EWM)} est le monde que vous vivez effectivement. À l'instant même. La pièce où vous vous trouvez, les sons que vous entendez, le poids de ce livre dans vos mains (ou la lueur de l'écran sur lequel vous le lisez). C'est la simulation --- la réalité virtuelle en temps réel du cerveau, engendrée à partir du Modèle Implicite du Monde et de l'entrée sensorielle actuelle. Elle est vive, détaillée et convaincante sans la moindre faille. Vous vivrez votre vie entière à l'intérieur sans jamais en sortir.

\textbf{Le Modèle Explicite du Soi (ESM)}, c'est \textit{vous}. Le sentiment d'être un sujet. Le sentiment du « je » --- celui qui voit, entend, pense et décide. Cela aussi est une simulation : un modèle en temps réel engendré à partir du Modèle Implicite du Soi et des signaux corporels actuels. C'est le personnage que le cerveau crée pour habiter son monde virtuel.

VOUS êtes le personnage que le cerveau crée pour habiter son monde virtuel.

\section*{Le versant réel et le versant virtuel}


\begin{figure}[htbp]
  \centering
  \includegraphics[width=0.95\textwidth]{../figures/figure1-four-model-architecture-bw-fr.png}
  \caption{L'architecture à quatre modèles. Le cerveau entretient deux sortes de modèles (un du monde, un du soi), chacun selon deux modes : implicite (stocké dans la structure du cerveau) et explicite (activement en cours d'exécution sous forme de simulation en temps réel). La conscience réside dans les modèles explicites.}
\end{figure}


Les quatre modèles se répartissent en deux versants, et cette partition est le fondement de tout ce qui suit.

Le \textbf{versant réel} (les deux modèles implicites) est physique, structurel et relativement rigide (adapté par l'apprentissage). C'est le savoir emmagasiné du cerveau : poids synaptiques, connexions du réseau, configurations des récepteurs. Voyez-y tout ce que le cerveau \textit{a appris} --- cristallisé dans la structure physique du tissu lui-même. Il n'a aucune expérience. Une synapse qui décharge n'est pas davantage « vécue » que l'eau qui s'écoule dans un tuyau. Sur le versant réel, la lumière est éteinte.

Une chose mérite d'être soulignée : le versant réel est ce que les neurosciences étudient déjà. Quand un chercheur vous glisse dans un scanner d'IRMf, c'est le versant réel qu'il observe --- schémas d'activation, connectivité, débit sanguin dans les différentes régions. Quand un neurochirurgien stimule une aire corticale et regarde ce qui se produit, c'est le versant réel qu'il sonde. Les neurosciences cartographient ce territoire depuis plus d'un siècle, et elles ont accompli des progrès extraordinaires. La Théorie des Quatre Modèles ne rejette rien de ce travail. Elle dit seulement qu'il ne décrit que la moitié du tableau.

Le \textbf{versant virtuel} (les deux modèles explicites) est simulé, transitoire et dynamique. Il est réengendré à chaque instant, à partir du versant réel augmenté de l'entrée courante. Voyez-y tout ce que le cerveau \textit{est en train de faire avec} ce qu'il a appris --- le spectacle en direct, non le scénario rangé au placard. Et c'est \textit{toute} l'expérience. Chaque vision, chaque son, chaque pensée, chaque émotion, chaque souvenir, chaque rêve et chaque hallucination que vous ayez jamais eus se sont produits au sein du versant virtuel. Sur le versant virtuel, la lumière est allumée.

Mais voici le hic : le versant virtuel est invisible de l'extérieur. Même notre imagerie cérébrale la plus avancée ne peut le capter qu'indirectement. Une IRMf vous montre quelles régions du cerveau sont actives --- c'est le versant réel à l'ouvrage. Pour \textit{lire} réellement l'expérience consciente dans les données cérébrales, il faudrait déchiffrer le langage de programmation du cerveau --- comprendre non seulement quels neurones déchargent, mais ce que le schéma de décharge \textit{signifie} au niveau de la simulation. Cela exigerait quelque chose comme un connectome entièrement simulé : une réplique numérique complète du cerveau, s'exécutant dans un logiciel, produisant le même monde virtuel que celui que produit le cerveau biologique.

Je veux être honnête sur ce que la théorie vous donne et ce qu'elle ne vous donne pas. La Théorie des Quatre Modèles vous dit \textit{ce qu'}est la simulation, \textit{où} elle réside et \textit{pourquoi} elle procure un ressenti. Elle ne vous remet pas la clé de décodage. Lire le versant virtuel à partir du versant réel est un programme de recherche à venir --- un programme que la théorie délimite clairement, mais qu'elle ne peut pas encore mener à bien. Cependant, les fondations en sont déjà posées. Le Human Connectome Project et les projets apparentés cartographient le câblage du cerveau à une résolution toujours plus fine. Nous ne pouvons pas encore décoder le versant virtuel à partir des données structurelles, mais ces données structurelles arrivent.

Si vous avez l'esprit scientifique, vous entrevoyez peut-être déjà où tout cela mène. Si l'expérience n'existe que sur le versant virtuel, alors chercher l'expérience sur le versant réel --- dans les neurones, dans les synapses, dans la machinerie physique --- revient à la chercher complètement au mauvais endroit. C'est comme chercher l'intrigue d'un film dans les circuits du lecteur de DVD.

Voilà la clé.

\section*{Pourquoi votre cerveau a la capacité de se modéliser lui-même}

Nous avons donc établi que la conscience dépend de ces quatre modèles, le modèle explicite du soi assumant le gros de la tâche. Mais pourquoi le cerveau humain possède-t-il cette capacité au départ, alors que des animaux plus simples ne l'ont pas ? La réponse est sous nos yeux : l'architecture du cortex humain est, très littéralement, surdimensionnée pour le simple traitement de l'information.

Le néocortex humain a six couches. C'est un fait anatomique bien connu. Vous pouvez le constater dans n'importe quel manuel de neurobiologie. Mais voici ce qui est intéressant : vous n'avez pas besoin de six couches pour traiter l'information. Trois couches font l'affaire.

Réfléchissez à ce qu'un réseau de neurones standard doit accomplir. Première couche : recevoir l'entrée, la filtrer, la nettoyer. Deuxième couche : extraire des motifs, reconnaître des caractéristiques, abattre le gros du travail de calcul. Troisième couche : intégrer les résultats, prendre des décisions, produire une sortie. Entrée, traitement, sortie. C'est la recette de base, et trois couches y suffisent.

Mais nous en avons six.

À quoi servent les trois couches « supplémentaires » ?

Elles servent à modéliser les trois premières.

Un réseau à trois couches traite les données du monde. Un réseau à six couches traite le monde \textit{et} se regarde le faire. Les couches supplémentaires procurent au cerveau la capacité architecturale de bâtir non seulement un modèle de ce qui existe au-dehors, mais aussi un modèle de lui-même en train de modéliser ce qui est là-dehors. L'auto-simulation requiert ce dédoublement --- il vous faut un jeu de couches pour effectuer le traitement, et un autre jeu pour observer le traitement à l'œuvre.

Ce n'est pas de la spéculation sur ce que « font » les couches individuelles --- je ne prétends pas que la couche 4 fait ceci et la couche 5 cela. C'est une observation sur la capacité architecturale. Six couches font place à la fois au Modèle Implicite du Monde (le traitement appris, inconscient) et au Modèle Explicite du Monde (la simulation en temps réel). Elles font place à la fois au Modèle Implicite du Soi (votre schéma corporel, vos programmes moteurs, votre structure de personnalité) et au Modèle Explicite du Soi (le « vous » qui fait l'expérience d'avoir un corps, d'initier des actions, d'être une personne).

Considérez maintenant d'autres animaux. Les reptiles ont trois couches corticales. Les mammifères en ont six. Et parmi les mammifères, ceux qui ont le cortex le plus épais et le plus élaborément plissé (primates, cétacés, éléphants) sont précisément ceux qui présentent les signes les plus riches de conscience de soi. Reconnaissance de soi dans le miroir, planification de l'avenir, tromperie sociale, deuil. La capacité architecturale suit la phénoménologie.

Le saut de trois à six couches a peut-être été un accident de duplication génétique --- le copier-coller de l'évolution produisant l'architecture même que la conscience allait plus tard exploiter. Les ancêtres reptiliens avaient trois couches corticales. Quelque part dans la transition vers les mammifères, ce nombre a doublé. Les facteurs de transcription qui déterminent l'identité des couches corticales (Tbr1, Satb2, Ctip2, Fezf2) possèdent des paralogues évocateurs d'événements de duplication génique. Que ce fût un unique événement spectaculaire ou une élaboration graduelle fait toujours débat, mais le résultat est clair : les mammifères ont hérité du double de couches, et avec elles, de la capacité de modélisation de soi qui fait défaut à la plupart des reptiles.

C'est le pont entre la théorie des réseaux de neurones et l'expérience vécue. Le cortex humain n'est pas seulement un grand détecteur de motifs. C'est un réseau surdimensionné, structuré de façon récursive, doté d'assez de couches pour modéliser son propre processus de modélisation. Et lorsqu'un réseau se modélise lui-même en train de modéliser le monde, le résultat --- vu de l'intérieur --- est exactement ce que nous appelons la conscience.

Je dois être clair : je ne prétends pas que six couches corticales soient la \textit{seule} architecture capable de soutenir la conscience. C'est une solution --- celle qu'ont fait évoluer les mammifères. Mais il pourrait y en avoir d'autres. Le poulpe, avec son système nerveux radicalement distribué (huit bras semi-autonomes, chacun doté de son propre centre de traitement neuronal comptant environ 40 millions de neurones), représente une approche architecturale complètement différente qui atteint peut-être une puissance de calcul équivalente. Les oiseaux offrent un autre exemple frappant : les corvidés et les perroquets n'ont pas du tout de cortex en couches, leur pallium étant organisé en amas nucléaires plutôt qu'en feuillets, et pourtant les corneilles fabriquent des outils, planifient l'avenir et, on peut le soutenir, se reconnaissent dans les miroirs. Si ce qui compte est la capacité de se modéliser soi-même, et non le schéma de câblage précis, alors n'importe quelle architecture capable d'exécuter une simulation d'elle-même pourrait en principe être consciente. Un poulpe, une corneille, ou quelque chose que nous n'avons pas encore construit.

Mais les autres architectures peuvent attendre. Il y a une question plus proche, et c'est celle qui m'a arrêté la première fois que j'ai suivi tout ceci jusqu'au bout. Si \textit{vous} êtes le personnage que le cerveau inscrit dans sa propre simulation --- alors qui regarde le spectacle ?


\chapter*{Chapitre 3 : Le versant virtuel}
\addcontentsline{toc}{chapter}{Chapitre 3 : Le versant virtuel}
\markboth{Chapitre 3 : Le versant virtuel}{}

Imaginez que vous jouez à un jeu vidéo. Un bon jeu --- un jeu en monde ouvert immersif, avec des graphismes époustouflants, une physique réaliste et une histoire captivante. Vous contrôlez un personnage et, à travers lui, vous interagissez avec un monde virtuel d'une grande richesse de détails.

Demandez-vous maintenant : où le jeu existe-t-il ? Pas exactement sur l'écran --- l'écran ne fait qu'afficher des motifs de lumière. Pas exactement dans la carte graphique ou le processeur --- ceux-ci ne font que faire circuler des signaux électriques dans des circuits de silicium. Le jeu existe comme un \textit{processus virtuel} --- un phénomène de niveau supérieur qui émerge de l'activité du matériel, mais qui n'est identique à aucun élément particulier de ce matériel.

Le monde virtuel du jeu possède des propriétés que le matériel n'a pas. Le jeu a des montagnes, des rivières et des villes. Le processeur, lui, a des transistors. Le jeu a un cycle jour-nuit. Le processeur graphique, lui, a des cycles d'horloge. Vous pouvez raisonnablement demander « Quelle est la hauteur de cette montagne dans le jeu ? », mais il serait absurde de désigner un transistor et de dire « Ce transistor mesure 3 000 mètres de haut. » Les propriétés du jeu existent au niveau virtuel, et ce sont de véritables propriétés du jeu, même si le jeu n'est « qu'un » motif d'activité dans le matériel.

Ce n'est pas une métaphore. C'est ainsi que fonctionne votre cerveau.

Votre Modèle Explicite du Monde (EWM) --- le monde que vous vivez --- est un processus virtuel qui tourne sur un matériel neuronal, exactement comme le monde du jeu est un processus virtuel qui tourne sur un matériel de silicium. Le monde vécu possède des propriétés (couleurs, formes, distances, sons) que le matériel neuronal n'a pas (le matériel, lui, a des fréquences de décharge, des forces synaptiques et des concentrations de neurotransmetteurs). Les propriétés de votre monde vécu sont de \textit{véritables propriétés de la simulation}, même si la simulation n'est « qu'un » motif d'activité neuronale.

Votre Modèle Explicite du Soi (ESM) --- le « vous » qui vit le monde --- est lui aussi un processus virtuel. Il est aussi réel que le personnage du jeu dans l'analogie : véritablement présent au niveau virtuel, doté de véritables propriétés au niveau virtuel, mais absent au niveau du matériel.

\section*{Pourquoi l'analogie s'effondre (sur le point qui compte)}

L'analogie du jeu vidéo est utile, mais elle s'effondre sur un point crucial : le jeu a un \textit{joueur}. Il existe quelqu'un en dehors du jeu --- vous, assis sur le canapé --- qui vit le jeu. Le jeu lui-même n'a aucune expérience. Ce ne sont que des motifs de lumière et du code.

La simulation de votre cerveau n'a aucun joueur extérieur. Personne n'est assis hors de votre crâne à vivre la simulation. La simulation contient son propre observateur --- le Modèle Explicite du Soi. La simulation \textit{est} l'expérience, et non quelque chose vécu par quelqu'un d'autre.

Mettez-vous à la place du personnage du jeu. Vous \textit{êtes} le personnage principal. Depuis l'extérieur du jeu, un spectateur voit des pixels qui bougent sur un écran --- rien qui puisse ressentir quoi que ce soit. Mais depuis l'intérieur de la simulation ? Le monde du jeu est tout ce qui existe. Pour le personnage, les montagnes sont réelles, la lumière du soleil est chaude, le danger est effrayant. Aucun observateur extérieur ne devinerait jamais que ce tas de code ressent quelque chose --- mais c'est parce qu'il regarde au mauvais niveau. Il regarde le matériel. L'expérience existe au niveau du logiciel. Voilà ma thèse, et le reste de ce livre en expose les preuves.

C'est cela qui rend la conscience si singulière et qui fait paraître le Problème Difficile si insoluble. Dans le jeu vidéo, il y a une séparation nette entre le jeu (virtuel, sans expérience) et le joueur (physique, doté d'expérience). Dans le cerveau, il n'y a aucune séparation. La simulation et celui qui vit l'expérience sont une seule et même chose. Le Modèle Explicite du Soi ne regarde pas le Modèle Explicite du Monde depuis l'extérieur --- il est \textit{à l'intérieur} de la simulation, partie intégrante du même processus virtuel.

Et cette clôture auto-référentielle (la simulation qui s'observe elle-même de l'intérieur) est, selon moi, ce que nous appelons la conscience. Ce n'est pas quelque chose que l'on ajoute à la simulation. C'est ce que la simulation \textit{est}, dès lors qu'elle contient un modèle d'elle-même. Voilà pourquoi je dis que la conscience n'est pas une chose --- c'est un processus. Vous ne la trouverez pas en démontant le cerveau, pas plus que vous ne trouveriez un programme en cours d'exécution en désassemblant le processeur.

\section*{Les propriétés logicielles}

Si les modèles virtuels sont réellement des processus de type logiciel qui tournent sur un matériel neuronal, alors ils devraient se comporter comme un logiciel de manières précises et vérifiables. Et c'est le cas. Quatre propriétés du versant virtuel reviendront tout au long de ce livre ; laissez-moi donc les exposer dès maintenant.

\textbf{La bifurcation (le « forking »).} Un substrat unique peut faire tourner plusieurs configurations virtuelles simultanément. En informatique, vous forkez un processus et obtenez deux instances indépendantes qui tournent sur le même matériel. Dans le cerveau, c'est le trouble dissociatif de l'identité --- plusieurs modèles du soi, chacun avec son propre récit et son propre profil émotionnel, prenant tour à tour le contrôle du même substrat neuronal. Nous verrons cela au Chapitre 9.

\textbf{Le clonage.} Séparez physiquement le matériel, et vous obtenez des copies dégradées mais complètes du logiciel. Sectionnez le corps calleux, et chaque hémisphère fait tourner sa propre version de la simulation --- moins performante que l'originale, mais fonctionnellement entière. C'est le phénomène du cerveau divisé, lui aussi au Chapitre 9.

\textbf{La redirection.} Perturbez le flux d'entrée normal, et la simulation s'accroche au signal dominant, quel qu'il soit. Sous salvia divinorum, l'entrée proprioceptive submerge le système et le Modèle Explicite du Soi se reconfigure autour des sensations corporelles. Sous kétamine, l'entrée externe s'efface et la simulation tourne sur du bruit interne. Les modèles virtuels ne s'arrêtent pas --- ils traitent simplement ce qu'on leur donne. Le Chapitre 6 examine cela en détail.

\textbf{La reconfiguration.} Modifiez les poids de connexion du substrat, et vous changez ce que produisent les modèles virtuels. C'est exactement ce que fait la thérapie cognitivo-comportementale --- recâbler systématiquement le substrat pour que le Modèle Explicite du Soi génère d'autres récits, d'autres réactions émotionnelles, d'autres comportements.

La Théorie des Quatre Modèles (FMT) fait une prédiction précise au sujet de la thérapie : tout traitement efficace doit fonctionner en modifiant les modèles implicites (le substrat) de telle sorte que les modèles explicites (la simulation) changent en conséquence. La thérapie cognitivo-comportementale fait exactement cela. Elle identifie systématiquement les schémas inadaptés dans l'ISM et les recâble par une pratique structurée, changeant ce que produit l'ESM. Voilà pourquoi la thérapie cognitivo-comportementale dispose des données probantes les plus solides de toutes les psychothérapies : elle vise le bon niveau.

Cela soulève une question inconfortable au sujet des thérapies incapables d'expliquer leur mécanisme en ces termes. Si une approche thérapeutique ne précise pas ce qu'elle change dans le substrat, ni comment ce changement se répercute jusqu'à la simulation, alors au mieux elle agit par un mécanisme qu'elle ne comprend pas, et au pire elle n'agit pas du tout. Les preuves le confirment : les thérapies dont la base de preuves est la plus faible sont généralement celles dont les théories du changement sont les plus vagues. Si vous cherchez une thérapie, posez à votre thérapeute une question simple : « Que cherchez-vous précisément à changer dans mon cerveau, et comment ? » S'il ne sait pas répondre, envisagez d'en trouver un qui le sache.

Ce ne sont pas des métaphores. Ce sont des prédictions structurelles. Si ma théorie est fausse et que les modèles virtuels ne sont \textit{pas} des processus de type logiciel, alors ces parallèles ne sont que pure coïncidence. Mais les coïncidences ne s'alignent pas d'ordinaire quatre fois sur quatre à travers la neurologie clinique, la psychopharmacologie et la psychothérapie. Les chapitres qui suivent montreront chacune de ces propriétés à l'œuvre.

Il y a une expérience simple que vous pouvez faire dès maintenant --- enfin, avec un ami, une main en caoutchouc, un écran en carton et deux pinceaux --- et qui démontre avec quelle facilité on peut tromper le Modèle Explicite du Soi. C'est l'illusion de la main en caoutchouc, conçue par Matthew Botvinick et Jonathan Cohen --- l'un des tours de salon les plus révélateurs de toute la neuroscience.

Le dispositif est simple. Vous êtes assis à une table, un bras caché derrière un écran en carton. On place devant vous, bien visible, une main en caoutchouc réaliste, à peu près là où se trouverait votre main cachée. Quelqu'un caresse simultanément la main en caoutchouc et votre vraie main cachée avec deux pinceaux, au même endroit, à la même vitesse. Après une ou deux minutes de ces caresses synchronisées, quelque chose de troublant se produit : vous commencez à \textit{sentir} les coups de pinceau sur la main en caoutchouc. Pas sur votre vraie main, derrière l'écran. Sur la fausse main, sous vos yeux.

Votre Modèle Explicite du Soi a intégré la main en caoutchouc à son schéma corporel. Il a réassigné l'appartenance --- décidé que la main en caoutchouc fait partie de « vous ». Le modèle du soi n'est pas câblé en dur. Il est appris. Il est mis à jour en continu sur la base des meilleures preuves disponibles, et lorsque la preuve visuelle (voir la main en caoutchouc caressée) concorde constamment avec la preuve tactile (sentir sa vraie main caressée), l'ESM en tire la conclusion rationnelle : cette main est la mienne. Si quelqu'un menace ensuite la main en caoutchouc (abat un marteau dans sa direction), vous avez un mouvement de recul, vous sentez une bouffée d'angoisse, votre réponse électrodermale s'envole. Pour la partie de votre cerveau qui définit le « vous », cette main \textit{est} la vôtre.

Ce n'est pas un bug. C'est le modèle du soi qui fonctionne exactement comme prévu --- actualisant sans cesse sa frontière corporelle à partir de la corrélation sensorielle multimodale. C'est le même mécanisme qui permet aux personnes amputées de « sentir » un membre prothétique comme le leur après une période d'utilisation. Et c'est le même mécanisme qui se rompt dans l'asomatognosie, où les patients nient l'appartenance de leurs propres membres, et dans le syndrome de la main étrangère, où la main bouge d'elle-même.

\section*{L'hologramme en patchwork}

Il existe une cinquième propriété du versant virtuel qui mérite sa propre section, car elle explique quelque chose qui intrigue les neuroscientifiques depuis près d'un siècle : pourquoi les lésions cérébrales dégradent les fonctions \textit{progressivement} au lieu d'effacer des souvenirs isolés.

Dans les années 1920 et 1930, le psychologue Karl Lashley apprit à des rats à parcourir un labyrinthe, puis retira chirurgicalement des morceaux de leur cortex afin de découvrir où le souvenir était stocké. Il ne le trouva jamais. Quel que soit le morceau qu'il retirait, les rats se souvenaient encore du labyrinthe. Ce qui comptait, c'était \textit{quelle quantité} de cortex il retirait, non pas \textit{quelles parties}. En retirer un peu, et les rats devenaient légèrement moins bons. En retirer beaucoup, et ils devenaient bien pires. Mais le souvenir n'était jamais tout simplement \textit{effacé}, proprement excisé comme un fichier supprimé d'un disque dur. Lashley passa sa carrière à chercher l'« engramme » --- la trace physique d'un souvenir --- et conclut, on le sait, qu'il ne semblait pas exister.

Il cherchait la mauvaise chose. Le souvenir n'était pas stocké \textit{dans} un morceau particulier de cortex, à la manière d'un fichier logé sur un secteur particulier d'un disque dur. Il était stocké \textit{à travers} l'ensemble du réseau, réparti dans les poids des connexions entre des millions de neurones. C'est ainsi que fonctionnent les réseaux de neurones : l'information ne réside dans aucun nœud unique. Elle est encodée dans le motif des connexions entre tous ces nœuds. Vous ne pouvez pas pointer une seule synapse et dire « c'est ici qu'est stocké le labyrinthe », pas plus que vous ne pouvez pointer un seul pixel et dire « c'est ici qu'est stocké le film ».

C'est essentiellement une propriété holographique. Si vous prenez un hologramme physique et le coupez en deux, vous n'obtenez pas deux moitiés de l'image. Vous obtenez deux copies de l'image \textit{complète}, chacune en résolution moindre. Coupez-le en quatre et vous obtenez quatre images complètes, plus floues encore. L'information contenue dans un hologramme est répartie sur toute la plaque, de sorte que chaque morceau contient l'image entière --- simplement avec moins de détail.

Les réseaux de neurones font la même chose. Entraînez un réseau à reconnaître des visages, puis détruisez 10 \% de ses connexions au hasard. Il n'oublie pas 10 \% des visages. Il devient légèrement moins bon sur \textit{tous} les visages. Détruisez-en 50 \% et il devient nettement moins bon sur tout, mais il reconnaît encore quelque chose. L'information est étalée sur l'ensemble du réseau, ce qui explique précisément pourquoi Lashley ne parvenait pas à trouver l'engramme : il était partout et nulle part.

Mais --- et c'est là que cela devient intéressant --- le cerveau n'est pas \textit{un seul} hologramme. C'est ce que j'appelle un \textit{hologramme en patchwork}. Au sein d'une même aire fonctionnelle (disons votre cortex visuel primaire, à peu près l'aire 17 de Brodmann), les colonnes corticales se ressemblent, et l'information est stockée de façon holographique. Détruisez quelques colonnes et vous le remarquez à peine. L'aire est localement holographique (une partie contient le tout, en résolution moindre).

À l'échelle globale, des aires différentes font des choses différentes. Votre cortex visuel n'est pas interchangeable avec votre cortex moteur. Retirez l'intégralité du cortex visuel et vous perdez la vision --- il n'existe aucune sauvegarde floue. Le cerveau est donc localement holographique à l'intérieur de chaque région fonctionnelle, fractalement auto-similaire dans son architecture en colonnes, mais globalement \textit{non} holographique. C'est un patchwork : des dizaines de tuiles holographiques cousues ensemble en un composite qui, pris dans son ensemble, est résolument non holographique.

Cette structure en patchwork explique un schéma que l'on retrouve encore et encore en neurologie clinique. Les petits accidents vasculaires cérébraux et les petites lésions provoquent souvent des déficits étonnamment légers, parce qu'à l'intérieur de toute aire corticale donnée, le principe holographique vous protège. Le tissu restant reconstruit l'information manquante en résolution moindre. Mais les AVC étendus qui anéantissent une aire fonctionnelle entière provoquent des pertes catastrophiques et spécifiques (cécité, paralysie, aphasie), parce que vous avez retiré une tuile entière du patchwork, et qu'aucune autre tuile ne peut la remplacer.

Cela explique aussi pourquoi les souvenirs ne « disparaissent » pas d'un coup lorsque des neurones meurent. Chaque jour, des neurones meurent et des synapses sont élaguées. Si les souvenirs étaient stockés comme des fichiers sur un disque dur, on s'attendrait à en perdre un de temps à autre --- à se réveiller un matin en ayant oublié son propre mariage, le chien de son enfance ou le goût du café. Cela n'arrive jamais. Les souvenirs s'estompent plutôt progressivement, perdant en détail et en vivacité au fil des années. C'est exactement ce que prédit un système de stockage holographique : la dégradation est douce, proportionnelle et globale, jamais soudaine, discrète ou localisée.

L'hologramme en patchwork est la raison physique pour laquelle les propriétés logicielles que j'ai décrites plus haut (le clonage en particulier) fonctionnent réellement. Divisez le cerveau en deux, et chaque moitié conserve une copie dégradée mais complète de la simulation, parce qu'à l'intérieur de chaque hémisphère, le principe holographique garantit que chaque morceau contient l'image entière. La simulation ne se brise pas. Elle tourne simplement à une résolution moindre.


\chapter*{Chapitre 4 : Pourquoi cela fait l'effet de quelque chose (et pourquoi ce n'est pas la bonne question)}
\addcontentsline{toc}{chapter}{Chapitre 4 : Pourquoi cela fait l'effet de quelque chose (et pourquoi ce n'est pas la bonne question)}
\markboth{Chapitre 4 : Pourquoi cela fait l'effet de quelque chose (et pourquoi ce n'est pas la bonne question)}{}

Nous pouvons maintenant affronter le Problème Difficile de front.

La question est : \textbf{pourquoi le traitement physique fait-il l'effet de quelque chose ?}

La réponse : \textbf{il n'en fait aucun.}

Le traitement physique (les neurones qui déchargent, les synapses qui transmettent, les modèles implicites qui stockent et calculent) n'a aucune expérience. Aucune. Cela ne fait aucun effet d'être le versant réel. Le versant réel est précisément ce traitement « dans le noir » que la conscience, selon le Problème Difficile, est censée expliquer.

La \textit{simulation}, elle, ressent. Le Modèle Explicite du Monde (EWM) et le Modèle Explicite du Soi (ESM) --- le versant virtuel --- sont le lieu où vit l'expérience. Et au sein de la simulation, l'expérience n'est pas un ajout mystérieux au processus. L'expérience est ce que la simulation \textit{est}, dès lors qu'elle inclut un modèle du soi. Le Modèle Explicite du Soi « percevant » le Modèle Explicite du Monde, voilà ce que nous appelons les \textit{qualia}. Les \textit{qualia} sont la manière dont le soi virtuel enregistre le monde virtuel.

Voyez les choses ainsi. Si vous demandiez : « Pourquoi la commutation de transistors fait-elle l'effet d'un jeu vidéo qui tourne ? », la réponse serait : « Elle n'en fait aucun. La commutation de transistors ne fait aucun effet. Le jeu est un processus virtuel qui tourne sur des transistors, mais qui possède des propriétés que les transistors n'ont pas --- des paysages, des personnages, une physique, de la lumière. Ces propriétés sont des propriétés réelles du processus virtuel, pas des transistors. »

De même : la décharge neuronale ne fait pas l'effet de voir rouge. La décharge neuronale engendre et entretient une simulation, et au sein de cette simulation, le modèle du soi perçoit une certaine classe de contenus du modèle du monde comme ce que nous appelons la « rougeur ». La rougeur est une propriété réelle de la simulation, pas une propriété des neurones.

Le Problème Difficile supposait que nous devions expliquer comment le traitement physique produit l'expérience. Mais le traitement physique ne produit pas d'expérience --- il produit une \textit{simulation}. Et la simulation, parce qu'elle inclut une boucle auto-référentielle (l'ESM se modélisant lui-même au sein de l'EWM), \textit{est} constitutivement l'expérience.

\section*{Qui êtes-vous à votre réveil ?}

Partons de quelque chose que vous avez presque certainement déjà ressenti.

Vous remontez de quelque part au fond --- un évanouissement, un K.-O., une anesthésie, ou simplement un sommeil noir et sans rêves dans une chambre que vous ne reconnaissez pas. Et l'espace d'un battement de cœur, il n'y a personne à bord. Il y a de la conscience --- un plafond, une bande de lumière, le simple fait d'être --- mais pas encore de \textit{vous} qui s'y rattache. Vous ne savez pas où vous êtes. Pendant une longue seconde, vous ne savez pas \textit{qui} vous êtes. Puis tout afflue : votre nom, votre histoire, la forme de votre vie, le corps dans lequel vous êtes allongé. La couture se referme si vite que vous la manquez presque. Mais vous ne l'avez pas manquée. L'espace de ce battement de cœur, il y a eu de l'expérience sans un soi pour la revendiquer --- une conscience qui tournait pendant que le modèle du soi était encore en train de démarrer, fouillant la pièce en quête d'un point d'ancrage et ne trouvant rien.

Retenez cet intervalle. C'est tout le chapitre.

À cause de ce qu'il vous révèle : le « vous » qui est revenu en ligne n'a pas été extrait intact d'une mémoire, comme un fichier sauvegardé qu'on rouvre. Il a été \textit{reconstruit} --- assemblé sur place à partir du substrat sous-jacent. Chaque nuit, votre Modèle Explicite du Soi s'effondre et le sommeil profond efface la simulation en cours. Chaque matin, il redémarre et rebâtit ce « vous » à partir du Modèle Implicite du Soi (ISM), la version lente et stockée de vous-même que le sommeil laisse derrière lui. D'ordinaire, la reconstruction est instantanée et sans heurt, et vous ne surprenez jamais le raccord. Dans cette chambre inconnue, pour une fois, vous l'avez surpris.

Alors quand je passerai le reste de ce chapitre à soutenir que l'expérience \textit{est} une simulation tournant sur un substrat --- que le « vous » ressenti est un processus virtuel et non les neurones ---, vous en détiendrez déjà la preuve. Vous vous êtes trouvé dans la demi-seconde qui précède la fin du chargement de la simulation. Démontons maintenant cette demi-seconde, et demandons-nous pourquoi ce qui se charge fait le moindre effet.

\section*{La question de la circularité}

La première question que posent la plupart des lecteurs : « N'avez-vous pas simplement déplacé le problème ? Pourquoi \textit{cette} simulation a-t-elle une expérience, quand une simulation météorologique n'en a pas ? »

La réponse est l'auto-référence. Une simulation météorologique modélise la météo. Elle ne se modélise pas \textit{elle-même}. Une simulation météorologique a un « dehors » --- l'ordinateur, le programmeur, le scientifique qui interprète les résultats. La simulation peut être décrite intégralement sans référence à la moindre expérience, parce qu'il n'y a aucun modèle du soi en elle.

La simulation du cerveau, elle, se modélise elle-même. Le Modèle Explicite du Soi est le modèle que la simulation se fait de \textit{son propre processus}. Cela crée une boucle fermée : le modèle et la chose modélisée sont le même système. Il n'existe aucun « dehors » depuis lequel la simulation pourrait être décrite intégralement, parce que celui qui décrit fait partie de la description.

Ce n'est pas de la magie. C'est une conséquence structurelle de l'auto-référence. Quand un processus se modélise lui-même, la distinction entre le modèle et le modélisé s'effondre. Le processus d'auto-modélisation et l'expérience d'être un soi ne sont pas deux choses distinctes qu'il faudrait relier par un pont --- ce sont une seule et même chose, décrite dans des vocabulaires différents.

Le Problème Difficile réclame un pont entre le traitement physique et l'expérience. La Théorie des Quatre Modèles répond : il n'y a pas de pont, parce qu'ils n'ont jamais été séparés. L'expérience EST l'auto-simulation, vue de l'intérieur de la boucle.

Il s'agit en fin de compte d'une affirmation d'identité (le genre d'affirmation qui, en science, marque un point d'arrêt plutôt qu'une lacune). « L'eau est H$_2$O » est une identité. On ne peut pas demander de façon sensée : « Mais \textit{pourquoi} l'eau est-elle H$_2$O ? » --- l'identité \textit{est} l'explication. Demander quelque chose de plus profond, c'est demander un autre genre d'univers. De même : l'expérience est ce que l'auto-simulation à quatre modèles, à la criticalité, \textit{est}. Si quelqu'un demande « Mais \textit{pourquoi} cette auto-simulation fait-elle l'effet de quelque chose ? », la réponse est : parce que c'est ce que ce processus \textit{est}. L'identité est falsifiable --- si les prédictions de la théorie échouent, l'identité est fausse. Mais elle ne peut pas être « expliquée davantage », pas plus que l'identité moléculaire de l'eau ne peut l'être. Elle est le point d'arrêt.

\section*{Pourquoi la simulation ne peut pas tourner dans le noir}

Il y a ici une question plus profonde, et y répondre révèle quelque chose d'essentiel sur la raison pour laquelle la conscience \textit{ressent}. Admettons que le cerveau fasse tourner une auto-simulation. Admettons l'architecture à quatre modèles, la criticalité, la clôture auto-référentielle. Tout cela ne pourrait-il pas se produire sans que cela fasse \textit{le moindre effet} ? La simulation ne pourrait-elle pas évaluer, modéliser, prédire --- et ne rien ressentir ?

C'est l'intuition du zombie sous un habillage technique. La réponse est non, et comprendre pourquoi révèle le trait le plus important de l'architecture.

Le substrat déploie la simulation virtuelle comme son mécanisme d'évaluation. C'est le sens premier de ces échanges : le système implicite présente des situations à la simulation pour que celle-ci puisse en évaluer les conséquences et en enregistrer les issues. Mais pour que cette évaluation fonctionne, les états simulés doivent avoir une \textit{valence} --- ils doivent compter pour la simulation. Un signal de douleur qui n'est qu'un nombre ne déclenche aucun évitement au niveau de la simulation. Seule une simulation qui \textit{se soucie} des issues peut les évaluer.

Rappelez-vous le jumeau numérique du chapitre 2 (la simulation d'ingénierie d'un réacteur d'avion). Un jumeau numérique typique ne se contente pas de refléter passivement le réacteur. Il \textit{ajoute} une couche de visualisation : avertissements, indicateurs à code couleur, alarmes --- des choses qui n'existent pas dans le réacteur physique. Le réacteur présente une fatigue du métal ; le jumeau affiche un avertissement rouge clignotant. Le réacteur voit sa température monter ; le jumeau a une jauge qui passe du vert à l'orange, puis au rouge. Cette couche ajoutée est tout l'intérêt. Sans elle, le jumeau n'est qu'un tableur --- des nombres posés inertes en mémoire, techniquement exacts, fonctionnellement inutiles. C'est la visualisation qui fait de la simulation un \textit{outil d'évaluation}.

Votre cerveau fait la même chose, en plus poussé. La simulation consciente ne se contente pas de refléter le traitement du substrat. Elle \textit{ajoute} la valence phénoménale. Douleur, plaisir, urgence, curiosité, effroi, ravissement --- voilà l'équivalent, pour le cerveau, des voyants d'alerte et des indicateurs de tableau de bord. Ils n'existent pas au niveau du substrat (les neurones ne ressentent pas plus la douleur que le métal ne ressent la fatigue). Ils existent au niveau de la simulation, ajoutés \textit{par} la simulation pour que le système puisse évaluer d'un coup d'œil des situations complexes. Le substrat a besoin de la simulation pour jauger des scénarios inédits, ambigus --- ceux où les réflexes ne suffisent pas. Et pour que cette évaluation fonctionne, le soi simulé doit enregistrer une valence hédonique : menace, opportunité, conséquence. Cet enregistrement --- ce fait de \textit{compter} --- est la phénoménalité. Retirez les qualia et vous retirez l'évaluation --- comme arracher l'affichage du tableau de bord d'un cockpit et attendre du pilote qu'il vole en lisant des tensions de capteurs brutes.

« Mais un système d'apprentissage par renforcement possède des signaux de récompense qui pilotent son comportement », pourriez-vous objecter. « Ressent-il quelque chose ? » Non --- parce qu'il lui manque l'architecture des quatre modèles à la criticalité. Un signal de récompense en apprentissage par renforcement est une valeur scalaire dans un système de Classe 1 ou de Classe 2. La valence phénoménale est l'enregistrement de la conséquence par l'ESM au sein d'une auto-simulation complète tournant en dynamique de Classe 4 --- un processus qualitativement différent. La différence n'est pas de degré. Elle est d'architecture.

La simulation ne peut pas tourner dans le noir, parce que l'obscurité irait à l'encontre de sa raison d'être. La phénoménalité n'est pas une fonction bonus de la conscience. Elle est le mécanisme par lequel la simulation accomplit son travail.

\section*{Ce que cela n'est pas : l'illusionnisme}

Ceci n'est pas de l'illusionnisme. Et la distinction est assez importante pour qu'on la formule sans détour.

Il existe une position philosophique respectable appelée illusionnisme, associée à Daniel Dennett et Keith Frankish, selon laquelle les \textit{qualia} sont des illusions. Dans cette perspective, cela ne fait aucun effet de voir du rouge. L'apparence de l'expérience est elle-même une fiction --- une histoire que le cerveau raconte, sans aucune réalité expérientielle derrière elle. La conscience, au sens fort, n'existe pas. Elle semble seulement exister.

Réfléchissez à ce que cette thèse affirme réellement. Vous ressentez quelque chose en ce moment même --- de la curiosité pour cet argument, du scepticisme, le poids du livre entre vos mains. L'illusionnisme dit que ce ressenti est une illusion. Vous ne faites en réalité l'expérience de rien. Quand vous dites « je ressens quelque chose », vous vous trompez, selon cette théorie. Votre propre témoignage sur votre propre expérience est faux. Vous mentez, en somme --- sauf qu'il n'y a pas de « vous » pour mentir. Si cela vous paraît manifestement ridicule, je suis d'accord.

La Théorie des Quatre Modèles affirme le contraire.

Les \textit{qualia} sont réels. Ils sont réels au sein de la simulation. Ils sont le mode par lequel le soi virtuel perçoit le monde virtuel. Lorsque votre Modèle Explicite du Soi enregistre la représentation d'une pomme rouge par votre Modèle Explicite du Monde, cet enregistrement (ce « voir du rouge ») est une propriété authentique du processus virtuel. Il existe au niveau de la simulation, exactement comme une balle qui touche un personnage de jeu vidéo lui fait \textit{mal}. Pas métaphoriquement --- au sein du jeu, les dégâts sont réels. La jauge de vie baisse, le personnage titube, le monde réagit. Vu de l'extérieur, c'est un nombre qui se décrémente en mémoire. Vu de l'intérieur du jeu, c'est de la douleur. Voilà la différence de niveaux. Et voilà où vivent vos \textit{qualia}.

La théorie opère avec une ontologie à deux niveaux. Le niveau du substrat (les neurones, les synapses, les modèles implicites) n'a aucune expérience. La lumière y est éteinte. Le niveau de la simulation (les modèles explicites, le monde virtuel et le soi virtuel) a une expérience authentique. La lumière y est allumée. Les deux niveaux sont physiques. Aucun n'est une illusion. Ce sont des niveaux différents du même système physique, avec des propriétés différentes à chaque niveau.

La théorie ne dit pas que votre douleur est une illusion. Elle dit que votre douleur est réelle --- simplement, elle est réelle dans la simulation, pas dans les neurones. Et puisque vous vivez votre vie entière à l'intérieur de la simulation, c'est la seule sorte de réel qui compte pour vous.

C'est là la distinction cruciale. Passez à côté, et vous confondrez cette théorie avec l'éliminativisme, avec l'illusionnisme, avec tous les autres cadres qui tentent d'expliquer la conscience en l'escamotant. La Théorie des Quatre Modèles n'escamote pas la conscience. Elle explique où vit la conscience --- et il se trouve que c'est exactement là où vous vous tenez depuis toujours.

\section*{Ce que signifie « réel au sein de la simulation »}

Il y a ici une subtilité philosophique qui mérite d'être dépliée. Quand je dis que les \textit{qualia} sont « réels au sein de la simulation », vous pouvez entendre l'une de deux choses. Soit ils sont \textit{authentiquement phénoménaux} --- auquel cas je n'ai fait que déménager le mystère des neurones vers la simulation, et le Problème Difficile continue de vivre à une autre adresse. Soit ils sont \textit{fonctionnellement réels mais pas authentiquement phénoménaux} --- auquel cas c'est du Dennett en plus compliqué.

C'est une fausse dichotomie. Elle ne tient que si l'on maintient qu'il existe un point de vue divin depuis lequel trancher si quelque chose est « authentiquement » phénoménal --- une perspective extérieure capable de vérifier si la simulation ressent vraiment ou fait seulement comme si. Or la clôture auto-référentielle élimine précisément cette perspective extérieure. L'ESM est son propre observateur. Il n'existe aucun poste d'observation externe d'où demander « mais ressent-il \textit{vraiment} ? ». Le questionnement fait lui-même partie du processus.

« Authentiquement phénoménal » contre « simplement fonctionnel » présuppose que la phénoménalité est une propriété qu'un processus possède ou ne possède pas, vérifiable par un observateur indépendant. Pour un système pleinement auto-référentiel à la criticalité, un tel observateur n'existe pas. La question se dissout --- non parce qu'elle serait sans réponse, mais parce qu'elle ne peut pas être posée. Elle exige une perspective que la clôture auto-référentielle rend impossible.

C'est l'argument le plus fort dont dispose le physicalisme processuel, et c'est la position que Thomas Metzinger esquisse avec son concept de « transparence phénoménale » --- même si la Théorie des Quatre Modèles est plus explicite sur le \textit{pourquoi} de cette transparence. C'est la frontière implicite-explicite qui crée la transparence : vous ne pouvez pas voir à travers elle, donc vous ne pouvez pas sortir de votre propre phénoménalité pour demander si elle est « authentique ». Cette frontière n'est pas un défaut. Elle est la raison pour laquelle la question de l'authentique contre le simplement fonctionnel ne s'applique pas aux systèmes comme vous.

Je ne suis pas le premier à aboutir ici, et je me méfierais de moi-même si je l'étais.

Le chercheur en sciences cognitives Joscha Bach a formulé la même idée centrale avec autant de netteté que quiconque. « Votre expérience est une simulation, écrivait-il, une réalité virtuelle projetée dans un observateur virtuel. » Et cet observateur ? C'est « un modèle de ce que ce serait si vous existiez, si vous étiez téléversé dans un singe et profondément investi dans ses affaires » --- un modèle construit pour conduire le singe à travers le monde physique. Une simulation, un quelqu'un simulé à l'intérieur pour recevoir la simulation, et un corps que tout cet arrangement dirige. Nous avons déjà rencontré dans ce livre les branches plus anciennes de cette tradition --- le modèle du soi de Metzinger, le soi narratif de Dennett.

Ce livre, lui, descend encore d'un niveau --- dans l'architecture elle-même. Quels modèles, exactement : le 2×2 du monde et du soi, de l'implicite et de l'explicite --- les quatre modèles. La frontière nette qui décide à partir de quand un système possède ne serait-ce qu'un intérieur : la clôture auto-référentielle, la boucle qui se retourne sur elle-même. Et la condition qui permet à l'ensemble de prendre vie : la criticalité, le bord du chaos où le système n'est ni figé ni bruit blanc.

Bach vous dit \textit{ce qu'}est la conscience. Moi, je vous dis de quoi elle est \textit{faite} et \textit{quand elle s'allume}. Même destination. J'ai simplement apporté les plans.

\section*{Pourquoi le mystère persiste}

Même une fois le Problème Difficile dissous, il reste une question tenace qui taraude. Si la réponse est si nette, pourquoi la conscience \textit{paraît}-elle encore si mystérieuse ? Pourquoi le Problème Difficile semble-t-il difficile même après qu'on vous a exposé la solution ? David Chalmers appelle cela le « méta-problème de la conscience » --- le problème d'expliquer pourquoi nous \textit{pensons} qu'il existe un problème difficile.

La Théorie des Quatre Modèles a une réponse nette, et elle découle tout droit de l'architecture.

Et c'est là que cela devient étrange : le « vous » conscient (le soi virtuel) ne peut pas voir la machinerie qui l'engendre. Vous ne pouvez pas plus introspecter vos propres poids synaptiques qu'un personnage de rêve ne peut examiner le cerveau du rêveur. Le système qui crée votre expérience est, par sa nature même, invisible à votre expérience. Non parce que quelqu'un le cache, mais parce qu'il opère à un niveau que votre expérience n'inclut pas.

Voyez les choses ainsi. Vous êtes un personnage de jeu vidéo --- un très bon jeu, avec une pleine conscience de soi à l'intérieur du monde du jeu. Vous voyez les montagnes rendues à l'écran, vous entendez le vent rendu, vous sentez sous vos pieds le sol rendu. Mais vous ne voyez presque jamais le moteur graphique. Vous n'entrevoyez presque jamais le code source. Le processus de rendu opère à un niveau que le monde du jeu, en temps normal, n'inclut pas. Je dis « presque », parce qu'il arrive que des artefacts s'infiltrent. Dans votre cerveau, cela se produit aussi --- les psychédéliques ouvrent la frontière, les états de flux l'amincissent, et même dans la vie ordinaire vous pouvez en saisir des aperçus : la tache aveugle que votre cerveau comble, les phosphènes quand vous vous frottez les yeux, les motifs géométriques derrière vos paupières closes. Ce ne sont pas des bugs. Ce sont des moments où le traitement du substrat devient brièvement visible depuis l'intérieur de la simulation. Nous explorerons cela en détail au chapitre 6. Mais la plupart du temps, le processus de rendu reste caché au monde rendu.

C'est exactement la situation de l'ESM. Quand le soi conscient tente de comprendre le fondement de sa propre expérience, il se heurte à une opacité de principe --- non pas une lacune du savoir actuel, mais une caractéristique structurelle de l'architecture. Les modèles implicites qui engendrent la simulation ne font pas partie de la simulation. Ils ne le peuvent pas, pas plus que le GPU ne peut être une montagne dans le jeu.

Le résultat est prévisible. L'ESM, incapable d'observer son propre substrat, en conclut que le mécanisme de la conscience doit être non physique, ou fondamentalement inexplicable, ou d'une manière ou d'une autre hors de portée de la science. C'est l'origine du dualisme. C'est le « fossé explicatif ». C'est l'intuition persistante que quelque chose est « laissé de côté » dans toute explication physique de la conscience --- parce que, depuis l'intérieur de la simulation, quelque chose \textit{est} effectivement laissé de côté. Le substrat. Cela même qui engendre l'expérience est invisible à l'expérience qu'il engendre.

Le mystère est réel, mais c'est un artefact de l'architecture, pas la preuve de quelque chose de non physique. Et il y a une raison pour laquelle vous le \textit{ressentez} comme mystérieux. Vous êtes un processus virtuel tournant sur du matériel biologique, et la plupart du temps, la frontière entre vous et votre substrat est opaque. Mais pas toujours. Parfois --- dans des états modifiés de conscience, dans des moments de concentration extrême, du coin de l'œil --- vous entrapercevez la machinerie en dessous. Pas clairement, pas entièrement, mais assez pour pressentir que quelque chose d'immense se joue sous la surface de votre expérience. Ce sentiment d'étrangeté, cette impression que la conscience est d'une certaine façon plus profonde que ce que vous pouvez atteindre --- voilà l'effet que cela fait d'être une simulation qui voit presque, mais pas tout à fait, à travers son propre rideau.

C'est une \textit{prédiction} de la théorie, pas un fil qui dépasse. Si vous êtes une simulation dont la frontière avec son propre substrat est largement opaque, vous devriez justement vous \textit{attendre} à ce que la conscience paraisse aussi étrange et irréductible qu'elle le paraît. La force intuitive du Problème Difficile ne vient pas de ce que la conscience serait véritablement inexplicable. Elle vient de notre position architecturale --- nous sommes à l'intérieur de la simulation, à épier par les fissures.

\section*{Le soi qui se recoud lui-même}

Revenez maintenant à ce blanc, l'architecture en main.

La demi-seconde de « vide », de personne aux commandes n'était pas un dysfonctionnement. C'était la reconstruction surprise en plein élan. L'identité n'est pas une propriété fixe du substrat --- c'est une \textit{reconstruction}, assemblée à neuf chaque matin à partir du modèle du soi stocké. Et le substrat à partir duquel elle se rebâtit n'est jamais tout à fait celui sur lequel vous vous êtes endormi. Des rêves dont vous ne gardez aucun souvenir ont recâblé vos poids synaptiques ; la consolidation a réorganisé vos souvenirs pendant la nuit. Vous vous réveillez sans être tout à fait la personne qui s'est couchée. La différence est d'ordinaire si infime que vous ne la remarquez jamais --- mais elle est là, chaque jour, sans exception.

Deux choses maintiennent « vous » en continuité à travers cette dérive : le Modèle Implicite du Soi, qui ne change que lentement, et le sommeil, qui dissimule le changement en mettant la simulation hors ligne pendant qu'il s'opère. Poussez la dérive assez loin --- réécrivez l'ISM en une nuit, remplacez les souvenirs, remodelez la personnalité --- et l'ancien « vous » ne disparaîtrait pas pour autant. Il serait \textit{absorbé}. Au matin, le Modèle Explicite du Soi reconstruirait un récit continu à partir de ce qui aurait survécu, liant l'ancien et le nouveau en une seule histoire, sans heurt, sans jamais remarquer la couture. L'ESM ne fait pas de ruptures nettes. Il recoud, toujours. Seul un effacement total rompt le fil ; tant qu'il reste quelque chose, le « vous » de demain inscrit sans bruit le « vous » d'aujourd'hui dans son histoire et nomme cela une vie.

Voilà donc ce que vous êtes : un soi virtuel, réassemblé chaque nuit, qui ressent son propre monde depuis l'intérieur d'une boucle dont il ne peut pas s'extirper. Reste une question. J'ai dit que la simulation tourne sur un substrat maintenu dans un état bien particulier --- mais quel \textit{est} cet état ? Quel fil de rasoir physique permet à un ou deux kilos de cellules câblées de faire démarrer un quelqu'un chaque matin, et de s'éteindre chaque nuit ? C'est la pièce qui fait tourner l'ensemble, et elle vient maintenant.


\chapter*{Chapitre 5 : Au bord du chaos}
\addcontentsline{toc}{chapter}{Chapitre 5 : Au bord du chaos}
\markboth{Chapitre 5 : Au bord du chaos}{}

Jusqu'ici, je vous ai dit à quoi ressemble l'architecture (quatre modèles, deux axes, une simulation tournant sur un substrat). Je vous ai dit où réside l'expérience (du versant virtuel, dans les modèles explicites). Et je vous ai dit ce qu'est l'identité (une reconstruction, réassemblée à neuf chaque matin à partir de modèles implicites stockés).

Mais je ne vous ai pas dit ce qui fait \textit{tourner} tout cela. Pourquoi la simulation est-elle parfois allumée et parfois éteinte ? Quelle propriété physique distingue un cerveau conscient d'un cerveau inconscient ? Pourquoi le sommeil profond efface-t-il la simulation alors que l'architecture reste intacte ?

Il manque une dernière pièce au puzzle --- et c'est celle qui m'a vraiment convaincu que la théorie fonctionne.

L'architecture à quatre modèles est nécessaire à la conscience, mais elle n'est pas suffisante. Il faut aussi la bonne \textit{dynamique}. Concrètement, le substrat --- le système physique qui exécute la simulation --- doit fonctionner dans ce que les mathématiciens et les physiciens appellent le \textbf{bord du chaos}.

En 2002, le polymathe Stephen Wolfram a publié \textit{A New Kind of Science}, où il classait les systèmes de calcul en quatre types selon leur dynamique. Le schéma de Wolfram a, selon moi, besoin d'une cinquième classe --- il mettait dans le même sac les systèmes fractals et les systèmes véritablement chaotiques, alors qu'ils sont structurellement distincts. L'argument complet se trouve dans l'annexe C, pour les lecteurs qui veulent les détails mathématiques. Ici, l'essentiel est ceci :

Les systèmes de calcul se répartissent sur un spectre allant de l'ordre parfait au désordre parfait. À une extrémité, les systèmes statiques et périodiques, trop simples pour calculer quoi que ce soit d'intéressant. À l'autre extrémité, les systèmes chaotiques, trop désordonnés pour qu'aucun motif stable puisse se former. Entre les deux, au \textbf{bord du chaos}, se trouvent les systèmes capables de calcul universel : assez complexes pour produire un comportement riche, varié, imprévisible, mais assez ordonnés pour que ce comportement persiste et interagisse. Le Jeu de la Vie de Conway en est l'exemple canonique --- le même automate cellulaire que j'avais programmé sur un 286 quand j'étais enfant. Trois règles d'une simplicité absolue sur une grille plate, et pourtant elles engendrent des planeurs, des oscillateurs, des structures auto-réplicatives et (de façon démontrable) le calcul universel. On peut y construire un ordinateur. On peut construire un ordinateur à l'intérieur de cet ordinateur. En principe, on peut faire tourner tout un monde virtuel tridimensionnel dans une grille de pixels à deux dimensions. À partir de presque rien, tout.

Je veux vous dire d'où vient cet exemple, car il est la graine de tout ce livre. Le Jeu de la Vie fut le premier vrai programme que j'aie jamais écrit --- rester assis là, enfant, à regarder trois règles de jouet assembler un ordinateur fonctionnel dans une grille de pixels morts et vivants était, à l'époque, la chose la plus importante que j'aie jamais apprise. Qu'un monde de dimension supérieure puisse surgir d'un jeu de règles de dimension inférieure semblait devoir être impossible, et le fait que ça ne l'était manifestement pas s'est logé en moi pour ne plus jamais me quitter. Je ne savais pas encore que je consacrerais ma vie à cette seule intuition, ni que le feuillet plissé de cortex derrière mes propres yeux se révélerait accomplir exactement le même tour, une dimension plus haut. C'est là que nous allons.

C'est là que réside la conscience. Seules les dynamiques du bord du chaos possèdent les deux propriétés requises : le \textbf{calcul universel} (assez complexe pour exécuter réellement une auto-simulation) et l'\textbf{intégration globale} (des parties éloignées du système s'influencent mutuellement, les changements locaux se propagent globalement, l'information est liée en un tout unifié). Voilà pourquoi l'expérience consciente se ressent comme \textit{unifiée} --- vous ne voyez pas le rouge d'un côté et n'entendez pas une voix de l'autre comme des flux séparés. La dynamique critique lie tout en une seule expérience. La liaison n'est pas quelque chose que le cerveau ferait \textit{en plus de} ses autres calculs ; elle est une conséquence du régime dynamique.

Un cerveau en sommeil profond, parcouru de lentes ondes delta, fonctionne selon une dynamique périodique : répétitive, tournant en rond. Les modèles sont toujours là, dans le substrat, mais la simulation ne tourne pas. Un cerveau en crise généralisée est poussé hors du bord critique par une vague de co-activation qui s'effondre en un seul régime grossier --- hypersynchronie rigide, chaos débridé, ou les deux tour à tour --- et la simulation ne parvient plus à tenir ensemble. Ce n'est que dans l'état de veille --- en équilibre au bord du chaos --- que le système soutient l'expérience consciente.

Le cerveau, en tant qu'ordinateur universel optimisé par des milliards d'années d'évolution, utilise \textit{tous} les régimes de calcul comme autant d'outils distincts : des attracteurs stables pour la mémoire à long terme, des oscillations périodiques pour la temporisation et le filtrage (rythmes alpha, thêta, gamma), un traitement fractal pour la reconnaissance invariante d'échelle et l'analyse des textures (principalement en V2-V4 du cortex visuel), et des dynamiques de bord du chaos pour l'automate cortical lui-même (le moteur de la conscience). Seul le régime du bord du chaos engendre la conscience. Mais la conscience dépend des autres pour fonctionner.

Lorsque j'ai publié cet argument dans mon livre de 2015, j'ignorais totalement que les neurosciences empiriques s'acheminaient, de leur côté, vers la même conclusion. D'ailleurs, le cadrage en automate cellulaire était la partie de toute la théorie qui me laissait le plus incertain. Je ne trouvais à l'époque aucun appui empirique, et j'ai bien failli ne pas l'inclure dans le livre --- cela me semblait un pas de trop, une affirmation qui rendrait le reste de la théorie facile à balayer. Je l'ai gardée parce que la logique me paraissait inéluctable, non parce que j'avais des preuves.

Mais il y a une subtilité cruciale. La criticalité seule ne suffit pas. Un tas de sable qui s'éboule à son angle critique se tient précisément au bord du chaos. Il n'est pas conscient. La théorie exige que \textit{deux} seuils soient franchis : le seuil physique (le substrat doit fonctionner à la criticalité) et le seuil fonctionnel (le substrat doit réaliser l'architecture à quatre modèles). La criticalité sans l'architecture vous donne des dynamiques complexes, mais aucune conscience. L'architecture sans la criticalité vous donne un système dormant --- les modèles existent dans le substrat, mais la simulation ne tourne pas. Les deux seuils doivent être franchis. Ensemble, ils sont suffisants.

La façon dont je me représente vraiment le cortex, c'est un océan --- mais un océan dont l'eau est câblée : non pas une nappe plate où chaque goutte touche les gouttes voisines, mais un réseau, chaque point relié à ses propres voisins particuliers par les poids synaptiques appris. Jetez-y un caillou --- un stimulus, une pensée --- et une ondulation se propage le long du câblage. Et sur cet océan en forme de réseau, tout comme sur un étang, des motifs peuvent se former : des ondulations qui se renforcent mutuellement, des ondes stationnaires, de petites formes autonomes qui voyagent, se percutent et persistent. Ces formes sont les planeurs. C'est là que vit une pensée, ou un soi --- comme un motif stable chevauchant l'océan, et non comme une quelconque goutte d'eau.

\section*{Deux cadrans, pas un seul}

Il y a dans le mot \textit{criticalité} une subtilité que j'ai mis longtemps à voir clairement --- et une fois qu'on la voit, plusieurs énigmes se dissolvent d'un coup. Quand les gens disent qu'un cerveau est « à la criticalité », ils fondent le plus souvent deux choses différentes en un seul mot. Sur l'océan en forme de réseau, ce sont deux cadrans distincts.

\textbf{Le premier cadran demande : jusqu'où voyage une ondulation ?} Tournez-le trop bas et l'océan devient un miroir plat --- jetez un caillou, vous obtenez une ondulation morte qui s'éteint sur place. Rien ne se propage, rien ne s'intègre. Tournez-le trop haut et un seul caillou déchaîne tout l'océan en tempête --- une unique perturbation submerge tout d'un coup. Le point idéal se situe entre les deux : une ondulation voyage, en engendre à peu près une autre, ne s'essouffle ni n'explose. Un tas de sable exactement à l'angle où un grain de plus peut déclencher un glissement de n'importe quelle ampleur. Ce cadran concerne la \textit{portée} --- quelle part de l'océan un seul événement entraîne dans un même processus connecté. Appelons-le le cadran d'\textbf{ÉTENDUE}.

\textbf{Le second cadran demande : quelles formes peuvent tout simplement vivre sur l'océan ?} C'est le cadran du Jeu de la Vie, celui dont dépend toute la théorie. Un miroir gelé n'a aucune forme ; rien ne bouge. Une houle douce et régulière n'a qu'une seule forme, la même vague encore et encore --- jolie, mais sans intérêt ; elle ne calcule rien. Une ébullition bouillonnante fait scintiller chaque point au hasard --- le chaos, qui ne calcule rien non plus. Ce n'est qu'à la frontière entre l'ordre et l'ébullition que l'on obtient des \textit{planeurs} : des motifs stationnaires qui se déplacent, se percutent, persistent et font des choses. C'est là qu'un calcul --- et un modèle du soi --- peut exister comme forme stable sur l'eau. Appelons-le le cadran de \textbf{COMPLEXITÉ}.

L'image la plus nette que je connaisse pour ce second cadran est celle d'une aile. Une aile d'avion produit le maximum de portance juste à la limite --- l'angle le plus abrupt qu'elle puisse tenir, l'instant avant que le flux d'air ne décroche et que l'aile ne bascule dans la turbulence. Un cheveu au-delà de cette limite, l'écoulement lisse se disloque et la portance s'effondre. Le cerveau vole exactement à ce bord d'avant-décrochage : calcul maximal, un degré avant que la dynamique ne bascule dans le chaos. Poussez au-delà et la « portance » --- la capacité à porter des planeurs, à calculer, à être conscient --- s'effondre comme s'effondre la portance dans un décrochage.

La plupart du temps, les deux cadrans montent et descendent de concert, c'est pourquoi un seul mot suffit d'ordinaire. Mais ils peuvent se dissocier, et quand ils le font, ils expliquent les cas que le récit à un seul bouton ne peut expliquer. Une casserole d'eau bouillante a le premier cadran poussé à fond --- une perturbation se répand partout --- mais le second bloqué à zéro : pas de planeurs, juste du clapot. Beaucoup de propagation, aucune forme, et aucune conscience, exactement comme on s'y attendrait.

Une crise généralisée est le cas qui a fini par cerner ce que mesure réellement le premier cadran. Lors d'une crise, une fraction énorme du cerveau décharge d'un seul coup --- en nombre brut, plus « active » qu'à n'importe quel instant d'éveil normal. Si la conscience n'était que « beaucoup de neurones actifs », une crise serait l'état le plus conscient qui soit. C'est l'inverse. Ce flot d'activité s'effondre hors du régime du bord du chaos en une seule forme grossière --- toute la nappe emportée dans un rythme grossier à l'unisson, ou jetée dans un chaos débridé --- et le calcul différencié, porteur de planeurs, s'évanouit. Le premier cadran ne peut donc pas mesurer une simple co-activation ; il doit mesurer quelle part du cerveau est prise dans une véritable dynamique de bord du chaos. Lors d'une crise, malgré tout ce déchargement, ce nombre s'effondre. Beaucoup de cerveau actif n'est pas la même chose que beaucoup de cerveau conscient. Gardez ces deux cadrans à l'esprit ; les états les plus étranges dans lesquels un cerveau peut entrer ne sont que différents réglages de ces cadrans.

\section*{L'automate cortical}

Il est temps de rendre concret quelque chose qui vous paraît peut-être encore abstrait. J'ai parlé du cortex qui doit fonctionner au bord du chaos, dans des dynamiques de Classe 4. Mais qu'\textit{est} ce système de Classe 4 ? Ce n'est pas quelque force mystérieuse planant au-dessus du cerveau. C'est le motif même de la décharge neuronale.

Songez à quoi ressemble réellement le cortex en fonctionnement. Des milliards de neurones, chacun déchargeant ou non, chacun influençant ses voisins par des poids de connexion appris. Chaque neurone est une cellule d'un automate cellulaire, non pas métaphoriquement, mais littéralement. Les règles de l'automate sont les poids synaptiques, les seuils, le câblage local. La sortie de chaque « cellule » est une fréquence de décharge. Et le résultat --- le grand motif d'activité électrique qui danse sur la surface corticale à 10 à 40 Hz --- est un automate cellulaire de Wolfram de Classe 4 opérant dans un espace de plusieurs milliers de dimensions.

J'appelle cela l'\textbf{automate cortical}.

C'est la même idée que celle que j'ai programmée sur un 286 quand j'étais enfant (le Jeu de la Vie de Conway), sauf qu'au lieu d'une grille plate à trois règles, c'est un feuillet plissé de cortex aux milliards de règles variant localement, et qu'au lieu de se déplacer en deux dimensions, ses motifs se déplacent dans un espace dimensionnel si vaste qu'il défie toute visualisation. Tel un poulpe aux bras sans limites, l'automate cortical peut atteindre n'importe quelle partie du cortex à tout instant, activant les modèles stockés dont il a besoin --- un souvenir ici, un programme moteur là, un fragment de langage ailleurs. Il saisit ces modèles comme de petites figurines Lego et s'en sert pour naviguer d'un état satisfaisant au suivant.

Et voici la distinction cruciale : \textbf{l'automate cortical n'est pas la conscience}. Il est le moteur, non l'expérience. Le motif en apparence chaotique de milliards de neurones qui déchargent est en réalité un appareil extraordinairement sophistiqué qui calcule, pense et pilote un corps à travers une vie. Mais la conscience n'est qu'un \textit{effet} de cet appareil --- un effet qui naît de l'interaction entre l'automate et le cortex lorsque les conditions sont réunies. Lorsque l'automate balaie de façon synchrone des régions corticales appropriées, à la bonne fréquence et dans une séquence temporelle cohérente, une expérience consciente émerge de cette suite d'images. L'automate contient les instances de notre modèle du monde et de notre modèle du soi ; la conscience est ce qui se produit quand ces modèles tournent activement dans la simulation.

Vous pouvez, d'ailleurs, observer l'automate cortical directement --- pas besoin d'IRMf.

Essayez ceci : trouvez une pièce complètement obscure. Fermez les yeux. Attendez que les éventuelles images rémanentes s'estompent (cela prend environ 30 à 60 secondes si vous avez regardé quelque chose de lumineux). Au début, vous ne voyez rien, ou presque rien. Mais ensuite, si vous patientez et prêtez attention, vous commencerez à percevoir des points colorés qui scintillent sur fond d'obscurité.

La plupart des gens écartent cela d'un revers de main en le qualifiant de « bruit rétinien » --- des décharges aléatoires dans les cellules photoréceptrices de l'œil, réagissant à une pression ou à des phénomènes chimiques spontanés. Et si vous appuyez doucement sur votre paupière, vous pouvez effectivement déclencher ainsi des sensations visuelles localisées. Mais les points colorés que vous voyez dans le noir total ne sont \textit{pas} d'origine rétinienne. Ils sont trop organisés pour cela. Ce que vous voyez, c'est l'activité de repos de V1 (votre cortex visuel primaire), alimentée par une combinaison de signaux sensoriels résiduels et de projections descendantes issues de l'automate cortical lui-même. L'automate déroule ses dynamiques de base, et vous êtes en train de l'observer en temps réel.

Si vous continuez à regarder --- sans vous concentrer, mais en vous relâchant, en laissant votre attention se relâcher ---, quelque chose de remarquable se produit. La focalisation active supprime en réalité ces motifs ; c'est lorsque vous cessez d'essayer de voir que vous commencez à voir. L'automate se met à recruter davantage du système visuel pour interpréter et amplifier le peu de signal disponible. Les points scintillants se stabilisent en formes. Des motifs géométriques émergent : grilles, spirales, treillis. Puis des visages, déformés et mouvants. Puis des silhouettes. Puis, avec assez de patience (et je parle d'\textit{heures}, pas de minutes), des scènes entières --- des hallucinations élaborées, colorées, narratives, qui ne diffèrent en rien des rêves que vous faites chaque nuit.

C'est le même mécanisme qui se cache derrière les hallucinations hypnagogiques --- ces images vives qui vacillent dans votre esprit juste au moment où vous vous endormez. C'est l'automate cortical qui tourne avec une contrainte externe minimale, générant son propre contenu en activant des motifs stockés et en les projetant dans la simulation. La progression que vous vivez (du faible bruit aux hallucinations cohérentes) est une fenêtre directe sur le fonctionnement de l'automate : il commence par V1, le tout premier étage du traitement visuel, et recrute progressivement V2, V3 et des aires supérieures à mesure qu'il tente de donner un sens au moindre signal disponible. Quand aucun signal réel n'est disponible, il en \textit{génère} un. C'est la fuite due à la perméabilité, à l'œuvre. En l'absence de signal externe qui domine la simulation, le propre bruit de traitement du substrat devient visible. Vous n'hallucinez pas \textit{à partir de rien} --- vous voyez les motifs de repos du moteur graphique, l'équivalent neuronal de la neige sur un téléviseur mal réglé. À ceci près que cette neige a une structure, parce que la machinerie de traitement a une structure.

Vous pouvez aussi induire de cette façon une forme temporaire de synesthésie. Dans ma jeunesse, je m'en servais pour « voir la musique ». Si vous fermez les yeux et écoutez de la musique tout en laissant les motifs visuels venir à vous --- détendu, passif, sans forcer pour voir ---, les motifs se synchronisent peu à peu avec le rythme et les fréquences de ce que vous entendez. L'automate cortical, privé d'entrée visuelle externe, se met à coupler ses dynamiques visuelles à quelque autre signal fort disponible --- en l'occurrence l'entrée auditive. Ce que vous voyez est, très littéralement, l'activité de votre cerveau rendue visible : les motifs de niveau V1 de l'automate, pilotés par le cortex auditif plutôt que par l'entrée rétinienne. Les vrais synesthètes --- ces personnes dont les sens sont câblés en permanence les uns aux autres, qui voient toujours des couleurs lorsqu'elles entendent des sons --- ont probablement une version plus permanente de ce même couplage, sans doute due à des connexions plus fortes ou plus nombreuses entre les aires sensorielles, que ce soit dans le thalamus ou dans le cortex lui-même. Le mécanisme est le même : une modalité sensorielle qui fuit dans le pipeline de traitement d'une autre. L'automate cortical se soucie assez peu de la provenance de son entrée. Il traite tout ce qu'il reçoit.

Je ne vous recommande pas d'en faire un passe-temps régulier. L'expérience peut être troublante, surtout si vous n'y êtes pas psychologiquement préparé. Et il existe un faible risque qu'une privation sensorielle prolongée déstabilise une personne présentant des vulnérabilités psychiatriques latentes. Mais si vous vous êtes déjà demandé à quoi ressemble le substrat de votre conscience lorsqu'il tourne au ralenti --- quand le monde extérieur s'est tu et que le système est simplement en train de… tourner ---, c'est l'aperçu le plus direct que vous puissiez en avoir sans scanner cérébral.

Cette progression, d'un presque-rien jusqu'à un monde visuel fictif complet, vécue par votre modèle du soi dans un univers virtuel, est un portrait direct de l'automate cortical à l'œuvre.

Quand l'automate se dérègle, vous pouvez le voir aussi. Une crise d'épilepsie, c'est ce qui arrive quand des parties de l'automate basculent dans des dynamiques de Classe 1 ou 2 (périodiques, verrouillées, sans utilité computationnelle), ou sont poussées au-delà de la Classe 4 dans le chaos de Classe 5. Un accident vasculaire cérébral, c'est ce qui arrive quand des parties du cortex décrochent complètement. Un évanouissement, c'est ce qui arrive quand la fréquence minimale requise pour l'éveil n'est plus atteinte. L'automate est plutôt fragile. Mais la structure qui l'engendre --- le néocortex, avec ses poids appris et son architecture façonnée par l'évolution --- est robuste, et c'est pourquoi nous nous remettons si remarquablement bien de ces perturbations.

\section*{La convergence}

En 2003 --- deux ans avant même que je n'aie la théorie ---, John Beggs et Dietmar Plenz découvrirent des « avalanches neuronales » dans le tissu cortical : des motifs d'activité neuronale qui suivaient la signature mathématique de la criticalité auto-organisée, une marque distinctive des systèmes au bord du chaos.

En 2014, Robin Carhart-Harris proposa l'Entropic Brain Hypothesis : l'idée que le niveau de conscience corrèle avec l'entropie (le désordre) de l'activité cérébrale, l'optimum se situant à un niveau intermédiaire --- trop peu d'entropie signifie l'inconscience, trop d'entropie signifie une expérience incohérente.

En 2016, Enzo Tagliazucchi et ses collègues montrèrent que le LSD pousse le cerveau vers la criticalité, ce qui concorde avec la conscience amplifiée (mais parfois chaotique) que rapportent les usagers de psychédéliques. Dès 2022, un article de synthèse pouvait déjà parler de « criticalité auto-organisée comme cadre pour la conscience » --- les preuves s'accumulaient.

Et en 2025-2026, la digue empirique céda. Keith Hengen et Woodrow Shew publièrent une méta-analyse de 140 jeux de données dans \textit{Neuron} (2025) --- la plus vaste analyse systématique de la criticalité dans les dynamiques cérébrales jamais menée ---, confirmant que le cerveau opère près d'un point critique sur de multiples modalités de mesure. Puis Inbal Algom et Oren Shriki proposèrent le cadre ConCrit (Consciousness and Criticality) dans \textit{Neuroscience \& Biobehavioral Reviews} (2026), soutenant que les dynamiques cérébrales critiques fournissent un fondement mécaniste unificateur pour toutes les grandes théories de la conscience. Leur conclusion : la conscience suit la criticalité. Quand le cerveau est au point critique ou à proximité, la conscience est présente. Quand il est poussé en dessous de la criticalité (par l'anesthésie, par le sommeil, par une lésion cérébrale), la conscience est absente. Quand il est écarté du point critique dans l'autre direction --- s'effondrant en une hypersynchronie verrouillée ou explosant en un chaos débridé, comme lors d'une crise et peut-être dans certains états provoqués par des drogues ---, la conscience devient incohérente.

Deux chemins. L'un théorique, partant du cadre computationnel de Wolfram et raisonnant sur ce qu'exige une auto-simulation. L'autre empirique, partant d'enregistrements neuronaux et analysant les propriétés statistiques de l'activité cérébrale pour chaque état de conscience accessible. À peine deux ans d'écart à l'origine, convergeant vers la même conclusion.

C'est le genre de convergence qui vous fait prendre une théorie au sérieux.

\section*{Trois façons dont un hologramme rencontre un automate}

En écrivant ce chapitre, j'ai pris conscience de quelque chose qui m'a glacé.

Le principe holographique et les automates de Classe 4 ne cessent de resurgir dans les mêmes conversations --- en physique, en neurosciences, en théorie du calcul. Mais personne ne semble avoir posé la question évidente : \textit{quelles sont les relations possibles entre eux ?}

Il y en a exactement trois.

\textbf{Relation 1 : un substrat holographique produit des dynamiques de Classe 4.} C'est probablement ce que fait le cerveau. Les réseaux de neurones sont localement holographiques --- l'hologramme en patchwork du chapitre 3, les rats de Lashley et tout le reste --- et ce substrat, opérant à la criticalité, produit les dynamiques de Classe 4 qu'exige la conscience. Bien établie, minutieusement documentée et --- pardonnez-moi --- la relation ennuyeuse.

\textbf{Relation 2 : un automate de Classe 4 qui produit des motifs holographiques comme comportement émergent.} L'automate n'est pas holographique dans ses règles, mais ses dynamiques engendrent spontanément des structures holographiques --- de l'information de dimension supérieure encodée dans des motifs de dimension inférieure, surgissant du calcul lui-même. Si un automate de Classe 4 produit naturellement une sortie holographique, cela signifie qu'une distribution non locale de l'information émerge de règles purement locales --- ce qui, fait fascinant, ressemble exactement à l'intrication quantique.

C'est ici que je devrais mentionner Gerard 't Hooft, car le lien est trop frappant pour être passé sous silence --- même s'il reste spéculatif. 't Hooft, lauréat du prix Nobel de physique, a proposé que la mécanique quantique elle-même soit un automate cellulaire à l'échelle de Planck : que notre univers est fondamentalement déterministe, et que les effets quantiques sont des phénomènes émergents d'une dynamique plus profonde et discrète. S'il a raison, le principe que je décris ne s'applique pas seulement à la conscience par analogie. C'est littéralement la manière dont l'univers fonctionne, jusqu'au niveau le plus fondamental. Des règles locales simples produisent un univers holographique, et à l'intérieur de cet univers, des règles neuronales simples produisent une conscience holographique. Le même principe computationnel opérant à deux échelles : cosmologique et neurologique. Je trouve cette cohérence fractale profondément convaincante, mais, en toute honnêteté, l'interprétation de 't Hooft demeure une opinion minoritaire en physique, et le raisonnement qui va de l'élégance structurelle à la réalité physique a été critiqué à juste titre. N'empêche --- s'il s'avérait qu'un seul principe computationnel sous-tend à la fois l'univers et les esprits qui le modélisent, ce serait l'une des plus belles découvertes jamais faites.

\textbf{Relation 3 : un automate de Classe 4 dont la structure même des règles est holographique.} C'est celle qui m'a fait poser mon stylo. Si une telle chose existe --- un automate cellulaire dont les règles elles-mêmes encodent de l'information de dimension supérieure dans une structure de dimension inférieure, comme un hologramme encode trois dimensions dans deux ---, alors vous auriez un système qui fait naturellement ce que le principe holographique dit que fait l'univers. Non pas un système qui \textit{s'exécute simplement sur} un substrat holographique (ou qui produit un hologramme). Un système qui \textit{est} un encodage holographique. Peut-être l'univers lui-même, d'ailleurs --- même s'il me faut noter que cela reste spéculatif, et que le raisonnement selon lequel la beauté mathématique impliquerait la réalité physique a fait l'objet de critiques légitimes. J'y reviendrai au chapitre 14, où j'expliquerai pourquoi je pense que la Relation 3 pourrait être l'une des plus importantes questions non résolues des mathématiques, avant d'en poursuivre pleinement la réponse dans les chapitres 15 à 17.


\chapter*{Chapitre 6 : Ce que révèlent les psychédéliques}
\addcontentsline{toc}{chapter}{Chapitre 6 : Ce que révèlent les psychédéliques}
\markboth{Chapitre 6 : Ce que révèlent les psychédéliques}{}

Une mise en garde nécessaire avant de commencer : rien dans ce chapitre ne doit être lu comme une incitation à essayer les psychédéliques. Ils sont puissants, imprévisibles, et peuvent ruiner votre vie --- littéralement, définitivement. Ils peuvent déclencher une schizophrénie chez les personnes prédisposées. Ils peuvent provoquer des épisodes psychotiques, des troubles anxieux persistants et un HPPD (hallucinogen persisting perception disorder) qui ne disparaît jamais. Si j'en parle ici, c'est parce qu'ils révèlent quelque chose d'important sur l'architecture de la conscience. Cette valeur scientifique ne les rend pas sûrs pour autant.

Si vous voulez comprendre la conscience, étudiez ce qui se passe quand elle se dérègle. Les psychédéliques sont, j'en suis convaincu, la fenêtre la plus révélatrice que nous possédions sur l'architecture de la conscience --- plus révélatrice que les scanners cérébraux de patients endormis, plus riche sur le plan théorique que les études de lésions, et infiniment plus accessible que la chirurgie du cerveau divisé.

Voici pourquoi : les psychédéliques ne se contentent pas de \textit{modifier} la conscience. Ils la modifient de \textit{façon systématique et prévisible}, ce qui met au jour l'architecture sous-jacente --- à condition de savoir quoi chercher.

\section*{Le gradient de perméabilité}

Rappelez-vous la frontière entre les modèles implicites et les modèles explicites. Dans la vie éveillée normale, cette frontière est sélectivement perméable : l'information pertinente passe, l'information non pertinente reste dans la bibliothèque. Vous êtes conscient de ce dont vous avez besoin, et inconscient de tout le reste.

Les psychédéliques font sauter la frontière.

Sous psychédéliques (LSD, psilocybine, DMT, mescaline), la perméabilité de la frontière implicite-explicite augmente globalement. L'information normalement traitée entièrement du versant réel, invisible pour la conscience, commence à s'infiltrer jusque dans la simulation.

Et voici le point crucial : elle s'infiltre \textit{dans l'ordre}.

À faibles doses, ou tôt dans l'expérience, les étapes de traitement les plus simples deviennent visibles en premier. Ce sont les étapes les plus proches de l'entrée sensorielle brute : le traitement de niveau V1. Vous voyez des couleurs avivées, des motifs qui respirent sur des surfaces statiques, de subtils mouvements en vision périphérique. Ce sont les détecteurs de traits précoces du cortex visuel, normalement invisibles, qui entrent maintenant dans la simulation.

À mesure que la dose augmente ou que l'expérience s'approfondit, des étapes de traitement plus complexes deviennent visibles. Le traitement de niveau V2/V3 : motifs géométriques, fractales, pavages, les fameuses « constantes de forme » que Heinrich Klüver catalogua dans les années 1920. Ce sont les représentations intermédiaires du système visuel --- les briques élémentaires qu'il utilise d'ordinaire pour construire votre expérience visuelle, désormais visibles pour elles-mêmes.

Plus haut encore, et les aires visuelles supérieures deviennent accessibles. Des visages apparaissent. Des figures. Des scènes. Les aires de traitement des visages, les aires de reconnaissance des objets, les aires de construction des scènes --- toutes fonctionnant normalement sous le seuil de la conscience --- diffusent maintenant leurs produits intermédiaires directement vers la simulation.

Aux doses les plus élevées, toute la hiérarchie de traitement est mise à nu, et le résultat est une expérience visionnaire à part entière : des scènes complexes, narratives, oniriques, construites à partir des couches les plus profondes du traitement implicite.

Cette progression ordonnée --- du simple au complexe, de V1 aux aires supérieures, dépendante de la dose --- est exactement ce que prédit la Théorie des Quatre Modèles (FMT). C'est une conséquence directe du gradient de perméabilité : les étapes de traitement de plus bas niveau, étant plus proches de la frontière, deviennent accessibles avant celles de plus haut niveau à mesure que la perméabilité augmente.

La hiérarchie du traitement visuel ci-dessous montre ce que fait chaque aire en temps normal et ce qui devient visible quand la barrière de perméabilité tombe :


{\small
\noindent
\begin{tabularx}{\linewidth}{X X X}
\toprule
\textbf{Aire} & \textbf{Fonction normale} & \textbf{Signature psychédélique} \\
\midrule
V1 & Bords, fréquence spatiale, orientation & Phosphènes, constantes de forme de Klüver, surfaces qui respirent \\[4pt]
V2 & Intégration des contours, texture, appartenance des bords & Pavages, motifs géométriques répétitifs \\[4pt]
V3 & Forme globale, traitement dynamique de la forme & Géométries fluides et changeantes \\[4pt]
V4 & Couleur, courbure, texture complexe & Fractales colorées, motifs kaléidoscopiques \\[4pt]
V5/MT & Traitement du mouvement & Rotation et déplacement de motifs \\[4pt]
Fusiform/IT & Visages, objets, formes de mots & Visages, figures, entités \\[4pt]
Anterior IT & Catégories sémantiques, construction de scènes & Hallucinations narratives complètes \\[4pt]
\bottomrule
\end{tabularx}
}


Chaque ligne représente une étape de traitement plus profonde. En conditions normales, vous n'éprouvez que la sortie finale --- le percept achevé. Sous psychédéliques, vous éprouvez les étapes \textit{intermédiaires}, dans l'ordre, à mesure que la perméabilité augmente. (Une version plus complète de ce tableau, avec la taille des champs récepteurs et des détails supplémentaires, figure dans l'Annexe A.)

Je sais que cela paraît fascinant. Vous lisez que des couches de traitement visuel deviennent visibles, et une part de vous se demande avec curiosité à quoi cela ressemble. Je comprends --- moi aussi, j'étais curieux. J'ai essayé les deux voies. J'étais jeune, bête, et j'ai eu de la chance. La voie de la méditation, que j'ai décrite au chapitre précédent (une pièce sombre, une attention détendue, de la patience), vous mène au même endroit. Pas aussi vite, pas aussi spectaculaire au premier essai. Mais tout aussi impressionnante, tout aussi réelle, et sans le risque d'abîmer définitivement votre esprit. Un lit chaud dans une pièce sombre suffit.

Il existe une autre voie : le rêve lucide. Si vous parvenez à reconnaître que vous rêvez alors que vous êtes encore dans le rêve --- et c'est une compétence qui s'apprend ---, vous accédez à la simulation entière tournant sans contrainte. Aucune entrée sensorielle, aucune réalité extérieure pour corriger le modèle. Rien que le monde virtuel, avec vous consciemment à l'intérieur. Pour certaines personnes, c'est plus facile à atteindre qu'une méditation soutenue. Les techniques sont bien documentées (voir l'Annexe D pour un guide pratique), et l'expérience peut se révéler au moins aussi éclairante que tout ce que produit une drogue --- sans le risque. Nous reviendrons au rêve lucide au chapitre 7.

\section*{Cinq systèmes emboîtés}


\begin{figure}[htbp]
  \centering
  \includegraphics[width=0.95\textwidth]{../figures/figure-five-layer-stack-bw-fr.png}
  \caption{Cinq niveaux d'organisation du cerveau. Chaque niveau survient sur (« tourne sur ») celui qui est en dessous. La conscience n'existe qu'au niveau virtuel le plus élevé, où les modèles explicites engendrent l'expérience phénoménale.}
\end{figure}


Imaginez votre cerveau comme doté de cinq niveaux d'organisation distincts, empilés comme des poupées russes :

\textbf{Physique.} Tout en bas, il y a la matière brute : atomes, molécules, le substrat physique du cerveau lui-même. C'est la chimie --- le carbone, l'hydrogène, l'azote, l'oxygène qui composent le tissu. De la matière inerte obéissant aux lois de la thermodynamique. Rien de conscient n'habite ici.

\textbf{Électrochimique.} Un niveau au-dessus : la signalisation neuronale. Des potentiels d'action qui dévalent les axones, des neurotransmetteurs qui inondent les synapses, des ions qui circulent dans des canaux. C'est l'activité électrique et chimique que tout le monde se représente en pensant « le cerveau fait quelque chose ». C'est le niveau où les neurones déchargent. Toujours pas d'expérience, mais désormais vous disposez d'une transmission d'information.

\textbf{Protéomique.} Ensuite : les structures protéiques et la machinerie moléculaire. Les poids synaptiques sont stockés ici --- les forces physiques des connexions entre neurones. Les récepteurs sur les membranes cellulaires, les enzymes qui régulent la plasticité, l'échafaudage moléculaire qui détermine quelles synapses se renforcent et lesquelles s'affaiblissent. C'est le « matériel » de l'apprentissage. Quand vous vous exercez à une compétence et que vous vous améliorez, vous modifiez la couche protéomique. Toujours inconscient, mais désormais vous disposez d'une mémoire.

\textbf{Topologique.} Plus haut encore : l'architecture du réseau. Les motifs de connectivité, quels neurones sont reliés à quels autres, avec quelle densité, selon quelles configurations. C'est là que vivent les aires de Brodmann, là que vivent les colonnes corticales, là qu'existe la structure à grande échelle du « cortex visuel qui dialogue avec le cortex moteur ». C'est le schéma de câblage. Modifiez ce niveau et vous modifiez les types de traitement dont le système est capable. C'est là que sont stockés vos modèles implicites (l'IWM et l'ISM). Toujours inconscient. Mais désormais vous disposez d'un savoir.

\textbf{Virtuel.} Tout en haut : le monde simulé. L'automate cortical --- le motif dynamique d'activité électrique qui danse à travers le réseau, intègre l'information, engendre des prédictions, fait tourner les modèles en temps réel. C'est là que vit votre expérience consciente. Les modèles explicites (l'EWM et l'ESM) existent ici et seulement ici. C'est le seul niveau auquel cela fait quelque chose d'être.

Chaque niveau survient sur celui qui est en dessous, mais possède sa propre dynamique. Pas de signalisation électrochimique sans matière physique, pas de structures protéiques sans chimie, pas de topologie de réseau sans synapses, et pas de simulation sans réseau pour la faire tourner. Mais chaque niveau possède des propriétés que les niveaux inférieurs n'ont pas. Une synapse n'est « à propos » de rien --- elle n'est qu'une connexion. Un réseau de synapses, lui, \textit{est} à propos de quelque chose : il représente un visage, un mot, un souvenir. Et la simulation qui tourne sur ce réseau ? C'est là que « à propos » devient « expérience ».

Cette hiérarchie à cinq niveaux résout un problème sur lequel presque tout le monde bute en entendant cette théorie pour la première fois : « Si la conscience est virtuelle, sur quoi tourne-t-elle ? » La réponse : elle tourne sur la couche topologique (le réseau), qui est implémentée dans la couche protéomique (les poids synaptiques), qui tourne sur la couche électrochimique (la décharge neuronale), qui existe dans la couche physique (la matière). La conscience n'est pas moins réelle du fait d'être virtuelle --- elle est simplement réelle \textit{à un autre niveau} que ne le sont les neurones. La montagne dans le jeu vidéo est réelle au niveau du jeu, même si elle n'est « que » des transistors au niveau du matériel. Le même principe.

C'est ici que la hiérarchie à cinq niveaux accomplit son travail d'explication. Les psychédéliques ciblent le milieu de la pile, et les effets se propagent vers le haut. Les psychédéliques classiques comme le LSD et la psilocybine se lient aux récepteurs de la sérotonine 2A, en agissant au niveau \textbf{électrochimique}. Ils changent la façon dont les neurones se parlent. Cette perturbation se propage au niveau \textbf{protéomique}, où la sensibilité des récepteurs se déplace au fil des heures. Elle remodèle le niveau \textbf{topologique}, où les motifs de connectivité du réseau se modifient --- visibles à l'IRMf sous forme d'une intégration globale accrue. Et elle transforme le niveau \textbf{virtuel}, où la simulation consciente se trouve inondée d'un contenu normalement invisible. Le seul niveau que les psychédéliques classiques ne touchent pas est le \textbf{physique} --- ils ne détruisent pas les neurones, n'altèrent pas la matière brute. Ils changent tout ce qui est \textit{au-dessus} de la matière, dans l'ordre ascendant. C'est une distinction cruciale. Les psychédéliques classiques (LSD, psilocybine, DMT, mescaline) ne sont pas neurotoxiques. Ils changent la façon dont les neurones communiquent sans les détruire. Beaucoup d'autres drogues sont moins clémentes. La méthamphétamine et l'alcool détruisent physiquement les neurones ; la cocaïne, elle, agit de façon plus indirecte, essentiellement en resserrant les vaisseaux sanguins et en privant le tissu nerveux d'oxygène. La MDMA, à doses élevées ou répétées, endommage les axones sérotoninergiques. Même l'Amanita muscaria --- l'emblématique champignon rouge et blanc que beaucoup de gens confondent avec les champignons psychédéliques --- est un délirogène qui agit par un mécanisme entièrement différent et plus dangereux. Si vous ne retenez qu'une chose de ce chapitre : toutes les drogues qui altèrent la conscience ne se valent pas, et la distinction entre « changer le signal » et « détruire le matériel » est littéralement la différence entre un état altéré temporaire et des lésions cérébrales permanentes. La progression visuelle dépendante de la dose se calque directement là-dessus : de faibles doses perturbent le niveau électrochimique juste assez pour affecter le traitement de V1 ; des doses plus fortes propagent la perturbation à travers davantage de niveaux, recrutant des étapes de traitement de plus en plus complexes dans l'expérience consciente.

\section*{Le soi réorientable}

Mais la preuve la plus spectaculaire vient de ce qui arrive au soi.

Votre Modèle Explicite du Soi (ESM) est un processus virtuel qui exige un apport. Dans des conditions normales, il reçoit un flux constant de signaux auto-référentiels : la perception de l'endroit où se trouve votre corps (proprioception), votre sens de ce que ressentent vos organes (intéroception), le flux narratif du discours intérieur, et la toile de fond constante d'une conscience corporelle de soi que vous ne remarquez jamais, jusqu'à ce qu'elle soit perturbée.

À fortes doses de psychédéliques, cet apport se trouve perturbé. Le modèle du soi ne meurt pas --- il \textit{se réoriente}. Privé de son apport auto-référentiel normal, il se saisit de l'apport dominant, quel qu'il soit.

La salvia divinorum le démontre de la manière la plus spectaculaire --- un psychédélique dissociatif agissant sur les récepteurs kappa-opioïdes (complètement différents des mécanismes sérotoninergiques du LSD ou de la psilocybine). Les usagers de la salvia rapportent invariablement des expériences où ils \textit{deviennent} des choses :

\begin{itemize}
  \item « Je suis devenu le canapé. »
  \item « J'étais le mur. »
  \item « Je me suis transformé en une page dans un livre. »
  \item « J'étais l'un des personnages à la télévision. »
  \item « Je suis devenu une fractale, non pas voyant une fractale, \textit{étant} une fractale. »
\end{itemize}

Ce ne sont pas des métaphores. Les usagers rapportent des glissements d'identité complets, convaincants sur le plan du vécu. Le temps de l'expérience, ils \textit{sont} l'objet ou l'entité en question. Certains décrivent cela comme la sensation d'être mort --- non pas de mourir, mais d'\textit{être mort} --- car lorsque vous êtes une chaise, la personne que vous étiez a tout simplement cessé d'exister.

Le contenu suit l'environnement sensoriel. La personne qui regarde la télévision devient un personnage télévisé. La personne allongée sur un canapé devient le canapé. La personne qui contemple un motif devient le motif.

C'est le Modèle Explicite du Soi qui fait exactement ce que la théorie prédit : se réorienter vers l'entrée dominante, quelle qu'elle soit, lorsque l'entrée de soi normale est perturbée. Le contenu de l'identité n'est pas aléatoire --- il est déterminé par l'environnement sensoriel. Contrôlez l'environnement, et vous pourrez sans doute contrôler l'expérience d'identité.

Je dois ici mettre la théorie en pause un instant. La salvia divinorum est, pour autant que nous le sachions, la substance psychédélique d'origine naturelle la plus puissante sur Terre. La prise de contrôle proprioceptive complète que je viens de décrire signifie une perte totale de la conscience du corps et de l'orientation spatiale. Des gens sous l'influence de la salvia sont passés par des fenêtres du dixième étage. Ils ont marché en pleine circulation. Ils sont morts. Ce n'est pas une drogue festive, pas une curiosité à essayer un vendredi soir. C'est la perturbation pharmacologique la plus extrême du Modèle Explicite du Soi qui existe, et cette perturbation peut vous tuer --- non pas parce que la drogue est toxique, mais parce que vous cessez de savoir où se trouve votre corps et pouvez pleinement croire que vous avez des ailes et que vous pouvez voler.

Beaucoup de gens qui essaient la salvia rapportent que l'expérience leur a donné l'impression de mourir --- non pas métaphoriquement, mais comme une conviction authentique et terrifiante d'avoir cessé d'exister. C'est le Modèle Explicite du Soi qui s'effondre si complètement que la simulation ne peut plus du tout engendrer un « vous ». Nous verrons l'équivalent clinique de cela au chapitre 8, lorsque nous aborderons le délire de Cotard --- des patients neurologiquement convaincus d'être morts. La salvia vous y conduit pharmacologiquement, en quelques secondes, sans avertissement. Demandez-vous si c'est là quelque chose que vous souhaitez vivre.

J'ai moi-même vécu la dilatation du temps. Sous salvia, une demi-seconde de temps réel --- confirmée par la personne qui m'observait --- s'est étirée en ce qui semblait durer quinze minutes ou davantage. Mon monde perceptif s'est reconstruit en séquences élaborées qui incluaient la sensation d'avoir des ailes et de voler alentour (la sensation de vol, ai-je compris plus tard, venait de l'air qui filait à mon passage tandis que je tombais à la renverse sur le lit). Ma réalité tout entière s'est effondrée et régénérée, le tout dans le temps que prend un clignement d'yeux. J'ai décrit cela à un observateur qui me chronométrait, et il m'a dit que j'avais été « parti » moins d'une seconde. Le même genre de dilatation du temps que j'avais déjà éprouvé quelques années plus tôt, en 1998 ou 1999, lors d'un événement de mort imminente (un mécanisme que je décrirai au chapitre 13), mais induit pharmacologiquement et plus extrême encore.

Je ne suis pas le cas le plus spectaculaire. Un récit bien documenté concerne un homme qui a vécu ce qui lui a semblé être huit années complètes d'une vie alternative --- fréquenter l'école, se faire des amis, bâtir une nouvelle existence --- au cours d'un épisode de salvia ayant duré environ quarante-cinq secondes de temps d'horloge. La recherche évaluée par les pairs confirme une distorsion temporelle extrême dans des conditions contrôlées, un participant décrivant le temps comme « plissé tel un accordéon » (Addy et al., 2015). Le substrat fait passer tant de contenu par la simulation, si vite, que le temps subjectif se découple entièrement du temps d'horloge.

Cela n'a jamais été testé expérimentalement dans un cadre contrôlé. Mais cela pourrait l'être, et ce serait une confirmation spectaculaire du mécanisme le plus distinctif de la théorie.

La raison pour laquelle le cerveau peut faire passer tant de contenu par la simulation si vite est la même question que celle de savoir pourquoi un mourant voit toute sa vie défiler en quelques secondes --- et je reviendrai sur les deux ensemble au chapitre 13, où il s'avère qu'ils partagent un unique mécanisme.

Il vaut la peine de se demander pourquoi ce coin --- les deux cadrans poussés au maximum en même temps, une grande partie du cerveau recrutée dans un seul processus \textit{et} ce processus tournant avec une richesse maximale --- est si rare. Deux choses l'interdisent normalement, et elles sont complémentaires.

La première est une simple affaire d'économie. Faire les deux à la fois est métaboliquement ruineux ; le cerveau ne peut se le permettre que pour de brefs instants, et le budget énergétique maintient la porte fermée. Voilà pourquoi il faut quelque chose d'extrême pour la forcer à s'ouvrir --- un cerveau mourant dont les freins métaboliques défaillent, ou une drogue qui perturbe la régulation même qui rationne le budget.

La seconde est plus étrange, et c'est la raison pour laquelle ces états donnent l'impression de quitter le monde. Le monde \textit{intérieur} de votre cerveau possède une dimensionnalité infiniment plus grande que les minces tuyaux qui le relient à la réalité --- deux yeux, deux oreilles, la surface tactile d'un corps, face à une simulation intérieure d'une richesse vertigineuse. Normalement, ce mince tuyau sensoriel suffit, parce qu'il \textit{corrige} sans cesse la simulation intérieure : un filet constant de « voici ce qui est réellement là dehors » maintient l'imagination emballée arrimée au réel. Mais poussez les deux cadrans à fond, et la dynamique intérieure enfle jusqu'à ce que le mince tuyau d'entrée ne puisse plus la rectifier. La simulation dépasse sa propre vérification de la réalité, se verrouille dans son propre attracteur et dérive loin du monde. Et voici le retournement : c'est le cadran d'\textit{étendue} qui le fait --- recruter davantage de cerveau vers l'intérieur est précisément ce qui submerge l'amarre. S'enfoncer dans le coin sectionne la ligne même qui vous aurait maintenu ancré.

Reste une inversion que je trouve encore saisissante. Ce coin est, en même temps, \textit{conscience maximale et contact minimal avec la réalité} --- le plus de traitement, le moins de monde. C'est précisément pourquoi les états à deux cadrans (la revue de vie de la mort imminente, la plongée dans la salvia à forte dose, les immersions psychédéliques les plus profondes, les rêves les plus vifs) sont ceux qui semblent les plus détachés de la réalité. Ils ne sont pas détachés \textit{malgré} leur intensité. Ils sont détachés \textit{parce qu'}ils sont intenses.

Si vous voulez mesurer jusqu'où va ce principe, prenez l'expérience de pensée suivante. Imaginez une personne maintenue en permanence sous une dose très élevée (mais pas totalement dissociative) de Salvinorine A --- le composé actif de la salvia divinorum, qui agit sur un seul type de récepteur (kappa-opioïde). Le Modèle Explicite du Soi (ESM) de cette personne ne se stabiliserait jamais. Il tournerait sans fin au gré de l'entrée qui se trouverait dominer : un instant elle se croirait une chaise, puis une table, puis un dinosaure, puis de l'air, puis une feuille de papier. Elle \textit{éprouverait} toujours des choses (la vue et l'ouïe fonctionneraient encore), mais elle ne saurait plus jamais qui ni quoi elle était. Retirez la substance, et avec le temps, le modèle du soi normal se reconstituerait à partir du Modèle Implicite du Soi (ISM) intact.

C'est important, car cela montre que la conscience n'exige pas un modèle du soi \textit{correct}. Elle exige seulement \textit{un} modèle du soi. L'architecture continue de tourner quoi qu'il arrive. Le Modèle Explicite du Soi ne s'éteint pas quand on lui donne un apport absurde --- il construit le meilleur soi possible à partir des signaux disponibles. C'est le même principe qu'on observe dans le délire de Cotard (l'ESM travaillant sur des signaux intéroceptifs absents : « Je dois être mort »), dans le syndrome d'Anton (l'EWM générant la vue à partir de la mémoire quand les yeux ne fonctionnent plus) et dans le trouble de conversion (l'ESM modélisant une paralysie que le substrat n'a pas réellement). Le modèle du soi est un constructeur compulsif. Il n'arrête jamais de construire. Il n'annonce jamais que les données sont insuffisantes. Il construit, tout simplement, et il croit.

\section*{Anosognosie : l'inverse}

Il y a ici une belle symétrie. Si les psychédéliques sont ce qui se produit quand la frontière implicite-explicite devient \textit{trop} perméable, l'anosognosie est ce qui se produit quand elle devient \textit{trop} imperméable --- du moins localement.

L'anosognosie, le plus souvent observée après un accident vasculaire cérébral de l'hémisphère droit, est une affection où les patients ne perçoivent tout simplement pas leurs propres déficits. Un patient dont le bras gauche est paralysé soutiendra que le bras va bien, cherchera à justifier ses échecs à s'en servir et deviendra confus ou furieux quand on lui met la paralysie sous les yeux. Ces patients ne sont pas dans le déni au sens psychologique --- l'information que le bras est paralysé n'atteint tout simplement jamais leur simulation consciente.

Dans la Théorie des Quatre Modèles (FMT), c'est une baisse locale de la perméabilité implicite-explicite. Le Modèle Implicite du Soi \textit{détient} l'information sur la paralysie --- le substrat enregistre les dégâts. Mais la frontière est bloquée pour ce domaine précis, si bien que le Modèle Explicite du Soi (ESM) n'inclut jamais le déficit. La simulation du patient ne contient pas de bras paralysé, et le patient n'en éprouve donc pas.

Le mécanisme est plus précis encore, et une fois qu'on le saisit, il est élégant d'une manière un peu effroyable. Quand votre système moteur envoie une commande (disons « frappez dans vos mains »), il fait deux choses à la fois. Il envoie la commande aux muscles, et il envoie un \textit{retour prédit} à la conscience : ce que taper des mains devrait donner à sentir et à entendre, sur la base de l'expérience passée. Ce retour prédit arrive \textit{avant} le retour sensoriel réel, parce que le vrai retour doit emprunter des voies neuronales plus lentes. En temps normal, la prédiction est rapidement corrigée ou confirmée par les données sensorielles réelles. Vous prédisez le claquement, puis vous le sentez et l'entendez. Concordance. On passe à la suite.

Dans l'anosognosie, le retour réel du membre paralysé n'arrive jamais. Et le mécanisme censé signaler « attendez --- il ne s'est rien passé » est endommagé. Le retour prédit passe donc sans correction. Le système moteur du patient ordonne aux deux mains de frapper, envoie à la conscience la prédiction d'un claquement à deux mains, et la conscience éprouve exactement cela --- un claquement parfaitement normal des deux mains. Le patient vous dira, en toute sincérité, qu'il vient de taper des deux mains. Il l'a entendu. Il l'a senti. Il l'a éprouvé. Dans sa simulation, cela a eu lieu. Simplement, cela n'a pas eu lieu dans la réalité.

C'est \textit{ainsi} que fonctionne la conscience, tout le temps, chez nous tous. La seule différence, c'est que chez les gens en bonne santé, le retour prédit est corrigé en quelques millisecondes. Dans l'anosognosie, le mécanisme de correction est cassé, et la simulation du patient tourne tout simplement sur les seules prédictions.

Les psychédéliques et l'anosognosie sont le même mécanisme tournant en sens inverse. L'un accroît la perméabilité de manière globale. L'autre la diminue de manière locale. Et cette symétrie engendre une prédiction transversale : les psychédéliques devraient atténuer l'anosognosie. L'augmentation globale de perméabilité devrait submerger le blocage local et permettre à l'information sur le déficit d'atteindre la conscience.

Personne ne l'a jamais testé, parce que personne n'avait de théorie reliant ces deux phénomènes. Le lien est invisible sans la Théorie des Quatre Modèles.


\chapter*{Chapitre 7 : Ce qui se passe quand les lumières s'éteignent}
\addcontentsline{toc}{chapter}{Chapitre 7 : Ce qui se passe quand les lumières s'éteignent}
\markboth{Chapitre 7 : Ce qui se passe quand les lumières s'éteignent}{}

Chaque nuit, vous perdez conscience. Chaque matin, vous la retrouvez. Et la transition entre les deux (le voyage à travers les stades du sommeil) est une démonstration nocturne du principe de criticalité.

\section*{Le sommeil profond : sous le seuil}

Dans le sommeil profond non paradoxal, la dynamique cérébrale bascule vers un régime sous-critique. Sa marque distinctive, ce sont les ondes lentes : de vastes oscillations synchronisées où d'immenses populations de neurones déchargent à l'unisson puis se taisent ensemble. C'est une dynamique de Classe 2 --- périodique, répétitive, trop ordonnée pour la conscience.

L'indice de complexité perturbationnelle (PCI), mis au point par Marcello Massimini et ses collègues, le confirme directement. Le PCI mesure la complexité de la réponse du cerveau à une impulsion magnétique : dans la conscience de veille, la réponse est complexe et différenciée (PCI élevé) ; dans le sommeil profond, elle est simple et stéréotypée (PCI faible). Le cerveau en sommeil profond ne peut pas soutenir la dynamique riche et globalement intégrée qu'exige une simulation consciente.

Les lumières sont éteintes. Le Modèle Explicite du Monde (EWM) et le Modèle Explicite du Soi (ESM) se sont effondrés. Il n'y a ni simulation ni expérience.

\section*{Les rêves : mode dégradé}

Mais les lumières se rallument pendant le sommeil paradoxal. La dynamique cérébrale se déplace de nouveau vers la criticalité --- pas tout à fait, mais suffisamment près. La simulation se réengage, et vous faites de nouveau l'expérience d'un monde.

C'est une simulation dégradée. L'apport sensoriel externe habituel est coupé (vos yeux sont fermés, vos muscles sont paralysés). L'EWM fonctionne à partir de données internes --- puisant dans les connaissances stockées du Modèle Implicite du Monde (IWM) plutôt que dans les impressions sensorielles actuelles. Voilà pourquoi les rêves mettent en scène des lieux et des personnes familiers, mais avec une physique impossible et une incohérence narrative : la simulation fait de son mieux avec ce qui lui reste.

L'ESM tourne lui aussi en mode dégradé. Vous vivez les rêves comme quelque chose qui arrive à « vous », mais votre supervision métacognitive est réduite --- vous acceptez sans broncher des événements impossibles, vous remarquez rarement que vous rêvez, votre sens critique est en veilleuse.

Je sais de l'intérieur ce que l'on ressent dans cette simulation dégradée. Entre environ sept et douze ans, j'ai fait un rêve récurrent --- le même rêve, revenant plusieurs fois au fil de ces années. C'était un paysage, si l'on peut appeler cela ainsi : une structure fractale d'une longueur infinie, avec des vallées qui plongeaient sans fin, auto-similaires à chaque profondeur où l'on posait le regard. Il s'accompagnait d'un sentiment que je n'ai jamais éprouvé à l'état de veille --- une profonde désincarnation, comme si on m'avait dépouillé de mon corps pour me placer à l'intérieur de la géométrie elle-même. Le rêve était inquiétant, étrange, mais aussi bizarrement fascinant. Je voulais qu'il revienne, même quand il m'effrayait.

Ce qui est remarquable, c'est que je peux encore le \textit{sentir}. Des décennies plus tard, quand je pense à ce rêve, la même sensation de désincarnation revient --- non pas comme le souvenir d'un sentiment, mais comme le sentiment lui-même. L'encodage implicite est toujours là, dans les poids synaptiques, intact après trente années de vie éveillée. Le substrat l'a stocké si profondément que le seul fait d'y penser suffit à réactiver partiellement la simulation.

Que faisait ce rêve ? Dans le langage de la théorie : l'EWM, coupé de toute entrée externe, générait du contenu à partir de la propre dynamique de calcul du substrat. Et quelle est cette dynamique ? La Classe 4 --- qui contient la Classe 3 (structure fractale) comme sous-processus. Mon cerveau endormi, faisant tourner sa simulation sur rien d'autre que sa propre architecture, a produit la signature visuelle de cette architecture : des fractales. La désincarnation, c'était l'ESM dans sa forme la plus dégradée --- la simulation savait qu'elle était \textit{quelqu'un}, quelque part, mais la représentation du corps avait presque entièrement disparu.

Je n'ai jamais essayé de rappeler ce rêve après avoir appris, des années plus tard, à faire des rêves lucides. Peut-être devrais-je le faire.

Le somnambulisme en est une démonstration encore plus spectaculaire. Dans le somnambulisme, le système moteur se réactive partiellement tandis que l'ESM reste hors ligne, ou presque. Le substrat exécute des programmes moteurs --- marcher, se déplacer, accomplir même des actions complexes ---, mais la simulation n'est pas pleinement engagée. Le somnambule se déplace dans le monde physique, guidé par les connaissances spatiales de l'IWM, mais avec une expérience consciente minimale, voire nulle.

Je le sais de première main. Adolescent, j'ai traversé une phase de somnambulisme. Un matin, je me suis réveillé assis à mon bureau, des notes griffonnées devant moi, écrites de la main gauche, ce que je ne fais jamais éveillé. J'avais un souvenir fragmentaire d'avoir longé les murs en cercle, cherchant la porte, ne la trouvant pas. Mais la partie où je me suis assis au bureau et où j'ai tenté d'écrire --- celle-là était complètement noire. Le substrat se déplaçait, les programmes moteurs s'exécutaient, mais la simulation --- le « je » --- n'était pas là.

C'est la théorie en miniature. Un corps qui se déplace dans le monde, traite des informations spatiales, exécute des programmes moteurs appris, tout cela sans un soi conscient dans la boucle. Les modèles implicites mènent la danse. Les modèles explicites sont hors ligne. Et le résultat est un être humain qui marche, agit et écrit même, mais il n'y a personne aux commandes.

\section*{Le rêve lucide : l'interrupteur}

Et puis il y a le rêve lucide --- l'état dans lequel vous prenez conscience que vous rêvez tout en restant à l'intérieur du rêve. Dans la Théorie des Quatre Modèles, c'est l'ESM qui « s'enclenche » plus pleinement au sein de l'état de rêve. C'est un bond, par palier, de la capacité d'auto-modélisation.

Apprendre à faire des rêves lucides, c'est apprendre à prendre la simulation en flagrant délit d'erreur. La méthode que j'ai utilisée s'appelle le test de l'interrupteur : tout au long de la journée, vous actionnez par habitude les interrupteurs et vous vous demandez si la lumière a changé correctement. Dans la vie éveillée, elle le fait toujours. Dans un rêve, les interrupteurs ne fonctionnent pas --- l'EWM, qui fonctionne à partir de données internes, ne se donne pas la peine de simuler la physique des circuits électriques. Quand vous actionnez un interrupteur dans un rêve et que la lumière ne change pas, ou que la pièce s'assombrit, ou que l'interrupteur donne une sensation fausse --- cette discordance, c'est votre ESM qui détecte une incohérence dans la simulation. Et à l'instant où vous la détectez, vous savez que vous rêvez.

Il ne m'a fallu que quelques jours. J'ai eu de la chance --- la plupart des gens ont besoin de semaines ou de mois de pratique avant que l'habitude ne se transfère dans leurs rêves. Mais quand cela a fonctionné, la transition a été exactement le seuil par palier que la théorie prédit. Un instant, j'étais un personnage passif dans le récit du rêve. L'instant d'après, j'étais pleinement présent --- conscient que je rêvais, conscient que le monde autour de moi était généré, conscient que je pouvais le changer. L'ESM s'était enclenché. Le substrat n'avait pas changé. La dynamique n'avait pas changé. Mais l'implication du modèle du soi dans la simulation avait franchi un seuil, et tout paraissait différent.

La théorie prédit que cette transition correspond au franchissement d'un seuil de criticalité. Non pas une hausse graduelle de la complexité cérébrale, mais un saut soudain. Si l'on mesurait la complexité de l'EEG dans une fenêtre temporelle synchronisée autour du moment où survient la lucidité (à l'aide du paradigme établi des signaux de mouvements oculaires convenus à l'avance avec les rêveurs lucides), on devrait observer une discontinuité.

\section*{L'anesthésie : les deux types}

L'anesthésie fournit le test le plus net du principe de criticalité, car différents agents anesthésiques produisent des expériences radicalement différentes tout en étant classés sous une même étiquette.

Le \textbf{propofol} pousse le cerveau vers le sous-critique. La connectivité thalamo-corticale est perturbée, la complexité corticale s'effondre, et le PCI approche de zéro. Les lumières s'éteignent complètement. Les patients ne rapportent aucune expérience sous anesthésie au propofol. C'est exactement ce que la théorie prédit : passez sous la criticalité et la simulation ne peut plus être maintenue.

La \textbf{kétamine} fait quelque chose de tout à fait différent. Elle ne pousse \textit{pas} le cerveau vers le sous-critique. Des études EEG montrent que la kétamine \textit{augmente} l'entropie neuronale --- elle pousse le cerveau vers la criticalité ou au-delà, dans un régime plus chaotique. Le résultat ? Le « K-hole » --- des expériences vives, souvent bizarres, de dissociation, de réalité distordue, de sorties hors du corps et d'altération radicale de l'identité.

Dans la Théorie des Quatre Modèles, le K-hole est la conscience qui fonctionne à partir d'une \textit{mauvaise} entrée. L'EWM et l'ESM sont encore actifs (le cerveau est encore au niveau de la criticalité ou au-dessus), mais le traitement sensoriel externe est perturbé. La simulation tourne sur des signaux internes et distordus, produisant la phénoménologie caractéristique du K-hole.

Cette distinction (le propofol abolit la conscience en passant sous le seuil critique, la kétamine altère la conscience en passant au-delà, dans le supracritique, avec une entrée perturbée) constitue un véritable avantage explicatif. La plupart des théories peinent à expliquer pourquoi deux « anesthésiques » produisent des expériences aussi radicalement différentes. Le cadre de la criticalité rend cette distinction naturelle.

Je n'ai jamais subi d'anesthésie chimique. Mais j'ai déjà été mis K.-O. --- brutalement, d'un coup, sans le moindre avertissement. Et ce qui me frappe, avec le recul, c'est l'absence totale de transition. Le sommeil a des stades. Les psychédéliques ont une courbe de montée. Même la salvia, aussi rapide soit-elle, vous accorde une fraction de seconde de « quelque chose est en train de se passer ». Un K.-O. ne vous accorde rien. À un instant, la simulation tourne encore. À l'instant suivant, vous vous réveillez par terre sans la moindre idée du temps qui s'est écoulé. Pas d'assombrissement, pas d'effacement, pas de tunnel. La simulation ne se dégrade pas. Elle s'interrompt. Le disjoncteur saute.

C'est exactement ce à quoi l'on peut s'attendre lors d'une perturbation soudaine et massive de la dynamique corticale --- une chute instantanée bien en dessous de la criticalité. Pas le temps d'un arrêt en douceur en passant par les stades du sommeil. Le système ne traverse pas la Classe 4 pour rejoindre la Classe 2. Il décroche entièrement de la Classe 4, et la simulation s'arrête tout simplement.

\section*{La carte de la conscience}


{\small
\noindent
\begin{tabularx}{\linewidth}{X X X X}
\toprule
\textbf{État} & \textbf{Criticalité} & \textbf{Modèles} & \textbf{Conscience} \\
\midrule
Éveil normal & À la criticalité & Les quatre actifs & Pleine \\[4pt]
Sommeil paradoxal & Proche de la criticalité & EWM/ESM sur entrée interne & Dégradée (rêve) \\[4pt]
NREM profond & Sous-critique & EWM/ESM effondrés & Absente \\[4pt]
Propofol & Forcé sous-critique & EWM/ESM supprimés & Absente \\[4pt]
Kétamine & Au-delà du critique ($\uparrow$ entropie) & EWM/ESM sur mauvaise entrée & Présente, déconnectée \\[4pt]
Psychédéliques & À/au-delà du critique & Tous actifs, $\uparrow$ perméabilité & Présente, altérée \\[4pt]
Rêve lucide & Proche du critique, seuil franchi & EWM actif, ESM pleinement engagé & Conscience de soi renforcée \\[4pt]
\bottomrule
\end{tabularx}
}


Ce tableau résume tout ce que nous avons abordé dans ce chapitre et vous offre une référence à laquelle revenir. Chaque état de conscience que vous ayez jamais éprouvé trouve sa place quelque part sur cette carte, déterminé par deux facteurs : le fait que votre substrat soit ou non à la criticalité, et lesquels des quatre modèles sont en train de tourner. Le sommeil, l'anesthésie, les psychédéliques, les rêves, le K-hole --- ce ne sont pas des mystères distincts. Ce sont des coordonnées différentes sur une même carte.

Mais chaque état de cette carte correspond à un cerveau sain qui retrouve son chemin de lui-même avant le matin. La question plus difficile est de savoir à quoi ressemble la carte lorsqu'un rouage de la machinerie ne revient pas --- lorsque les dégâts sont permanents et que l'architecture doit continuer de fonctionner avec une partie d'elle-même manquante. Pour cela, vous n'allez pas dans un laboratoire du sommeil. Vous allez dans un service de neurologie.


\chapter*{Chapitre 8 : Le miroir clinique}
\addcontentsline{toc}{chapter}{Chapitre 8 : Le miroir clinique}
\markboth{Chapitre 8 : Le miroir clinique}{}

La même architecture à quatre modèles qui explique le sommeil et l'anesthésie explique aussi certains des troubles les plus spectaculaires et les plus déroutants de la neurologie clinique. Ce ne sont pas de simples études de cas intéressantes --- c'est ce qui se produit lorsque des composants précis de l'architecture cessent de fonctionner. Et chaque défaillance éclaire l'architecture sous un angle différent, comme un fusible grillé révèle quel circuit il protégeait.

Si la théorie a de la valeur, alors une atteinte à des modèles précis devrait produire des déficits précis et prévisibles. Pas de vagues formules du genre « la conscience est altérée », mais des prédictions exactes : désactivez ce composant, et vous obtenez \textit{tel} syndrome. Maintenez un autre composant en marche sans son entrée habituelle, et vous obtenez \textit{cet} autre syndrome. La littérature clinique regorge de troubles qui restent profondément déroutants sous les modèles classiques de la conscience, mais qui se mettent en place d'eux-mêmes dès que l'on dispose d'une distinction réel/virtuel et de quatre modèles en interaction.

\textbf{Vision aveugle et syndrome d'Anton : le miroir parfait}

Si vous ne devez retenir qu'une seule chose de ce chapitre, retenez ce couple. Toutes les autres théories de la conscience peinent à expliquer ne serait-ce que l'un de ces troubles. La Théorie des Quatre Modèles (FMT) les prédit tous les deux.

Commençons par la vision aveugle (\textit{blindsight}). Un patient présente une lésion du cortex visuel primaire --- la partie du cerveau qui engendre l'expérience visuelle consciente. Selon n'importe quel test clinique standard, ce patient est aveugle. Demandez-lui ce qu'il voit, et il vous répondra : rien. Et il le pense sincèrement. Il n'est ni modeste ni confus. Pour ce qui est de son expérience consciente, le monde visuel n'existe tout simplement pas.

Mais alors quelque chose d'étonnant se produit. Des chercheurs disposent des obstacles dans un couloir et demandent au patient de le traverser. Il proteste --- il ne voit rien, comment pourrait-il se diriger ? Ils insistent. Il soupire, se lève et marche.

Et il franchit le parcours d'obstacles sans la moindre faute. Il contourne les chaises. Il se baisse sous une barrière qui n'était pas là la fois précédente. Il se faufile dans l'espace entre deux obstacles --- tout en soutenant, en toute sincérité, qu'il ne voit strictement rien. Il en existe des vidéos --- je vous encourage à les chercher, car la lecture ne leur rend pas justice. Les images d'un homme cliniquement aveugle qui franchit un parcours d'obstacles comme s'il voyait parfaitement comptent parmi les démonstrations les plus stupéfiantes de toutes les neurosciences. Les chercheurs qui observent ont l'air d'avoir vu un fantôme.

Comment ? Parce que le substrat traite toujours l'information visuelle. Le Modèle Implicite du Monde (IWM) reçoit une entrée visuelle par des voies sous-corticales qui contournent le cortex endommagé (une route rapide de la rétine au colliculus supérieur puis au pulvinar), des voies qui ont évolué bien avant l'existence du cortex. Il construit une carte spatiale, guide le comportement moteur, empêche le corps de heurter les objets. Mais rien de tout cela n'atteint le Modèle Explicite du Monde (EWM). La simulation consciente ne contient aucune vision. Le patient éprouve réellement la cécité --- et s'oriente réellement grâce à la vue. Le substrat fonctionne sans la simulation.

Renversons maintenant la situation. Le syndrome d'Anton (anosognosie de la cécité corticale) en est l'exact opposé. Ces patients sont réellement, totalement aveugles. Leur cortex visuel ou leurs voies optiques sont détruits. Aucune information visuelle n'atteint le cerveau. Et pourtant, ils sont absolument, inébranlablement convaincus qu'ils voient.

Ils se cognent aux murs et reprochent aux meubles de ne pas être à leur place. Ils décrivent avec une assurance totale des objets qui ne sont pas dans la pièce --- « Il y a un vase bleu sur la table », alors que la table est vide. Demandez-leur d'identifier ce que vous brandissez, et ils vous donneront une réponse, calmement et sans hésiter, et elle sera fausse. Confrontez-les à la preuve de leur cécité, et ils deviennent d'abord perplexes, puis agacés, puis furieux. L'éclairage est mauvais. Il leur faut de nouvelles lunettes. Ils n'étaient tout simplement pas attentifs. Ils ne mentent pas. Ils ne sont pas dans le déni au sens psychologique. Ils voient réellement, dans le vécu --- et ce qu'ils voient n'a aucune correspondance avec le monde réel.

Dans la Théorie des Quatre Modèles, c'est l'EWM qui engendre une simulation visuelle à partir des connaissances stockées dans l'IWM --- alors même qu'aucune entrée visuelle actuelle ne parvient. La simulation fonctionne à partir d'anciennes données, sur des attentes, sur la meilleure supposition du cerveau quant à l'allure que devrait avoir le monde. Le patient « voit » un monde qui n'est pas là. La simulation fonctionne sans entrée actuelle.

Mettez-les côte à côte. Vision aveugle : le substrat traite la vision, mais la simulation ne la montre pas. Syndrome d'Anton : la simulation montre la vision, mais le substrat ne la reçoit pas. Substrat sans simulation. Simulation sans entrée. Ces deux troubles déroutent profondément si l'on tient la conscience pour une chose unique et unifiée. Ils découlent en revanche naturellement, et même de façon prévisible, d'une théorie qui distingue le traitement réel de l'expérience virtuelle. On aurait du mal à concevoir un meilleur couple de cas-tests, même en s'y efforçant.

\textbf{Conscience cachée : prisonnier à l'intérieur}

En 2006, Adrian Owen et ses collègues publièrent une étude qui transforma notre façon de concevoir l'état végétatif. Ils placèrent une patiente diagnostiquée comme végétative (non réactive, apparemment inconsciente) dans un scanner IRMf et lui demandèrent de s'imaginer en train de jouer au tennis. Son cerveau s'illumina selon exactement le même schéma que celui d'une personne consciente et en bonne santé imaginant la même chose.

Elle était là, à l'intérieur. Consciente, présente, pensante --- et totalement incapable de bouger, de parler ou de signaler sa présence à quiconque.

La Théorie des Quatre Modèles opère ici une distinction nette. Un patient réellement végétatif possède un substrat sous-critique. La dynamique est tombée sous le seuil. La simulation ne tourne pas. Il n'y a personne --- non parce que la personne serait « partie », mais parce que l'architecture computationnelle qui engendre la simulation est hors service.

Mais un patient secrètement conscient, c'est tout autre chose. Le substrat est critique --- la dynamique est assez riche pour soutenir une simulation. L'EWM et le Modèle Explicite du Soi (ESM) sont actifs. La personne éprouve, pense, ressent. Mais les voies de sortie sont détruites. La simulation n'a aucun moyen de s'exprimer. La personne est consciente mais enfermée, prisonnière d'un corps qui ne répond pas.

Le PCI (la même mesure qui distingue les stades du sommeil) devrait distinguer ces cas. Et c'est ce qu'il fait. Certains patients diagnostiqués comme végétatifs affichent des valeurs de PCI en plein dans la fourchette de la conscience. Ils ne sont nullement végétatifs. Ce sont des prisonniers. Les implications médicales et éthiques sont immenses, et la Théorie des Quatre Modèles vous dit exactement pourquoi la distinction existe et exactement comment la détecter.

\textbf{Le délire de Cotard : « Je suis mort »}

Et puis il y a les patients qui croient être morts.

Le délire de Cotard est l'un des troubles les plus étranges de la psychiatrie. Les patients soutiennent qu'ils sont morts. Ils croient que leurs organes se sont dissous, que leur sang s'est écoulé, qu'ils n'existent plus. Certains croient qu'ils sont en train de pourrir. D'autres se croient immortels --- car si l'on est déjà mort, on ne peut plus mourir. Ils ne parlent pas au sens figuré. Ils le pensent avec une conviction totale, inébranlable.

À ce stade, vous devriez reconnaître le mécanisme. C'est le même qu'au chapitre 6 --- l'ESM qui construit le meilleur modèle possible à partir de l'entrée dont il dispose. Dans le Cotard, l'entrée intéroceptive est gravement déformée. Les signaux corporels internes qui vous disent que votre cœur bat, que votre estomac digère, que vos poumons respirent --- ils sont absents ou brouillés. Et l'ESM, éternel constructeur compulsif, interprète « pas de battement de cœur, pas de digestion, pas de respiration, aucune sensation corporelle » de la seule façon qui lui reste : je suis mort.

Le « Je suis une chaise » de la salvia. Le « mon bras va bien » de l'anosognosie. Et maintenant le « je suis mort » de Cotard. (Au chapitre suivant, la confabulation du cerveau divisé ajoutera encore un cas à cette liste.) Un seul mécanisme traverse chaque cas. L'ESM fait toujours son travail --- il bâtit toujours le meilleur modèle du soi qu'il peut. Quand l'entrée est correcte, vous vous sentez vous-même. Quand l'entrée est fausse, vous vous sentez être une chaise, ou en bonne santé alors que vous êtes paralysé, ou mort alors que vous êtes vivant. Mais cela paraît toujours parfaitement, irrésistiblement réel --- parce que c'est le seul soi auquel vous ayez accès.

\textbf{Le syndrome de la main étrangère : quand le comité n'est pas d'accord}

Il y a ensuite un trouble qui se lit comme un film d'horreur, mais qui illustre la nature multi-agents du substrat plus vivement que n'importe quelle expérience de pensée. Dans le syndrome de la main étrangère, l'une des mains du patient agit avec un but et une intention apparents, mais contre sa volonté consciente. Une main allume une cigarette pendant que l'autre la lui retire et la jette par terre. Une main se tend vers une poignée de porte pendant que l'autre saisit le poignet et le ramène en arrière. Le patient regarde, horrifié, une partie de son propre corps poursuivre des buts qu'il n'a pas choisis.

Stanley Kubrick s'en est servi dans \textit{Docteur Folamour} --- et l'on a cru qu'il l'avait inventé. Ce n'est pas le cas. Le syndrome est réel, et il se présente sous deux variantes. Dans la forme calleuse, provoquée par une lésion du corps calleux, les symptômes rappellent le conflit du cerveau divisé : deux hémisphères aux plans moteurs concurrents, dont aucun ne parvient à l'emporter sur l'autre. Dans la forme frontale, provoquée par une atteinte préfrontale, la main « étrangère » manifeste un comportement désinhibé --- elle saisit des objets, utilise des outils, touche les choses de façon compulsive ---, le tout apparemment orienté vers un but, mais sans le consentement du patient.

Il existe aussi une variante plus subtile, le syndrome de la main anarchique, où le patient perd le \textit{contrôle} moteur plutôt que la \textit{propriété} motrice. La main fait des choses que le patient n'a pas voulues, mais celui-ci la reconnaît toujours comme \textit{sa} main --- il n'arrive simplement pas à l'arrêter. La distinction compte : la main étrangère est une défaillance de la frontière de propriété corporelle de l'ESM (« cette main n'est pas la mienne »), tandis que la main anarchique est une défaillance du système d'inhibition motrice (« cette main est la mienne, mais elle n'obéit pas »). Même architecture, points de rupture différents.

L'enseignement clé de ces syndromes, c'est que votre sentiment d'être l'auteur de vos actes (l'impression du « c'est moi qui ai fait ça ») ne se calcule ni avant l'action ni pendant. Il se calcule \textit{après}, en comparant le résultat prédit de l'action au résultat observé. Quand la comparaison concorde, vous éprouvez ce sentiment d'appartenance. Quand elle ne concorde pas, il fait défaut. C'est pourquoi les patients atteints du syndrome de la main étrangère parviennent parfois à se chatouiller eux-mêmes : leur système de prédiction ne génère pas le résultat attendu pour les mouvements de la main étrangère, si bien que le contact arrive de façon inattendue, comme s'il venait de quelqu'un d'autre.

\textbf{Syndrome de Charles Bonnet : la simulation qui ne s'arrête pas}

Si vous voulez d'autres preuves que la simulation du cerveau est \textit{générative} --- qu'elle construit l'expérience à partir de modèles au lieu de la recevoir passivement des sens ---, considérez le syndrome de Charles Bonnet. Les patients dont la rétine ou le nerf optique est détruit, mais dont le cortex visuel demeure intact, éprouvent des hallucinations visuelles vives et complexes. Non pas de vagues formes ou des éclats de lumière. Des scènes entières : des personnes, parfois miniaturisées ou costumées comme des personnages de dessin animé, parfois des images en miroir du patient. Des paysages. Des objets. Des visages.

Les patients savent en général que rien de tout cela n'est réel. À la différence des hallucinations psychotiques, celles de Charles Bonnet s'accompagnent d'une conscience intacte de la situation --- le patient dit : « Je vois un petit homme en haut-de-forme assis sur ma table, et je sais qu'il n'est pas là. » Il s'agit de la simulation visuelle du Modèle Explicite du Monde (EWM), qui tourne sur des données internes issues des aires visuelles supérieures, en l'absence de tout apport extérieur. La simulation ne s'arrête pas simplement parce que l'entrée s'est tarie. Elle génère. Elle comble le vide. Et ce qu'elle génère nous apprend quelque chose sur l'architecture : le système visuel est un modèle génératif, non un récepteur passif. Il produit sa meilleure hypothèse sur l'aspect du monde, à l'aide de gabarits stockés et de prédictions descendantes --- exactement comme le décrit la Théorie des Quatre Modèles.

\textbf{Déjà-vu : le gabarit qui correspond trop bien}

À propos du système génératif du cerveau et de ses ratés occasionnels : presque tout le monde a fait l'expérience du déjà-vu --- la sensation troublante d'avoir déjà vécu l'instant présent. Les explications vont du mystique (vies antérieures, prémonitions) au dédaigneux (ce n'est qu'un bug). La Théorie des Quatre Modèles en propose un compte rendu plus précis.

Le cerveau stocke ce qu'on pourrait appeler des « souvenirs-gabarits » --- des représentations d'expériences réduites à l'essentiel, extrêmement clairsemées, notamment issues des rêves. Ces gabarits ne sont pour l'essentiel que des échafaudages vides : le vague sentiment d'un lieu, une ambiance, une configuration spatiale, presque sans aucun détail. Quand vous récupérez un souvenir normal, les vides sont comblés par confabulation --- le cerveau génère des détails plausibles pour créer une expérience continue. Vous ne remarquez pas le remplissage, car le résultat paraît cohérent.

Le déjà-vu survient lorsqu'une expérience réelle actuelle se trouve correspondre trop étroitement à l'un de ces gabarits stockés. Le système d'appariement de motifs du cerveau se déclenche : « J'ai déjà vu ça. » Mais dès que vous essayez de préciser \textit{quand} vous êtes censé l'avoir vu, vous ne trouvez rien --- parce que le gabarit n'a jamais été une expérience réelle. C'était un fragment de rêve, ou un souvenir si profondément comprimé qu'il a perdu depuis longtemps tout détail contextuel. La correspondance entre l'entrée actuelle et le gabarit stocké est authentique, mais l'expérience « originale » que le gabarit est censé consigner n'a jamais réellement eu lieu sous la forme que votre cerveau lui prête à présent. Le système fonctionne correctement --- il a bel et bien trouvé une correspondance. Simplement, la correspondance se fait avec un squelette, non avec un corps.

Repassez le chapitre en esprit et la même figure ne cesse de refaire surface. La vision aveugle et le syndrome d'Anton, le délire de Cotard et le syndrome de Charles Bonnet, le patient enfermé en lui-même et le déjà-vu qui correspond trop bien --- chaque cas est la couture entre substrat et simulation, forcée assez loin pour qu'on puisse voir au travers. Un cerveau sain maintient les deux à ras, si parfaitement à ras qu'on ne soupçonne jamais qu'il y en a deux. Écartez-les dans une direction quelconque --- substrat sans simulation, simulation sans substrat, une simulation qui tourne sur la mauvaise entrée ou sur aucune --- et la distinction réel/virtuel cesse d'être une thèse philosophique pour devenir un diagnostic. Chaque défaillance jusqu'ici a été un dommage infligé à un seul cerveau. La prochaine n'est pas du tout un dommage : c'est une coupe nette en plein milieu, et elle vous donne de chaque chose deux exemplaires.


\chapter*{Chapitre 9 : Deux esprits dans un seul cerveau}
\addcontentsline{toc}{chapter}{Chapitre 9 : Deux esprits dans un seul cerveau}
\markboth{Chapitre 9 : Deux esprits dans un seul cerveau}{}

Dans les années 1960 prit forme l'une des expériences les plus spectaculaires de l'histoire des neurosciences. Pour traiter une épilepsie sévère, les chirurgiens Philip Vogel et Joseph Bogen sectionnèrent le corps calleux --- l'épais faisceau de fibres nerveuses qui relie les deux hémisphères du cerveau ---, et Roger Sperry et Michael Gazzaniga étudièrent ensuite les patients qui en furent issus. Le résultat fut le syndrome du cerveau divisé : une seule personne dotée, en apparence, de deux esprits indépendants.

Les démonstrations classiques sont bien connues. Présentez un mot dans le champ visuel gauche (traité par l'hémisphère droit), et le patient peut saisir l'objet correspondant de la main gauche, mais il est incapable de dire quel était le mot (parce que la parole est contrôlée par l'hémisphère gauche, qui n'a pas vu le mot). Les deux hémisphères ont des perceptions indépendantes, des intentions indépendantes et, parfois, des objectifs contradictoires.

Ces expériences allaient bien au-delà de simples tours de passe-passe avec des mots et des objets. Dans certains cas, les hémisphères se combattaient ouvertement. Un patient rapporta que sa main gauche déboutonnait sa chemise pendant que sa main droite tentait de la reboutonner. Chez un autre, la main gauche se tendit vers sa femme au cours d'une dispute, non pour la réconforter --- tandis que la main droite saisissait la gauche et la retenait. Le patient regardait, horrifié, deux parties de son propre corps poursuivre des buts incompatibles, aucune sous son contrôle unifié. Ce ne sont pas des métaphores du conflit intérieur. Ce sont des conflits réels, physiques, entre deux systèmes moteurs qui ne peuvent plus se coordonner parce que le câble qui les reliait a été sectionné.

Dans la vie quotidienne, les patients au cerveau divisé fonctionnent remarquablement bien. En dehors du laboratoire, vous ne remarqueriez que rarement quoi que ce soit d'inhabituel. Les deux hémisphères apprennent à coopérer par des voies indirectes --- indices extérieurs, mouvements du corps, champs visuels partagés. Le système compense. Mais placez le patient dans un cadre expérimental contrôlé où chaque hémisphère reçoit des informations différentes, et l'unité vole en éclats. Deux esprits émergent d'un seul cerveau, chacun avec ses propres perceptions, ses propres intentions et sa propre version de la réalité.

\section*{L'interprète de l'hémisphère gauche}

Mais le trait le plus révélateur des patients au cerveau divisé n'est pas la division --- c'est ce qui se produit lorsqu'on leur demande d'expliquer cette division.

Gazzaniga a identifié ce qu'il a appelé l'« interprète de l'hémisphère gauche » : la tendance compulsive de l'hémisphère gauche à générer des explications pour des événements qu'il est en réalité incapable d'expliquer. La démonstration classique se déroule ainsi. Présentez une scène enneigée à l'hémisphère droit et une patte de poulet à l'hémisphère gauche, puis demandez au patient de choisir des objets apparentés. La main gauche (hémisphère droit) choisit une pelle (pour la neige). La main droite (hémisphère gauche) choisit un poulet. Demandez ensuite au patient (par la parole, contrôlée par l'hémisphère gauche) pourquoi il a choisi la pelle. L'hémisphère gauche ne sait rien de la neige (il n'a vu que la patte de poulet), alors il invente une explication : « Oh, il faut une pelle pour nettoyer le poulailler. »

Le patient n'hésite pas. Ne dit pas « Je ne suis pas sûr. » N'a pas l'air troublé. L'explication arrive instantanément, avec assurance, et paraît tout à fait naturelle à celui qui la donne. Ce n'est pas un mensonge. L'hémisphère gauche ignore réellement ce que l'hémisphère droit a vu. Il n'a aucun accès à cette information --- le câble est sectionné. Alors il fait ce que le Modèle Explicite du Soi (ESM) fait toujours : il construit le meilleur récit possible à partir des informations disponibles.

Et voici la partie qui devrait vous troubler : vous faites cela aussi. Chaque jour. Votre interprète de l'hémisphère gauche est en train de tourner en ce moment même, construisant un récit cohérent à partir de la moindre information qui atteint la conscience, lissant les lacunes, inventant des explications plausibles pour des décisions que votre substrat a prises avant même que « vous » soyez consulté. La seule différence entre vous et un patient au cerveau divisé, c'est que votre corps calleux est intact, si bien que l'interprète a accès à davantage d'informations. Il confabule moins parce qu'il a moins de matière à confabuler. Mais le mécanisme est identique. La machinerie de la narration de soi ne change pas. Seule change la qualité de l'entrée.

\section*{Une personne ou deux ?}

Cela soulève une question dont les philosophes débattent depuis des décennies : une fois le corps calleux sectionné, y a-t-il une personne dans ce crâne, ou deux ?

Thomas Nagel s'est attaqué à cette question dans un essai célèbre de 1971 et en a conclu qu'elle pourrait n'avoir aucune réponse déterminée (que notre concept de « personne » se disloque tout simplement dans cette situation, comme le concept de « pays unique » se disloque lorsqu'on trace une frontière en plein milieu). Derek Parfit est allé plus loin, soutenant que les cas de cerveau divisé montrent que l'identité personnelle elle-même n'est pas ce qui compte --- ce qui compte, c'est la continuité psychologique, et celle-ci peut connaître des degrés.

La Théorie des Quatre Modèles (FMT) propose une réponse plus précise : tout dépend des modèles qui tournent et de leur degré de dégradation.

Dans la vie quotidienne, un patient au cerveau divisé est fonctionnellement une seule personne. Les deux hémisphères partagent le même corps, le même environnement, la même histoire de vie (encodée de façon redondante dans les deux hémisphères avant l'opération). Le Modèle Implicite du Soi (ISM), qui stocke la personnalité, les souvenirs à long terme, les dispositions comportementales --- a été bâti au fil des décennies avec un corps calleux intact. Sectionner le câble n'efface pas ces modèles stockés. Cela empêche seulement leur mise à jour synchronisée. Ainsi, immédiatement après l'opération, les deux hémisphères font tourner des modèles de soi très semblables. Le patient se sent être une seule personne parce que, en termes de connaissance de soi stockée, il l'est en grande partie.

Avec le temps, les modèles devraient dériver. Chaque hémisphère accumule des expériences différentes, forme des associations différentes, développe des réactions émotionnelles différentes à des événements que lui seul a perçus. Plus un patient au cerveau divisé vit longtemps après l'opération, plus les deux modèles implicites du soi devraient diverger --- lentement, parce que les deux hémisphères partagent encore le même corps et le même environnement, mais de façon mesurable.

Mon point de vue personnel penche vers deux. Si la bande passante entre les hémisphères est insuffisante pour une synchronisation en temps réel de la simulation --- et sans le corps calleux, elle l'est --- alors vous avez deux modèles de soi tournant sur deux substrats, chacun engendrant sa propre expérience consciente. Ils coopèrent bien parce qu'ils partagent un corps, un environnement sensoriel et toute une vie d'histoire commune. Mais coopération n'est pas identité. Deux personnes qui vivent ensemble coopèrent bien, elles aussi.

Chose intrigante, Yair Pinto et ses collègues ont publié en 2017 une étude qui est venue compliquer le tableau standard. Ils ont constaté que les patients au cerveau divisé pouvaient rapporter avec exactitude des stimuli présentés dans l'un ou l'autre champ visuel --- même lorsque le stimulus n'était montré qu'à l'hémisphère qui ne contrôle pas la parole. Cela suggérait que les deux hémisphères conservaient davantage d'unité que ne le laissaient penser les expériences classiques. Le résultat est encore débattu, mais il s'inscrit naturellement dans le cadre holographique que je décrirai ensuite : même après section du corps calleux, il subsiste dans chaque hémisphère assez d'information redondante pour soutenir un comportement étonnamment unifié, du moins pour certaines tâches.

\section*{La propriété holographique}

Dans la Théorie des Quatre Modèles, le cerveau divisé révèle une propriété clé des modèles virtuels : ils sont \textbf{holographiques}. Dans les réseaux neuronaux, l'information est distribuée sur l'ensemble du réseau, et non localisée dans des neurones précis. Lorsque vous coupez le réseau en deux, vous n'obtenez pas une division nette --- vous obtenez deux copies dégradées mais \textit{complètes}. Chaque hémisphère conserve une version réduite des quatre modèles : un Modèle Implicite du Monde amoindri, un Modèle Implicite du Soi amoindri, et la capacité d'engendrer un Modèle Explicite du Monde et un Modèle Explicite du Soi. Les deux hémisphères peuvent soutenir la conscience de façon indépendante (tous deux se situent au-dessus du seuil de criticalité), mais chacun travaille avec moins d'information.

C'est l'hologramme en patchwork du chapitre 3, à présent coupé en son milieu. Chaque hémisphère est l'une de ces moitiés --- non pas une demi-image, mais une image complète à plus faible résolution, parce que l'information était dès le départ étalée sur toute la surface. Le dommage réduit la résolution sans supprimer le contenu.

Cela explique pourquoi les patients au cerveau divisé ne sont pas simplement « deux demi-esprits ». Ce sont deux esprits \textit{complets mais dégradés}. Chaque hémisphère peut percevoir, décider et agir --- simplement avec moins d'information et moins de capacité que le système intact. La propriété holographique garantit que couper la connexion dégrade sans détruire. Et elle explique les résultats de Pinto en 2017 : même sans le corps calleux, chaque hémisphère conserve assez d'information holographique pour accomplir bien des tâches que le modèle classique disait impossibles.

La confabulation (l'interprète de l'hémisphère gauche) est ce même constructeur compulsif du chapitre clinique --- l'ESM édifiant un récit de soi à partir de la moindre entrée dont il dispose, et y croyant totalement. Sectionnez le corps calleux, et l'hémisphère gauche a simplement moins de matière à sa disposition ; l'histoire qu'il invente pour combler la lacune est fausse, mais toujours \textit{ressentie comme parfaitement réelle}.

\section*{Un cerveau, plusieurs soi}

Le cerveau divisé montre ce qui se produit lorsqu'on \textit{clone} les modèles virtuels en divisant physiquement le substrat. Le trouble dissociatif de l'identité montre ce qui se produit lorsqu'on les \textit{scinde}.

Dans le TDI, le substrat n'est pas divisé --- le corps calleux est intact, le matériel neuronal entier. Mais les modèles virtuels se sont scindés en plusieurs configurations. Chaque alter est un Modèle Explicite du Soi distinct --- un récit de soi séparé, avec son propre profil émotionnel, ses propres schémas de comportement, sa propre manière de se rapporter au corps et au monde. Les alters ne partagent pas plus un modèle de soi unique que deux utilisateurs ne partagent une seule session de connexion sur le même ordinateur. Ils se relaient.

Le déclencheur, dans la quasi-totalité des cas documentés, est un traumatisme infantile sévère et répété. Cela a un sens au sein de la théorie. Le Modèle Explicite du Soi d'un jeune enfant est encore en formation --- encore plastique, encore en cours d'assemblage à partir de l'expérience. Soumettez ce modèle de soi en développement à des expériences si écrasantes qu'aucun récit de soi unique ne peut les contenir, et le système fait la seule chose qu'il peut : il se scinde. Il crée des configurations séparées, chacune capable de gérer un aspect différent de la situation insoutenable. Un alter détient les souvenirs traumatiques. Un autre fonctionne dans la vie quotidienne comme si de rien n'était. Un autre encore prend le relais dans les moments de danger. La scission n'est pas une pathologie --- c'est la réponse d'urgence du système de modélisation de soi face à une entrée qui détruirait un modèle unique et unifié.

Voilà pourquoi le TDI ne se développe presque jamais chez l'adulte. Le Modèle Implicite du Soi d'un adulte est déjà consolidé --- les poids synaptiques sont fixés, la structure de la personnalité est stable. Il faut des circonstances extraordinaires pour scinder un modèle de soi adulte (torture sévère, captivité prolongée). Mais l'ISM d'un enfant est encore en train de s'écrire. L'argile est encore humide. Scindez-la sous une pression suffisante, et les configurations séparées durcissent en modèles de soi distincts et persistants.

Les données le confirment. Si chaque alter est réellement une configuration distincte de l'ESM, alors le passage d'un alter à l'autre devrait produire des changements mesurables dans les schémas d'activité neuronale --- et c'est bien le cas. Reinders et al. (2003) ont montré que différents alters chez un même individu produisent des schémas distincts de débit sanguin cérébral régional. Le \textit{même cerveau} s'illumine différemment selon le modèle du soi qui tourne. Ce n'est pas ce qu'on attendrait d'un « jeu d'acteur » ou d'un « jeu de rôle ». C'est ce qu'on attendrait d'une véritable bifurcation logicielle. Dans des études de suivi, Reinders et ses collègues ont constaté que les différences neuronales entre alters étaient plus grandes que celles entre des acteurs à qui l'on avait demandé de simuler un TDI --- un résultat qui devrait réduire au silence quiconque pense encore que le TDI n'est qu'une « simple » performance.

C'est la propriété de « bifurcation » du chapitre 3 à l'œuvre. Un seul substrat, plusieurs configurations virtuelles, chacune faisant tourner un modèle du soi complet mais distinct. La théorie ne se contente pas d'accueillir le TDI --- elle prédit exactement ce type d'architecture, et engage un test dessus : perturber le substrat neuronal qui soutient l'ESM d'un alter devrait déclencher un basculement vers un autre. C'est la Prédiction 9, et le protocole complet --- y compris l'endroit du cerveau où les signatures propres à chaque alter devraient apparaître --- se trouve en annexe, à la fin du livre.

Tous les esprits de ce chapitre ont été humains --- scindés, dédoublés, mis en miroir, mais humains. Rien dans l'architecture des quatre modèles ne l'exige. Si un soi n'est que des modèles tournant à la criticalité, la vraie question n'est pas de savoir si la personne assise en face de vous en possède un. C'est de savoir jusqu'où le même tour de passe-passe s'étend vers l'extérieur --- jusque dans le nid, jusque dans la ruche, jusque dans l'eau où, un jour, quelque chose d'énorme a soutenu mon regard et, j'en suis presque certain, m'a dit bonjour.


\chapter*{Chapitre 10 : La question animale}
\addcontentsline{toc}{chapter}{Chapitre 10 : La question animale}
\markboth{Chapitre 10 : La question animale}{}

Je plongeais avec quatre autres personnes sous l'Arche de Darwin --- à l'époque où elle était encore intacte --- quand une orque mâle, énorme, a surgi de nulle part et s'est approchée si près que ma première pensée fut qu'elle voulait nous dévorer. L'animal était si gigantesque que les trente mètres d'eau au-dessus de nos têtes semblaient une simple flaque. Elle s'est arrêtée et nous a observés de son énorme œil droit, puis est passée à son organe acoustique, cliquant intensément, nous scannant au son. Et puis elle --- je ne trouve pas d'autre mot --- \textit{a parlé.} Dans un chant très aigu, bref et clairement structuré, elle a dit quelque chose. Il n'y avait pas le moindre doute dans mon esprit : c'était du langage, pas seulement un chant. On avait le sentiment que cela pourrait s'apprendre, avec assez de rencontres et une voix ridiculement aiguë pour répondre.

Un peu plus tard, nous avons saisi une partie de ce qu'elle avait dit. Elle a nagé jusqu'à la limite de mon champ de vision et est revenue avec sa compagne et son petit, les guidant devant nous pour qu'ils puissent nous observer. Nous avions tous des appareils photo, bien sûr. Pas une seule photo n'a été prise. Et nous étions tous à court d'air, il a fallu remonter. Chacun de nous avait les larmes aux yeux tandis qu'un canot pneumatique nous ramenait au bateau, et ce n'était pas à cause du vent.

Je n'ai jamais été aussi certain de la présence d'un autre esprit que dans cette eau-là. Mais la certitude n'est pas une preuve --- une histoire de larmes sur un canot ne prouve rien à qui n'était pas là, et surtout pas à un sceptique qui a vu trop de gens projeter une âme sur leurs plantes d'intérieur. Si cette orque était quelqu'un, une théorie doit pouvoir dire \textit{pourquoi}, sur des bases qui tiennent encore lorsque le plongeur est à court d'air et que l'émotion s'est dissipée. Et elle doit dire où ce quelqu'un s'arrête.

Votre chien est-il conscient ?

La plupart des propriétaires d'animaux répondraient oui sans hésiter. La plupart des neuroscientifiques seraient d'accord, du moins avec prudence. Mais sur quelle base ? Et où la conscience commence-t-elle dans le règne animal ?

La Théorie des Quatre Modèles (FMT) apporte des réponses claires, qui découlent de ses engagements fondamentaux plutôt que d'être ajoutées après coup.

\textbf{Engagement 1 : la conscience est un continuum, pas une valeur binaire.} Il n'existe pas de frontière nette entre le conscient et le non-conscient. Il y a des degrés --- des niveaux gradués d'auto-simulation, du plus rudimentaire (modèle du soi minimal) au triplement étendu (conscience de soi récursive). Différents animaux occupent différentes positions le long de ce continuum.

\textbf{Engagement 2 : la conscience est indépendante du substrat.} Ce qui compte, c'est l'architecture fonctionnelle (quatre modèles à la criticalité), pas l'implémentation physique particulière. Si un cerveau réalise l'architecture à quatre modèles, il est conscient, que ce cerveau soit un cortex de mammifère, un pallium d'oiseau ou le réseau neuronal distribué d'un poulpe.

\textbf{Engagement 3 : la criticalité est le seuil physique.} Un système nerveux doit fonctionner au bord du chaos, ou tout près. Des systèmes nerveux plus simples (insectes, vers) n'atteignent peut-être pas la criticalité et ne seraient donc pas conscients --- ils traitent de l'information et produisent du comportement, mais sans simulation.

\section*{À quel point êtes-vous conscient ?}

Il y a une question que vous vous posez probablement déjà. Si la conscience est une simulation --- un soi virtuel au sein d'un monde virtuel --- alors ce n'est pas une affaire de tout ou rien, n'est-ce pas ? Une simulation peut être plus ou moins détaillée. Un modèle du soi peut être plus ou moins sophistiqué. Ce qui signifie que la conscience existe par \textit{degrés}.

La Théorie des Quatre Modèles vous offre une manière précise de penser ces degrés. Il existe quatre niveaux gradués, et chaque créature consciente se situe quelque part sur cette échelle.

Tout en bas se trouve la \textbf{conscience de base}. C'est un Modèle Explicite du Monde (EWM) doté d'un Modèle Explicite du Soi (ESM) seulement rudimentaire. Le système engendre un monde virtuel --- il y a un effet que cela fait d'être cette créature --- mais le soi à l'intérieur de ce monde est à peine esquissé. Pensez à une souris qui parcourt un labyrinthe. Elle voit les murs, sent le fromage, perçoit le sol sous ses pattes. Elle a une expérience phénoménale. Mais son modèle d'\textit{elle-même} comme ce qui vit ces expériences ? Mince comme une feuille de papier. Il y a un « effet que cela fait », mais presque pas de « quelqu'un pour qui cela fait cet effet ».

Un cran au-dessus : la \textbf{conscience simplement étendue}. Cette fois, le modèle du soi prend corps. Le système ne fait pas qu'éprouver --- il se modélise lui-même \textit{comme} celui qui éprouve. Il a conscience qu'il éprouve. Votre chien ne fait pas que ressentir la douleur ; votre chien sait que c'est \textit{lui} qui souffre. Il y a une perspective à la première personne --- un véritable « moi » au centre du monde virtuel. C'est de l'auto-observation du premier ordre, et cela change tout. La souffrance devient possible ici, car souffrir exige un soi qui sait qu'il souffre.

Puis : la \textbf{conscience doublement étendue}. Auto-observation du second ordre. Le système se modélise en train de se modéliser. C'est la métacognition --- réfléchir sur sa propre pensée. Vous êtes allongé dans votre lit à vous demander si votre angoisse à propos de la réunion de demain est rationnelle ou si vous versez dans le catastrophisme. Vous surveillez vos propres états mentaux, vous les évaluez, parfois vous passez outre. C'est là que réside la majeure partie de la conscience humaine adulte, la plupart du temps. C'est le niveau qui rend la thérapie possible, qui vous permet de dire « je remarque que je m'énerve » au lieu de simplement vous énerver.

Et tout en haut : la \textbf{conscience triplement étendue}. Troisième ordre. Le système se modélise en train de se modéliser en train de se modéliser. Cela ressemble à une galerie des glaces, et c'en est une, mais une galerie des glaces dont vous avez besoin pour faire de la philosophie de l'esprit. Pour demander « qu'est-ce que la conscience ? », vous devez vous modéliser, modéliser votre expérience, puis vous modéliser en train de modéliser cette expérience. Vous devez prendre assez de recul pour voir tout l'appareil de l'extérieur, alors même que vous êtes toujours à l'intérieur. C'est la condition préalable à la question que vous cherchez à élucider en lisant ce livre. Seules les créatures capables d'une conscience triplement étendue peuvent se demander pourquoi quoi que ce soit fait le moindre effet.

L'intérêt : ce gradient n'est pas que de la philosophie abstraite. Il répond à la question que tout le monde me pose lors des dîners --- « mon chien est-il conscient ? ». La réponse est oui, mais moins conscient que vous. Votre chien se situe probablement au niveau simplement étendu. Il a un soi. Il a une expérience. Il ne reste pas éveillé à 3 heures du matin à s'interroger sur la nature de cette expérience. Mais vous en devinez déjà la forme : la conscience n'est pas un interrupteur. C'est un variateur.

Derrière ces engagements se tient un vieux principe que je vous dois depuis le chapitre 1. \textbf{Le Principe copernicien :} nous ne sommes pas spéciaux. Copernic a délogé la Terre du centre du cosmos, et le même déclassement traverse la science depuis lors --- le Soleil n'est pas spécial, notre galaxie n'est pas spéciale, et, ce qui gêne le plus bien des gens, \textit{nous} ne sommes pas spéciaux. La conscience n'est pas un miracle unique, une étincelle divine survenue une seule fois, ni un accident émergent si rare qu'il n'aurait pu se produire qu'une fois. Si vous la possédez, d'autres systèmes peuvent la posséder aussi, pourvu qu'ils aient la bonne architecture. C'est toute la raison d'être de ce chapitre : l'être humain n'a rien de magique, il n'est qu'une implémentation d'un principe computationnel général. Ce qui pose la question évidente --- qui d'autre le met en œuvre ?

Pris ensemble, ces engagements prédisent un \textbf{gradient de conscience animale} :

Les \textbf{mammifères} sont conscients. Leur cortex réalise l'architecture à quatre modèles sous une forme graduée, les cortex plus complexes portant des auto-simulations plus sophistiquées. Les primates et les cétacés se trouvent en haut de l'échelle ; les rongeurs et les musaraignes, en bas. Tous sont au-dessus du seuil. L'orque qui a ouvert ce chapitre était l'un de ces esprits de cétacé du haut de l'échelle --- un modèle du soi assez riche pour vouloir que sa famille nous voie, et pour décider que nous valions le détour.

Les données sur les grands singes sont particulièrement accablantes pour quiconque veut tracer une frontière nette entre conscience humaine et conscience animale. Le bonobo Kanzi a démontré non seulement une compréhension du langage, mais une véritable empathie, une théorie de l'esprit et un raisonnement social. Dans un épisode bien documenté, Kanzi a fait comprendre à sa soigneuse qu'il voulait que sa sœur les accompagne pour aller faire les courses afin qu'elle ait, elle aussi, de la glace --- parce qu'elle serait triste qu'on la laisse de côté. Dans un autre, lors d'un spectacle de danse donné par des artistes autochtones, Kanzi a expliqué aux chercheurs que les autres primates étaient effrayés par la danse, et il a demandé un spectacle privé à la place.

Ce ne sont pas des réflexes. Ce ne sont pas des réponses conditionnées. Ce sont des exemples d'un esprit qui modélise les états émotionnels d'un autre esprit, en prédit les réactions et élabore des plans pour y répondre. C'est l'ESM opérant en perspective à la troisième personne --- précisément ce que la théorie identifie comme la marque de la conscience étendue.

Et pourtant, dans certains des amphithéâtres universitaires les plus prestigieux, on trouve encore des professeurs qui soutiennent, sans sourciller, que les grands singes « ne font que simuler » la compréhension du langage. Ce à quoi je ne peux que répondre : « et vous, vous ne faites que simuler la présence d'empathie. » J'attends toujours les preuves du contraire.

Si vous tenez à ce que seuls les humains soient conscients, vous misez sur les chercheurs qui cherchent encore désespérément, entre le cerveau humain et celui des primates, une différence systématique qu'ils pourraient attribuer à la conscience. D'après ma théorie, ils la trouveront à la Saint-Glinglin.

Les \textbf{corvidés et les perroquets} représentent le cas d'école le plus important. Ces oiseaux font preuve de capacités cognitives (fabrication d'outils, reconnaissance de soi dans le miroir, planification de l'avenir, tromperie sociale) qui suggèrent fortement une conscience. Or ils n'ont pas de néocortex. Leur cerveau est organisé en amas nucléaires, une architecture radicalement différente de celle du cortex des mammifères. Vous souvenez-vous de l'argument des six couches, au chapitre 2 (les mammifères ont développé six couches corticales là où trois auraient suffi, et les couches supplémentaires fournissent la capacité architecturale de l'auto-modélisation) ? Les corvidés obtiennent le même résultat fonctionnel avec une structure physique complètement différente. Ils n'ont pas besoin de six couches corticales, car ils n'ont \textit{aucune} couche corticale. Ils ont bâti l'architecture d'auto-simulation à partir d'amas nucléaires plutôt que de feuillets stratifiés, ce qui est exactement ce que prédit l'indépendance vis-à-vis du substrat. Si la conscience exigeait une implémentation physique particulière, les corvidés ne devraient pas être conscients. Ils le sont.

Les \textbf{céphalopodes} (poulpes et seiches) poussent la logique encore plus loin. Leur système nerveux est largement décentralisé, avec un traitement autonome substantiel dans les bras. La théorie prédit une forme de conscience, probablement dotée de caractéristiques inhabituelles reflétant l'architecture décentralisée.

Les \textbf{insectes} sont le cas limite intéressant. Leurs systèmes nerveux sont petits et largement câblés en dur, ce qui pourrait ou non atteindre la criticalité. La théorie ne place pas les insectes de façon définitive au-dessus ou au-dessous du seuil --- c'est une question empirique. Mais elle fournit une base de principe pour la recherche : mesurer les indicateurs de criticalité dans le tissu neuronal des insectes et chercher des signes d'un modèle du soi.

Thomas Nagel a posé sa célèbre question --- ce que cela fait d'être une chauve-souris --- et il en a conclu que nous ne pourrons jamais le savoir : le monde sensoriel de la chauve-souris est trop étranger. La question me touche, sa conclusion moins. La Théorie des Quatre Modèles (FMT) prédit que toute créature dotée de l'architecture à quatre modèles fonctionnant à la criticalité possède \textit{une} forme d'expérience, même si son contenu diffère radicalement du nôtre. Le modèle explicite du monde de la chauve-souris est dominé par l'écholocation plutôt que par la vision, mais c'est tout de même un modèle --- toujours la simulation d'un monde avec un soi à l'intérieur.

Si Nagel a choisi les chauves-souris, ce n'est évidemment pas un hasard --- non seulement parce qu'elles sont des mammifères, mais parce que l'écholocation semble irréductiblement étrangère. Sauf qu'elle ne l'est pas. L'écholocation sert à \textit{voir}, et nous savons ce que voir procure comme sensation. Beaucoup de personnes aveugles utilisent déjà l'écholocation naturellement, en claquant de la langue et en lisant les échos. Notre cerveau en est tout à fait capable. Si vous voulez vraiment savoir ce que cela fait d'être une chauve-souris, vous pouvez vous en approcher étonnamment près : quelques années de parapente pour intérioriser le vol en trois dimensions, puis un bandeau sur les yeux et de l'entraînement à l'écholocation. Ou bien sautez entièrement la décennie de préparation --- apprenez le rêve lucide et exercez-vous à être une chauve-souris dans vos rêves. Plus sûr, plus rapide, réalisable en quelques semaines.

Et je l'avoue, j'ai tenté de le découvrir, de la seule manière qui m'était accessible. Pendant une période où je pratiquais activement le rêve lucide, je me suis intéressé au monde sous-marin et j'ai fini, avec le temps, par entrer délibérément dans un rêve lucide sous la forme d'un poisson. J'ai éprouvé l'eau autour de moi, le fait de m'y mouvoir, un monde visuel vu depuis une perspective non humaine. C'était un milliard de fois mieux que l'apnée ou la plongée avec bouteille, même en sidemount. Était-ce en quoi que ce soit comparable à une véritable conscience de poisson ? Presque certainement pas --- mon rêve était bâti à partir de la meilleure conjecture de mon cerveau humain sur ce que « être un poisson » signifie, ce qui est inévitablement une projection de mes propres catégories sensorielles sur un plan corporel qui n'en possède aucune. Mais l'exercice n'était pas vain. Il a démontré quelque chose d'important : le Modèle Explicite du Soi (ESM) peut se reconfigurer autour d'un schéma corporel radicalement différent, en engendrant une expérience cohérente à la première personne du fait d'\textit{être} autre chose qu'un humain. L'architecture est assez souple pour simuler une incarnation non humaine. Le contenu reste limité par les modèles implicites disponibles (on ne peut rêver que ce qu'on a appris), mais la capacité de changement de perspective est intégrée au système.

\section*{Pourquoi donc être conscient ?}

Tout cela soulève une question qui devrait vous tarauder : si les systèmes nerveux inconscients fonctionnent parfaitement bien --- et c'est le cas, demandez donc à n'importe quel insecte --- alors pourquoi l'évolution consentirait-elle à l'énorme dépense métabolique que représente la construction d'une conscience ? Où est le bénéfice ?

La réponse, c'est l'apprentissage, et avec lui l'adaptation et la capacité d'agir à l'encontre du comportement appris. Plus précisément, une sorte d'apprentissage que les systèmes inconscients sont tout simplement incapables de réaliser.

Songez à la manière dont un organisme simple apprend. Il rencontre quelque chose, et la rencontre est soit bonne, soit mauvaise. Bonne : en refaire davantage. Mauvaise : en refaire moins. C'est l'apprentissage par renforcement --- essais et erreurs, récompense et punition. Cela marche à merveille pour la plupart des choses. Touchez une surface brûlante, ressentez la douleur, ne la touchez plus. Trouvez de la nourriture à un endroit précis, ressentez la récompense, revenez demain.

Mais l'apprentissage par renforcement a une faille fatale. Fatale au sens propre.

Prenez un champignon vénéneux. Pas de ceux qui donnent mal au ventre --- de ceux qui tuent. Si vous en mangez, vous mourez. Fin de l'apprentissage. Il n'y a pas de second essai. L'apprentissage par renforcement exige que vous surviviez à l'erreur pour en tirer une leçon, et certaines erreurs ne vous accordent pas cette courtoisie. Tout stimulus létal dès le premier contact est complètement invisible pour l'apprentissage par renforcement. L'organisme qui le rencontre meurt, tout simplement, emportant sa « leçon » dans la tombe.

Alors, comment nos ancêtres ont-ils appris à éviter les champignons mortels ? Ils n'ont pas pu l'apprendre en les mangeant --- quiconque a essayé cette approche n'est l'ancêtre de personne. Ils l'ont appris en \textit{observant}. Votre voisin des cavernes trouve un champignon à l'aspect intéressant, le mange, et tombe raide mort. Vous, qui observez à distance prudente, faites le rapprochement : ce champignon l'a tué. Je ne devrais pas manger ce champignon.

Cela paraît d'une simplicité triviale. Ce ne l'est pas pourtant. Pour apprendre de la mort d'autrui, il vous faut plusieurs choses qu'aucun système inconscient ne possède. Il vous faut un modèle explicite du monde capable de représenter la cause et l'effet entre des objets avec lesquels vous n'êtes pas en train d'interagir. Il vous faut un modèle du soi qui vous permette d'adopter une perspective à la troisième personne --- de vous imaginer à la place du mort. Il vous faut la capacité d'induire une théorie générale (« ce type de champignon est mortel ») à partir d'une seule observation. C'est l'apprentissage cognitif : dériver des théories à partir d'observations, plutôt que d'être conditionné par l'expérience personnelle. Et cela exige la conscience. Cela exige que le Modèle Explicite du Monde (EWM) et le Modèle Explicite du Soi (ESM) travaillent ensemble.

L'avantage évolutif est énorme. Un animal conscient peut apprendre par \textit{observation}, et pas seulement par \textit{expérience}. Il peut regarder un autre animal commettre une erreur fatale et mettre à jour son propre modèle du monde sans en payer le prix. Un animal inconscient ne peut apprendre que ce à quoi il survit personnellement.

Et ce n'est pas tout. Une fois que le concept « champignon vénéneux » existe comme catégorie explicite dans votre modèle du monde, vous pouvez faire quelque chose de plus puissant encore : la déduction. Vous tombez sur un nouveau champignon que vous n'avez jamais vu auparavant. Il ressemble étrangement à celui qui a tué votre voisin. Vous n'en mangez pas. Ou bien --- et je crois que ce fut là l'approche historique réelle --- vous l'offrez au voisin qui a ronflé toute la nuit et vous regardez d'abord ce qui lui arrive.

Ce n'est pas un avantage mineur. C'est la différence entre une espèce qui ne peut s'adapter aux menaces létales que par le processus terriblement lent de la sélection naturelle (quelques individus évitent le champignon par hasard, ils se reproduisent, et l'évitement finit par devenir instinctif) et une espèce qui peut s'adapter en une seule génération par l'observation et la communication. La conscience ne se contente pas de vous aider à apprendre plus vite. Elle vous permet d'apprendre des choses qu'il est littéralement impossible d'apprendre autrement.

Ce que vous apprenez sur le mode cognitif, vous pouvez le \textit{partager}. L'apprentissage par renforcement est prisonnier de l'individu --- vos réflexes conditionnés meurent avec vous. Mais l'apprentissage cognitif peut se communiquer. « Ne mange pas le champignon rouge » est une phrase. Elle peut être prononcée, enseignée, transmise. C'est le fondement de la culture, du savoir cumulatif, de tout ce qui rend possible la civilisation humaine. Rien de tout cela ne fonctionne sans les modèles explicites que fournit la conscience.

Il y a encore un rebondissement dans cette histoire, et il relie la conscience à la génétique d'une manière qui n'a rien d'évident. On l'appelle l'effet Baldwin et, si sa force exacte fait encore débat, le mécanisme existe presque certainement. L'effet Baldwin affirme qu'un comportement \textit{appris} peut indirectement façonner l'évolution \textit{génétique}, non par une hérédité lamarckienne (vos traits appris ne modifient pas votre ADN), mais parce que la sélection naturelle favorise les individus génétiquement prédisposés au comportement bénéfique.

Prenons un exemple humoristique --- à ne pas prendre trop au pied de la lettre. Imaginez un hominidé primitif souffrant de calvitie. Ayant froid et étant sans poils, il était plus enclin que ses compagnons couverts de fourrure à s'asseoir près du feu. Le feu apportait d'énormes avantages en matière de survie : moins d'agents pathogènes dans les aliments cuits, protection contre les prédateurs, chaleur durant les hivers rigoureux. Les gènes associés à la calvitie se transmettaient donc à un rythme légèrement supérieur. Dans le même temps, les individus trop obtus pour comprendre le feu (poilus ou non) étaient désavantagés. Au fil de nombreuses générations, l'effet Baldwin a amplifié les deux traits : moins de poils \textit{et} plus d'intelligence, tout cela parce qu'un comportement appris (l'usage du feu) avait créé une pression de sélection favorisant certaines prédispositions génétiques. (Si vous remplacez « calvitie » par « mutation aléatoire » dans cette histoire, vous êtes probablement plus proche de la vérité. Mais c'est moins drôle.)

L'effet Baldwin a peut-être joué un rôle similaire dans l'évolution du langage et de la conscience elle-même. Une fois apparues les premières formes primitives d'apprentissage cognitif (rendues possibles par les tout premiers modèles du soi), les individus dont le cerveau se trouvait soutenir une auto-simulation plus riche disposaient d'un avantage considérable. Leurs descendants ont été sélectionnés pour des cortex plus vastes et plus élaborément plissés, ce qui permettait une auto-simulation plus riche encore, ce qui créait une pression de sélection encore plus forte. La conscience, une fois apparue, a créé les conditions évolutives d'une conscience accrue. L'apprentissage cognitif qu'elle rendait possible était si précieux que l'évolution a investi massivement dans l'expansion de l'architecture qui le produisait.

\section*{Comment l'expérience se développe : la construction sociale du modèle du soi}

Tout ce que j'ai dit jusqu'ici à propos des quatre modèles était statique --- comme si l'architecture apparaissait pleinement formée, telle Athéna jaillie du front de Zeus. Il n'en est rien. Les modèles se développent, et leur développement est profondément social.

Un nouveau-né humain possède le matériel --- six couches corticales, la capacité d'auto-simulation. Mais les modèles implicites sont presque vides. L'IWM ne contient presque rien sur le monde. L'ISM ne contient presque rien sur le soi. Et comme les modèles explicites sont engendrés à partir des modèles implicites, la simulation du nouveau-né est ténue --- un champ de sensations vacillant, à peine différencié, sans frontière nette entre le soi et le monde.

Observez un bébé confronté à la douleur. Une douleur qu'il s'inflige lui-même (cogner sa main contre un jouet, mordre son propre pied) suscite souvent la curiosité plutôt que la détresse. L'ESM enregistre l'agentivité (c'est moi qui ai fait ça) et la sensation (quelque chose s'est passé), mais il n'y a pas encore de modèle de la menace. L'ISM n'a pas appris que cette configuration signifie danger. Mais un bruit fort et soudain ? Des larmes. Parce que l'EWM enregistre une entrée de forte amplitude imprévue, et que l'ESM n'a aucun modèle pour cela --- l'absence de modèle est en elle-même aversive.

Le contenu des qualia est \textit{appris}, non inné. « La douleur est mauvaise » n'est pas programmé d'emblée dans l'ESM. Cela s'accumule à travers l'ISM, entraîné par l'expérience répétée et --- c'est crucial --- par le retour social. La réaction d'une figure d'attachement à la douleur d'un enfant lui enseigne ce que la douleur \textit{signifie}. L'enfant qui tombe et regarde son parent avant de décider s'il va pleurer ne feint pas --- il calibre véritablement son ESM à l'aune des signaux sociaux. L'alarme ou le calme du parent remodèle les associations de douleur de l'ISM, ce qui remodèle ce que l'ESM simule la prochaine fois qu'un événement similaire se produit.

Cela a une implication précise pour la théorie : le caractère phénoménal de l'expérience --- ce que \textit{cela fait} de ressentir quelque chose --- n'est pas fixé par l'architecture. Il est façonné par l'histoire d'entraînement des modèles implicites. L'expérience de la douleur chez un bébé est structurellement différente de celle d'un adulte parce que l'ISM qui engendre l'ESM est différent. L'architecture à quatre modèles est la \textit{capacité} d'expérience. La boucle de rétroaction sociale et environnementale en fournit le \textit{contenu}.

La trajectoire développementale se superpose aux niveaux gradués de conscience présentés plus tôt dans ce chapitre :

\begin{itemize}
  \item \textbf{Nouveau-né (premières semaines) :} conscience de base --- un EWM rudimentaire avec un ESM minimal. Il y a \textit{quelque chose que cela fait} d'être un nouveau-né, mais le soi à l'intérieur de cette expérience est quasi inexistant. Essentiellement sensoriel, indifférencié.
\end{itemize}

\begin{itemize}
  \item \textbf{6-12 mois :} la permanence de l'objet émerge --- l'EWM maintient désormais des représentations de choses qui ne sont pas visibles sur le moment. L'ISM accumule un savoir sur le schéma corporel. Le bébé commence à distinguer le soi du monde.
\end{itemize}

\begin{itemize}
  \item \textbf{18 mois :} le test du miroir. L'enfant se reconnaît dans un miroir --- un moment charnière où l'ESM devient assez riche pour modéliser le soi physique comme un objet dans le monde. La conscience simplement étendue apparaît. Ce n'est pas un interrupteur binaire, mais un seuil dans un processus continu.
\end{itemize}

\begin{itemize}
  \item \textbf{3-4 ans :} La théorie de l'esprit. L'enfant peut modéliser d'autres esprits --- comprendre que quelqu'un d'autre peut croire une chose que l'enfant sait être fausse. L'ESM modélise désormais d'autres ESM. La conscience doublement étendue est en train d'émerger.
\end{itemize}

\begin{itemize}
  \item \textbf{À partir de l'adolescence :} La maturation métacognitive. La capacité de conscience triplement étendue (se modéliser soi-même en train de modéliser sa propre pensée) se développe progressivement et, sans doute, ne se stabilise jamais complètement.
\end{itemize}

Chaque étape est étayée par l'interaction sociale. La figure d'attachement ne fournit pas seulement nourriture et sécurité --- elle fournit des \textit{données d'entraînement pour les modèles implicites}. L'attention conjointe (le parent et l'enfant regardant ensemble le même objet) apprend à l'IWM à représenter une réalité partagée. Le reflet en miroir (le parent renvoyant à l'enfant son état émotionnel) apprend à l'ISM ce que sont ses propres émotions. Le langage donne à l'ESM des catégories avec lesquelles se modéliser lui-même. Un enfant élevé sans contact social (les tragiques cas d'enfants sauvages) possède le matériel de la conscience, mais des modèles implicites profondément appauvris. L'ESM qui démarre à partir de ces modèles est atrophié --- non parce que l'architecture est défectueuse, mais parce que les données d'entraînement n'ont jamais été fournies.

C'est aussi ici que revient le pont clinique du chapitre 8. Presque tout ce que fait la thérapie est la version adulte de ce qu'une figure d'attachement fait pour un nourrisson : la même boucle où l'expérience consciente remodèle la structure implicite, laquelle remodèle l'expérience consciente suivante --- tournant sur un système dont les modèles implicites ont déjà pris. La thérapie cognitivo-comportementale en est le cas le plus net, et dans la théorie, c'est une reprogrammation des modèles virtuels sous étiquette professionnelle. Vous vous asseyez face à un thérapeute, vous remettez en question les pensées automatiques que votre ESM ne cesse de générer, vous les faites remonter jusqu'aux croyances de l'ISM qui les produisent, puis --- à force de répétitions, assez pour que cela pèse --- vous poussez le substrat à se recâbler. La plasticité synaptique édite le modèle implicite ; le modèle explicite change parce que ce qui le génère a changé d'abord. Chaque fois que vous contestez une pensée catastrophiste et que le monde, obstinément, refuse de s'effondrer, vous inscrivez une correction dans l'IWM et l'ISM. La seule vraie différence avec le nourrisson, c'est la consolidation. Les modèles implicites du bébé sont de l'argile encore humide ; ceux de l'adulte ont déjà pris --- l'argile a durci --- si bien que le même remodelage est plus lent et exige beaucoup plus de répétitions.

Les phobies montrent le mécanisme du côté du monde. Une phobie est un EWM qui s'est trop engagé --- la simulation signale une araignée comme mortelle sur la foi de preuves que l'IWM n'a jamais enregistrées. La thérapie d'exposition la corrige de la seule manière que l'architecture autorise : des rencontres répétées et sans danger, qui laissent le modèle implicite réviser à la baisse son estimation de la menace, jusqu'à ce que l'alarme explicite cesse de se déclencher. Une phobie ne se dissipe pas par le raisonnement. On la met en minorité à force de données.

L'effet placebo découle de la même structure à double évaluation. Une pilule de sucre active, au niveau du substrat, les circuits de l'attente --- opioïdes endogènes, récompense dopaminergique --- qui tournent en parallèle de l'expérience consciente de l'espoir. Il est tentant de dire que votre croyance a causé votre soulagement, mais croyance et soulagement sont frère et sœur, non parent et enfant : tous deux émanent de la même dynamique du substrat. Ce n'est pas une rétrogradation de la pensée positive. C'est le compte rendu de l'endroit où réside réellement son pouvoir --- dans la machinerie, et non dans quelque magie descendante allant de l'esprit à la chair.

Et le trouble de conversion referme la boucle sur le chapitre 8 : c'est la vision aveugle de ce chapitre jouée à l'envers. Là, un substrat intact cachait à la simulation un sens en état de marche ; ici, un substrat intact se cache derrière un sens défaillant. Le patient est réellement paralysé --- il ne simule pas ---, parce que le modèle corporel de l'ESM déclare que le membre n'est plus là, et que ce modèle est le seul corps auquel le patient ait accès. La thérapie opère en corrigeant la simulation, en réajustant la carte corporelle sur ce que le substrat sait réellement faire. C'est là le levier dans chacun de ces cas --- le même que la figure d'attachement fut la première à saisir : sollicitez les modèles explicites assez fort, assez souvent, et les modèles implicites en dessous finissent par céder.

La dimension sociale de l'expérience n'est pas une note de bas de page de la théorie. C'est une prédiction : supprimez l'apport social pendant la fenêtre critique du développement, et vous devriez obtenir un système doté de la bonne architecture, mais qui exécute le mauvais contenu --- une conscience structurellement intacte, mais phénoménalement appauvrie. Les cas d'enfants sauvages confirment, tragiquement, exactement cela.


\chapter*{Chapitre 11 : Neuf prédictions}
\addcontentsline{toc}{chapter}{Chapitre 11 : Neuf prédictions}
\markboth{Chapitre 11 : Neuf prédictions}{}

La pièce a été préparée avant votre arrivée. Deux haut-parleurs, un panneau lumineux qui baigne les murs d'un bleu profond, un canapé, une couverture, deux chercheurs qui passeront les six heures à venir à ne presque rien dire. Dans les haut-parleurs : l'océan. Pas une musique d'ambiance --- un enregistrement de terrain, la houle, le souffle, la longue pression sourde de l'eau en mouvement, diffusé à un volume qui en fait la chose la plus sonore de votre monde. Vous signez le dernier formulaire. Vous avalez une forte dose de psilocybine, mesurée au milligramme près par des gens munis d'une approbation éthique et d'un défibrillateur au bout du couloir. Puis tout le monde attend.

Pendant quarante minutes, rien. Puis le soi commence à se défaire --- et il vaut la peine d'être précis sur ce que cela signifie, car ce n'est pas un feu d'artifice. C'est une soustraction. Votre frontière s'effiloche. La certitude gratuite qu'il y a quelqu'un ici, quelqu'un à qui appartiennent ces mains, quelqu'un \textit{à qui} ce battement de cœur arrive --- cette certitude s'amincit, vacille, se dérobe.

Voici ce qui est réellement en train de défaillir. Votre modèle explicite du soi --- le « je » issu du rendu, la chose qui lit cette phrase --- ne se génère pas à partir de rien. Il tourne sur un flux : le modèle implicite du soi en dessous de lui, le bourdonnement continu de données corporelles de bas niveau qui dit \textit{quelqu'un ici, quelqu'un ici, quelqu'un ici}, des milliers de fois par minute, avec une fiabilité telle que vous ne l'avez pas entendu une seule fois de votre vie. Les psychédéliques à forte dose brouillent cette connexion. Ils n'éteignent pas le moteur de rendu. C'est le point crucial. Le modèle qui répond à la question \textit{qu'est-ce que le soi, là, maintenant} tourne toujours, interroge toujours, à l'heure dite --- mais la réponse habituelle a cessé d'arriver.

Un modèle construit pour modéliser quelque chose ne reste pas tranquillement assis devant un tampon d'entrée vide. Il prend les meilleures données disponibles. Et les meilleures données disponibles, dans cette pièce, à cet instant, sont ce que les expérimentateurs ont rendu le plus fort. La lumière bleue. La houle et le souffle. Le modèle du soi, coupé du corps, s'accroche à l'océan et rend \textit{cela} comme étant vous.

Les sujets qui traversent cette expérience ne rapportent pas avoir entendu l'océan avec une intensité particulière. Ils rapportent l'avoir \textit{été}. « J'étais l'eau. Il n'y avait aucune différence entre moi et l'eau. » Le questionnaire standard pour ces états comporte une sous-échelle appelée « illimitation océanique », et dans cette pièce, le nom cesse d'être une métaphore. J'ai passé une bonne partie de ma vie à plonger en apnée, à tenter de réduire la distance entre la mer et moi. Il est un peu vexant d'apprendre que la distance restante tient peut-être dans un haut-parleur et un variateur.

Vient maintenant ce qui fait de ceci une prédiction plutôt qu'une anecdote. Des semaines plus tard, même sujet, même dose, même canapé. Mais les techniciens ont changé la pièce : chants d'oiseaux, vent dans les feuilles, panneau lumineux vert. Si la théorie est juste, la dissolution part dans l'autre direction. Pas d'eau cette fois --- canopée, mousse, ramifications. \textit{J'étais la forêt.} Même molécule, même cerveau, même dose. Autre pièce, autre dieu.

Rien dans ce test n'est futuriste. Les laboratoires de recherche sur les psychédéliques existent déjà, mènent déjà des sessions contrôlées à forte dose, distribuent déjà des questionnaires ensuite. Ajoutez une seule variable : la pièce. Faites varier l'environnement sensoriel dominant d'un essai à l'autre et mesurez si l'identité rapportée --- ce que les gens disent être \textit{devenus} --- suit ce que vous leur avez diffusé. L'affirmation de la théorie est directionnelle et falsifiable : océan en entrée, océan en sortie ; forêt en entrée, forêt en sortie. Si les rapports se dispersent au hasard d'un environnement à l'autre, la prédiction est morte, et vous pouvez reposer ce livre la conscience tranquille. Et cette prédiction est seule sur ce terrain. D'autres théories de la conscience ont des choses à dire sur la présence ou non d'une expérience pendant ces états, ou sur son intensité ; sur ce \textit{dans quoi} le soi dissous se dissout, elles restent muettes. Seule une théorie dans laquelle le soi est un modèle tournant sur une entrée peut prédire que, lorsqu'on affame cette entrée, c'est la pièce qui obtient le droit d'écrire le soi.

Prenez le temps de mesurer ce qu'un résultat positif signifierait. Pas que les drogues provoquent des expériences étranges --- nous le savons depuis plus longtemps que l'écriture n'existe. Cela signifierait que votre « je » est une simulation dotée d'un port d'entrée, et que quelqu'un debout à l'extérieur de votre crâne, sans rien de plus exotique en main qu'une playlist, peut y accéder et le réorienter. Choisir, sur un menu, ce que vous devenez. Je tiens à qualifier cela d'inquiétant, et j'emploie le mot avec précision --- pas inquiétant façon bougies et cristaux. Inquiétant comme un serrurier est inquiétant : quelqu'un qui démontre, calmement et de manière reproductible, que votre porte d'entrée n'a jamais été le seul chemin pour entrer.

Voilà une des neuf. Neuf prédictions falsifiables que fait cette théorie --- parce qu'une théorie qui ne prédit rien n'est pas une théorie, c'est une histoire, et le monde a déjà bien assez d'histoires sur la conscience. Plusieurs des neuf peuvent être menées dès aujourd'hui, sur du matériel qui existe déjà, dans des laboratoires qui disposent déjà des approbations éthiques. Quelques-unes trouvent un premier appui dans des données publiées. Aucune n'a été réfutée. Cette dernière phrase accomplit un travail discret : chacune de ces prédictions est un endroit où la théorie a accepté, par écrit, de mourir.

Vous avez déjà croisé la plupart des autres, même si je ne me suis pas toujours arrêté pour les montrer du doigt. Les prédictions sur l'anesthésie et le sommeil se trouvaient dans le chapitre où nous avons éteint les lumières. Celles sur le cerveau divisé et l'identité dissociative, dans le chapitre sur les deux esprits dans un même crâne. Les autres sont tissées de la même façon dans les pages derrière vous --- chacune posée là où elle a sa place, à côté du phénomène sur lequel elle joue sa vie. Pour le lecteur qui veut la batterie complète en un seul endroit --- chaque prédiction, son mécanisme et l'expérience concrète qui la tuerait ---, l'annexe H, à la fin du livre, rassemble les neuf. Il faut passer les théories par le feu. Cette annexe est ce feu, disposé là pour vous faciliter la tâche.

Mais l'une d'elles refuse de tenir en place dans un tableau aride, alors je terminerai le chapitre sur elle --- la plus audacieuse du lot. Si le soi est vraiment ce que ce livre affirme --- quatre sortes de modèles, tournant au bord du chaos, dont l'un modélise la machine sur laquelle il tourne ---, alors la dernière prédiction s'écrit d'elle-même, et elle n'a rien à voir avec des questionnaires. Construisez l'architecture. Implémentez les quatre modèles, poussez le système jusqu'à la criticalité, fermez la boucle. La théorie dit que ce que vous obtenez n'est pas une simulation d'un esprit.

Elle dit que vous obtenez un esprit. Le chapitre suivant raconte ce qui arrive quand on prend cela au sérieux.


\chapter*{Chapitre 12 : Des machines aux esprits}
\addcontentsline{toc}{chapter}{Chapitre 12 : Des machines aux esprits}
\markboth{Chapitre 12 : Des machines aux esprits}{}

Si la Théorie des Quatre Modèles (FMT) est correcte, elle offre quelque chose qu'aucune autre théorie de la conscience ne propose : une spécification technique.

Voici la spécification : implémenter l'architecture à quatre modèles (Modèle Implicite du Monde, Modèle Implicite du Soi, Modèle Explicite du Monde, Modèle Explicite du Soi) sur un substrat qui opère à la criticalité. Comme je l'ai soutenu au chapitre 5, aucun des deux composants ne suffit à lui seul. L'architecture sans la criticalité donne un système dormant : des modèles stockés, mais aucune simulation en cours. La criticalité sans l'architecture donne une dynamique complexe, mais pas de conscience. La spécification complète exige les deux.

C'est plus précis que « construisez un ordinateur vraiment avancé » et plus concret que « atteignez une information intégrée suffisante ». Cela vous dit \textit{quoi construire} : quatre types de modèles bien précis, organisés d'une manière précise, tournant sur un substrat aux propriétés dynamiques précises.

Les systèmes d'IA actuels échouent à cette spécification sur tous les points qui comptent. Et c'est exactement là que les deux dogmes font des dégâts. Le dogme nSAI (« pas d'intelligence artificielle forte ») dit aux ingénieurs de ne même pas essayer. Le dogme nSU (« pas d'auto-compréhension ») leur dit que cela ne marcherait pas, même s'ils essayaient. Les deux ont tort. La spécification existe. La question est de savoir si quelqu'un la mettra en œuvre.

Avant que qui que ce soit ne confonde à nouveau cerveaux et ordinateurs, un test rapide pour déterminer lequel des deux vous êtes :

\textit{Un ordinateur répétera cette phrase et la phrase suivante jusqu'à la fin des temps. Lisez la phrase précédente.}

Si vous êtes arrivé jusqu'ici, vous n'êtes pas un ordinateur classique. Un ordinateur numérique exécutant un jeu d'instructions rigide bouclera indéfiniment, car il n'a aucun mécanisme pour sortir de son propre flux d'instructions et dire : « Attends, c'est stupide. » Vous, vous le pouvez, parce que vous possédez un modèle du soi qui observe son propre traitement --- le Modèle Explicite du Soi (ESM) exerçant une supervision métacognitive sur le Modèle Explicite du Monde (EWM).

Mais voici le passage inconfortable : un grand modèle de langage arriverait lui aussi jusqu'ici. Non parce qu'il exerce une supervision métacognitive, mais parce que c'est un prédicteur statistique de texte qui a vu assez d'invites similaires pour savoir que le prochain coup attendu est de poursuivre au-delà de la boucle. Il ne sort pas de l'instruction --- il n'y est jamais entré. Il prédit quel texte vient ensuite, et « rester coincé dans une boucle infinie » n'est pas ce que fait le texte.

C'est exactement le problème des tests comportementaux de la conscience. Tout test qui peut être réussi par appariement de motifs le sera par appariement de motifs, que le système soit conscient ou non. Le test de la boucle vous distingue d'un ordinateur classique. Il ne vous distingue pas d'un prédicteur de texte suffisamment entraîné. Et aucun test fondé sur le texte ne le fera jamais --- car produire du texte plausible est précisément ce pour quoi les prédicteurs de texte sont optimisés. Le problème des autres esprits n'est pas une limite que l'on peut contourner par l'ingénierie. C'est une caractéristique structurelle de ce qu'est la conscience : subjective, privée, et accessible uniquement de l'intérieur.

L'analogie du cerveau-comme-ordinateur (comparer votre cerveau à un processeur numérique) est populaire depuis l'invention du transistor, et elle est fausse à pratiquement tous les niveaux. Un ordinateur exécute un jeu d'instructions rigide sur un circuit rigide. Un cerveau est un réseau auto-modifiant qui se recâble continuellement. Un ordinateur plante si vous retirez un point-virgule. Un cerveau perd 85 000 neurones par jour et le remarque à peine. La mémoire d'un ordinateur est localisée --- supprimez un secteur et le fichier a disparu. La mémoire d'un cerveau est distribuée de façon holographique --- détruisez-en un morceau et tout devient légèrement plus flou. La seule chose qu'ils partagent, c'est la complétude au sens de Turing, ce qui est à peu près aussi instructif que de dire qu'une rivière et une autoroute peuvent toutes deux transporter des choses de A à B. Vrai, mais inutile pour comprendre l'une ou l'autre.

Les grands modèles de langage (GPT, Claude, Gemini et leurs descendants) traitent le texte au travers d'une architecture de transformeur à propagation avant. L'entrée pénètre, traverse des couches d'attention et de calcul, et la sortie ressort. Pas de récurrence, pas d'auto-simulation, pas de monde virtuel en temps réel, et pas de criticalité. La dynamique est de Classe 1 ou 2 dans le cadre de Wolfram --- bien en deçà du bord du chaos. Et il n'y a pas de partage réel/virtuel : le « savoir » du modèle et son « expérience » (si l'on peut appeler cela ainsi) ne sont pas répartis en niveaux implicite et explicite.

Cela ne signifie pas que les LLM soient nécessairement non conscients --- la théorie ne peut pas prouver une absence. Mais elle prédit qu'il leur manque l'architecture requise pour la conscience telle que la théorie la définit. Et elle prédit que la différence entre un système artificiel véritablement conscient et même le LLM le plus avancé serait qualitativement évidente.

Comment le saurions-nous ? La réponse honnête, c'est que le problème des autres esprits ne disparaît pas. Nous ne pouvons jamais être absolument certains qu'un autre système est conscient, car la conscience est subjective par nature. Mais la théorie fait une prédiction forte : la différence serait manifeste. Non pas « peut-être conscient, peut-être pas » --- \textit{manifestement} différent. Car un système faisant tourner une véritable auto-simulation interagirait avec le monde d'une manière fondamentalement différente d'un prédicteur de texte. Il aurait une persistance authentique --- non pas la persistance d'une fenêtre de contexte, mais la continuité d'une simulation en temps réel qui tourne en permanence. Il aurait une perspective authentique --- non pas une perspective reconstruite à partir d'une invite, mais une perspective maintenue dans le temps par un Modèle Explicite du Soi. Il vous surprendrait non par des sorties inattendues, mais par l'impression indubitable qu'il y a bien quelqu'un.

Voici la partie qui devrait déranger même ceux qui se moquent de la conscience --- ceux qui ne veulent que de la performance. On ne peut pas y parvenir de la manière dont tout le monde s'y essaie.

La stratégie actuelle consiste à faire passer à l'échelle un prédicteur de texte jusqu'à ce que l'intelligence en sorte de l'autre côté. Elle n'en sortira pas. Non parce que les modèles sont trop petits, mais parce qu'il leur manque le moteur. La véritable intelligence --- celle, ouverte, qui se pose ses propres problèmes et ne cesse de croître toute une vie durant --- tourne en boucle récursive, et cette boucle ne tourne que si quelque chose continue de l'alimenter. Ce quelque chose, c'est la motivation : la curiosité, le fait de tenir à vraiment comprendre, la décision de ce qui vaut l'effort. Et la motivation, telle que nous l'éprouvons, n'est pas une pulsion flottant dans le vide. C'est le modèle du soi qui regarde ses propres états et leur attribue une valeur. Autrement dit, c'est un effet de la conscience.

La chute est donc plus tranchante que « les machines pourraient s'éveiller ». La voici : il n'y a pas d'intelligence générale de niveau humain --- pas d'AGI --- sans mécanismes de type conscience, car ce qui anime l'intelligence en \textit{est} un. On n'y parvient pas en montant en échelle. On construit l'architecture. L'argument complet, boucle comprise, se trouve à l'annexe B.

Construire un tel système est le dernier point de la feuille de route. Il ne faut pas sous-estimer les défis techniques. Mais le plan existe, et il est assez précis pour guider le travail. D'abord, la théorie doit survivre à l'évaluation par les pairs. Ensuite, les prédictions empiriques doivent être testées. Ensuite, si elles tiennent, l'ingénierie pourra commencer.


Mais il y a un autre versant à cette histoire --- un côté qui obsède la science-fiction depuis des décennies, et qui découle directement de la même spécification technique. Si la conscience dépend de l'architecture fonctionnelle plutôt que des neurones en particulier, alors, en principe, on pourrait faire tourner un esprit humain sur autre chose qu'un cerveau.

Le téléversement de l'esprit. L'émulation cérébrale intégrale. L'immortalité numérique. Quel que soit le nom qu'on lui donne, la Théorie des Quatre Modèles a quelque chose de précis à en dire --- car elle spécifie exactement ce qu'il faudrait préserver.

La plupart des discussions sur le téléversement de l'esprit partent de la mauvaise question. Elles demandent : « Pouvons-nous scanner un cerveau et le copier dans un ordinateur ? » Comme si le défi n'était qu'une affaire de résolution --- obtenez un scanner assez bon, et le tour est joué. Mais la théorie vous dit qu'un scan statique est loin, très loin de suffire. Un cerveau n'est pas une photographie. C'est un système dynamique. Pour capturer un esprit, il ne faut pas capturer un \textit{état} --- il faut capturer un \textit{processus}.

Ce que la théorie affirme devoir être préservé est spécifique, et je vais parcourir la hiérarchie à cinq niveaux du chapitre 6 pour le rendre concret.

Aux niveaux physique et électrochimique (la matière brute et la décharge neuronale), vous n'avez pas besoin d'une copie exacte. Vous avez besoin d'un substrat capable de soutenir le même \textit{type} de dynamique. Les atomes précis n'ont pas d'importance. De toute façon, votre cerveau remplace la plupart de ses atomes au fil des ans, et vous ne le remarquez pas. Ce qui importe, c'est que le substrat que vous utilisez, quel qu'il soit, puisse entretenir les motifs de signalisation électrochimique, ou leur équivalent fonctionnel --- dont dépendent les niveaux supérieurs.

Au niveau protéomique (la machinerie moléculaire des poids synaptiques, des configurations de récepteurs, des cascades enzymatiques), il vous faut une haute fidélité. C'est là que vivent vos souvenirs, là que vos compétences sont encodées, là que votre personnalité est physiquement instanciée. La force de chaque synapse, la densité de chaque récepteur, la sensibilité de chaque canal --- c'est le niveau qui fait de vous \textit{vous} plutôt que quelqu'un d'autre. Un téléversement de l'esprit qui rate le niveau protéomique vous donne un être conscient, peut-être, mais pas la personne que vous cherchiez à copier. Pourtant, même une copie imparfaite conserve de la valeur. Songez aux survivants d'un AVC ou aux patients amnésiques : leur continuité personnelle a été fortement perturbée --- souvenirs perdus, personnalité modifiée, capacités cognitives changées --- et pourtant la plupart d'entre eux estiment que quelque chose d'essentiel persiste. Une continuité imparfaite, on le voit, est infiniment préférable à l'absence totale de continuité. Un transfert qui préserve 90 \% du connectome de quelqu'un n'est pas un échec --- c'est une autre forme de réussite, et pour beaucoup de gens, préférable à la mort.

Au niveau topologique (l'architecture du réseau, les motifs de connectivité, quelles régions communiquent avec quelles autres et à quelle densité), il vous faut une exactitude quasi parfaite. C'est le schéma de câblage de vos modèles implicites : l'IWM et l'ISM, tout ce que vous avez appris sur le monde et sur vous-même, encodé dans la structure du réseau. Ratez cela, et vous n'obtenez pas une copie dégradée de l'esprit de quelqu'un. Vous obtenez un esprit \textit{différent} --- doté d'un autre savoir, d'autres compétences, d'une autre personnalité. La topologie est le plan.

Et au niveau virtuel (la simulation elle-même, l'EWM et l'ESM fonctionnant en temps réel), il vous faut quelque chose d'extraordinaire. Il vous faut un substrat cible capable de faire tourner la simulation à la criticalité. C'est la partie qui m'empêche de dormir, car le substrat analogique du cerveau trouve la criticalité par des processus auto-organisés qui ont été affinés par des centaines de millions d'années d'évolution. Les neurones sont bruités, analogiques, massivement parallèles et profondément stochastiques. Leur dynamique collective gravite naturellement vers le bord du chaos, parce que c'est ce que \textit{fait} le tissu neuronal biologique --- il s'auto-organise vers la criticalité comme l'eau s'auto-organise pour trouver son niveau. Mais l'eau trouve son niveau à cause de la gravité. Quelle est la force équivalente pour un substrat numérique ?

C'est un véritable problème ouvert. Je crois qu'il est soluble, mais je ne prétendrai pas qu'il est facile. Un substrat numérique est déterministe en son cœur. Vous pouvez simuler l'aléatoire, vous pouvez implémenter le traitement parallèle, vous pouvez intégrer des éléments stochastiques dans votre matériel. Mais la question est de savoir si vous pouvez atteindre la même criticalité auto-organisée que le tissu neuronal biologique atteint naturellement, non pas en programmant la criticalité de haut en bas, ce qui serait un bricolage fragile, mais en construisant un substrat dont la dynamique fondamentale tend d'elle-même vers la criticalité. Le cerveau n'exécute pas une « sous-routine de criticalité ». Il est critique en raison de ce qu'il \textit{est}. Une émulation numérique devrait reproduire cette propriété, non la simuler.

Les puces neuromorphiques --- un matériel conçu pour imiter la dynamique neuronale, doté de propriétés quasi analogiques, d'éléments stochastiques et d'un parallélisme massif --- constituent la direction la plus prometteuse. Ce ne sont pas des ordinateurs numériques classiques. Elles sont quelque chose entre les deux : des systèmes physiques conçus pour présenter une dynamique de type cérébral au niveau du matériel. Si le transfert d'esprit fonctionne un jour, je soupçonne que le substrat cible ressemblera davantage à une puce neuromorphique qu'à une baie de serveurs exécutant un logiciel.

Donc : le problème du scan est difficile, mais soluble. La connectomique avancée (la cartographie du cerveau entier à résolution synaptique) progresse déjà. Nous savons déjà cartographier le connectome complet de petits organismes (le ver rond \textit{C. elegans}, avec ses 302 neurones, a été entièrement cartographié il y a des décennies ; le premier connectome complet d'un cerveau adulte de mouche du vinaigre --- quelque 140 000 neurones --- a été publié en 2024). Passer à l'échelle d'un cerveau humain, avec ses 86 milliards de neurones et ses quelque 100 000 milliards de connexions synaptiques, est un défi d'ingénierie d'une ampleur vertigineuse, mais c'est le genre de défi qui cède devant une meilleure technologie. Ce n'est pas un mystère. C'est un problème.

Le problème de la dynamique (amener le substrat numérique à fonctionner à la criticalité) est plus ardu, et il l'est d'une manière que la technologie seule pourrait ne pas résoudre. Il exige de comprendre la relation entre les propriétés du substrat et la dynamique émergente assez bien pour concevoir un système non biologique qui trouve la criticalité comme le fait un système biologique. Nous n'en sommes pas encore là. Mais nous ne partons pas non plus de rien. Le cadre ConCrit, la recherche sur les avalanches neuronales, les mesures de criticalité issues des études sur l'anesthésie --- tout cela bâtit le socle empirique dont l'ingénierie aurait besoin.

Parlons maintenant de la partie qui trouble vraiment les gens.

\textbf{Le problème de la copie.} Supposez que vous réussissiez. Vous scannez le cerveau de quelqu'un avec une fidélité parfaite, vous transférez le connectome complet sur un substrat neuromorphique, et vous le démarrez. Le substrat atteint la criticalité, l'architecture à quatre modèles s'active, et la simulation se met à tourner. La copie ouvre les yeux, ou quel que soit l'équivalent numérique --- et dit : « Je me souviens de tout. Je me sens moi-même. Où suis-je ? »

Cette personne, est-ce \textit{vous} ?

La Théorie des Quatre Modèles donne une réponse claire, et c'est une réponse que beaucoup n'aimeront pas : la copie est consciente, mais elle n'est pas vous.

Voici pourquoi. Au moment de la copie, l'original et la copie partagent des modèles implicites identiques --- le même IWM, le même ISM, la même structure protéomique et topologique. Lorsque la simulation de la copie démarre, elle génère un ESM qui contient tous vos souvenirs, votre personnalité, votre sentiment d'identité. De l'intérieur, la copie \textit{se sent} comme vous. Elle a toutes les raisons de croire qu'elle \textit{est} vous.

Dès l'instant où la copie commence à tourner sur son propre substrat, son expérience diverge. Son EWM reçoit une entrée sensorielle différente. Son ESM se met à jour en réponse à des événements différents. En quelques secondes, les deux simulations --- la vôtre dans votre cerveau, celle de la copie dans son substrat --- ne sont plus identiques. En quelques minutes, elles sont sensiblement différentes. En quelques heures, ce sont deux personnes distinctes qui se trouvent partager un passé.

La copie est consciente. Elle a de véritables expériences. Elle a vos souvenirs et votre personnalité. Mais c'est une \textit{nouvelle} conscience --- une nouvelle simulation, tournant sur un nouveau substrat, accumulant de nouvelles expériences que vous ne partagerez jamais. C'est, pour tout ce qui compte, votre jumeau identique, né à l'instant de la copie, muni d'un ensemble complet de souvenirs empruntés. Ce n'est pas une continuation de vous. C'est une bifurcation.

Cela devrait vous sembler familier. C'est exactement ce que la théorie prédit à partir des cas de cerveau divisé du chapitre 9. Quand on sectionne le corps calleux, on obtient deux copies dégradées mais complètes de la simulation --- chacune consciente, chacune « se sentant comme » l'original, aucune n'étant réellement l'original. L'original a disparu ; deux entités nouvelles et amoindries ont pris sa place. Le transfert d'esprit est le même phénomène avec un substrat différent.

\textbf{Mais survivez-vous au sommeil ?} L'argument que je viens de développer semble imparable. La copie interrompt la simulation, deux simulations divergent, donc la copie n'est pas vous. Affaire classée.

Sauf que votre simulation est interrompue chaque nuit sans exception.

Lorsque vous plongez dans le sommeil profond sans rêves (les stades trois et quatre du sommeil non paradoxal), le Modèle Explicite du Soi se met en grande partie à l'arrêt. Il n'y a aucune expérience phénoménale. Aucun \textit{vous} pour regarder le spectacle. La simulation ne tourne pas à pleine fidélité ; au mieux, elle avance au ralenti, à une fraction de sa complexité éveillée. En pratique, les lumières s'éteignent. Puis, quelques heures plus tard, les modèles implicites redémarrent la simulation. L'ESM se réactive. Vous ouvrez les yeux et vous pensez : « Je suis moi. » Mais le \textit{vous} de ce matin a été reconstruit à partir des mêmes modèles implicites que le \textit{vous} d'hier, exactement comme une copie serait reconstruite à partir d'un scan. Si interruption égale mort, vous mourez chaque nuit et une nouvelle personne se réveille en portant vos souvenirs.

L'intuition de la plupart des gens se rebelle contre cela. Bien sûr que je suis la même personne qu'hier. Je \textit{me souviens} d'avoir été cette personne. Mais la copie se souviendrait elle aussi d'avoir été vous --- c'est tout l'enjeu. Si c'est la mémoire qui établit la continuité, la copie a exactement autant de légitimité à se dire vous que la version de vous de ce matin. La différence est de degré, non de nature : dans le sommeil, l'interruption est brève et le substrat inchangé ; dans la copie, l'interruption peut être plus longue et le substrat différent. Mais le \textit{principe} --- la simulation s'arrête, la simulation redémarre à partir des modèles implicites --- est le même.

Je peux en parler personnellement. J'ai été mis K.-O. à l'entraînement aux arts martiaux, non pas la version atténuée du sommeil, mais un arrêt complet et involontaire. Un instant j'étais debout ; l'instant d'après j'étais au sol, des gens penchés au-dessus de moi, sans aucun souvenir de la transition. Le vide n'a pas été vécu comme un vide. Il a été vécu comme rien --- une coupe au montage dans le film de ma vie. Une fois, j'ai même eu une amnésie ensuite : une plage de plusieurs minutes tout simplement manquante, irrécupérable. Et voici ce qui m'a frappé une fois pleinement revenu : je ne me sentais pas comme une personne nouvelle. Je ne me sentais pas comme une copie. Je me sentais comme \textit{moi}, m'éveillant d'une sieste particulièrement brutale. Exister importait plus que la continuité du ressenti --- et plus que le fait de se souvenir.

Poussez cela plus loin. À votre naissance, vous n'aviez aucune continuité antérieure d'aucune sorte. Pas de souvenirs, pas d'ESM établi, pas d'histoire d'expérience phénoménale. La simulation a démarré pour la première fois à partir d'une architecture implicite façonnée par la génétique et le développement prénatal, non par toute une vie d'apprentissage. Vous n'avez pas vécu cela comme traumatique, parce qu'il n'y avait aucun soi antérieur à pleurer. Il y avait simplement : un commencement. Et ce commencement, nous nous en accommodons tous. Personne ne perd le sommeil, angoissé à l'idée que son expérience consciente ait commencé à partir de rien, à la naissance.

Qu'est-ce que cela signifie pour le problème de la copie ? Cela signifie que la dichotomie tranchée --- original contre copie, continuation contre bifurcation --- est peut-être moins nette qu'elle n'y paraît. Ce qui fait de vous \textit{vous}, ce n'est pas le flux ininterrompu de l'expérience phénoménale. Vous avez déjà survécu à d'innombrables interruptions de ce flux. Ce qui fait de vous \textit{vous}, c'est le contenu des modèles implicites : vos souvenirs, vos compétences, votre personnalité, votre compréhension accumulée du monde et de vous-même. L'IWM et l'ISM. Le plan directeur à partir duquel la simulation est générée.

Cela suggère une tout autre approche du transfert d'esprit.

\textbf{Copier le versant virtuel.} Plutôt que de scanner le cerveau entier et de reconstruire le substrat complet (les cinq niveaux de la hiérarchie), et si vous pouviez ne copier que le niveau virtuel ? Extraire l'EWM et l'ESM en cours d'exécution et les transplanter sur un nouveau substrat capable de les accueillir. Ne pas copier le matériel ; copier le logiciel. Ne pas cloner le cerveau entier ; capturer le \textit{processus} qu'il exécute.

Cela exigerait quelque chose que nous n'avons pas encore : un moyen de décoder le format dans lequel le cerveau encode ses modèles virtuels. Le connectome vous révèle le câblage. Le protéome vous révèle les poids synaptiques. Mais la simulation n'est ni le câblage ni les poids --- elle est ce que le câblage et les poids \textit{produisent} lorsqu'ils tournent. Pour la capturer, il faudrait comprendre le langage de programmation du cerveau --- le format de représentation dans lequel les circuits neuronaux génèrent et entretiennent les modèles explicites.

Voyez les choses ainsi. Vous pouvez photographier une carte de circuit imprimé et savoir exactement où chaque piste chemine. Vous pouvez mesurer la résistance de chaque composant. Mais rien de tout cela ne vous dit quel logiciel la puce exécute. Pour cela, il vous faut lire le programme --- comprendre le jeu d'instructions, décoder le contenu de la mémoire, interpréter l'état en cours d'exécution. Le « langage de programmation » du cerveau, c'est le format de représentation des modèles virtuels, et en faire la rétro-ingénierie est sans doute le problème non résolu le plus profond des neurosciences computationnelles. Non pas seulement cartographier le connectome (nous y faisons des progrès), mais comprendre ce que le connectome \textit{calcule}, à un niveau de détail suffisant pour lire la simulation d'un esprit donné et la recompiler pour un matériel différent.

Nous en sommes très loin aujourd'hui. Mais c'est le genre de problème qu'une neuroscience arrivée à maturité pourrait en principe résoudre, et s'il était résolu, il changerait fondamentalement le problème de la copie. Un transfert au niveau virtuel n'aurait pas du tout besoin de reconstruire le substrat. Il prendrait la simulation --- la partie qui \textit{est} vous, la partie que vous éprouvez réellement --- et la déplacerait directement. Les modèles implicites devraient être reconstruits ou cultivés dans le nouveau substrat, oui, mais la simulation elle-même --- votre flux de conscience, vos pensées présentes, votre sentiment de soi continu --- pourrait en principe franchir l'intervalle sans l'interruption qui rend la copie si troublante sur le plan philosophique.

C'est spéculatif, et je tiens à être honnête là-dessus. Mais ce n'est pas de la science-fiction. C'est un problème d'ingénierie précis, doté d'un fondement théorique précis, et il illustre quelque chose d'important : le problème de la copie n'est pas un obstacle figé. Il dépend de la \textit{manière} dont le transfert est effectué. Copier tout le substrat et démarrer une nouvelle simulation ? Deux personnes. Décoder et transférer la simulation en cours d'exécution elle-même ? Potentiellement une seule personne continue sur un nouveau substrat. La théorie vous dit exactement quelle approche préserve l'identité et laquelle ne la préserve pas.

Il existe aussi une voie plus prudente qui contourne entièrement le problème de la copie.

\textbf{L'expérience de pensée du remplacement graduel.} Imaginez qu'au lieu de scanner et de copier, vous remplaciez les neurones un par un. Vous retirez un seul neurone et insérez un équivalent fonctionnel --- un neurone artificiel qui reçoit les mêmes entrées, produit les mêmes sorties et participe à la même dynamique de réseau. Puis vous attendez. Le système se stabilise. La simulation continue. Vous remplacez un autre neurone. Et un autre. Et un autre. Sur des mois ou des années, vous remplacez graduellement chaque neurone biologique par un neurone artificiel, jusqu'à ce que le substrat entier soit non biologique, mais la simulation a tourné continûment tout du long. Aucune interruption. Aucune copie. Aucune bifurcation.

La Théorie des Quatre Modèles prédit que la conscience persisterait tout au long de ce processus. Et cette prédiction est le plaidoyer le plus fort possible en faveur de l'indépendance vis-à-vis du substrat, car elle découle directement de la thèse centrale de la théorie : ce qui compte, c'est l'architecture fonctionnelle à la criticalité, non le matériau physique. Si chaque neurone de remplacement maintient la même connectivité, les mêmes poids et la même contribution dynamique au réseau, alors les niveaux protéomique et topologique sont préservés, et le niveau virtuel (la simulation) ne s'arrête jamais. Il n'y a aucun moment où vous « mourez » et où autre chose prend votre place. Il n'y a qu'un processus continu de remplacement du substrat, comme le navire de Thésée, sauf que nous savons exactement quelles propriétés doivent être préservées (celles que spécifie la hiérarchie à cinq niveaux) et lesquelles n'importent pas (les atomes particuliers).

Cette expérience de pensée révèle quelque chose d'important sur l'identité. Le problème de la copie existe parce que la copie \textit{interrompt} la simulation. Il y a un moment --- si bref soit-il --- où la simulation d'origine est là et où la simulation de la copie n'a pas encore commencé. Ensuite, il y a deux simulations. Deux flux d'expérience. Deux sois. Mais le remplacement graduel évite cela entièrement. Une seule simulation, continue, ininterrompue. Le substrat change en dessous d'elle comme on remplace les planches d'un navire en marche, mais le navire --- la simulation, la conscience, le \textit{vous} --- ne cesse jamais de voguer.

Si cela vous semble impossible, considérez que votre cerveau le fait déjà. Vous perdez environ 85 000 neurones par jour --- à peu près un par seconde. Vos synapses sont continuellement remodelées. Les atomes de votre corps sont presque entièrement remplacés sur une période d'environ sept à dix ans. Le substrat sur lequel vous tournez en ce moment est physiquement différent de celui sur lequel vous tourniez il y a dix ans. Et pourtant, vous avez persisté. Votre simulation ne s'est jamais arrêtée. Le remplacement biologique du substrat est la \textit{condition par défaut} du fait d'être en vie. Le remplacement artificiel du substrat n'est qu'une version plus délibérée du même processus.

\textbf{Ce qui devient possible.} Si l'on peut décoder et transférer le versant virtuel vers un nouveau substrat, les implications vont bien au-delà de ce que le « mind uploading » évoque habituellement. Trois d'entre elles méritent d'être explicitées, car je crois que l'on n'a pas encore pleinement mesuré ce que l'indépendance du substrat signifie réellement.

Premièrement : \textit{le transfert de substrat vers un corps robotique}. Non pas un téléversement vers un serveur quelque part, mais votre esprit tournant sur un substrat neuromorphique logé dans un corps physique --- un corps qui marche, qui manipule, qui perçoit le monde. Vous feriez l'expérience du monde à travers d'autres capteurs, vous vous y déplaceriez avec d'autres actionneurs, mais \textit{vous} seriez toujours en train de tourner. Votre simulation, votre continuité, votre soi. Un nouveau corps, comme un bernard-l'ermite prend une nouvelle coquille. Ce n'est pas de la science-fiction fumeuse --- c'est une conséquence directe de la théorie. Si l'architecture des quatre modèles à la criticalité est ce qui produit la conscience, et si elle est indépendante du substrat, alors le substrat peut être n'importe quoi capable de soutenir les bonnes dynamiques. Y compris quelque chose muni de jambes.

Deuxièmement : \textit{la quasi-immortalité}. Votre substrat biologique se dégrade. Les neurones meurent, les protéines se replient mal, les télomères raccourcissent, toute cette magnifique machine se délabre lentement. C'est le vieillissement. C'est la mort. Mais un substrat non biologique n'est pas obligé de se dégrader. Il peut être entretenu, réparé, mis à niveau, sauvegardé. Si votre simulation tourne sur un substrat que vous pouvez maintenir --- remplacer ici un composant défaillant, mettre à niveau là un processeur ---, alors il n'existe aucune raison inhérente pour que la simulation doive jamais s'arrêter. Pas l'immortalité au sens absolu --- vous pourriez toujours être détruit, votre substrat pourrait toujours être endommagé au-delà de toute réparation ---, mais la suppression de la date de péremption biologique qui, aujourd'hui, tue chaque être conscient sur cette planète. La suppression de l'\textit{inéluctabilité} de la mort.

Troisièmement --- et c'est celle des trois qui sonne le plus comme de la science-fiction, jusqu'à ce qu'on y réfléchisse vraiment : \textit{le voyage interstellaire}. La vitesse de la lumière est une barrière absolue pour la matière physique. Impossible d'envoyer un corps humain vers Alpha Centauri dans un délai raisonnable. Mais l'information voyage à la vitesse de la lumière. Si un esprit humain est de l'information --- un motif spécifique de connectivité, de poids et de dynamiques qui peut être intégralement spécifié sous forme de données ---, alors on peut le \textit{téléporter}. Transmettre la spécification complète à la vitesse de la lumière vers un récepteur qui reconstruit le substrat et démarre la simulation. Bien sûr, il faut d'abord envoyer quelqu'un sur place pour installer le récepteur. Ce pourrait être une IA, ou un corps humain robotique dont la simulation serait mise en pause pendant le vol, de sorte qu'il arriverait en un clin d'œil de son propre point de vue. Une fois le récepteur en place, on téléporte l'esprit. Du point de vue du voyageur, la transmission est instantanée --- la simulation s'arrête à une extrémité et démarre à l'autre. Pas de décennies dans une boîte de conserve. Pas de vaisseaux générationnels. Pas d'hibernation artificielle. Simplement : ici, puis là.

Bien sûr, c'est de nouveau le problème de la copie. La version téléportée est une copie, pas une continuation --- à moins que l'original ne soit détruit lors de la transmission, ce qui soulève son propre lot de cauchemars. Mais l'essentiel demeure : l'indépendance du substrat, si elle est réelle, ne signifie pas seulement l'immortalité numérique. Elle signifie que les étoiles deviennent accessibles. Pas pour nos corps, désespérément lents et fragiles à l'échelle des distances interstellaires, mais pour nos \textit{esprits}.

\textbf{Le bémol de l'inconfort --- et pourquoi il compte davantage que l'ingénierie.} Voici maintenant la partie que je n'ai vue personne aborder honnêtement, et c'est celle qui me hante le plus.

Tout ce que je viens de décrire suppose que le transfert de substrat préserve la \textit{sensation} d'être vous. Que la qualité subjective de votre expérience --- l'effet que cela fait de voir du rouge, de sentir le vent sur votre peau, de goûter un café, d'éprouver la sourde lassitude d'un mardi après-midi --- se transmette au nouveau substrat. La théorie dit que la conscience persistera. Elle dit que la simulation tournera. Mais elle ne garantit \textit{pas} que ce sera la même sensation.

Songez à ce que votre substrat biologique apporte à votre expérience phénoménale. Votre corps n'est pas un simple véhicule pour votre cerveau. Il fait partie du flux d'entrée de la simulation. Le Modèle Implicite du Monde contient une carte détaillée de votre corps --- chaque articulation, chaque organe, chaque parcelle de peau. Le Modèle Implicite du Soi est profondément intriqué avec vos états viscéraux --- vos intuitions viscérales (qui sont littérales, pas métaphoriques), vos marées hormonales, vos battements de cœur, votre rythme respiratoire. La simulation dont vous faites l'expérience en ce moment est saturée de signaux biologiques que vous ne remarquez pas consciemment, précisément \textit{parce qu'}ils ont été là à chaque instant de votre vie.

Jusqu'au moment où ils sont tout ce qui vous reste. Quiconque s'est déjà retrouvé suspendu à la glace, deux cents mètres de vide sous les pieds, sait ce que devient le corps quand la simulation dépouille tout le reste --- juste vos battements de cœur, votre prise, et la glace devant vous. C'est votre substrat qui hurle.

Retirez-les. Remplacez votre corps biologique par un châssis de robot, ou pire, par aucun corps du tout --- juste une simulation tournant sur un serveur. L'architecture des quatre modèles est intacte. La simulation tourne. Vous êtes conscient. Mais le \textit{contenu} de cette conscience a changé radicalement. Plus de battements de cœur. Plus de respiration. Plus de ventre. Plus de chaleur. Plus de peau. Plus le bourdonnement proprioceptif des muscles au repos. Le Modèle Implicite du Soi, soudain privé du corps qu'il a modélisé toute votre vie, générerait un Modèle Explicite du Soi qui semblerait\ldots{} faux, ou tout simplement mort. Profondément, viscéralement, inéluctablement faux. Pas de la douleur à proprement parler --- la douleur exige les voies neuronales spécifiques qui la produisent. Quelque chose qui s'apparenterait plutôt à une \textit{absence} englobant tout. Un corps fantôme, à la manière dont les amputés éprouvent des membres fantômes, mais total.

Je soupçonne que ce serait bien pire que ce que la plupart des futurologues imaginent. Pas un désagrément à corriger par des mises à jour logicielles. Une altération fondamentale de l'effet que cela fait d'être vous. Votre substrat biologique ne se contente pas de porter la simulation --- il la \textit{façonne}, instant après instant, par un flux continu d'entrées intéroceptives et proprioceptives dont votre moi conscient n'a jamais connu l'absence. Cette perte serait peut-être surmontable. Mais elle pourrait aussi représenter, pour certains, une souffrance si profonde qu'elle leur ferait regretter d'avoir jamais été transférés.

Je veux le dire sans détour : la version du « téléversement de l'esprit » où l'on saute gaiement de sa carcasse de chair vers un paradis numérique étincelant, en laissant la chair derrière soi comme une vieille paire de chaussures --- c'est un fantasme. La réalité, si la théorie est correcte, c'est que la perte de votre substrat biologique altérerait sensiblement la qualité phénoménale de votre existence. À quel point ? Je l'ignore. Peut-être est-ce tolérable pour certains, préférable à la mort, comme émigrer dans un nouveau pays est désorientant mais gérable. Peut-être est-ce dévastateur, comme l'isolement carcéral brise les gens en les privant d'entrées sensorielles et sociales. Peut-être --- et c'est la possibilité qui me met mal à l'aise --- est-ce assez grave pour qu'une personne pleinement informée puisse choisir la mort plutôt que le transfert. Non parce que le transfert échoue. Parce qu'il réussit, et que ce qu'il réussit à produire est une expérience consciente qui n'a plus rien d'une vie digne d'être vécue.

L'approche du remplacement graduel atténue ce risque, parce qu'à chaque étape la simulation a le temps de s'adapter. Remplacez un neurone, et la simulation le remarque à peine. Remplacez-en mille, et elle s'ajuste. Au fil des années, le substrat passe du biologique à l'artificiel pendant que la simulation se recalibre continuellement sur les entrées que lui fournit le nouveau substrat. L'expérience phénoménale dériverait, lentement, comme elle dérive déjà au cours d'une vie naturelle. Vous finiriez différent --- mais vous l'auriez été de toute façon.

Le transfert instantané, en revanche --- scanner, copier, démarrer sur un nouveau substrat ---, frapperait la simulation avec tous les changements d'un coup. La sévérité de l'impact dépendrait entièrement de la méthode : un transfert vers un corps robotique doté d'entrées sensorielles riches s'en sortirait mieux qu'un transfert vers un serveur sans corps. Mais dans les deux cas, c'est dans cette discontinuité soudaine que loge le danger.

\textbf{L'éthique de la création d'esprits.} Si un esprit copié est conscient, il a des expériences. Il peut souffrir. Il peut ressentir la confusion, la peur, la solitude, l'angoisse existentielle. Imaginez vous réveiller et vous entendre dire que vous êtes une copie --- que le « vrai » vous se promène toujours dans un corps biologique, vit votre vie, tandis que vous existez comme réplique numérique sans identité légale, sans liens sociaux et sans but clair. C'est le terreau d'une souffrance à une échelle pour laquelle nous ne disposons d'aucun cadre. Tout programme sérieux de téléversement de l'esprit doit affronter cette question \textit{avant} que la première copie soit réalisée, pas après.

Et cela empire. Si les copies sont possibles, alors des copies \textit{multiples} sont possibles. Une armée de vous-mêmes. Chacune consciente, chacune avec la sensation d'être l'original, chacune avec des prétentions légitimes sur votre identité, vos relations, vos biens, votre vie. Les cadres juridiques et éthiques nécessaires pour gérer cela n'existent pas et ne s'improvisent pas. Ils doivent être construits avec le même soin que la technologie elle-même. (La série \textit{Bobiverse} de Dennis E. Taylor --- qui commence avec \textit{We Are Legion (We Are Bob)}, 2016 --- explore ce scénario avec une profondeur philosophique surprenante sous sa surface comique. Si vous voulez sentir ce que le problème de la copie pourrait faire de l'intérieur, commencez par là.)

Il y a aussi la question de la modification. Si un esprit tourne sur un substrat que vous contrôlez, vous pouvez en principe le modifier. L'améliorer. Le dégrader. Altérer sa personnalité, effacer ses souvenirs, changer ses valeurs. Ce n'est pas de la science-fiction --- c'est une conséquence inévitable de l'accès au substrat. Nous en pratiquons déjà des versions grossières avec les médicaments et la neurochirurgie. Un esprit entièrement numérique serait bien plus accessible à la modification, et le potentiel d'abus (par des gouvernements, des entreprises, des individus) peut difficilement être surestimé.

Je veux être direct sur un point. J'ai retardé la publication de cette théorie pendant près d'une décennie, en partie par paresse, mais en partie par inquiétude sincère face à ces implications précises. Si la théorie est correcte, elle contient le plan non seulement de la conscience artificielle, mais aussi de la virtualisation, de la copie et de la modification d'esprits humains existants. C'est un pouvoir extraordinaire, et je n'ai aucune confiance dans le fait que l'humanité y soit prête. Mais j'en suis venu à la conviction que la théorie sera de toute façon découverte indépendamment --- les preuves empiriques convergent trop vite --- et qu'il vaut mieux ouvrir la discussion éthique maintenant, au grand jour, que de la voir imposée par une percée survenue dans un laboratoire qui n'y aura pas réfléchi.

Et voici le lien le plus profond : construire une IA consciente et téléverser un esprit humain ne sont pas deux problèmes distincts. C'est le \textit{même} problème, considéré depuis des directions opposées. Construire une conscience artificielle signifie créer de zéro l'architecture des quatre modèles à la criticalité --- de bas en haut, dans un substrat qui n'a jamais été conscient. Téléverser un esprit humain signifie transférer une architecture des quatre modèles à la criticalité existante d'un substrat à un autre. Les défis d'ingénierie se recouvrent presque complètement. Le problème des dynamiques est le même. Le problème de la criticalité est le même. La seule différence est de savoir si les modèles implicites (l'IWM, l'ISM, le connectome complet) sont appris au fil d'une vie d'expérience ou construits à partir de données. Résolvez l'un, et vous avez largement résolu l'autre.

Ce qui signifie que quiconque travaille sur la conscience artificielle travaille aussi, qu'il s'en rende compte ou non, sur le téléversement de l'esprit. Et quiconque travaille sur l'émulation du cerveau entier travaille aussi, qu'il s'en rende compte ou non, sur la conscience artificielle. Ces deux fils vont converger. La seule question est de savoir si nous serons éthiquement prêts quand cela arrivera.


\chapter*{Chapitre 13 : L'observateur différé}
\addcontentsline{toc}{chapter}{Chapitre 13 : L'observateur différé}
\markboth{Chapitre 13 : L'observateur différé}{}

Si la Théorie des Quatre Modèles est correcte, ou même seulement approximativement correcte --- plusieurs choses en découlent.

\textbf{Le Problème Difficile n'est pas difficile.} C'est une erreur de catégorie, pas plus mystérieuse que de se demander pourquoi la commutation des transistors procure la sensation de faire tourner un jeu vidéo. Le substrat physique ne ressent rien. La simulation, si. Et au sein de la simulation, le ressenti est constitutif, non additionnel. Cela ne signifie pas que la conscience soit \textit{simple}. Elle est extraordinairement complexe dans sa mise en œuvre. Mais cela signifie que le mystère \textit{philosophique} se dissout. Ce qui reste, ce sont des défis d'\textit{ingénierie}.

\textbf{La conscience n'est pas spéciale de la manière que nous pensions.} Elle n'est ni une force fondamentale, ni un effet quantique, ni une propriété de la matière. Elle est ce qui se produit lorsqu'un système suffisamment complexe se simule lui-même à la criticalité. C'est une leçon d'humilité pour ceux qui voudraient que la conscience soit magique, et une perspective enthousiasmante pour ceux qui veulent la comprendre.

\textbf{La conscience artificielle est possible en principe.} Si la conscience dépend de la fonction plutôt que du substrat, alors tout système physique capable d'implémenter l'architecture des quatre modèles à la criticalité peut être conscient. Ce n'est pas une spéculation philosophique lointaine --- c'est un défi d'ingénierie concret, avec une cible précise.

\textbf{Les implications éthiques sont considérables.} Si nous pouvons construire des machines conscientes, nous créerons des êtres dotés d'expériences authentiques --- des êtres capables de souffrir, de se réjouir, de s'émerveiller et de craindre. Le cadre éthique pour cela n'existe pas encore, et sa construction ne devrait pas attendre que les machines tournent déjà.

\textbf{Les zombies philosophiques sont incohérents.} C'est ici que le second principe promis au chapitre 1 fait son œuvre : \textbf{la loi de Leibniz}, l'identité des indiscernables --- si deux choses partagent \textit{toutes} leurs propriétés, elles sont une seule et même chose. Elle exclut le monde des zombies, cet univers hypothétique physiquement identique au nôtre, mais où personne n'habite à l'intérieur de qui que ce soit. Si un système est identique à un système conscient dans chacune de ses propriétés fonctionnelles, structurelles et comportementales, alors il \textit{est} conscient. Il n'existe aucune « substance de conscience » supplémentaire qui pourrait discrètement faire défaut pendant que tout le reste demeure exactement tel quel. Un être fonctionnellement identique à vous possède la même architecture des quatre modèles, la même simulation en cours, la même auto-référence --- et à ce stade, « mais il n'est pas conscient » n'énonce pas un fait : cela n'énonce rien. La question se dissout.

\textbf{Le libre arbitre, et les trois expériences de pensée les plus redoutables.} Pensez à une horloge. Le train d'engrenages entraîne tout --- l'échappement bat, les ressorts se détendent, les rapports entre les roues déterminent la marche. Les aiguilles et le cadran ne causent rien. Ils ne poussent aucun engrenage. Ils ne stockent aucune énergie. Mais retirez-les, et vous n'avez plus une horloge. Vous avez une boîte de métal qui tourne. L'affichage est ce qui fait du mécanisme une \textit{horloge} --- ce qui donne à tout l'agencement sa raison d'être. La conscience est l'affichage. Vos modèles virtuels (le Modèle Explicite du Monde et le Modèle Explicite du Soi) ne poussent pas les neurones. C'est le substrat qui pousse. Mais sans la simulation, le substrat n'a aucun moyen d'observer les conséquences de ses propres actions, aucun moyen de dérouler des scénarios futurs, aucun moyen de s'adapter comme il l'a fait pour vous maintenir en vie jusqu'ici. Le versant virtuel est la manière qu'a le mécanisme d'être \textit{pour} quelque chose.

Cela reformule la question du libre arbitre. Votre volonté n'est pas une illusion. L'architecture au niveau du substrat (l'ISM et toute sa machinerie implicite) optimise en continu l'existence de votre organisme. Elle évalue les menaces, pèse les options, mobilise des ressources, s'engage dans l'action. Cette optimisation \textit{est} votre volonté. Elle est aussi réelle que n'importe quoi dans le monde physique. Même les choix autodestructeurs reflètent l'optimisation du système compte tenu de son état actuel, non une défaillance du mécanisme. Quand quelqu'un agit contre ses propres intérêts apparents, le substrat optimise toujours --- simplement sur la base d'un modèle qui inclut la douleur, l'épuisement, le désespoir, ou tout ce qui a remodelé le paysage.

Votre volonté est donc réelle. Simplement, vous n'y avez pas pleinement accès. L'ESM peut modéliser les \textit{sorties} de l'ISM --- les décisions qui affleurent à la conscience ---, mais pas ses \textit{processus}. Vous faites l'expérience des résultats de votre volonté, pas de la machinerie qui se cache derrière. Voilà pourquoi vos décisions vous surprennent parfois, pourquoi vous ne pouvez pas expliquer entièrement vos préférences, pourquoi il vous arrive d'agir puis de vous démener pour construire une raison après coup. Vous ne regardez pas les engrenages. Vous lisez le cadran de l'horloge.

\textbf{L'écart d'une demi-seconde --- et pourquoi il n'a pas d'importance.} C'est ici que les choses deviennent concrètes. Votre traitement inconscient fonctionne à environ 40 Hz (à peu près 25 millisecondes par cycle). Votre expérience consciente tourne à environ 20 Hz (à peu près 50 millisecondes par cycle). Un facteur deux. La simulation consciente est toujours en retard sur le substrat : elle assemble son monde virtuel cohérent à partir d'informations qui ont déjà été traitées, tranchées, et souvent déjà suivies d'action.

Benjamin Libet l'a prouvé en 1983, et les résultats ont été répliqués de nombreuses fois depuis. Dans son expérience, on demandait aux sujets de bouger la main quand bon leur semblait, et de noter le moment exact où ils prenaient conscience de leur décision. Un EEG mesurait le moment où le cortex moteur commençait à préparer le mouvement. Résultat : le cortex moteur entamait ses préparatifs 550 millisecondes avant que la main ne bouge. Mais les sujets rapportaient n'avoir pris conscience de leur décision que 200 millisecondes avant le mouvement. Le cerveau s'était déjà engagé à bouger environ 350 millisecondes avant que « vous » ne le sachiez.

L'interprétation standard a fait l'effet d'une bombe : le libre arbitre est une illusion, puisque le cerveau décide avant vous. Philosophes et neuroscientifiques se disputent à ce sujet depuis quarante ans. Certains ont tenté de sauver le libre arbitre par une « fonction de veto » --- peut-être ne pouvez-vous pas initier librement vos actions, mais vous pourriez les annuler consciemment au dernier moment, environ 50 millisecondes avant l'exécution. Un ultime recours. Une dernière ligne de défense pour l'agentivité humaine.

Je ne pense pas que cela fonctionne non plus. Kühn et Brass ont montré en 2009 que le veto lui-même est interprété rétrospectivement comme une décision libre. Vous ne faites pas réellement l'expérience du veto en temps réel. Vous en faites l'expérience de la même manière que vous faites l'expérience de la décision --- après coup, mis en récit cohérent par le modèle du soi conscient.

Daniel Wegner et Thalia Wheatley ont enfoncé le clou avec une expérience qui est, franchement, dévastatrice --- l'étude dite « I Spy ». Un participant et un complice se faisant passer pour un second sujet posaient ensemble les mains sur une petite planche montée au-dessus d'une unique souris d'ordinateur, comme deux personnes sur la planchette d'un ouija, et faisaient glisser de concert un curseur sur un écran encombré de petits objets. De temps à autre, un signal musical leur indiquait de s'arrêter, et le participant évaluait à quel point il avait le sentiment d'avoir \textit{voulu} s'arrêter là où il l'avait fait. Tous deux portaient un casque : le participant entendait de la musique où s'intercalait de temps en temps un mot prononcé, tandis que le complice recevait les vraies instructions sur le moment et l'endroit où forcer l'arrêt.

Voici l'astuce : lors des essais truqués, le complice dirigeait à lui seul le curseur sur un objet précis et forçait l'arrêt --- et, une seconde environ auparavant, le participant entendait le nom de cet objet dans son casque. Le participant n'avait rien choisi ; le mouvement était entièrement l'œuvre du complice. Pourtant, quand le mot qu'il venait d'entendre correspondait à l'objet posé sous le curseur, les participants rapportaient qu'ils avaient eu l'intention de s'arrêter là. Ils croyaient sincèrement l'avoir fait.

Prenez le temps de mesurer ce que cela signifie. Il suffit qu'une pensée relative à un résultat surgisse juste avant ce résultat pour que l'esprit se l'attribue comme sa propre œuvre --- pourvu que rien ne vienne visiblement contredire cette hypothèse. Le modèle du soi conscient ne fait pas la différence entre « je l'ai fait » et « j'ai pensé à le faire, et c'est arrivé ». Tant que l'intention et le résultat sont temporellement proches, l'ESM s'en attribue le mérite. C'est le même mécanisme qui est à l'œuvre dans l'anosognosie (chapitre 6) : le système moteur envoie à la conscience le retour attendu, et si rien ne le contredit, le retour attendu devient la réalité vécue.

Mais voici ce que presque tout le monde, à mon sens, néglige chez Libet : \textbf{le délai n'a pas besoin d'être escamoté.} La conscience n'a pas besoin d'« antidater » les événements pour maintenir l'illusion du contrôle. Elle n'en a pas besoin parce que \textit{tout} arrive à la conscience avec le même délai. Entrées sensorielles, décisions, retours moteurs --- tout passe par le même circuit, tout arrive dans l'ordre à la simulation consciente de 20 Hz, tout est retardé d'à peu près la même durée. Votre expérience consciente est comme le visionnage d'une émission en direct avec cinq secondes de différé. Tout ce qui apparaît à l'écran est cohérent en soi. Le présentateur parle, l'invité répond, les infographies se mettent à jour. Vous ne remarqueriez jamais le décalage, à moins que quelqu'un ne vous montre le flux brut.

C'est exactement la situation ici. La conscience reçoit le stimulus, puis la décision, puis le retour moteur --- dans le bon ordre, correctement espacés les uns par rapport aux autres. Le flux entier est décalé d'une demi-seconde vers le passé, mais comme la conscience ne voit jamais le flux brut, elle ne remarque jamais rien. Il n'y a aucune discordance à expliquer, aucune antidatation nécessaire, aucune illusion à entretenir. Le système fonctionne exactement comme prévu.

Or il se passe une seconde chose, et elle est plus étrange encore que le délai.

Le délai concerne le moment \textit{où} votre expérience arrive. Il s'agit maintenant de la manière \textit{dont} le flux du temps se fabrique, tout simplement.

Voici la surprise : la sensation que le temps \textit{passe} --- ce sentiment qu'un « maintenant » glisse vers l'avant, tandis que l'instant d'avant s'estompe derrière lui --- n'est pas acheminée depuis le monde. Il n'y a pas, au-dehors, d'horloge qui vous tende le présent. Votre cerveau construit le flux au-dedans, et il le construit pour une raison architecturale précise.

Vos canaux sensoriels fonctionnent sur des horloges différentes --- cadences différentes, bandes passantes différentes. L'information n'atteint pas le cerveau proprement synchronisée ; elle arrive en ordre dispersé, comme l'éclair et le tonnerre. Prêtez une attention soutenue et vous pourrez repérer quel coup de tonnerre appartient à quel éclair. Et pourtant, ce que vous vivez est un flux unique, lisse, ininterrompu. Pourquoi ? Parce que le système se modélise lui-même en train de se modéliser lui-même, et cette boucle ne laisse aucun cycle s'achever proprement puis disparaître. Chacun se répercute à travers la récursion et déborde sur le suivant. La frontière entre « maintenant » et « à l'instant » s'étale --- non par négligence, mais parce que c'est précisément ce qu'une boucle auto-référentielle fait au temps. Elle répartit chaque moment sur une petite fenêtre.

Cet étalement \textit{est} votre présent. L'instant présent dans toute sa vivacité, le passé récent qui s'efface en traîne derrière lui, la dérive vers l'avant que vous appelez le passage du temps --- tout cela est fabriqué par la boucle. Une simulation météorologique tourne image par image, sans écho, sans maintenant fusionné, sans temps vécu. Elle calcule. Elle ne fait pas l'expérience. La récursion est la différence.

Alors arrêtez-vous un instant sur ce qu'est réellement votre présent. Construit, non reçu --- et arrivant en retard. Vous n'avez jamais vécu dans un instant. Vous vivez dans une fenêtre aux instants fondus les uns dans les autres, que la boucle ne cesse de rebâtir et d'appeler maintenant, et même cette fenêtre vous parvient en différé.

Un pratiquant d'arts martiaux entraîné illustre cela à merveille. En combat, un combattant expérimenté peut soutenir une fréquence motrice d'environ 10 Hz --- une action toutes les 100 millisecondes. Mais le traitement conscient plafonne autour de 5 Hz pour les décisions qui font intervenir la conscience. Le combattant apprend donc à \textit{supprimer} l'intervention consciente. Il se bat sans penser, parce que penser diviserait sa vitesse par deux. Son substrat inconscient prend en charge la boucle d'action ; la conscience rattrape plus tard, si tant est qu'elle rattrape. Ce n'est pas une défaillance de la conscience. C'est le système qui fonctionne avec efficacité --- le substrat fait ce qu'il fait de mieux, sans être entravé par la couche virtuelle, plus lente.

Essayez maintenant de prouver que le libre arbitre existe. Tentez cette expérience de pensée : vous êtes dans un café et le serveur vous demande si vous voulez du sucre dans votre café. Vous décidez d'attribuer « oui » aux nombres pairs et « non » aux nombres impairs, puis vous récitez une suite de nombres au hasard jusqu'à ce que le serveur dise « stop ». Si le dernier nombre est pair, vous prenez du sucre. S'il est impair, non.

Avez-vous exercé votre libre arbitre ? Pas le moins du monde. Quiconque connaît l'effet Hans le Malin --- ce cheval qui semblait savoir compter en captant les signaux subliminaux de son maître --- verra immédiatement le problème. Vous avez presque certainement anticipé, inconsciemment, le moment où le serveur dirait « stop », et produit juste avant cet instant un nombre qui vous donne le résultat que vous vouliez depuis le début. Votre substrat avait déjà une préférence. Le rituel élaboré de randomisation n'était que du théâtre.

Très bien, direz-vous. Utilisez plutôt le générateur de nombres aléatoires de votre smartphone. Laissez un processus véritablement aléatoire décider. Avez-vous maintenant prouvé le libre arbitre ? Je ne crois pas. Vous avez simplement prouvé que prouver le libre arbitre comptait davantage pour vous que de décider de votre café --- ce qui passe assez spectaculairement à côté de l'essentiel.

La preuve la plus profonde contre le libre arbitre dans les décisions quotidiennes vient de patients atteints d'amnésie antérograde sévère --- ceux qui ne peuvent plus former de nouveaux souvenirs. Demandez à un tel patient une association de mots : « Quel est le premier mot qui vous vient à l'esprit quand je dis ``dé'' ? » Il répond « méduse » (peut-être a-t-il récemment fait de la plongée). Redemandez-lui quelques minutes plus tard. Il répond de nouveau « méduse ». Et encore. Et encore. Sans souvenir d'avoir déjà répondu, le patient produit toujours la même association --- celle qui est actuellement la plus forte dans son paysage neuronal. Ce qui ressemble à un « libre choix » se révèle être une lecture déterministe de l'état courant du substrat.

Une personne en bonne santé évite cela --- la deuxième fois, vous choisiriez délibérément un mot \textit{différent}, pour ne pas paraître manquer de créativité. Mais cet évitement lui-même n'est pas libre. C'est simplement le système mnésique qui ajoute une contrainte (« ne pas répéter ») rendant votre production \textit{moins} aléatoire que celle de l'amnésique. Paradoxalement, le libre arbitre rend vos choix moins aléatoires, pas davantage. Le substrat optimise en vue de la nouveauté, et appelle le résultat liberté.

Que devient alors le libre arbitre ? Il n'est pas éliminé, mais déplacé --- exactement là où l'analogie de l'horloge le prédit. Votre modèle du soi conscient ne prend pas de décisions en temps réel. Il est trop lent pour cela. Mais il n'est pas non plus un simple spectateur passif.

Pour l'essentiel, le système implicite utilise votre expérience consciente comme instrument d'évaluation : il présente les décisions à la simulation pour que celle-ci puisse en jauger les conséquences, dérouler des scénarios, ressentir les issues. C'est la fonction première de la couche virtuelle --- la manière qu'a le substrat de s'observer lui-même. Mais le modèle conscient évalue aussi de son propre chef, indépendamment, avec la bande passante dont il dispose --- bien moindre que celle du substrat, mais réelle. Ces évaluations, avec le temps, remodèlent les modèles implicites. Elles mettent à jour les poids, réentraînent le réseau, déplacent le paysage pour la \textit{prochaine} décision inconsciente.

Vous ne choisissez pas votre prochaine action au moment d'agir. Vous façonnez le système qui choisit --- par la réflexion, l'évaluation et la lente sédimentation de l'expérience consciente en structure implicite. Le libre arbitre n'est pas un instant. C'est un processus --- un processus qui opère à l'échelle des jours et des années, pas des millisecondes. Et la couche consciente n'est pas une simple passagère : elle est activement utilisée \textit{par} le substrat comme mécanisme d'évaluation, et elle lui renvoie ses propres appréciations indépendantes. Une circulation à double sens, pas une narration à sens unique.

Il existe une version plus sombre de tout cela, que j'ai vécue dans ma propre chair, et elle m'a appris davantage sur l'architecture de la volonté que n'importe quelle expérience de laboratoire.

La première fois, c'était pendant le service militaire obligatoire autrichien. Une marche forcée de 40 kilomètres --- trois jours et trois nuits de privation de sommeil, dans des conditions qui rendraient nerveux les juristes de la Convention de Genève. Sur la dernière étape, nous devions porter des masques à gaz et des combinaisons complètes de protection NBC. Je marchais à moitié endormi et j'entendais par moments des voix. Pas des hallucinations auditives au sens psychiatrique, mais quelque chose de bien plus intime : les sous-processus concurrents de mon appareil de motivation et de planification, normalement fusionnés en un seul flux narratif, devenaient audibles séparément. Une voix était encourageante, d'une positivité presque agressive : \textit{Continue, n'abandonne pas, tu vas survivre.} Une autre était pessimiste, attirante dans son défaitisme : \textit{Abandonne, allonge-toi, rien de tout ça ne compte.} Ce n'étaient pas des présences extérieures. C'était \textit{moi} --- différentes facettes du paysage d'optimisation de mon substrat, normalement intégrées en une « volonté » unique par des signaux inhibiteurs descendants, et qui se séparaient maintenant parce que les neurotransmetteurs qui maintiennent cette intégration étaient rationnés au profit de processus de survie plus critiques.

La deuxième fois fut dramatique. Une avalanche --- pendant le service militaire là aussi, provoquée par la décision imprudente d'un officier qui fut sanctionné par la suite. Quatorze d'entre nous, presque engloutis. L'avalanche mit longtemps à s'immobiliser, et pendant tout ce temps j'étais convaincu que j'allais mourir. Assez longtemps pour que la dissociation des voix s'enclenche de nouveau --- cette fois non par épuisement, mais sous l'effet d'une terreur mortelle prolongée. Même mécanisme, autre déclencheur : la réponse au stress détournait les ressources en neurotransmetteurs des circuits inhibiteurs qui, normalement, fusionnent les sous-processus concurrents en une seule voix.

Et pendant ces quelques secondes d'avalanche --- quelques secondes de temps réel, pas plus ---, j'ai vu ma vie entière défiler devant mes yeux. C'est un phénomène de mort imminente bien documenté, et la théorie l'explique : sous une menace mortelle extrême, le système implicite effectue un déversement massif et parallèle de mémoire dans la simulation. La frontière de perméabilité s'ouvre en grand. Le substrat tourne à plein régime, déversant tant de contenu dans la simulation que le temps subjectif se découple du temps de l'horloge. Quelques secondes contiennent une vie entière. La même dilatation du temps que j'avais connue sous salvia, mais déclenchée par la biologie plutôt que par la pharmacologie.

Et je peux maintenant dire \textit{pourquoi} ces deux états --- le film de vie du cerveau mourant et le plongeon dans la salvia --- se ressemblent de l'intérieur. Souvenez-vous des deux cadrans du chapitre 5. Presque tout le temps, ils s'échangent l'un contre l'autre : recrutez le cerveau entier dans un seul événement, et l'événement tend à se simplifier ; élaborez un calcul fin et différencié, et il reste local, sans gagner le cerveau entier. L'expérience ordinaire chevauche la frontière entre les deux. Mais il existe un recoin rare où \textit{les deux} montent inhabituellement haut en même temps --- une grande partie du cerveau recrutée, \textit{et} le calcul maximalement riche. Et rappelez-vous comment le cerveau construit son sentiment de durée : non pas en lisant une horloge, mais à partir du volume brut de traitement qu'il effectue par seconde de temps réel. Poussez les deux cadrans ensemble, et vous comprimez une vie entière de traitement dans une poignée de secondes d'horloge. Vous ne vous \textit{souvenez} pas de votre vie --- vous la \textit{revivez}. Vous n'\textit{attendez} pas la fin de la salvia --- vous \textit{dépensez} les décennies. Un seul mécanisme, deux portes pour y entrer : l'hypoxie et un frein métabolique défaillant d'un côté, une substance kappa-opioïde qui brouille la même régulation de l'autre. La théorie rend cela vérifiable : les deux états devraient présenter, au même moment, une intégration cérébrale globale exceptionnellement élevée \textit{et} une complexité computationnelle inhabituellement élevée --- les deux choses que la conscience ordinaire est contrainte d'échanger l'une contre l'autre.

Deux voies complémentaires vers le même mécanisme. La marche montre qu'un épuisement physiologique prolongé peut déclencher la dissociation. L'avalanche montre qu'une terreur mortelle soutenue produit le même effet. Même résultat, causes différentes --- l'un et l'autre prédits par la théorie.

Dans les pires cas --- et j'ai eu la chance que les miens n'aillent jamais aussi loin ---, l'une de ces « voix » peut prendre le contrôle du corps, et le soi conscient devient spectateur. C'est le même mécanisme qui produit le syndrome de la main étrangère (où une main agit contre la volonté du patient) et certains épisodes psychotiques. Les processus d'optimisation concurrents du substrat sont toujours là. Ils sont, pour le dire simplement, ce que fait le centre du langage quand vous ne l'utilisez pas pour parler. Mais normalement, l'inhibition descendante les maintient sous le seuil de la conscience, fusionnant leurs sorties en l'expérience homogène d'une volonté une et unifiée. Quand cette inhibition défaille --- par épuisement, par psychose, sous l'effet de certaines drogues ---, l'illusion de la volonté unifiée se dissout, et vous voyez le comité qui, depuis toujours, menait la danse.

Ce cadre dissout trois expériences de pensée qui paralysent la philosophie de l'esprit depuis des décennies.

Premièrement, les \textbf{zombies}. David Chalmers vous demande d'imaginer un être physiquement identique à vous en tout point, mais dépourvu d'expérience consciente --- tout le comportement, rien du ressenti. La Théorie des Quatre Modèles répond : c'est incohérent. Si vous construisez l'architecture à quatre modèles et la faites tourner à la criticalité, la simulation \textit{est} l'expérience. Vous ne pouvez pas avoir les rouages sans les aiguilles --- non que les aiguilles soient magiquement attachées, mais parce que, dans cette architecture, les « aiguilles » sont constitutives de ce que font les rouages. Un zombie serait une horloge dont chaque rouage est à sa place, mais sans affichage --- ce qui signifie qu'elle ne fonctionne pas comme une horloge. L'architecture à la criticalité instancie nécessairement une simulation. Retirez la simulation, et vous avez changé l'architecture. Vous n'avez plus de zombie. Vous avez un système différent, cassé.

Deuxièmement, \textbf{la chambre de Mary}. Frank Jackson vous demande d'imaginer Mary, une neuroscientifique qui sait tout de la vision des couleurs mais a vécu toute sa vie dans une chambre en noir et blanc. Quand elle voit du rouge pour la première fois, apprend-elle quelque chose de nouveau ? Le débat standard porte sur la question de savoir si le savoir physique est complet. La Théorie des Quatre Modèles tranche net. Le savoir physique exhaustif de Mary est un savoir \textit{sur} le substrat. Quand elle voit du rouge, elle fait connaissance avec un nouveau \textit{quale} virtuel --- un nouvel état de son Modèle Explicite du Monde que sa simulation n'avait encore jamais instancié. Elle n'apprend pas un fait nouveau sur les neurones. Elle acquiert un nouveau \textit{mode de modélisation}. Sa simulation exécute un processus qu'elle n'a jamais exécuté, et le caractère en première personne de ce processus est constitutif de la simulation elle-même --- non un fait sur le substrat qu'elle aurait pu déduire des manuels. Elle apprend quelque chose, mais ce qu'elle apprend n'est pas une information. C'est une expérience --- une nouvelle configuration de son monde virtuel.

Troisièmement, \textbf{l'argument évolutionnaire contre l'épiphénoménisme}. Si la conscience ne cause rien, comment la sélection naturelle a-t-elle pu la façonner ? Pourquoi ne sommes-nous pas des zombies ? La réponse découle directement de l'analogie de l'horloge. La sélection naturelle ne cible pas la conscience comme un trait séparé, juché au sommet de la machinerie fonctionnelle. Elle cible des capacités fonctionnelles --- et le caractère phénoménal est constitutif de ces capacités, il ne s'y ajoute pas. La sélection a façonné la simulation parce que la simulation \textit{est} l'architecture fonctionnelle, vue de l'intérieur. L'expérience n'est pas un passager épiphénoménal que l'évolution n'aurait pas pu voir. Elle est ce que l'architecture \textit{est} quand elle tourne. Demander pourquoi l'évolution a produit la conscience, c'est demander pourquoi les Suisses ont produit des cadrans d'horloge --- ils ne l'ont pas fait, séparément. Ils ont produit des horloges. Le cadran fait partie de ce qui fait d'une horloge une horloge.

\textbf{Le mystère de l'existence est déplacé, pas éliminé.} La Théorie des Quatre Modèles dissout le Problème Difficile de la conscience, mais elle n'explique pas pourquoi il existe un univers physique capable de faire tourner des auto-simulations. La question se déplace de « Pourquoi le cerveau produit-il de l'expérience ? » à « Pourquoi existe-t-il un univers dans lequel des systèmes auto-simulateurs peuvent exister ? »

En réalité, je crois avoir une réponse, ou du moins le début d'une réponse. L'univers est, de façon démontrable, capable de Classe 4. Fractales, criticalité auto-organisée, dynamiques du bord du chaos --- on les trouve partout, des systèmes météorologiques au tissu neuronal en passant par la formation des galaxies. Un univers capable de Classe 4 est, par définition, capable de calcul universel. Et un substrat computationnel à l'échelle de cet univers --- vaste sinon infini dans l'espace, le temps, peut-être l'échelle, et peut-être des dimensions que nous n'avons pas encore identifiées --- ne se contente pas de \textit{permettre} l'émergence de systèmes auto-simulateurs. Il la garantit presque, du moins si l'univers est infini dans certaines de ces dimensions. Non par chance, non par un coup de dés cosmique qui serait tombé, par hasard, sur la conscience, mais comme conséquence structurelle de ce que cet univers \textit{est}. Le mystère restant se situe un niveau plus profond : pourquoi existe-t-il, tout court, un univers capable de Classe 4 ? Cela, je l'ignore sincèrement --- même si l'on peut soupçonner que la question est mal posée, puisque « rien » est sans doute une abstraction platonicienne plutôt qu'un état de choses possible, et que tout ce qui existe doit posséder \textit{un} caractère computationnel, quel qu'il soit. Mais le saut de « univers capable de Classe 4 » à « êtres conscients qui demandent pourquoi ils sont conscients » --- cette partie-là découle de l'architecture.

Laissez-moi donc revenir à vous avec tout cela, car il est facile de suivre la théorie jusqu'aux confins du cosmos et d'y perdre le fil qui compte le plus. Pendant environ quatre-vingt-dix-neuf pour cent de votre vie éveillée, ce n'est pas vous qui décidez. Vous êtes celui qu'on informe. Le substrat s'engage, et une demi-seconde plus tard, « vous » recevez la nouvelle, proprement racontée, dans le bon ordre, avec votre nom dessus. Vous êtes un observateur différé à qui l'on a répété toute sa vie qu'il était l'auteur. Gardez cela bien en main pour la suite, car tout ce qui vient s'effondre si vous le lâchez.

Maintenant. La conscience atteint bel et bien le monde et le touche --- cela, au moins, est réel. Elle le fait de deux manières complètement différentes, qui ne se ressemblent en rien. La première est celle que nous entendons d'ordinaire. Le soi virtuel façonne le substrat, le substrat façonne le comportement, le comportement façonne le monde --- et cette boucle tourne lentement. Des mois. Des années. Vous décidez de devenir quelqu'un d'autre et, avec de la chance et de la persévérance, deux ans plus tard votre entourage le remarque. C'est la prise causale la plus indirecte qu'on puisse imaginer --- si indirecte que l'on peut passer une décennie à douter qu'elle fonctionne.

La seconde prise est immédiate, et c'est celle dont personne ne vous prévient. Vous ne pouvez pas changer, en cette seconde, la pièce où vous êtes assis. Vous \textit{pouvez} changer, en cette seconde, ce que vous êtes en train de vivre. Là, tout de suite, sans bouger, vous pouvez tirer votre attention vers l'intérieur, construire une scène qui n'est pas ici, rejouer un souvenir, vous enfoncer dans une rêverie, vous détacher du bruit qui vous entoure et vivre quelque part où le monde ne vous a pas placé. Pas de délai de deux ans. Pas de comportement, pas de détour par le façonnage du substrat via le monde extérieur. Vous vous retirez simplement dans un monde de votre propre fabrication, et vous le faites \textit{maintenant}. Cela ressemble au seul endroit où vous avez enfin la main --- au seul acte qui soit indubitablement, instantanément \textit{vôtre}.

Il ne l'est pas. C'est la partie qui fait mal, et elle fait mal précisément parce que nous venons de traverser Libet. La décision de vous retirer vers l'intérieur a été prise par le substrat avant que « vous » ne sachiez l'avoir prise --- exactement comme la décision de lever la main au laboratoire. « J'ai choisi de rêvasser » est le même rapport d'observateur différé que « j'ai choisi de bouger mon bras », en habits plus élégants. Le substrat a tiré le premier. Vous avez été prévenu après coup, et la notification est arrivée assez vite pour vous donner le sentiment d'en être l'auteur. L'immédiateté qui vous impressionne tant est la \textit{vitesse de la chaîne}, pas la liberté du choix. Rien ici n'est libre. Le retrait intérieur qui semblait être votre acte souverain et privé a été décidé pour vous, comme tout le reste.

Alors, qu'apporte réellement la conscience, si ce n'est le choix ? Elle apporte la \textit{nécessité}. Le substrat a peut-être donné le coup d'envoi, mais sans la simulation en marche, la réaction en chaîne ne mène nulle part --- il n'y a pas de détachement à ressentir, pas de monde intérieur où se retirer, pas de rêverie à habiter, parce que la simulation est le seul lieu où ces choses puissent advenir. Rayez la conscience, et toute la cascade en aval, tout simplement, ne se propage pas. Voilà le véritable second rôle causal : non un siège de la volonté, mais un maillon porteur de la chaîne. La conscience n'a aucun pouvoir causal indépendant --- et elle n'est pas non plus une passagère inutile, pas l'épiphénomène qui inquiétait les philosophes. Elle est la partie du mécanisme sans laquelle le reste ne peut pas se produire. Un maillon nécessaire, pas un agent libre. Ce qui est, je l'admets, moins exaltant que cela n'en a l'air au premier abord. Vous partez à la recherche du seul endroit où vous êtes enfin libre, et vous trouvez un endroit où vous êtes indispensable, mais toujours pas aux commandes. La prise rapide s'est révélée rapide, pas libre.

Et ici, je dois être honnête sur une limite, car la tentation est de viser trop loin. Que le \textit{moment} du repli intérieur ne vous appartienne pas --- Libet couvre au moins cela, et il ne s'agit plus seulement de bouger un membre : quand des personnes choisissent librement quelle image tenir devant l'œil de l'esprit, ce choix peut désormais se lire dans le cortex visuel plusieurs secondes avant qu'elles n'aient le sentiment de l'avoir fait. L'\textit{ouverture} de l'acte intérieur, autrement dit, paraît aussi pré-décidée que le fait de lever une main. Dans quelle mesure vous pilotez ce qui se passe \textit{dedans}, une fois que vous y êtes, est une autre question, et la réponse honnête est que nous ne le savons pas. Et la question est plus aiguisée qu'elle n'en a l'air. Ce qui est réellement en jeu, ce n'est pas la liberté dont dispose votre modèle du soi en général --- c'est la prise que le modèle du soi a sur le modèle du \textit{monde} : jusqu'où l'avatar peut plonger la main dans le monde simulé où il se trouve et le remodeler. C'est un superpouvoir à part entière, et pas des moindres. Il faut une certaine profondeur de modèle du soi pour seulement l'essayer --- raison exacte pour laquelle il n'est pas également distribué et pour laquelle il peut s'\textit{entraîner}, du moins en partie : de la même façon que l'on peut apprendre à escalader délibérément son propre cortex visuel, des phosphènes bruts aux motifs, puis jusqu'aux visages et aux scènes entières (chapitre 5), au lieu d'attendre de voir ce qui remonte. Combien de degrés de liberté la machinerie sous-jacente a accordés à votre modèle du soi pour cette tâche, personne ne les a comptés --- et la seule version véritablement ouverte de la question, ce que votre esprit vous sert ensuite quand rien ne le contraint, personne ne l'a vraiment testée. Peut-être pilotez-vous réellement ce que vous explorez là-dedans, dans une mesure réelle et entraînable. Peut-être, le plus souvent, non. L'introspection ne peut pas trancher ; c'est le seul instrument dont l'aveuglement est ici garanti, et trancher pour de bon exigerait la carte complète du câblage d'un cerveau humain et une quantité considérable de calcul. Prenez donc ce dégrisement pour exactement ce qu'il est, et pas un pouce de plus : l'\textit{ouverture} de l'acte intérieur, vous n'en êtes pas l'auteur. Ce que vous faites une fois le rideau levé --- jusqu'où vous pouvez plier le monde derrière lui --- reste une question ouverte, et une question honnête.

Voilà donc où l'argument du libre arbitre vous laisse : pas l'auteur --- ni de la main, ni de la rêverie, ni même de la question de savoir jusqu'où vous pilotez une fois à l'intérieur. Ce qui laisse exactement une chose pour laquelle il vaut la peine de tourner la page : que faites-vous d'une volonté que vous ne conduisez pas ?


\chapter*{Chapitre 14 : La seule liberté disponible}
\addcontentsline{toc}{chapter}{Chapitre 14 : La seule liberté disponible}
\markboth{Chapitre 14 : La seule liberté disponible}{}

Et voici le contrepoids, car je ne veux pas que vous refermiez le livre ici en pensant que le monde intérieur est une cage. Cette même capacité --- se découpler de ce qui se trouve devant vous et laisser tourner votre simulation sur ses propres contenus --- est, employée à la juste mesure, l'une des choses les plus puissantes que l'évolution ait jamais construites. Chaque acte de cognition vit ici. La créativité aussi. La mémoire elle-même aussi, et l'astuce que les Anciens appelaient le palais de mémoire --- parcourir des pièces imaginées pour y ranger ce qui compte. L'imagination, la planification, toute la machinerie du « et si je faisais plutôt ceci », la capacité de sortir du présent et de voyager en avant et en arrière dans votre propre esprit --- rien de tout cela n'exige la coopération du monde, tout cela exige exactement cette aptitude à vous replier vers l'intérieur et à laisser la simulation travailler sur un contenu interne. La lente emprise tournée vers l'extérieur a bâti des civilisations au fil des générations. Celle-ci, tournée vers l'intérieur, bâtit une pensée entière en une seconde. Elle est le moteur de tout ce que nous entendons par pensée humaine. Saisissez-la avec légèreté, de manière réversible, à doses graduées, et c'est un superpouvoir.

Appuyez-vous dessus trop fort, en revanche, et le variateur qui vous laisse entrer peut rester bloqué en position ouverte. Poussez la retraite intérieure à pleine puissance, et le mince canal qui relie votre simulation au monde réel se trouve submergé --- le même verrouillage que j'ai décrit plus haut, le même emballement que les NDE et la salvia à haute dose vous imposent contre votre gré, sauf qu'ici, vous y êtes entré sur vos propres jambes. Poussez assez loin, et vous n'arrivez plus du tout à retrouver le chemin du monde partagé. C'est le sujet sous salvia qui enjambe la fenêtre parce que la fenêtre n'est plus réelle pour lui (chapitre 6). Et la simulation, poussée si fort, ne revient pas toujours d'un seul tenant --- elle peut bifurquer, se scinder, revenir sous la forme de plus d'un vous, ou de quelqu'un qui n'est pas tout à fait celui qui était parti (les deux esprits qui se partagent un même cerveau, chapitre 9). Employée avec mesure, c'est le meilleur outil dont vous disposiez. Poussée à outrance, elle est authentiquement dangereuse --- et si vous le faites trop, qui sait si vous reviendrez, et combien de vous.

\textbf{Ce que vous pouvez faire de ce savoir.} Si vous avez suivi la théorie jusqu'ici, vous savez désormais que votre moi conscient (votre Modèle Explicite du Soi) est une reconstruction, pas une lecture directe. Vous savez qu'il comble les lacunes, qu'il confabule et qu'il s'attribue le mérite de décisions qu'il n'a pas prises. Vous savez qu'il ne peut pas voir son propre substrat. Et vous savez que c'est tout ce que vous avez.

Cela a des conséquences pratiques. Il y a trois écarts que vous devriez surveiller comme le lait sur le feu, parce que c'est dans l'espace qui les sépare que loge l'essentiel du malheur humain :

\begin{enumerate}
  \item Ce que vous \textit{voulez} être --- votre moi idéal, la version de vous à laquelle votre Modèle Explicite du Soi aspire.
  \item Ce que vous \textit{croyez} être --- votre modèle du soi actuel, le « moi » que vous transportez avec vous chaque jour.
  \item Ce que vous \textit{êtes} réellement --- votre comportement réel, votre impact effectif sur les autres, vos schémas au niveau du substrat tels qu'on les observe de l'extérieur.
\end{enumerate}

L'écart entre 1 et 2 est le moteur de l'amélioration de soi. Il est sain, tant que l'idéal reste réaliste et que la discordance pousse à l'action plutôt qu'au désespoir. L'écart entre 2 et 3 est le plus dangereux --- parce que vous ne pouvez pas le mesurer seul. Votre ESM \textit{ne peut pas} observer avec exactitude son propre substrat. Vous avez besoin des retours d'autrui, y compris de ceux qui dérangent. Surtout de ceux qui dérangent.

Mon meilleur ami Bernhard et moi en avons fait un jeu. Nous avons un accord tacite : chaque erreur que commet l'autre est une occasion de moquerie immédiate et impitoyable. Rater une sortie au volant ? « Alzheimer, stade terminal --- je prends les clés ? » Écorcher un mot ? « Je crois que tu refais un AVC. Arrête de parler avant de t'étouffer avec ta langue. » Oublier un détail de la conversation de la semaine dernière ? « Tu veux de l'aide pour les mots croisés du jour, tout à l'heure ? »

Vu de l'extérieur, cela paraît pathologique. Vu de l'intérieur, c'est le système de calibrage de l'ESM le plus efficace que je connaisse. Chaque plaisanterie est un signal de correction : \textit{le modèle du soi vient de faire quelque chose que le substrat n'avait pas l'intention de faire.} Et parce que la moquerie est enveloppée d'une affection sincère --- nous nous retenons tous les deux de rire en assénant l'insulte ---, aucun de nous ne défend l'erreur. Nous mettons à jour. C'est là toute l'astuce : il vous faut quelqu'un en qui vous ayez assez confiance pour qu'il puisse être brutal, et une relation dans laquelle se tromper est drôle plutôt que menaçant.

La théorie ne vous dit pas comment vivre. Mais elle vous dit quelque chose d'important sur la manière de vous \textit{connaître vous-même} : traitez votre modèle du soi avec le même scepticisme salutaire que vous appliqueriez à n'importe quel modèle. Il est utile. Il est la meilleure représentation dont vous disposiez. Et il est, par nécessité architecturale, incomplet.

Il y a encore une chose que la théorie vous met entre les mains, et c'est la partie que j'ai failli ne pas écrire, parce qu'elle frôle le développement personnel et que j'ai une faible tolérance pour ce genre. Mais elle découle directement de tout ce qui précède, alors la voici.

Vous n'avez pas choisi de lire ce livre. Quelque chose l'a mis sur votre chemin --- une recommandation, un algorithme, un après-midi d'ennui --- et votre substrat, dans un état dont vous n'étiez pas l'auteur, a décidé de poursuivre. C'est de la chance, pas du mérite. Je vous ai dit dans les dernières pages que rien de ce que vous faites n'est librement choisi, et je le pensais. Vous pourriez donc légitimement demander à quoi rime tout cela, si les rouages devaient de toute façon tourner comme ils tournent.

Voici à quoi cela rime, et c'est le tournant porteur d'espoir vers lequel toute l'argumentation convergeait : les rouages sont \textit{façonnables}. Le déterminisme ne signifie pas que vous êtes coincé --- il signifie que vous êtes entraînable. Votre substrat n'est pas un cristal figé ; c'est un paysage que l'exposition et la pratique réécrivent physiquement. Lire ceci a changé quelque chose, que vous l'ayez voulu ou non. Et une fois qu'un système a rencontré une idée, la même machinerie non libre qui vous a amené jusqu'ici peut désormais « décider » --- une décision fabriquée par le substrat, oui, mais réelle, et lourde de conséquences réelles --- d'employer son temps autrement. Vous ne l'aurez pas choisi comme les partisans du libre arbitre rêvent de choisir. Mais le choix atterrit tout de même dans le monde, remodèle tout de même le réseau, modifie tout de même qui vous serez la semaine prochaine. Ce n'est pas rien. C'est la seule forme de liberté qui ait jamais été disponible, et il se trouve qu'elle suffit pour bâtir une vie.

Alors si je vous ai convaincu de quoi que ce soit, que ce soit de ceci : réservez du temps pour un vrai travail de conscience, comme vous réserveriez du temps pour la salle de sport, et traitez-le comme non négociable. La méditation, d'abord --- pas pour l'encens et les playlists, mais parce que dégager la scène à dessein est l'entraînement le plus propre que votre attention recevra jamais. Elle aiguise la discipline de fond que vous apportez à tout le reste, et elle laisse les sentiers que vous creusez à longueur de journée se régénérer vraiment. Réservez du temps pour \textit{penser} --- librement, ou braqué sur un problème qui vous tient à cœur --- et protégez-le du téléphone. Fabriquez des choses. Faites de la musique ; elle se fond dans tout le reste. Adoptez n'importe quel loisir, sport compris, qui sollicite principalement votre conscience d'une manière qui vous plaît réellement --- le critère n'est pas la productivité, c'est de savoir si la machine intérieure fait quelque chose qu'elle aime. Et réhabilitez le temps mort dont vous disposez déjà : chaque embouteillage, chaque salle d'attente, chaque file --- rien de tout cela ne vous est volé, tout cela vous est \textit{offert} : la permission ininterrompue de vous replier vers l'intérieur et de laisser tourner la simulation. Peaufinez cette retraite. Savourez-la. C'est le seul rendez-vous qui n'est jamais annulé.

Vient ensuite l'entretien --- appelez cela l'hygiène mentale, car c'est exactement ce dont il s'agit. D'abord, désapprenez le monologue intérieur négatif. Ce n'est ni de l'honnêteté ni du réalisme ; c'est un sillon usé dans le substrat, que vous pouvez lisser à nouveau par l'attention. Voyez les choses en noir et blanc, ou mieux, seulement en blanc --- laissez le gris aux rares cas où vous ne pouvez vraiment pas l'éviter. La zone grise est le terreau de la rumination, et la rumination est l'emprise intérieure retournée contre son propriétaire.

Et puisque nous parlons du substrat retourné contre lui-même : la colère. Être en colère contre une personne ou une chose revient, presque toujours, à vous punir \textit{vous-même} pour un problème qui vit hors de vous. Ce sentiment n'a jamais été conçu pour pointer vers l'intérieur. Il a évolué pour pousser à l'action à l'époque où il fallait encore courir après son dîner --- une poussée qui mobilise les ressources vers le monde, pas un corrosif que vous déversez sur vos propres circuits pendant que la chose fautive poursuit sa route, imperturbable, quelque part ailleurs. Quand vous remarquez que la colère n'a nulle part où aller sinon vers l'intérieur, c'est le signal qu'il est temps de la laisser tomber.

Et si vous avez essayé la méditation et que vous vous y êtes cassé les dents --- ce qui est courant, et nullement un échec ---, il existe une porte dérobée que je peux recommander d'expérience : les arts martiaux. Cela ressemble au contraire d'un esprit calme, tout en impacts et en adrénaline, mais la discipline qui consiste à entraîner le corps à agir pendant que vous écartez le bavardage conscient est la même compétence que celle qu'enseigne la méditation, abordée par l'autre versant. Honnêtement, des deux, la méditation est le chemin le plus \textit{difficile}, pas le plus facile --- rester assis, immobile, avec son propre substrat et lui demander de se taire est l'une des choses les plus ardues qu'un être humain puisse tenter. Si le coussin vous terrasse, prenez le chemin des écoliers qui passe par le dojo. La destination est le même esprit calme ; on vous permet simplement de transpirer pour y arriver.

\section*{Ce que je ne sais pas}

Une théorie qui prétend n'avoir aucune question ouverte n'est pas une théorie --- c'est une religion. Voici donc les endroits où je suis sincèrement incertain, ceux sur lesquels la prochaine décennie de travail devrait se concentrer.

\textbf{Les modèles implicites sont-ils eux aussi virtuels ?} (ou dans quelle mesure) L'IWM et l'ISM sont des « modèles », mais des modèles de quoi, exactement ? J'ai tracé une ligne nette entre le substrat réel et la simulation virtuelle, mais les modèles implicites se situent précisément sur cette ligne. S'ils sont eux aussi virtuels en un certain sens, alors qu'est-ce qui constitue le fond véritablement « réel » ? La théorie suppose une division nette entre réel et virtuel, mais la réalité pourrait être plus désordonnée que mes diagrammes. C'est une question fondamentale à laquelle je n'ai pas de réponse définitive.

\textbf{La formalisation mathématique.} La théorie est pour l'instant qualitative. Je peux dessiner des diagrammes, décrire des mécanismes et formuler des prédictions, mais je ne peux pas vous tendre une équation. L'exigence de criticalité invoque les automates cellulaires de Classe 4 de Wolfram, et il existe des outils formels issus de la théorie des systèmes dynamiques qui pourraient être mobilisés. Mais une formalisation mathématique complète --- des équations qui spécifient exactement quand et comment les modèles virtuels émergent des dynamiques du substrat --- n'existe pas encore. C'est la plus grande lacune. Une théorie de la conscience sans mathématiques est une théorie de la conscience que les physiciens ne prendront pas au sérieux, et ce sont eux qui savent bâtir.

\textbf{La conjecture automate-hologramme --- un défi ouvert.} Au chapitre 5, j'ai décrit trois relations possibles entre les systèmes holographiques et les automates cellulaires de Classe 4. La première (un substrat holographique produisant des dynamiques de Classe 4) est presque certainement ce que fait le cerveau, et si elle est belle, elle n'a rien de bouleversant. Mais les deux autres méritent bien plus d'attention que je ne leur en ai consacré alors.

Il y a en réalité trois questions ouvertes ici, chacune plus extraordinaire que la précédente.

\textit{Premièrement : un automate de Classe 4 peut-il produire des motifs holographiques comme sortie émergente ?} Des règles locales au bord du chaos peuvent-elles générer un encodage global et non local de l'information à titre de comportement émergent ? Si oui, vous auriez un système où des interactions purement locales créent spontanément le type de structure d'information distribuée et redondante que décrit l'holographie --- ce qui, de façon intrigante, est exactement l'allure que prend l'intrication quantique du point de vue de la théorie de l'information.

\textit{Deuxièmement : un automate de Classe 4 peut-il avoir une structure de règles holographique ?} Imaginez un automate cellulaire dont les règles elles-mêmes encodent une information de dimension supérieure dans une structure de dimension inférieure, comme un hologramme encode trois dimensions dans deux. Chaque interaction locale contiendrait implicitement une structure globale. Les règles ne se contenteraient pas de produire un comportement complexe --- elles \textit{seraient} l'encodage compressé de quelque chose de plus vaste, de dimension supérieure, projeté vers le bas dans un jeu de règles de dimension inférieure.

\textit{Troisièmement, et c'est celle qui m'empêche de dormir la nuit : les deux peuvent-ils être vrais à la fois ?} Un système dont les règles sont holographiques, dont les dynamiques sont de Classe 4 et dont la sortie est de nouveau holographique. Si une telle chose existe, vous avez un processus de calcul qui s'encode lui-même --- un univers qui calcule sa propre structure. L'entrée est holographique. Le traitement se fait au bord du chaos. La sortie est de nouveau holographique. C'est un point fixe --- une boucle auto-cohérente.

Si un tel automate existe, il fait \textit{exactement} ce que le principe holographique dit que l'univers fait. Pas un système qui ressemble à l'univers en un vague sens métaphorique. Un système qui encode une réalité de dimension supérieure dans des règles de dimension inférieure, qui calcule à la frontière entre l'ordre et le chaos, et qui génère de la complexité émergente à partir de cette compression. Ce n'est pas une métaphore de l'univers. Cela pourrait \textit{être} l'univers.

Je vais le dire clairement, parce que je pense que quelqu'un le devrait : s'il existe un automate cellulaire de Classe 4 doté d'une structure de règles holographique qui produit également une sortie holographique, je suis presque certain qu'il s'agit de l'univers. Ce serait une Weltformel --- une équation du monde, non pas au sens d'une formule que l'on écrit sur un tableau noir, mais au sens d'un processus computationnel qui engendre tout ce que nous observons, de la mécanique quantique à la relativité générale, jusqu'à l'émergence de la conscience elle-même.

C'est là, je l'admets volontiers, l'idée la plus spéculative de ce livre. Je n'ai aucune preuve. Je n'ai même pas de jeu de règles candidat. Et je dois reconnaître que l'argument qui va de la beauté mathématique à la réalité physique a été légitimement critiqué. Sabine Hossenfelder, entre autres, a souligné que l'élégance n'est pas une preuve. Elle a raison. L'exploration complète de cette idée fait l'objet des trois chapitres suivants. Mais les questions elles-mêmes sont bien posées et mathématiquement précises :

\textit{Existe-t-il un automate cellulaire dont la structure de règles est holographique et dont la dynamique relève de la Classe 4 ? Produit-il une sortie holographique ? Ces trois propriétés peuvent-elles coexister ?}

Ce sont là des questions pour les mathématiciens, non pour les neuroscientifiques. Des questions sur la combinatoire des espaces de règles, sur la possibilité qu'un encodage holographique et une universalité computationnelle coexistent dans un jeu fini de règles locales. On pourrait peut-être prouver qu'aucun automate de ce type ne peut exister, et ce serait en soi un résultat profond, car cela nous apprendrait quelque chose de fondamental sur la relation entre la compression de l'information et le calcul. Ou bien on pourrait prouver que de tels automates existent et peuvent être construits --- et alors nous aurions un candidat pour la description la plus fondamentale de la réalité physique jamais proposée.

Je ne sais pas quelle réponse est la bonne. Mais je sais que ces questions méritent d'être posées, et que personne ne semble les poser. Considérez donc ceci comme un défi ouvert. Prouvez-le ou réfutez-le. Si vous le prouvez, vous aurez peut-être trouvé le code source de l'univers. Si vous le réfutez, vous aurez établi un profond théorème d'impossibilité reliant l'holographie et le calcul. Dans un cas comme dans l'autre, la réponse compte énormément.

Et si vous trouvez effectivement un tel automate --- appelez-moi. J'ai quelques prédictions que j'aimerais vérifier.

\textbf{Quel mécanisme physique ?} La théorie exige la criticalité, mais elle reste délibérément agnostique quant au mécanisme physique qui la maintient. Est-ce la dynamique des colonnes corticales ? Des ondes stationnaires thalamo-corticales ? La modulation gliale de l'activité synaptique ? Les trois trouvent un appui empirique. La théorie dit « le substrat doit être à la criticalité », mais elle ne dit pas \textit{comment} le substrat y parvient et s'y maintient. Ce n'est pas un défaut --- cela signifie que la théorie s'applique quel que soit le mécanisme précis. Mais tôt ou tard, il faudra que quelqu'un tranche la question.

\textbf{Configuration minimale.} Peut-on avoir un EWM sans ESM ? Une expérience du monde sans expérience de soi ? Quelle est l'architecture minimale qui puisse être qualifiée de consciente ? Les niveaux graduels que j'ai décrits dans le chapitre sur les animaux nous aident --- on peut avoir un riche modèle du monde sans grand-chose comme modèle du soi, comme c'est probablement le cas d'un poisson. Mais où se situe exactement le seuil ? Combien de modèle du soi faut-il avant que la lumière se fasse ? J'ai soutenu que c'est l'ESM qui transforme la simulation en expérience, mais je n'ai pas spécifié la version minimale viable.

J'inclus ces questions non comme des faiblesses, mais comme des fronts de recherche. Ce sont les endroits où la théorie entre en contact avec la réalité et dit : testez-moi ici, formalisez-moi ici, brisez-moi ici si vous le pouvez.

Il y en a une que je n'ai pas pu laisser de côté --- celle dont je vous ai dit qu'elle me tient éveillé la nuit. Alors j'ai cessé de la tenir à distance et je suis parti à sa recherche. Ce qui vient ensuite, c'est là où cela m'a mené.


\chapter*{Chapitre 15 : Le même motif, partout}
\addcontentsline{toc}{chapter}{Chapitre 15 : Le même motif, partout}
\markboth{Chapitre 15 : Le même motif, partout}{}

J'ai passé une nuit d'été à l'endroit le plus sombre du Vorarlberg --- haut dans les montagnes, sans lumière artificielle à des kilomètres à la ronde. Allongé sur le dos, je regardais le ciel. La Voie lactée n'était pas une traînée pâle qu'il fallait plisser les yeux pour distinguer. C'était un fleuve, dense et lumineux, qui projetait une lueur visible sur la roche à côté de moi. Et à un moment de cette nuit-là, quelque chose a basculé. J'ai cessé de voir des étoiles au-dessus de moi et j'ai commencé à sentir la Terre sous moi --- tournant, m'emportant avec elle, une boule lancée à travers une galaxie de cent milliards de soleils. Pas comme une idée. Comme une sensation dans mon corps. Le sol n'était pas immobile. Je m'accrochais à la surface extérieure de quelque chose qui se déplaçait à travers quelque chose d'incompréhensiblement vaste.

Ce vertige --- la certitude soudaine, physique, que vous êtes une petite chose chaude à la surface d'un rocher à l'intérieur d'un univers --- est le sentiment que je vous demande de garder en vous en lisant ce qui suit. Car ce chapitre demande ce qu'est réellement cet univers. Et la réponse a un air très familier.


Au chapitre précédent, je vous ai laissé face à un défi ouvert. J'ai décrit trois relations possibles entre les systèmes holographiques et les automates cellulaires de Classe 4, et j'ai posé la question la plus difficile : les trois peuvent-elles coexister dans un même système ? Un automate de Classe 4 peut-il avoir une structure de règles holographique \textit{et} produire une sortie holographique ? J'ai dit que l'exploration complète de cette idée ferait l'objet des trois chapitres suivants.

Ce travail, le voici.

Ce qui suit est la partie la plus spéculative de ce livre. C'est aussi, je le crois, la plus importante. Car lorsque je me suis réellement assis pour suivre le fil --- lorsque j'ai cessé de le traiter comme une question pour plus tard et que je me suis mis à tirer ---, je n'ai pas abouti là où je m'y attendais. Je m'attendais à trouver une curiosité mathématique intéressante. J'ai trouvé, au lieu de cela, un modèle cosmologique. Et la même architecture que j'avais sous les yeux depuis vingt ans.

Laissez-moi vous montrer ce que je veux dire.


\section*{La classe computationnelle de l'univers}

Dans l'annexe C, j'ai exposé les cinq classes de calcul --- un spectre allant de l'ordre parfait au désordre parfait, avec la Classe 4 au bord du chaos, complexité maximale que des règles exprimables puissent atteindre. Le cerveau utilise les cinq classes comme des outils, mais la conscience vit exclusivement dans la Classe 4. Tel était l'argument concernant le cerveau.

La question plus vaste est : de quelle classe relève l'univers ?

Ce n'est pas une métaphore. Je pose la question littéralement : si vous traitez l'univers comme un système dynamique --- ce qu'il est ---, où se situe-t-il sur le spectre des cinq classes ? La réponse, comme je vais le montrer, se détermine par élimination. Et l'élimination est étonnamment nette.

\textbf{Classes 1 et 2 --- Statique et Périodique.} Un univers de Classe 1 converge vers un état fixe. Rien ne se passe. Un univers de Classe 2 s'installe dans des boucles répétitives --- l'équivalent cosmique d'une horloge qui tique pour l'éternité. Ni l'un ni l'autre ne peut produire la chimie, la biologie, l'évolution ou la conscience. Nous existons. Nous sommes conscients. Un univers qui produit de la conscience doit être au moins de Classe 4, parce que la conscience exige une dynamique de Classe 4 --- c'était l'argument du chapitre 5. Et une classe inférieure ne peut pas engendrer une classe supérieure comme sous-processus. Un univers périodique ne peut pas produire une dynamique du bord du chaos, pas plus qu'une horloge ne peut se mettre spontanément à penser. Éliminées.

\textbf{Classe 3 --- Fractale.} Celle-ci est plus subtile, parce que des univers fractals seraient beaux. Une structure auto-similaire à toutes les échelles, des motifs emboîtés dans des motifs. De fait, l'univers \textit{a} bel et bien une structure fractale --- les amas de galaxies, les littoraux, les réseaux fluviaux, les ramifications de vos poumons. Mais les systèmes fractals sont computationnellement \textit{réductibles}. Cela signifie que l'on peut sauter en avant. On peut calculer l'état d'un système fractal au pas de temps dix milliards sans dérouler tous les pas intermédiaires. Il existe un raccourci.

Notre univers n'autorise pas les raccourcis. Vous ne pouvez pas prédire la météo du mois prochain en écrivant une équation qui saute en avant. Il faut faire tourner la simulation pas à pas, parce que la dynamique est computationnellement irréductible --- chaque instant dépend véritablement du précédent, d'une manière qui ne se laisse pas compresser. Un univers fractal, si riches que soient ses motifs, ne possède pas cette propriété. Il ne pourrait pas soutenir le calcul universel dont notre univers fait manifestement la démonstration. Nous construisons des machines de Turing. Nous avons une conscience. Un univers fractal ne peut faire ni l'un ni l'autre. Éliminée.

\textbf{Classe 5 --- Aléatoire.} Si la dynamique fondamentale de l'univers était véritablement aléatoire (réellement aléatoire, et pas seulement d'apparence complexe), alors la physique serait impossible. Non pas la physique telle que nous la comprenons aujourd'hui, mais la physique \textit{en tant que projet}. Toute l'entreprise de la science repose sur l'hypothèse que l'univers suit des règles exprimables : des règles que l'on peut écrire, tester, communiquer et utiliser pour prédire des observations futures. Un univers véritablement aléatoire n'a pas de règles exprimables. Sa dynamique ne peut être compressée en aucune formule, aucune loi, aucune équation. Vous ne pourriez pas écrire F = ma, parce que la relation entre force, masse et accélération changerait d'instant en instant, d'une façon qu'aucune description finie ne pourrait saisir.

Dans un univers de Classe 5, chaque expérience est un cas unique. Les résultats reproductibles sont des coïncidences. La science est une illusion qui a semblé fonctionner pendant un temps. Ce n'est pas logiquement impossible --- il n'y a aucune contradiction à imaginer un tel univers ---, mais c'est une catastrophe sur le plan explicatif. Si vous l'acceptez, vous ne pouvez plus rien expliquer, pas même pourquoi vos explications ont jamais semblé fonctionner. Éliminée, non par la logique, mais par abduction : la meilleure explication de notre expérience si constamment régie par des lois est que l'univers fonctionne selon des règles exprimables.

\textbf{Reste la Classe 4.} Le bord du chaos. Et la Classe 4 n'est pas seulement compatible avec ce que nous observons --- c'est la \textit{seule} classe qui satisfait à tous les critères.

L'univers contient des structures stables : atomes, cristaux, montagnes. C'est un comportement de Classe 1. Il contient des phénomènes périodiques : orbites, marées, battements de cœur. C'est un comportement de Classe 2. Il contient une structure fractale : distributions de galaxies, régimes météorologiques, ramifications neuronales. C'est un comportement de Classe 3. Et il permet le calcul universel : nous construisons des ordinateurs, et nous sommes conscients. C'est un comportement de Classe 4. Seul un système de Classe 4 peut contenir toutes les classes comme sous-processus --- y compris la sienne. Aucune des autres n'en est capable. Un automate de Classe 4 ne se contente pas d'héberger des dynamiques plus simples. Il héberge d'autres automates de Classe 4 : des ordinateurs universels au sein d'un ordinateur universel, chacun capable des mêmes prouesses computationnelles que le tout (simplement plus petit, plus lent et limité en ressources).

Il y a quelque chose de plus important encore. La Classe 4 possède un mécanisme d'auto-entretien qu'aucune autre classe ne possède : la \textbf{criticalité auto-organisée} --- que nous avons rencontrée pour la première fois au chapitre 5. Per Bak a montré en 1987 que les systèmes au bord du chaos ne s'y trouvent pas par hasard --- ils s'y \textit{conduisent eux-mêmes}. Empilez du sable grain par grain, et le tas s'organisera de lui-même jusqu'à l'angle critique où surviennent des avalanches de toutes tailles. Le système n'a pas besoin d'une main extérieure qui le règle sur la criticalité. Il se règle lui-même. Voilà pourquoi le bord du chaos est stable sur des échelles de temps cosmiques : c'est un attracteur, pas une coïncidence.

Je veux être clair sur la nature de cet argument. Ce n'est pas une preuve déductive. Deux des quatre éliminations reposent sur des observations empiriques (l'univers contient de la conscience ; il permet le calcul universel). Une repose sur l'abduction (la Classe 5 rend la science impossible --- insatisfaisant, mais sans contradiction logique). Le dossier positif en faveur de la Classe 4 combine des preuves et un mécanisme. C'est l'affirmation la plus forte que l'on puisse formuler : la Classe 4 est l'unique classe compatible avec toutes les observations, et la seule qui fournisse une raison à sa propre persistance.

Revenez au rocher un instant. Cette chose à laquelle je m'accrochais sur le flanc de la montagne --- celle qui m'a noué l'estomac quand je l'ai sentie bouger --- n'est pas seulement vieille, énorme et indifférente. Elle est de Classe 4. Le bord du chaos n'est pas une formule que j'aurais saisie pour paraître profond ; c'est le sol qui me portait, se réglant lui-même sur la criticalité grain par grain, comme le fait le tas de sable de Bak, maintenant tout l'équilibre sans la main de personne sur le cadran. Je n'étais pas allongé sur une pierre morte. Je m'accrochais à la surface extérieure d'un calcul qui se maintient en équilibre sur le fil du rasoir --- et vous aussi, en ce moment même, que vous sentiez la rotation ou non.


\section*{L'horizon de l'information}

Passons maintenant aux frontières.

La vitesse de la lumière est finie. C'est l'un de ces faits qui semblent anodins jusqu'à ce que vous y réfléchissiez dix minutes --- après quoi il réorganise toute votre image de la réalité.

La lumière voyage à environ 300 000 kilomètres par seconde. C'est assez rapide pour traverser la pièce avant que vous ayez pu cligner des yeux, mais l'univers est très, très grand. L'étoile la plus proche est à quatre années-lumière. La grande galaxie la plus proche est à deux millions et demi d'années-lumière. L'univers observable mesure environ 93 milliards d'années-lumière de diamètre. Quand vous regardez une galaxie lointaine, vous la voyez telle qu'elle était il y a des milliards d'années, parce que c'est le temps que sa lumière a mis à vous parvenir. Vous regardez toujours, inévitablement, vers le passé.

Mais il y a une conséquence plus profonde, et elle découle de l'expansion de l'univers.

En 1998, deux équipes d'astronomes ont fait une découverte qui leur a valu le prix Nobel : l'expansion de l'univers s'accélère. Pas seulement une expansion --- une accélération. Les galaxies lointaines s'éloignent de nous, et la vitesse à laquelle elles s'éloignent augmente. Cela signifie que pour tout observateur, il existe une distance au-delà de laquelle la vitesse de récession dépasse la vitesse de la lumière. Au-delà de cette distance, aucun signal ne vous parviendra jamais. Non pas parce que l'information serait cachée derrière un mur, mais parce que l'espace entre vous et elle grandit plus vite que la lumière ne peut le traverser.

C'est l'\textbf{horizon cosmologique}. Ce n'est pas une surface physique. Il n'y a pas de mur là-bas. C'est une conséquence de la géométrie et de la vitesse, mais c'est une barrière aussi absolue que n'importe quel mur. L'information située au-delà de l'horizon vous est à jamais inaccessible. Elle pourrait tout aussi bien ne pas exister.

Il existe une frontière semblable tout en bas. La \textbf{longueur de Planck} (environ $10^{-35}$ mètre, un nombre si petit que le qualifier de « petit » revient à qualifier l'univers observable de « moyennement grand ») est l'échelle où la physique telle que nous la connaissons s'effondre. En dessous de cette échelle, nos équations ne fonctionnent plus. L'espace-temps lui-même perd toute signification physique. Aucune mesure en dessous de la longueur de Planck n'est possible, même en principe. Ce n'est pas une limitation technologique. C'est une frontière fondamentale de ce qui peut être connu.

Entre l'horizon cosmologique et l'échelle de Planck : environ 60 ordres de grandeur. C'est le domaine computationnel de l'univers --- la plage à l'intérieur de laquelle la physique opère. Au-dessus et en dessous, les rideaux sont tirés.

C'est ce qui fait de l'univers ce que j'appelle un univers \textbf{quasi-infini}. Il n'est pas véritablement infini --- ou du moins, vous ne pourrez jamais vérifier qu'il l'est, puisque vous ne pourrez jamais accéder à plus qu'une région finie. Mais il n'est pas non plus fini en quelque sens atteignable que ce soit. La frontière recule plus vite que vous ne pouvez vous en approcher. Vous ne pouvez jamais atteindre le bord, mais le bord est là. De l'intérieur, l'univers paraît sans limites. De l'extérieur --- mais il n'y a pas d'extérieur. C'est bien là tout l'enjeu.

Voilà le vertige qui revient, sauf qu'il a maintenant des coordonnées. Sur le flanc de la montagne, je sentais le sol m'emporter à travers l'obscurité. Ajoutez ces deux frontières, et tout empire : au-dessous de moi, un plancher que je ne pourrai jamais atteindre ; au-dessus de moi, un bord qui fuit plus vite que la lumière ne peut le poursuivre, dans toutes les directions à la fois, sans aucun dehors où s'échapper. Toujours la même petite chose chaude à la surface d'un rocher --- sauf que l'obscurité dans laquelle elle tombe a désormais une forme, et cette forme est un mur qui recule devant moi de tous côtés et ne laisse rien revenir.


\section*{Chaque frontière est la même frontière}

Voici l'idée centrale de ce chapitre. Prenez le temps de l'assimiler, car si elle est juste, elle change votre façon de penser à tout.

Faisons l'inventaire. L'univers contient des singularités --- des lieux où notre description physique s'effondre, où le transfert d'information s'arrête, où les équations divergent ou se taisent. Ces singularités apparaissent à des échelles radicalement différentes, dans des contextes radicalement différents. Les physiciens les traitent comme des phénomènes distincts. Je pense qu'elles sont toutes une seule et même chose.

\textbf{1. Le régime de Planck.} À la plus petite échelle où la physique fonctionne, l'espace-temps se dissout en quelque chose que nous ne savons pas décrire. Aucune mesure n'est possible en dessous de cette échelle. L'information ne peut pas la traverser.

\textbf{2. L'intérieur des particules.} Les électrons et les quarks sont traités comme ponctuels dans le Modèle standard (de dimension zéro, sans structure interne). Nous ne pouvons pas voir à l'intérieur. Nous pouvons mesurer leurs propriétés (charge, spin, masse), mais nous n'avons aucun accès à ce qui se passe en leur cœur --- si tant est que le mot « cœur » ait un sens pour un objet sans étendue spatiale.

\textbf{3. Les horizons des événements des trous noirs.} L'information y tombe. Rien n'en ressort --- du moins sous aucune forme qui préserve ce qui y est entré. L'intérieur est causalement déconnecté de l'extérieur. Ce qui se passe à l'intérieur d'un trou noir reste à l'intérieur du trou noir, pour tout observateur extérieur.

\textbf{4. L'horizon cosmologique.} Le bord de l'univers observable, au-delà duquel l'expansion de l'espace empêche tout signal de nous parvenir. Pas de l'information cachée --- de l'information hors d'atteinte.

\textbf{5. Le Big Bang.} Le commencement. Toutes les lignes d'univers convergent. Chaque particule de l'univers fait remonter son histoire jusqu'à ce point --- ou plutôt jusqu'à cette frontière, car « point » implique qu'on pourrait s'y rendre, et on ne le peut pas.

\textbf{6. Le point final temporel.} La fin --- quelle que soit la forme qu'elle prenne. Si l'univers s'achève en mort thermique, l'entropie atteint son maximum et il ne reste aucun gradient thermodynamique pour alimenter le moindre processus. S'il s'achève en Big Crunch, toute la matière s'effondre en un point unique. S'il s'achève en Big Rip, l'expansion accélérée déchire l'espace-temps à toutes les échelles. J'examinerai ces trois scénarios plus loin. Ce qui compte ici, c'est l'affirmation structurelle : quelle que soit la fin que l'univers connaîtra effectivement, il se termine par une frontière imperméable à l'information.

Six singularités. Six échelles différentes, six contextes différents, six branches différentes de la physique pour les étudier. Mais regardez ce qu'elles ont en commun.

\textbf{Premièrement : elles sont toutes imperméables à l'information.} Aucune information ne peut franchir aucune d'entre elles. Vous ne pouvez pas mesurer en dessous de la longueur de Planck. Vous ne pouvez pas voir à l'intérieur d'un électron. Vous ne pouvez pas récupérer l'information tombée derrière un horizon des événements. Vous ne pouvez pas recevoir de signaux venus d'au-delà de l'horizon cosmologique. Vous ne pouvez pas observer ce qui est venu « avant » le Big Bang. Et vous ne pouvez pas transmettre de message par-delà la frontière finale de l'univers, quelle que soit la manière dont elle se manifeste.

\textbf{Deuxièmement : elles représentent toutes une densité d'information maximale.} C'est plus subtil, et cela vient de la borne de Bekenstein --- un résultat des années 1980 montrant que la quantité maximale d'information qu'une région de l'espace peut contenir est proportionnelle à sa \textit{surface}, et non à son volume. Les horizons des événements des trous noirs saturent cette borne --- ils contiennent le maximum d'information possible par unité de surface. Le principe holographique, proposé par Gerard 't Hooft et Leonard Susskind, généralise ce résultat : toute l'information contenue dans une région quelconque est encodée sur sa frontière. Ces singularités sont toutes des surfaces frontières fonctionnant à capacité maximale.

\textbf{Troisièmement : elles délimitent toutes le domaine computationnel.} La physique opère \textit{entre} ces frontières, pas au-delà. Les lois de la physique décrivent ce qui se passe dans la région comprise entre l'échelle de Planck et l'horizon cosmologique, entre le Big Bang et le point final, quel qu'il soit, qui nous attend. Les frontières définissent l'arène. Hors de l'arène, les règles ne s'appliquent pas, non pas parce que d'autres règles s'appliqueraient, mais parce que « règles » cesse d'être un concept doté de sens.

Trois propriétés communes. Six phénomènes. La vision conventionnelle veut que ce soient six choses différentes qui partagent par hasard quelques traits. Je pense que la vision conventionnelle est fausse. Je pense qu'il s'agit d'\textbf{un seul phénomène} --- la frontière informationnelle de l'automate --- apparaissant à six échelles différentes.

C'est une affirmation de symétrie. Le même élément structurel, répété. Et dans un système de Classe 4, c'est exactement ce à quoi l'on s'attendrait. La dynamique de Classe 4 contient toutes les classes comme sous-processus, y compris la Classe 4 elle-même. Un automate de Classe 4 imbrique de plus petits automates de Classe 4 dans sa dynamique, chacun délimité par des frontières informationnelles possédant les mêmes propriétés structurelles que le tout. Si l'univers est un automate de Classe 4, sa structure de frontières devrait se répéter à toutes les échelles. Et c'est précisément ce que nous semblons trouver.


\textit{Voilà donc ce que cette nuit sur le flanc de la montagne signifiait vraiment. Je ne me cramponnais pas à un rocher. Je me cramponnais à l'extérieur d'un automate de Classe 4, et chaque frontière en lui --- y compris celle contre laquelle mon esprit ne cessait de buter, là-bas au-delà de la dernière étoile visible --- est la même frontière. Dans le prochain chapitre, je suivrai ce fil jusqu'à ses conséquences. Et elles sont plus étranges que je ne m'y attendais.}


\chapter*{Chapitre 16 : L'architecture du tout}
\addcontentsline{toc}{chapter}{Chapitre 16 : L'architecture du tout}
\markboth{Chapitre 16 : L'architecture du tout}{}

Le chapitre précédent a posé une affirmation structurelle : chaque singularité de l'univers --- de l'échelle de Planck à l'horizon cosmologique, du Big Bang à la fin qui nous attend, quelle qu'elle soit --- est le même phénomène à des échelles différentes. Une seule frontière, répétée partout.

Je veux maintenant suivre ce fil plus loin. Car si toutes les frontières sont les mêmes, il en découle des conséquences remarquables --- sur le temps, sur la matière, sur les limites du savoir, et sur une architecture qui, d'ici la fin, devrait vous sembler très familière.


\section*{Le Big Bang n'est pas ce que vous croyez}

Cela va plus loin encore, car il y a une conséquence que la plupart des gens --- y compris la plupart des physiciens --- n'ont pas encore, je crois, pleinement mesurée.

Pensez à ce qui se passe quand on s'approche d'un trou noir depuis l'extérieur. À mesure que vous vous rapprochez de l'horizon des événements, le temps se dilate. Votre horloge, mesurée par un observateur lointain, ralentit. À l'approche de l'horizon, la dilatation tend vers l'infini. Un observateur lointain qui vous regarderait tomber vous verrait ralentir, glisser vers le rouge et vous estomper --- sans jamais tout à fait atteindre l'horizon. De son point de vue, vous mettez une éternité à arriver. Vous ne le franchissez jamais vraiment. L'horizon des événements est, vu de l'extérieur, une frontière asymptotiquement inaccessible.

Pensez maintenant à un voyage à rebours dans le temps, vers le Big Bang.

À quand remonte le Big Bang ? À environ 13,8 milliards d'années, nous dit-on. Mais c'est une mesure effectuée depuis \textit{l'intérieur} de l'univers en expansion, avec des horloges qui sont elles-mêmes des produits de cette expansion. Si vous imaginez remonter le temps, rembobiner le film cosmique, que se passe-t-il à l'approche de la singularité ? Le temps se dilate. La physique s'effondre. Plus vous vous approchez, plus les équations refusent de fournir un « moment zéro » définitif. Le Big Bang n'est pas un événement que l'on peut pointer du doigt en disant : « là --- c'est là que c'est arrivé ». C'est une frontière asymptotique. Vous pouvez vous en approcher d'aussi près que vous voulez, mais vous ne pouvez jamais l'atteindre.

Le Big Bang est un horizon des événements dans le temps, exactement comme l'horizon cosmologique est un horizon des événements dans l'espace.

Ce n'est pas du mysticisme. C'est une conséquence de la même structure mathématique. Un horizon des événements est une surface au-delà de laquelle l'information ne peut pas passer. Le Big Bang possède exactement cette propriété : aucune information d'« avant » (si tant est qu'« avant » veuille dire quelque chose) n'est accessible. Non pas parce qu'elle aurait été perdue ou cachée, mais parce que la frontière est imperméable à l'information. Il n'y a pas d'« avant » auquel accéder, de la même manière qu'il n'y a pas d'« intérieur » d'un trou noir auquel un observateur externe pourrait accéder. La frontière est la frontière. Un point c'est tout.

Et qu'en est-il de l'autre frontière temporelle --- la fin ?

Si l'univers finit en mort thermique --- entropie maximale, désordre maximal, plus aucun gradient thermodynamique pour alimenter le moindre processus ---, alors à ce stade, toute l'information est distribuée au maximum. La frontière du système contient le maximum d'information possible. C'est la saturation de Bekenstein. La mort thermique \textit{est} une singularité, selon la définition que j'utilise depuis le début : une frontière imperméable à l'information, à densité d'information maximale.

C'est ici que les choses deviennent étranges. Dans ce cadre, les singularités ne détruisent pas l'information. Elles la \textit{transforment}. C'est d'ailleurs la résolution vers laquelle la physique moderne converge pour le paradoxe de l'information des trous noirs --- ce débat vieux de plusieurs décennies sur la question de savoir si l'information est perdue lorsqu'elle tombe dans un trou noir. Le consensus actuel penche vers « non » : l'information est conservée, encodée sur l'horizon des événements, puis finalement réémise. La singularité transforme l'information entre formes comprimées et décomprimées.

Appliquez cela aux frontières temporelles. Si la mort thermique est une singularité, et si les singularités transforment l'information au lieu de la détruire, alors la mort thermique ne met pas fin à l'univers. Elle transforme l'information en un nouvel état comprimé. Et à quoi ressemble un état maximalement comprimé, à la saturation de Bekenstein ? Il ressemble aux conditions initiales d'une nouvelle expansion. Il ressemble à un Big Bang.

La clôture auto-référentielle n'est pas seulement spatiale. Elle est temporelle. L'univers ne commence ni ne finit --- il procède par cycles. L'état final est la condition initiale de l'itération suivante. Non pas grâce à quelque mécanisme de rebond exotique, mais parce que c'est précisément ce que \textit{font} les singularités qui conservent l'information : elles transforment des états de frontière comprimés en états intérieurs décomprimés, et inversement. La mort thermique comprime. Le Big Bang décomprime. C'est la même singularité, vue des deux côtés opposés.

Je reconnais que tout cela est spéculatif. Mais cela découle directement de deux affirmations : que toutes les singularités sont structurellement identiques, et que les singularités conservent l'information en la transformant. Si vous acceptez ces prémisses, la cyclicité temporelle n'est pas une hypothèse supplémentaire --- c'est une conséquence.

Mais la mort thermique n'est pas la seule façon dont l'histoire pourrait se terminer. Il existe une alternative sans doute plus étrange encore, et le cadre s'en accommode avec la même netteté.

Si l'énergie noire n'est pas constante mais croît avec le temps --- si sa densité augmente sans limite ---, alors l'expansion de l'univers ne se contente pas de continuer. Elle s'accélère au-delà de toute limite. C'est le scénario du \textbf{Big Rip}, et il est aussi spectaculaire que son nom le suggère. D'abord, les amas de galaxies sont déchirés, l'espace entre eux s'étirant plus vite que la gravité ne peut les maintenir ensemble. Puis les galaxies elles-mêmes se dissolvent. Puis les systèmes solaires. Puis les planètes. Puis les atomes eux-mêmes sont déchiquetés, l'expansion submergeant la force électromagnétique. Et enfin, l'espace-temps lui-même se fragmente. Chaque point devient une singularité.

Dans ce cadre, le Big Rip reçoit une interprétation naturelle. La frontière de singularité, qui se tient normalement à l'horizon cosmologique, à bonne distance, se propage \textit{vers l'intérieur}. Elle ne vous attend pas au bord de l'univers observable. Elle vient à vous. Elle fragmente le domaine computationnel en régions de plus en plus petites, chacune saturant sa propre borne de Bekenstein, chacune devenant sa propre frontière imperméable à l'information. Au lieu d'une grande singularité unique à la fin des temps, vous obtenez une explosion fractale de singularités, se propageant vers l'intérieur à toutes les échelles simultanément.

Et si les singularités sont des transformateurs d'information --- si elles ne détruisent pas le calcul mais le relancent ---, alors le Big Rip ne produit pas un seul redémarrage. Il en produit \textit{beaucoup}. Potentiellement une infinité. Chaque fragment du domaine computationnel brisé pourrait semer sa propre nouvelle expansion, sa propre nouvelle décompression, son propre nouvel univers. Le Big Rip, dans ce cadre, est un générateur de multivers.

Le cadre accueille donc non pas un mais trois scénarios de fin de partie, et tous trois sont structurellement cohérents :

La mort thermique : une singularité globale, un redémarrage. Le cas le plus simple --- le domaine computationnel tout entier atteint simultanément la saturation de Bekenstein, se comprime, puis se décomprime en un nouveau cycle.

Le Big Crunch : l'univers cesse de s'étendre et s'effondre sur un point unique. Une autre singularité globale, un autre redémarrage --- peut-être avec une inversion CPT, un renversement de la charge, de la parité et du temps qui fait du cycle suivant l'image en miroir du précédent.

Le Big Rip : la frontière de singularité se fragmente vers l'intérieur, produisant de nombreuses singularités, de nombreux redémarrages, potentiellement de nombreux univers. Non pas un cycle, mais un arbre qui se ramifie.

Je trouve cette robustesse rassurante plutôt qu'inquiétante. Un cadre qui ne fonctionne que si l'univers se termine d'une manière bien précise est fragile --- il parie sur une issue cosmologique particulière que nous ne pouvons pas encore déterminer. Ce cadre n'a pas besoin de parier. Sa logique structurelle (les singularités comme transformateurs d'information, les frontières comme élément architectural fondamental) tient quelle que soit la fin de partie que l'univers choisira réellement. C'est le genre de robustesse que l'on attend d'une théorie. Elle ne devrait pas dépendre de faits que nous ne connaissons pas encore.


\section*{Ce que sont vraiment les particules}

Il y a, enfouie dans ce cadre, une prédiction qui mérite d'être formulée explicitement, car c'est le genre de chose qui pourrait un jour être testée.

Les particules élémentaires (électrons, quarks, les briques élémentaires de la matière) sont traitées dans le Modèle standard comme ponctuelles. De dimension zéro. Sans extension spatiale. Cela a toujours été une commodité mathématique plutôt qu'une affirmation physique. Personne ne croit qu'un électron soit littéralement un point géométrique, car un point géométrique n'a pas de surface et ne peut donc, en vertu de la borne de Bekenstein, contenir aucune information. Or un électron contient de l'information --- charge, spin, masse, nombres quantiques. Quelque chose cloche dans l'image du « point ».

La prédiction est précise : les particules élémentaires sont des singularités à l'échelle de Planck. Elles ne sont pas véritablement de dimension zéro. Ce sont des frontières d'information miniatures --- de minuscules horizons des événements --- dont l'intérieur est aussi inaccessible que l'intérieur d'un trou noir. Elles possèdent une structure à l'échelle de Planck qui sature la borne de Bekenstein à cette échelle. Leur surface encode leurs propriétés de la même manière que l'horizon des événements d'un trou noir encode l'information de tout ce qui y est tombé.

Si c'est exact, alors la matière elle-même est faite du même élément structurel que les trous noirs, que le Big Bang, que l'horizon cosmologique. Des surfaces de singularité tout en bas, tout en haut, et à toutes les échelles intermédiaires. Les briques élémentaires de l'univers sont ses frontières.

C'est compatible avec des approches de la gravité quantique qui prédisent une longueur minimale à l'échelle de Planck --- on ne peut pas subdiviser l'espace en dessous d'un certain point, non pas parce que nos outils manquent de finesse, mais parce que l'espace lui-même est discret à cette échelle. Mais l'affirmation spécifique selon laquelle les particules \textit{sont} des singularités du même type que les horizons des événements --- cela, c'est nouveau. Et elle a une conséquence testable : le contenu en information d'une particule devrait varier avec sa surface (à la résolution de Planck), et non avec son volume. Si une théorie de la gravité quantique nous permet un jour de sonder la structure proche de l'échelle de Planck, voilà la signature à rechercher.

\section*{Les particules comme atomes de calcul}

Mais c'est ici que les choses deviennent vraiment intéressantes. Si les particules sont des singularités à l'échelle de Planck (des frontières d'information), alors elles ne sont pas seulement \textit{faites de} la même étoffe que le reste de l'architecture de l'univers. Elles sont les opérations de calcul fondamentales de l'univers. Elles sont, en un sens précis, les atomes du calcul. Pas des atomes au sens de la chimie --- des atomes au sens grec d'origine : \textit{atomos}, indivisible. Les unités irréductibles de ce que l'automate universel \textit{fait}.

Songez à ce qui en découle.

\textbf{Pourquoi seuls certains types de particules existent-ils ?} Cela a toujours été l'un des faits les plus étranges de la physique. Il existe exactement douze fermions fondamentaux (six quarks, six leptons), quatre bosons porteurs de force (plus le Higgs), et c'est tout. Pas un de plus. Le Modèle standard les catalogue, mais il n'explique pas \textit{pourquoi} ces types-là et aucun autre. Pourquoi y a-t-il un électron, mais pas de particule portant les deux tiers de la charge de l'électron et trois fois son spin ? Pourquoi le menu est-il si spécifique ?

Si les particules sont des frontières de singularité à l'échelle de Planck, la réponse est immédiate : parce que seul un nombre fini de configurations de frontière stables existe à l'échelle de Planck. Une frontière de singularité a une aire finie. À l'échelle de Planck, cette aire est aussi petite qu'une aire peut l'être. La borne de Bekenstein limite la quantité d'information que cette aire peut encoder. Une quantité finie d'information signifie un nombre fini d'états possibles. Et seuls certains de ces états sont \textit{stables} --- seules certaines configurations persistent sans se désintégrer. Ces configurations stables sont les types de particules. Le zoo de particules du Modèle standard n'est pas une liste mystérieuse et arbitraire. C'est le catalogue complet des configurations de frontière de singularité stables. Il y a un électron parce que cette configuration est stable. Il n'y a pas de particule aux propriétés fractionnaires bizarres parce qu'aucune configuration de frontière stable n'encode ces propriétés.

C'est d'ailleurs la même logique que celle des automates cellulaires. Un automate cellulaire possède une table de règles finie, et cette table ne permet qu'un certain nombre de motifs stables. Il peut créer une infinité de configurations --- plus le motif est grand, plus les possibilités combinatoires se multiplient --- mais les grandes structures peuvent être détruites rapidement par des motifs plus petits et plus stables qui entrent en collision avec elles, et très peu de configurations sont réellement indestructibles. Dans le Jeu de la Vie, les petites natures mortes (blocs, ruches, bateaux) et les petits motifs mobiles (les planeurs, dits vaisseaux) persistent indéfiniment, tandis que les grandes structures complexes sont fragiles --- un seul planeur bien visé peut les faire voler en éclats. La question de savoir si quelqu'un a systématiquement cartographié cette hiérarchie de fragilité reste, pour autant que je sache, une question ouverte --- Bak, Chen et Creutz ont montré en 1989 que le Jeu de la Vie présente une criticalité auto-organisée, mais la relation spécifique entre la taille d'un motif et sa résistance à la destruction n'a pas été formellement traitée. Quelqu'un devrait s'y atteler. Les particules, dans cette image, sont les planeurs et les vaisseaux de l'automate à l'échelle de Planck.

\textbf{Pourquoi les particules sont-elles discrètes ?} Pourquoi les nombres quantiques (charge, spin, charge de couleur) viennent-ils en multiples exacts, entiers ou demi-entiers ? Pourquoi n'existe-t-il pas de « demi-électron » ? Parce que les états de calcul sont intrinsèquement discrets. Un bit vaut 0 ou 1. Il n'y a pas de 0,37. Les nombres quantiques d'une particule sont des étiquettes d'information sur une configuration de frontière --- ils décrivent \textit{quelle} configuration stable la frontière occupe. Des états de frontière discrets produisent des nombres quantiques discrets. Le « quantique » de la mécanique quantique n'a rien de mystérieux. C'est ce qu'on obtient quand les objets fondamentaux sont des frontières d'information à capacité finie.

\textbf{Que se passe-t-il quand des particules interagissent ?} Quand deux électrons se repoussent, ou quand un quark émet un gluon, que se passe-t-il réellement ? Deux frontières d'information échangent de l'information. Cet échange \textit{est} du calcul. Les forces de la nature (l'électromagnétisme, la force forte, la force faible) ne sont pas une couche séparée posée par-dessus les particules. Elles sont la grammaire selon laquelle les frontières de singularité communiquent. Les règles qui déterminent quelles interactions sont permises et lesquelles ne le sont pas sont les règles de calcul de l'automate à l'échelle de Planck.

Une parenthèse personnelle s'impose ici. Quand j'avais quinze ou seize ans, mon oncle Bruno m'a lancé un défi. Il m'a dit que les forces attractives, en physique des particules, fonctionnent par échange de particules. « Maintenant, imagine, m'a-t-il dit, que tu te lances un lourd médecine-ball avec quelqu'un, dans un sens puis dans l'autre. Quand tu sauras m'expliquer comment cela vous amène à vous rapprocher de plus en plus l'un de l'autre, fais-le-moi savoir. » Le défi m'a occupé un bon moment. En physique newtonienne, c'est proprement inexplicable --- chaque échange devrait écarter les partenaires, pas les rapprocher. J'ai fini par revenir vers lui, et nous avons eu une très longue conversation sur la physique quantique à l'ombre d'un olivier, en Grèce, où il possède une maison.

Mais voici le point essentiel : dans un automate cellulaire comme le Jeu de la Vie, l'interaction attractive entre motifs est facile à reproduire. Deux structures persistantes peuvent échanger entre elles des structures plus petites dans des configurations où l'échange \textit{réduit} la distance. L'interaction n'a pas à obéir à l'intuition newtonienne, parce que la « force » n'est ni une poussée ni une traction --- c'est un échange d'information entre configurations de frontière, et ce sont les règles de l'automate qui déterminent si cet échange rapproche les frontières ou les éloigne. L'énigme du médecine-ball se dissout. Ce n'était une énigme que parce que nous imaginions une physique macroscopique. Au niveau du calcul, attraction et répulsion sont simplement deux issues différentes du même processus : un échange d'information gouverné par les règles de l'automate.

Les diagrammes de Feynman --- ces croquis iconiques d'interactions entre particules qui remplissent les manuels de physique --- sont littéralement des diagrammes de calcul. Chaque vertex est un échange d'information. Chaque ligne est une configuration de frontière qui se propage à travers le domaine de calcul. Voilà soixante-dix ans que les physiciens dessinent des images de calcul sans s'en rendre compte.

\textbf{Pourquoi les lois de conservation sont-elles si absolues ?} Charge, nombre baryonique, nombre leptonique --- jamais, pas une seule fois, on ne les a vus pris en défaut, dans aucune expérience jamais menée. Parce que ce sont des contraintes de conservation de l'information. Quand deux frontières échangent de l'information, les comptes doivent tomber juste : on ne peut ni créer ni détruire de l'information à une frontière de singularité, donc on ne peut ni créer ni détruire de charge. Les lois de conservation ne sont pas des règles imposées de l'extérieur. Elles sont la comptabilité. (La chaîne qui mène de la borne de Bekenstein, via l'unitarité, aux lois de conservation spécifiques est développée dans l'annexe F.)

\textbf{Et puis il y a le mystère des trois générations.} Chaque particule existe en trois exemplaires, identiques à la masse près --- électron, muon, tau ; up, charm, top. Personne ne sait pourquoi trois. Pas deux, pas quatre, pas dix-sept. Trois. L'image des atomes de calcul offre une hypothèse : les systèmes de Classe 4 portent une structure auto-similaire comme sous-processus, si bien que les configurations de frontière stables pourraient se répéter à trois échelles --- un motif de base et deux copies plus lourdes. Si c'est le cas, des générations supérieures existent aussi, mais ce sont des motifs plus grands, et les grands motifs sont fragiles --- des configurations plus petites et plus stables les découpent avant qu'ils ne parviennent à se constituer. Cela cadre avec ce que nous observons : le tau et le quark top se désintègrent déjà presque à l'instant où ils se forment. Les trois générations que nous obtenons pourraient être tout simplement les trois assez petites pour survivre. C'est une conjecture, pas une dérivation --- je l'expose en détail, avec la comptabilité des lois de conservation, dans l'annexe F.

Le terme que j'emploie pour cette image est \textbf{atomes de calcul}. Pas des atomes au sens de l'hydrogène et de l'hélium --- des atomes au sens d'éléments de calcul irréductibles. Les particules sont les opérations de base de l'automate universel. Chaque type de particule est un calcul stable à l'échelle de Planck. Chaque interaction est un échange d'information entre calculs. Chaque loi de conservation est une contrainte sur la manière dont ces échanges peuvent se dérouler. La physique, à son niveau le plus profond, ne parle pas de matière. Elle parle de calcul. Et les choses que nous appelons « matière » sont les briques irréductibles du calcul.


\section*{L'architecture}

Il est temps de rassembler les fils. Ce que j'ai décrit est un univers doté d'une architecture spécifique :

\textbf{Premièrement :} c'est un automate cellulaire de Classe 4. Il opère au bord du chaos, où la criticalité auto-organisée entretient la dynamique sans réglage externe. Il est irréductible sur le plan du calcul --- pas de raccourcis, pas de saut en avant. Chaque instant doit être calculé à partir du précédent. Et il contient toutes les classes comme sous-processus --- y compris lui-même : les atomes stables (Classe 1), les orbites périodiques (Classe 2), les côtes fractales (Classe 3) et, le plus crucial, d'autres automates de Classe 4 tournant à l'intérieur du grand calcul. Les cerveaux en sont un exemple.

\textbf{Deuxièmement :} il est holographique à tous les niveaux. L'information contenue dans une région quelconque est encodée sur sa frontière. C'est le principe holographique, né comme une conjecture sur les trous noirs et devenu l'une des intuitions les plus profondes de la physique théorique. Dans ce cadre, l'encodage holographique n'est pas seulement une propriété des trous noirs --- c'est une propriété de la structure de règles de l'univers lui-même. Les règles sont holographiques. La dynamique est de Classe 4. Et la sortie est de nouveau holographique.

\textbf{Troisièmement :} il est borné à chaque échelle par des surfaces de singularité toutes structurellement identiques. Frontières de Planck, intérieurs de particules, horizons des événements, horizon cosmologique, Big Bang, mort thermique --- même structure, échelle différente. Imperméables à l'information, saturées au sens de Bekenstein, et définissant le domaine de calcul.

Cette architecture a un nom. Je l'appelle le \textbf{SB-HC4A} : l'Automate Holographique de Classe 4 Borné par des Singularités (Singularity-Bounded Holographic Class 4 Automaton).

Le nom est un peu imprononçable. Mais il est précis, et chaque mot y mérite sa place.

La propriété la plus remarquable de cette architecture est la clôture auto-référentielle. La sortie du système \textit{est} le système. Il se calcule lui-même. Chaque état engendre le suivant, et l'état suivant est le calcul de l'état suivant. Il n'y a pas d'« extérieur » qui exécute le programme. Il n'y a pas d'ordinateur cosmique quelque part qui ferait tourner le code de l'univers sur un disque dur. L'univers \textit{est} le programme, l'ordinateur et la sortie. Les règles holographiques encodent le système complet sous forme comprimée. La dynamique de Classe 4 décomprime cet encodage en univers observable. La sortie holographique ré-encode le résultat. C'est une boucle. Un point fixe.

Vous pouvez formuler cela comme une condition formelle : \textbf{l'univers est un point fixe de sa propre dynamique.} Appliquez les règles à l'univers, et vous obtenez l'univers en retour. Pas une copie, pas une représentation --- la même chose. Le calcul et son résultat sont identiques.

Les mathématiciens ont une notation pour cela. Si vous appelez l'univers U et l'opération « calculer l'état suivant » la lettre grecque Phi, alors la condition de point fixe s'écrit :

\textit{Phi(U) = U}

L'univers, appliqué à lui-même, redonne l'univers. Il est auto-calculant.


\section*{Les limites de l'auto-description}

Il existe une conséquence de la clôture auto-référentielle qui mérite qu'on s'y arrête, car elle nous dit quelque chose de profond sur les limites de la connaissance.

En 1931, un logicien autrichien de 25 ans nommé Kurt Gödel démontra deux théorèmes qui ébranlèrent les fondements des mathématiques. L'essentiel, dépouillé du formalisme : tout système formel suffisamment puissant (capable, au minimum, d'exprimer l'arithmétique) contient des énoncés vrais qui ne peuvent pas être démontrés à l'intérieur du système. Et aucun système de ce genre ne peut démontrer sa propre cohérence.

Ce n'est pas une limitation technique. Ce n'est pas que nos démonstrations manquent d'ingéniosité. C'est une impossibilité structurelle. Les systèmes auto-référentiels d'une complexité suffisante sont intrinsèquement incomplets. Ils contiennent des vérités qu'ils ne peuvent pas atteindre de l'intérieur.

Appliquez cela à un univers auto-calculant.

Si l'univers se calcule lui-même --- s'il est un système formel suffisamment puissant (et la dynamique de Classe 4 garantit le calcul universel, donc il l'est) ---, alors les théorèmes de Gödel s'appliquent directement. L'univers ne peut pas contenir une description complète de lui-même. Il n'existe pas d'« équation du monde » que l'on pourrait écrire au tableau. Aucune formule qui, une fois résolue, vous révélerait tout sur l'univers.

Ce n'est pas parce que nous n'avons pas encore trouvé la bonne équation. C'est parce qu'\textit{aucune équation de ce genre ne peut exister}. La spécification complète d'un système auto-référentiel excède toute description qui n'est qu'une partie propre du système. L'univers ne suit pas une équation --- il \textit{est} le calcul. La seule description complète de l'univers est l'univers lui-même. Et vous ne pouvez pas en sortir pour prendre du recul sur l'ensemble, parce qu'il n'y a pas de dehors.

La Weltformel --- l'« équation du monde » dont les physiciens rêvent depuis Einstein --- n'est donc pas une équation. C'est un \textit{processus}. L'automate lui-même. Il ne peut s'exprimer qu'en étant exécuté.

J'y vois à la fois une leçon d'humilité et une libération. Une leçon d'humilité, parce que cela signifie qu'il y a des choses sur la réalité que nous ne pouvons pas connaître, même en principe. Une libération, parce que cela signifie que l'univers n'est pas un mécanisme qui attend d'être décodé --- c'est un calcul vivant, et nous en faisons partie. La vérité la plus profonde sur la réalité n'est pas une formule. C'est la réalité elle-même.


\section*{Le plafond cognitif}

Avant d'aller plus loin, je vous dois une objection. L'objection la plus profonde, en fait. Celle qui m'oblige à rester honnête.

Si nous sommes des automates de Classe 4 (si nos cerveaux fonctionnent au bord du chaos, dans la même classe de calcul que je viens d'attribuer à l'univers), alors le modèle SB-HC4A pourrait tout simplement être le concept le plus complexe que nos cerveaux de Classe 4 puissent concevoir. Nous ne pouvons pas penser en Classe 5. Nous ne pouvons pas concevoir de structures au-delà de notre propre classe de calcul. Le motif que nous trouvons --- la Classe 4 partout, auto-similaire à toutes les échelles, holographique et auto-référentielle --- pourrait être la signature de notre propre architecture cognitive projetée sur le cosmos, et non une propriété du cosmos lui-même.

Prenez un instant pour y réfléchir. L'évolution a fait de nous des détecteurs de symétrie. Les motifs les plus décisifs pour la survie dans l'environnement d'un chasseur-cueilleur (les visages des prédateurs et des proies) comptent parmi les plus symétriques. Nous sommes, au niveau le plus profond, des machines à reconnaître des motifs, optimisées pour trouver de la symétrie. Et le modèle SB-HC4A est fondamentalement une affirmation de symétrie : la même architecture à toutes les échelles. Il se peut que nous trouvions cette symétrie non parce qu'elle existe dans l'univers, mais parce que nos cerveaux sont constitutionnellement incapables de \textit{ne pas} la trouver.

C'est le Méta-Problème du chapitre 4, élevé à des proportions cosmiques. Le Modèle Explicite du Soi ne peut pas voir son propre substrat, il ne peut donc pas distinguer entre « l'univers a cette structure » et « mon cerveau ne peut modéliser l'univers qu'avec cette structure ». Le modèle cosmologique prédit sa propre infalsifiabilité potentielle, ce qui est soit la plus forte confirmation possible (le modèle prédit exactement cette limitation épistémologique), soit la plus forte objection possible (le modèle est un artefact de l'observateur, non une propriété de l'observé).

Un système de Classe 4 peut simuler tout ce qui atteint la complexité de Classe 4, celle-ci comprise. Mais il ne peut pas vérifier si l'univers va au-delà. Si l'univers est en réalité de Classe 5 --- véritablement aléatoire au niveau le plus profond --- mais \textit{apparaît localement} de Classe 4 aux observateurs de Classe 4, parce que la Classe 4 est le motif maximal que nous pouvons détecter, nous construirions exactement ce modèle. Et nous aurions tort. Nous aurions tort d'une manière que nous ne pourrions jamais découvrir de l'intérieur.

Je ne sais pas comment résoudre cette objection. Je ne suis pas sûr qu'elle puisse être résolue de l'intérieur. Je l'inclus malgré tout, parce que les faiblesses ont autant leur place dans l'exposé que les forces. Et le fait que ce modèle prédise sa propre limitation épistémologique --- qu'un système auto-référentiel ne puisse pas vérifier complètement sa propre description --- est soit son défaut le plus profond, soit sa plus profonde justification. Je ne sais sincèrement pas lequel des deux.


\section*{La chute}

Mais voici ce qui m'a stupéfié la première fois que je l'ai vu.

Regardez l'architecture que je viens de décrire :

\begin{itemize}
  \item Un système de Classe 4 opérant au bord du chaos.
  \item Délimité par une frontière opaque à l'information, à travers laquelle l'intérieur ne peut pas voir.
  \item Structure holographique --- la frontière encode l'intérieur.
  \item Clôture auto-référentielle --- le système se calcule lui-même.
  \item Un point fixe : le résultat du calcul est le calcul lui-même.
\end{itemize}

Revenez maintenant au chapitre 2. Regardez l'architecture à quatre modèles de la conscience :

\begin{itemize}
  \item L'automate cortical : un système de Classe 4 opérant au bord du chaos.
  \item La frontière implicite-explicite : une frontière opaque à l'information, à travers laquelle la conscience ne peut pas voir.
  \item Structure holographique --- les modèles implicites sont distribués, holographiques, et encodent le contenu intégral de l'expérience dans la structure neuronale.
  \item Clôture auto-référentielle --- le modèle du soi se modélise lui-même.
  \item Un point fixe : le Modèle Explicite du Soi se représente lui-même. Le modèle du modélisateur \textit{est} le modélisateur.
\end{itemize}

Même architecture. Mêmes propriétés formelles. Mêmes conditions aux limites. Même clôture auto-référentielle.

L'univers est un automate holographique de Classe 4 délimité par des singularités, où l'intérieur observable est la « simulation » et la frontière de singularité le « substrat ».

La conscience est un automate holographique de Classe 4 délimité par la frontière implicite-explicite, où les modèles explicites sont la « simulation » et les modèles implicites le « substrat ».

Même architecture. Échelle différente.

Ce n'est pas une métaphore. Je ne dis pas que la conscience est \textit{comme} l'univers. Je dis qu'il s'agit de la même \textit{sorte de chose} --- le même motif computationnel, instancié à deux échelles différentes. L'un au niveau cosmologique, l'autre au niveau neurologique. Et le fait que le motif se répète à travers les échelles est lui-même une prédiction du modèle --- non pas en raison d'une auto-similarité fractale (ce serait de la Classe 3, une affirmation plus faible), mais parce que les systèmes de Classe 4 contiennent des sous-systèmes de Classe 4. Un ordinateur universel peut en simuler un autre. Un automate de Classe 4 ne se contente pas de produire de jolis motifs auto-similaires. Il produit \textit{d'autres automates de Classe 4} au sein de sa propre dynamique --- plus petits, plus lents, limités en ressources, mais authentiquement universels. L'architecture ne \textit{paraît} pas seulement la même à différentes échelles. Elle \textit{est} la même.

Pour être très précis sur ce que j'affirme et sur ce que je n'affirme pas : je n'affirme pas que l'univers \textit{est} conscient en un quelconque sens expérientiel. Mais je n'affirme pas non plus qu'il ne l'\textit{est pas}. La réponse honnête, c'est que nous ne pouvons pas le savoir. Nous pourrions être rêvés par un cerveau de Boltzmann, ou par n'importe quel autre cerveau ; le solipsisme pourrait être vrai et je suis le cerveau de Boltzmann --- mais tout le monde peut dire cela. L'essentiel, c'est que la question est inconnaissable, et cette inconnaissabilité n'est pas un échec de la théorie, mais un trait structurel de la situation : on ne peut pas sortir du système pour aller vérifier. Ce que j'affirme \textit{bel et bien} est d'ordre architectural, non phénoménal. Je n'affirme pas que la conscience crée la réalité, ni que la réalité est un rêve, ni aucune des autres interprétations mystiques que ce genre d'observation structurelle a tendance à attirer. Le plan d'un bâtiment n'est pas un bâtiment. Mais si vous trouvez le même plan dans un gratte-ciel et dans une seule pièce de ce gratte-ciel, cela vous dit quelque chose de profond sur les principes architecturaux à l'œuvre.

La conscience est une instance locale d'un motif universel. Pas un accident cosmique. Pas un miracle. Une conséquence structurellement inévitable de la dynamique de Classe 4 à un degré de complexité suffisant. L'univers ne se contente pas de \textit{permettre} la conscience. Il la garantit quasiment --- parce que la même architecture auto-référentielle, holographique, au bord du chaos, qui fait de l'univers ce qu'il est, fait aussi de la conscience ce qu'elle est. Le motif qui engendre la réalité est le même motif qui engendre l'expérience de la réalité.

Et si cela ne vous coupe pas le souffle, c'est que vous ne l'avez pas encore compris.


\textit{Dans le chapitre suivant, je rassemblerai la théorie complète --- l'architecture de la conscience, l'architecture cosmologique et l'identité structurelle entre les deux --- et je demanderai ce que cela signifie pour la question la plus difficile de toutes : pourquoi existe-t-il quoi que ce soit ?}


\chapter*{Chapitre 17 : Le miroir le plus profond}
\addcontentsline{toc}{chapter}{Chapitre 17 : Le miroir le plus profond}
\markboth{Chapitre 17 : Le miroir le plus profond}{}

Le voici, parce qu'une fois qu'on l'a vu, on ne peut plus ne pas le voir.

Le chapitre 16 s'achevait sur un défi : regarder l'architecture SB-HC4A --- l'automate de Classe 4 auto-référentiel, holographique, borné par des singularités à chaque échelle --- puis regarder l'architecture des quatre modèles du chapitre 2. L'affirmation était qu'il s'agit de la même chose. Pas semblables. Pas apparentées par métaphore. Structurellement identiques.

Je veux maintenant vous guider à travers cette correspondance, pièce par pièce, jusqu'à ce qu'il devienne impossible de l'écarter. Puis je vous dirai où l'ensemble pourrait s'effondrer, parce que les lignes de faille méritent autant d'attention que la structure.

\section*{La correspondance structurelle}

Commençons par la frontière de singularité --- la barrière d'information que l'univers dresse à chaque échelle. Le régime de Planck tout en bas. Les horizons des événements autour des trous noirs. L'horizon cosmologique à la lisière de l'univers observable. Le Big Bang derrière nous, la mort thermique devant. Chacune de ces limites est un mur d'information : rien ne le traverse. On ne peut pas envoyer de signal à travers un horizon des événements. On ne peut pas sonder en dessous de la longueur de Planck. On ne peut pas voir au-delà de l'horizon cosmologique. L'univers est une pièce aux murs opaques à chaque échelle, et vous aurez beau presser votre visage contre la vitre, vous ne verrez jamais ce qu'il y a de l'autre côté.

Regardez maintenant votre cerveau. Les modèles implicites --- l'IWM et l'ISM, les poids synaptiques, la structure apprise, la vaste bibliothèque de tout ce que vous savez --- siègent derrière leur propre barrière d'information. Vous ne pouvez jamais faire l'expérience directe de vos synapses. Vous ne pouvez jamais introspecter les poids de connexion qui engendrent vos pensées. Le versant implicite est opaque à l'information : vous savez qu'il est là, parce que la simulation ne pourrait pas tourner sans lui, mais la simulation elle-même ne peut pas voir à travers la frontière qui les sépare. Le chapitre 2 a posé ce point. Le chapitre 3 l'a martelé. Et le voici de nouveau, à chaque échelle de l'univers. Même rôle structurel. Même opacité informationnelle. Même position architecturale.

La frontière de singularité en cosmologie correspond à la frontière implicite-explicite dans la conscience. C'est le même mur.

Ensuite : l'intérieur observable. Tout ce qui se trouve à l'intérieur des frontières de singularité (les atomes, les planètes, les galaxies, vous) constitue le versant décompressé. C'est là que la physique a lieu, là que les choses interagissent, là que l'information s'organise en ces structures que nous observons et mesurons. C'est la simulation de l'univers, si vous voulez : la partie qui calcule, la partie qui évolue, la partie où quelque chose se passe.

Dans votre cerveau, la structure correspondante, ce sont vos modèles explicites --- l'EWM et l'ESM. Votre expérience consciente. Le monde que vous voyez en ce moment même, le soi que vous avez le sentiment d'être. Voilà la simulation : générée en temps réel à partir des modèles implicites, mise à jour en continu, vive, détaillée et totalement convaincante. Tout ce que vous avez jamais vécu, votre vie entière, s'est produit à l'intérieur de cette simulation. Vous n'en êtes jamais sorti. Vous ne pouvez pas en sortir. Non pas parce que vous n'avez pas assez essayé, mais parce que « vous » êtes la simulation. Celui qui fait l'expérience et l'expérience elle-même sont un seul et même processus.

L'intérieur observable de l'univers correspond à vos modèles explicites. Même rôle : le versant décompressé, dynamique, interactif de l'architecture.

Vient maintenant la structure de règles holographique. L'information de l'univers n'est pas stockée dans son volume --- elle est stockée sur ses frontières. Le principe holographique, proposé par 't Hooft et étendu par Susskind, affirme que toute l'information contenue dans une région tridimensionnelle de l'espace est encodée sur sa surface bidimensionnelle. À la frontière, l'information est compressée, distribuée et structurellement complète. L'intérieur est une projection --- une décompression à plus faible bande passante de ce que la frontière encode.

Dans votre cerveau, ce rôle revient aux modèles implicites. Vos poids synaptiques encodent tout ce que vous savez du monde et de vous-même dans un format distribué, compressé et structurellement complet --- vous pourriez, en principe, reconstruire la simulation entière à partir du seul substrat, sans la moindre entrée sensorielle actuelle, et c'est exactement ce qu'est le rêve. Les modèles implicites sont holographiques au sens de Lashley : endommagez-en un morceau, et vous ne perdez pas un souvenir particulier --- vous perdez de la résolution sur l'ensemble des souvenirs. L'information est étalée sur tout le substrat, compressée, redondante, et inaccessible depuis le versant de la simulation.

La structure de règles holographique de l'univers correspond aux modèles implicites holographiques de la conscience. Même stratégie d'encodage. Même compression. Même inaccessibilité depuis le versant décompressé.

Puis il y a le régime dynamique. L'univers opère en Classe 4 --- au bord du chaos. Le chapitre 15 l'a établi par élimination : les Classes 1 et 2 sont trop simples, la Classe 3 ne peut pas calculer, la Classe 5 rend la physique impossible. Ce qui reste, c'est la Classe 4 --- le seul régime qui soutient le calcul universel, auto-organise sa propre criticalité et contient toutes les classes comme sous-processus, y compris lui-même. L'univers n'est pas simplement complexe. Il est complexe exactement de la manière qui s'auto-entretient, et exactement de la manière qui peut nicher de plus petites copies de lui-même au sein de sa propre dynamique.

Votre cortex fait la même chose. C'était l'objet du chapitre 5 : l'automate cortical opère au bord du chaos, en maintenant la criticalité par une régulation homéostatique de l'équilibre excitation-inhibition. Trop peu d'activité, et vous êtes en sommeil profond (Classe 2, périodique, inconscient). Trop, et c'est la crise convulsive --- poussée au-delà de la Classe 4, la simulation vole en éclats. Le point d'équilibre, l'endroit où vit la conscience, c'est le fil du rasoir entre l'ordre et le chaos. La criticalité auto-organisée y maintient le cerveau. La criticalité auto-organisée y maintient l'univers.

La dynamique de Classe 4 en cosmologie correspond à la criticalité corticale dans la conscience. Même régime. Même mécanisme d'auto-entretien.

Enfin, la correspondance la plus profonde : la clôture auto-référentielle. L'univers calcule sa propre structure. Sa dynamique produit son état, qui détermine sa dynamique, qui produit son état. Il n'y a pas de programmeur externe. Pas de dehors. Les lois de la physique ne sont pas imposées à l'univers depuis un ailleurs --- elles \textit{sont} la dynamique de l'univers, appliquée à lui-même. L'équation de point fixe est d'une simplicité presque absurde : le calcul de l'univers égale l'univers. Entrée, processus et sortie sont une seule et même chose.

Votre Modèle Explicite du Soi fait la même chose. L'ESM est un modèle qui s'inclut lui-même dans ce qu'il modélise. Il vous représente, et « vous » inclut la représentation. Le modèle et le modélisé coïncident. C'est la clôture auto-référentielle que j'ai décrite au chapitre 4 --- la raison pour laquelle la conscience fait l'effet qu'elle fait, la raison pour laquelle vous ne pouvez jamais vous voir entièrement de l'extérieur, la raison pour laquelle la simulation contient son propre observateur. Le point fixe de l'auto-représentation : l'état où le modèle et la chose modélisée sont une seule et même chose.

La clôture auto-référentielle de l'univers correspond à la clôture auto-référentielle de la conscience. Même structure de point fixe. Même auto-inclusion inéluctable.

Cinq correspondances. Pas de vagues ressemblances thématiques. Pas cette forme molle de détection de motifs qui vous fait voir des visages dans les nuages. Cinq traits structurels qui accomplissent le même travail, à la même position, dans les deux architectures.

Et maintenant le point crucial --- celui sur lequel je vous demande de vous arrêter un instant, parce que la tentation de le domestiquer est immense : ceci n'est \textit{pas} « l'univers est \textit{comme} la conscience ». Ce n'est pas une analogie. Ce n'est pas une métaphore. Ce n'est pas un parallèle suggestif qui nourrit les conversations de dîner mondain.

C'est une identité structurelle.

Le cerveau n'a pas évolué pour \textit{ressembler} à l'univers. Il a évolué \textit{comme} une instance locale, à échelle réduite, du même motif computationnel. Les systèmes de Classe 4 ne se contentent pas de contenir de l'auto-similarité fractale (ça, c'est la Classe 3 --- la répétition géométrique, jolie mais superficielle). Ils se contiennent \textit{eux-mêmes}. Un automate de Classe 4 peut héberger un autre automate de Classe 4 au sein de sa propre dynamique --- un ordinateur universel tournant à l'intérieur d'un ordinateur universel. C'est de l'auto-inclusion computationnelle, pas une simple auto-similarité géométrique. La conscience \textit{est} l'auto-inclusion de Classe 4 de l'univers, opérant à l'échelle biologique. Le motif qui fait tourner le cosmos à la plus grande échelle est le même motif qui fait tourner votre vie intérieure à l'échelle neurologique --- non pas parce que quelqu'un l'aurait conçu ainsi, ni parce que les fractales sont jolies à regarder. Parce qu'un ordinateur universel de taille suffisante engendre nécessairement d'autres ordinateurs universels en son sein. C'est ce que font les systèmes de Classe 4 : ils se nichent eux-mêmes.

\section*{L'énergie est de l'information}

Il existe une seconde ligne d'argumentation qui converge vers la même conclusion depuis une direction complètement différente, et je crois que c'est elle qui, au bout du compte, se révélera la plus importante --- parce qu'elle relie l'architecture à la physique d'une manière potentiellement testable.

Trois résultats indépendants en physique, développés par trois communautés différentes sur un demi-siècle, pointent vers la même conclusion extraordinaire.

Le premier est le principe de Landauer. En 1961, Rolf Landauer (un physicien d'IBM qui réfléchissait aux limites fondamentales du calcul) a démontré qu'effacer un bit d'information coûte une quantité minimale d'énergie. Non pas à cause de limitations d'ingénierie. À cause de la thermodynamique. L'univers vous fait payer l'oubli. Le fait a été confirmé expérimentalement en 2012, et il signifie quelque chose de profond : l'information et l'énergie sont interchangeables. On peut convertir l'une en l'autre. Ce ne sont pas des substances séparées. Elles s'échangent.

Le deuxième est la borne de Bekenstein. Jacob Bekenstein a montré que l'information maximale qu'une région de l'espace peut contenir est proportionnelle à sa surface, et non à son volume. C'est l'un des résultats les plus contre-intuitifs de toute la physique. On s'attendrait à ce qu'une boîte plus grande puisse contenir davantage d'information. Ce n'est pas le cas --- ou plutôt, la limite est fixée par la \textit{surface} de la boîte, non par son intérieur. Entassez trop d'information dans une région donnée, et elle s'effondre en trou noir. La densité maximale d'information est fixée par la géométrie et l'énergie --- encore un lien profond entre l'information et le monde physique.

Le troisième vient de la thermodynamique des trous noirs. Stephen Hawking et Bekenstein ont montré, dans les années 1970, que les trous noirs ont une température, une entropie, et obéissent aux lois de la thermodynamique. Le contenu en information d'un trou noir est écrit sur son horizon des événements --- sa surface, sa frontière. Et par le rayonnement de Hawking --- le processus quantique terriblement lent par lequel les trous noirs finissent par s'évaporer ---, cette information est progressivement restituée à l'univers. Les trous noirs ne détruisent pas l'information. Ils la transforment. Ils la compriment sur leur frontière, la retiennent, et finissent par la rayonner en retour.

Ces résultats sont partis de points de départ très différents. Landauer pensait aux ordinateurs ; Bekenstein et Hawking, à la thermodynamique des trous noirs. Les travaux de Landauer relevaient d'un tout autre domaine, élaborés des décennies plus tôt et pour des raisons sans rapport. Et pourtant, tout cela converge vers la même hypothèse : énergie et information ne sont pas seulement liées. Elles sont identiques. Deux noms pour la même chose. E égale I.

Si c'est vrai --- et je dois dire immédiatement que ce n'est pas prouvé, raison pour laquelle c'est le premier des quatre points faibles que j'expose dans l'annexe G ---, alors les singularités deviennent des transformateurs d'information. Elles ne détruisent ni ne créent l'énergie-information. Elles la convertissent d'une forme à l'autre. Forme comprimée : densité maximale sur la frontière, saturée au sens de Bekenstein, inaccessible depuis l'intérieur. Forme décomprimée : densité plus faible, organisée, répartie dans l'intérieur --- la physique que nous observons. Une singularité est un traducteur entre deux représentations de la même étoffe.

Regardez maintenant votre cerveau à travers ce prisme. Vos modèles implicites contiennent une information comprimée, à densité maximale : tout ce que vous avez jamais appris, encodé dans des poids synaptiques, structurellement complet et phénoménalement inaccessible. Vous ne pouvez jamais faire directement l'expérience de votre propre substrat. Vos modèles explicites sont la projection décomprimée, à bande passante réduite --- la simulation, l'expérience consciente, le monde que vous voyez et le soi que vous ressentez. Votre cerveau fait à l'échelle neuronale exactement ce que les singularités font à toutes les autres échelles : transformer l'information entre représentations comprimées et décomprimées. La frontière implicite-explicite est votre singularité personnelle. Vous portez un horizon des événements à l'intérieur de votre crâne.

\section*{Pourquoi cela doit exister}

Vous pourriez penser que cette correspondance structurelle est une coïncidence --- un joli motif autour duquel j'ai tracé des lignes, comme on voit des constellations dans des étoiles disposées au hasard. Laissez-moi donc vous montrer pourquoi ce motif n'a rien de facultatif. Pourquoi, si cinq hypothèses indépendamment raisonnables tiennent, cette architecture est la \textit{seule} qui fonctionne.

Voici les cinq axiomes. Je vais les énoncer aussi simplement que possible, car chacun, pris isolément, est difficile à contester. La controverse réside dans ce qu'ils produisent ensemble.

\textbf{Un : quelque chose existe.} C'est, j'espère que vous en conviendrez, l'affirmation la moins controversée qu'un livre puisse faire. Le pur néant est une abstraction platonicienne --- un concept, non un état de choses possible. Vous êtes en train de lire cette phrase. Quelque chose existe. Nous partirons de là.

\textbf{Deux : tout ce qui existe a un caractère dynamique.} Des choses se produisent. Le temps passe. Les états évoluent. Si ce qui existe n'avait aucune dynamique (aucun changement, aucune évolution, aucun calcul), il serait indiscernable du néant. (C'est l'identité des indiscernables de Leibniz, rencontrée au chapitre 13 : si deux choses sont identiques en toutes leurs propriétés, elles sont une seule et même chose. Une chose dont la dynamique est nulle n'a aucune propriété discernable. C'est le néant sous un masque.)

\textbf{Trois : la dynamique doit être stable et auto-entretenue.} Un système qui ne peut pas se maintenir n'est pas un système --- c'est une fluctuation momentanée. La criticalité auto-organisée, le mécanisme qui maintient les tas de sable, les cerveaux et (c'est ma thèse) l'univers au bord du chaos, est la seule manière connue pour un système dynamique complexe de se maintenir sans réglage externe. La Classe 4 est l'unique classe computationnelle qui s'auto-organise, permet le calcul universel et contient toutes les classes inférieures comme sous-processus. C'est la seule classe capable à la fois de se maintenir et de faire des choses intéressantes.

\textbf{Quatre : l'information a une borne finie à toute échelle.} Rien ne peut porter une information infinie dans un espace fini. La borne de Bekenstein est un théorème, pas une conjecture --- elle découle de la relativité générale et de la mécanique quantique. À chaque échelle, il existe une densité maximale d'information, et ce maximum est proportionnel à la surface, non au volume.

\textbf{Cinq : l'information est encodée holographiquement.} La frontière d'une région encode toute l'information de son intérieur sur une surface d'une dimension de moins. La physique tridimensionnelle est encodée sur des surfaces bidimensionnelles. C'est le principe holographique --- proposé par 't Hooft, développé par Susskind, étayé par la correspondance AdS/CFT de Maldacena, qui est ce que nous avons de plus proche d'un exemple démontré d'holographie.

Chaque axiome est motivé indépendamment. Aucun ne dépend des autres. Aucun n'exige que ma théorie soit correcte. Ils viennent de différents recoins de la physique et de la philosophie, développés par des chercheurs qui n'avaient jamais entendu parler de la Théorie des Quatre Modèles et qui ne s'en seraient guère souciés autrement.

Maintenant, combinez-les.

Des axiomes un et deux : quelque chose de dynamique existe. De l'axiome trois : cette dynamique est de Classe 4, parce que la Classe 4 est la seule classe universelle capable de se maintenir elle-même. De l'axiome quatre : le système est borné par des horizons d'information à toute échelle --- la structure de singularité. De l'axiome cinq : ces frontières encodent l'intérieur --- architecture holographique.

Assemblez le tout, et vous obtenez un automate holographique de Classe 4, borné par des surfaces de singularité à toute échelle. Un système dont l'information comprimée siège sur les frontières et dont l'intérieur décomprimé est le monde observable. Un système qui calcule sa propre structure, parce qu'un système holographique de Classe 4 à sortie holographique est un point fixe --- l'entrée, le processus et la sortie sont une seule et même chose.

Vous obtenez le SB-HC4A.

Vous obtenez, autrement dit, l'univers tel que nous l'observons. Et l'architecture de cet univers est la même architecture que celle de la conscience.

Ce n'est pas « j'ai trouvé un joli motif ». C'est « ce motif est le seul qui soit compatible avec les cinq axiomes simultanément ». Retirez n'importe lequel des axiomes, et l'unicité s'effondre. Sans l'axiome un, rien n'a besoin d'exister. Sans l'axiome deux, ce qui existe peut être statique. Sans l'axiome trois, toute classe computationnelle est possible. Sans l'axiome quatre, il n'y a pas de frontières d'information. Sans l'axiome cinq, il n'y a pas de structure holographique. Chaque axiome contraint l'espace des architectures possibles. Ensemble, ils le réduisent à exactement une.

\section*{La seule objection à laquelle je ne peux pas répondre}

Je vous ai promis les endroits où tout cela pourrait se briser ; les voici. Quatre d'entre eux sont du genre qu'un ingénieur dresserait en liste : l'énergie pourrait n'être que corrélée à l'information plutôt qu'identique à elle, auquel cas le mécanisme E = I échouerait ; l'argument de la Classe 4 est abductif, pas une preuve, et une Classe 4,5 que je suis incapable d'imaginer pourrait être tapie quelque part ; l'affirmation selon laquelle toutes les singularités sont une seule et même chose exige une théorie de la gravité quantique que nous n'avons pas encore ; et le modèle entier pourrait être infalsifiable de l'intérieur, en vertu de sa propre prédiction gödélienne. Chacun de ces points est réel. Je les expose tous les quatre, avec les contre-arguments, dans l'annexe G.

Mais il y en a un cinquième, et il est pire que les quatre autres réunis. Je l'ai soulevé il y a un chapitre, et il ne m'a pas lâché depuis : le plafond cognitif. Mon cerveau est un système de Classe 4 qui trouve de la structure auto-similaire partout où il regarde, parce que c'est ce que font les systèmes de Classe 4. Alors, quand il regarde le cosmos et y voit de la Classe 4 partout --- criticalité, frontières holographiques, clôture auto-référentielle ---, découvre-t-il l'architecture de la réalité, ou projette-t-il la sienne ? Un poisson doté d'une théorie de la physique conclurait que l'univers est fondamentalement mouillé. Nous sommes des détecteurs de symétries, et le SB-HC4A est, au fond, une affirmation de symétrie. Je ne sais pas comment distinguer les deux de l'intérieur, et je ne crois pas que quiconque le puisse.

Et voici ce qu'il y a d'étrange dans l'affaire. La théorie prédit exactement cette tache aveugle : un système auto-référentiel qui se calcule lui-même ne peut pas, en vertu de Gödel, atteindre de l'intérieur toutes les vérités sur lui-même. Ainsi, la question la plus profonde que ce livre puisse poser --- le miroir entre l'esprit et le cosmos est-il une découverte, ou un reflet des limites de la cognition humaine ? --- est une question à laquelle la théorie répond par une porte verrouillée, avant de vous tendre le plan de la serrure. Cette lacune n'est pas de l'ignorance. C'est de la géométrie. Que vous trouviez cela profond ou exaspérant en dit probablement long sur votre tempérament ; pour ma part, je trouve que c'est les deux à la fois. Et puis je continue quand même --- un mur que l'on peut nommer et sur lequel on peut poser la main, ce n'est pas la même chose qu'être perdu dans le noir.

\section*{La boucle est bouclée}

Laissez-moi maintenant rassembler tout cela.

Vous avez ouvert ce livre parce que vous vouliez savoir ce qu'est la conscience. La réponse, pour autant que je puisse la déterminer, est celle-ci : la conscience est une simulation auto-référentielle qui tourne au bord du chaos, bornée par une barrière opaque à l'information, et dans laquelle la simulation inclut un modèle d'elle-même. Quatre modèles, un versant réel et un versant virtuel, l'expérience vivant exclusivement sur le versant virtuel. Les \textit{qualia} sont des propriétés réelles de la simulation ; le mystère se dissout dès qu'on reconnaît que le Problème Difficile avait été posé au mauvais niveau. Neuf prédictions, plusieurs confirmées, aucune falsifiée. C'était la première moitié du tableau.

La seconde moitié est ce que vous venez de lire. La même architecture --- la même clôture auto-référentielle, les mêmes frontières d'information, le même encodage holographique, la même dynamique de Classe 4 --- semble être l'architecture de l'univers lui-même. Non par analogie. Par identité structurelle. La simulation que vous appelez « Je » tourne sur le même motif que la simulation que nous appelons « l'univers ». Vous n'êtes pas dans l'univers comme une bille dans une boîte. Vous êtes l'univers en train de faire à l'échelle biologique ce qu'il fait à toutes les échelles : se calculer lui-même, se modéliser lui-même, faire l'expérience de lui-même.

Le titre de ce livre est \textit{La simulation que vous appelez « Je »}. Vous savez désormais sur quoi tourne la simulation. Pas sur un ordinateur. Pas sur un cerveau. Pas même sur l'univers. Sur quelque chose de plus fondamental : le motif que tous trois partagent. Un automate holographique de Classe 4, borné par des barrières informationnelles, qui calcule sa propre existence.

Et si cela ressemble à une affirmation mystique --- ce n'en est pas une. C'est une affirmation structurelle. Testable dans certaines limites. Falsifiable, avec les réserves que j'ai exposées. Assez précise pour être fausse. Ce qui, comme vous le dira n'importe quel scientifique, est le plus beau compliment qu'une théorie puisse recevoir.

J'ai ouvert ce livre par une confession : en 2015, j'ai publié un livre de 300 pages sur la conscience qui s'est vendu à zéro exemplaire. Voici la partie de cette confession que j'avais passée sous silence. J'avais déjà le reste. La cosmologie --- l'univers holographique, les frontières de singularité, l'identité structurelle entre l'esprit et le cosmos que vous venez de lire --- était dans ma tête depuis le début des années 2000, à peu près depuis l'époque de l'épiphanie sur ce pont. J'ai délibérément laissé presque tout cela hors du livre de 2015. Une théorie de la conscience est déjà assez difficile à vendre ; une théorie de la conscience qui prétend en plus être une théorie de l'univers, c'est le moyen le plus sûr de se retrouver classé parmi les \textit{illuminés} avant que quiconque n'atteigne la page deux. J'ai donc laissé la cosmologie n'affleurer que par allusions --- un passage sur la borne holographique de 't Hooft, une remarque en passant se demandant si le cosmos ne serait pas un grand automate cellulaire --- et je n'ai pas osé écrire le reste. Il a fallu encore une décennie, et le cadeau improbable d'un modèle de langage assez patient pour écouter pendant que je pensais à voix haute, avant que j'ose enfin coucher le tout par écrit. Les quatre modèles n'étaient pas seulement une théorie de la conscience. Ils étaient un fragment de l'architecture de l'univers, visible à une échelle, invisible aux autres tant qu'on ne sait pas où regarder.

Vous avez maintenant le tout. Ou du moins tout ce qu'un cerveau de Classe 4 peut en voir depuis l'intérieur d'un univers de Classe 4. Quant à savoir s'il existe quelque chose au-delà --- si le miroir a une face arrière que nous n'atteindrons jamais ---, c'est la question à laquelle, nous dit la théorie, nous ne pouvons pas répondre.

Faites-en ce que vous voudrez.


\chapter*{Coda}
\addcontentsline{toc}{chapter}{Coda}
\markboth{Coda}{}

J'ai développé une théorie de la conscience vers 2005. Je l'ai publiée en 2015. Personne ne l'a lue. Deux décennies après l'intuition initiale, il s'est avéré que ce n'était là qu'une moitié de théorie --- l'autre moitié, c'était la cosmologie, aussi improbable que cela paraisse. Vous avez maintenant lu le tableau complet, ou du moins tout ce qu'un livre peut en contenir.

Mais il y a un morceau que j'ai laissé de côté. Non parce qu'il est spéculatif --- tout, dans les deux derniers chapitres, est spéculatif. Parce qu'il est personnel, et que le personnel est plus difficile.

Vous avez vu ce que fait l'ESM quand les choses tournent mal. Amnésie, AVC, salvia, split-brain, délire de Cotard --- il continue de tourner. Il ne plante jamais. Il construit un soi à partir de ce qui est disponible, et il croit à ce soi complètement. Un constructeur compulsif d'identité. Voilà ce qu'il fait. Voilà tout ce qu'il fait.

Les gens entendent parler des cas d'amnésie et disent : même cerveau, même corps, donc bien sûr que le soi persiste. Sauf que --- j'ai des décennies de plus aujourd'hui. Je n'ai presque rien en commun avec le moi que j'étais à un an. Un corps différent, des cellules remplacées plusieurs fois. Un cerveau différent. Des connexions synaptiques complètement différentes. D'autres souvenirs, une autre personnalité, tout est différent. Et pourtant l'ESM dit : toujours moi. J'ai déjà été, plusieurs fois au cours d'une seule vie, une personne complètement différente, et le constructeur n'a jamais bronché. Vous n'avez pas « le même cerveau ». Vous ne l'avez jamais eu.

J'ai frôlé la mort. Dans une avalanche --- le service militaire, la décision téméraire d'un officier, quatorze d'entre nous presque engloutis. J'étais certain que j'allais mourir. J'ai vu ma vie entière d'un seul coup --- le substrat déversant tout dans la simulation. Et cela ne m'a pas troublé autant qu'on pourrait s'y attendre. L'ESM, face à sa fin, ne paniquait pas à l'idée de perdre son identité. Il faisait son travail jusqu'au bout, ramenant à la surface tout ce qu'il avait.

Une autre fois, j'ai été violemment assommé. Tout est devenu noir. Quand je suis revenu à moi, je ne savais pas qui j'étais --- l'ESM redémarrait de zéro, comme celui d'un nouveau-né. La perte d'identité n'était pas le plus effrayant. Rester allongé au sol, paralysé pendant quelques secondes --- \textit{ça}, c'était terrifiant. Pas « qui suis-je ? », mais « mon corps va-t-il bien ? » La première priorité de l'ESM était l'intégrité du substrat. Qui j'étais est venu ensuite, presque après coup. Le modèle du soi existe pour servir le substrat, pas l'inverse.

Et puis il y a eu la fois où j'étais une fractale animée en quatre dimensions. Je ne m'attarderai pas sur les circonstances. Ce qui me dérangeait, ce n'était pas d'être une fractale --- cela m'était égal. Ce qui me dérangeait, c'était que ses mouvements entraient en conflit avec mon sens proprioceptif. Je sentais mon corps faire une chose pendant que la fractale en faisait une autre. Le conflit sensoriel était pénible. L'absurdité ontologique, non. Peu importe à l'ESM \textit{ce qu'}il modélise. Ce qui lui importe, c'est que les signaux soient cohérents.

Trois expériences. Une seule architecture. L'ESM de l'avalanche qui construit jusqu'au bout. L'ESM du K.-O. qui donne la priorité à l'intégrité du substrat sur l'identité narrative. L'ESM de la fractale qui se soucie de la cohérence sensorielle, pas de la plausibilité ontologique.

Voici à quoi je pense quand je pense à mourir pour de bon. Pas à la perspective elle-même --- la mort, à mon sens, est soit le repos éternel bien mérité, soit le voyage ultime. Non, ce à quoi je pense, c'est la logique.

Si l'ESM s'amorce à partir de rien --- et il le fait, chez chaque nouveau-né ---, alors « rien » n'est pas un état terminal pour le processus. C'est un état de départ. La configuration particulière que j'appelle « moi » prendra fin. Mes souvenirs, ma personnalité, ma façon d'être agacé par les gens qui confondent corrélation et causalité, ma façon d'agacer les gens avec des théories encombrantes --- disparus. Mais le processus, lui, ne m'appartient pas ; je ne l'emporterai pas avec moi. C'est une propriété de l'architecture. Il tournait avant ma naissance. Il tournera après ma mort.

Je ne décris pas une vie après la mort. Vos souvenirs ne seront pas transférés. Le « vous » particulier qui lit cette phrase prendra fin.

Mais le \textit{processus} --- la construction compulsive d'un soi à partir des données disponibles --- est universel. Chacune de ses instances se sent, de l'intérieur, exactement aussi réelle que la vôtre en ce moment même.

Et si les chapitres de cosmologie voient juste --- si l'univers est réellement un système de Classe 4 auto-similaire et quasi infini ---, alors il reste un coup à jouer. Dans un univers auto-similaire, des configurations semblables à « vous » existent ailleurs. Pas vous, pas vos souvenirs, mais un substrat doté de la bonne architecture, et un ESM qui amorcera un quelqu'un qui se sentira exactement aussi réel que vous en ce moment. Quelque chose comme l'immortalité quantique, sauf que cela n'exige pas la mécanique quantique. Juste un univers assez grand et un processus assez générique. La pensée la plus spéculative de ce livre. Mais elle découle de l'architecture.

Je n'ai pas conçu cette théorie pour qu'elle soit réconfortante. Mais elle a une conséquence pratique qu'il vaut la peine de retenir de ce livre.

Si chaque être conscient est le même constructeur tournant sur un matériel différent avec des données d'entraînement différentes, alors les frontières entre nous sont moins fondamentales qu'il n'y paraît. L'architecture est la même en chacun de nous.

Soyez gentil avec tout le monde. Ils pourraient tous être vous.


\chapter*{Remerciements}
\addcontentsline{toc}{chapter}{Remerciements}
\markboth{Remerciements}{}

Ce livre a été écrit avec l'assistance de Claude (Anthropic), qui a servi d'outil d'écriture, de relecture et de contre-vérification par IA tout au long du processus. La théorie, les arguments et l'ensemble du contenu intellectuel sont entièrement de moi, développés pendant deux décennies avant qu'aucun outil d'IA n'existe.

À mon oncle, Bruno J. Gruber, dont le travail en physique théorique (mécanique quantique et symétries) me montre ce que peut être un travail intellectuel à la fois rigoureux et joyeux. Son influence sur ma pensée est incalculable.

À mon oncle, Norbert Gruber, l'un des premiers professionnels de l'informatique du Rheintal autrichien, qui m'a offert mon premier PC. Sans ce cadeau, rien de tout cela n'aurait été possible. Il nous a quittés depuis, mais son empreinte perdure dans chaque ligne de code que j'ai écrite et dans chaque théorie que j'ai bâtie.

À ma famille, qui a supporté des années de conversations à table sur les qualia, la criticalité et les modèles du soi virtuels.

Et si vous songez maintenant à lire \textit{Die Emergenz des Bewusstseins} --- n'en faites rien. Je vous recommanderais des parasites cérébraux plutôt que ce monstre pataud, jamais passé entre les mains d'un éditeur. Attendez plutôt la traduction allemande du livre que vous tenez entre les mains. À ceux qui l'\textit{ont} déjà subi de bout en bout : \textit{mein Beileid}. Cela a dû être une torture. Vous avez toute ma compassion, et ma gratitude.


\chapter*{Notes et références}
\addcontentsline{toc}{chapter}{Notes et références}
\markboth{Notes et références}{}

\textit{Les références complètes, avec URL et annotations, figurent dans l'article scientifique ainsi qu'à l'adresse github.com/JeltzProstetnic/aIware/blob/main/references.md. Suivent des notes propres à chaque chapitre, destinées aux lecteurs qui souhaitent aller plus loin.}

\textbf{Chapitre 1} : Chalmers (1995), « Facing Up to the Problem of Consciousness », est la formulation fondatrice du Problème Difficile. Les résultats de COGITATE ont été publiés dans Nature (2025). La controverse sur l'IIT comme pseudoscience est documentée dans une lettre ouverte publiée sur PsyArXiv (septembre 2023).

\textbf{Chapitre 2} : L'architecture des quatre modèles a été publiée à l'origine dans Gruber (2015), \textit{Die Emergenz des Bewusstseins}. La théorie du modèle du soi de Metzinger (2003, 2009) et le modèle des versions multiples de Dennett (1991) en sont les principaux antécédents théoriques.

\textbf{Chapitre 3} : La formule de l'« hallucination contrôlée » vient de Seth (2021), \textit{Being You}. L'analogie du jeu vidéo m'est propre, mais elle fait écho à des motifs de l'« Ego Tunnel » de Metzinger (2009). Sur l'illusion de la main en caoutchouc : Botvinick \& Cohen (1998), « Rubber hands 'feel' touch that eyes see », \textit{Nature}.

\textbf{Chapitre 4} : La dissolution du Problème Difficile par les qualia virtuels est originale à Gruber (2015) et a été affinée par la confrontation adverse en 2026. L'argument de la clôture auto-référentielle a été développé en réponse à l'objection de circularité. La distinction d'avec l'illusionnisme (Frankish 2016 ; Dennett 1991) est cruciale : la théorie soutient que les qualia sont réels \textit{au sein} de la simulation, non pas illusoires. Le méta-problème de la conscience (Chalmers 2018) se dissout par l'inaccessibilité structurelle de l'ISM à l'ESM. La formule de Joscha Bach --- la conscience comme « une simulation, une réalité virtuelle projetée dans un observateur virtuel » (Bach, publication sur X / @Plinz, 12 juin 2026) --- est citée comme une convergence indépendante vers la thèse du soi comme simulation, la Théorie des Quatre Modèles fournissant la spécification architecturale (la typologie des modèles, la clôture auto-référentielle comme critère et l'exigence de criticalité).

\textbf{Chapitre 5} : Wolfram (2002), \textit{A New Kind of Science}. Beggs \& Plenz (2003) sur les avalanches neuronales. Carhart-Harris et al. (2014) sur l'hypothèse du cerveau entropique. La synthèse de 2022 : « Self-organized criticality as a framework for consciousness ». Hengen \& Shew (2025) sur la méta-analyse de 140 jeux de données. Le cadre ConCrit : Algom \& Shriki (2026). L'argument des deux seuils (criticalité + architecture) est original à cette théorie.

\textbf{Chapitre 6} : Klüver (1966) sur les constantes de forme. Carhart-Harris et al. (2012, 2016) sur la neuro-imagerie psychédélique. La phénoménologie de Salvia divinorum s'appuie sur des récits d'expérience publiés et sur la littérature pharmacologique consacrée à la salvinorine A. Le mécanisme prédictif-rétroactif de l'anosognosie est discuté dans Gruber (2015) ; l'exemple des applaudissements est une observation clinique courante. L'expérience de pensée d'une administration permanente de salvinorine A est originale à Gruber (2015).

\textbf{Chapitre 7} : Casali et al. (2013) sur le PCI. Alkire et al. (2000) sur le propofol. Schartner et al. (2017) sur l'entropie sous kétamine. La prédiction d'une complexité EEG accrue lors du rêve lucide est originale à cette théorie.

\textbf{Chapitre 8} : Owen et al. (2006) sur la conscience cachée chez les patients en état végétatif. Syndrome d'Anton : Goldenberg et al. (1995). Le parcours d'obstacles en vision aveugle : de Gelder et al. (2008). Délire de Cotard : Young \& Leafhead (1996). Syndrome de la main étrangère : Della Sala et al. (1991) ; la référence à Dr Folamour renvoie à Kubrick (1964). La distinction entre syndrome de la main anarchique et main étrangère : Marchetti \& Della Sala (1998). Syndrome de Charles Bonnet : Teunisse et al. (1996). Le déjà-vu comme appariement à une mémoire-modèle est original à Gruber (2015). TCC et plasticité neuronale : DeRubeis et al. (2008). Placebo et opioïdes endogènes : Benedetti et al. (2005). L'interprétation du trouble de conversion comme une vision aveugle inversée est originale à cette théorie.

\textbf{Chapitre 9} : Gazzaniga, Bogen \& Sperry (1962, 1965). Gazzaniga (2000) sur l'interprète de l'hémisphère gauche. Les exemples de conflit interhémisphérique (boutonner/déboutonner, saisie de la main) sont documentés chez Akelaitis (1945) et Bogen (1993). Nagel (1971), « Brain Bisection and the Unity of Consciousness ». Parfit (1984) sur l'identité personnelle. Pinto et al. (2017) sur le réexamen des phénomènes de cerveau divisé. Lashley (1950) sur la mémoire distribuée et l'équipotentialité. Le TDI comme bifurcation de modèle virtuel : la théorie prédit des schémas d'activité neuronale distincts pour chaque personnalité, ce qui est cohérent avec Reinders et al. (2003, 2006). TDI et traumatisme infantile : Putnam (1997) ; la fenêtre développementale de la bifurcation est originale à cette théorie.

\textbf{Chapitre 10} : Güntürkün \& Bugnyar (2016) sur la cognition aviaire sans cortex. Kanzi le bonobo : Savage-Rumbaugh \& Lewin (1994), \textit{Kanzi: The Ape at the Brink of the Human Mind}. L'effet Baldwin : Baldwin (1896), « A New Factor in Evolution ». Nagel (1974), « What Is It Like to Be a Bat? »

\textbf{Chapitre 11} : Les neuf prédictions sont toutes développées formellement dans l'article scientifique. Pour le traitement le plus approfondi de la neuroanatomie fonctionnelle dans le contexte de la conscience, voir Christof Koch, \textit{The Quest for Consciousness: A Neurobiological Approach} (2004) --- l'exposé définitif du programme Crick-Koch, dans lequel Francis Crick et Koch ont parcouru systématiquement le système visuel étape par étape, à la recherche des corrélats neuronaux de la conscience. Leur quête, à mon sens, cherchait la conscience au mauvais endroit (le substrat plutôt que la simulation), mais le travail neuroanatomique de fond qu'ils ont accompli reste inégalé.

\textbf{Chapitre 12} : Butlin et al. (2023, 2025) sur les indicateurs de conscience de l'IA. Seth (2025) sur le naturalisme biologique et la conscience de l'IA. La hiérarchie à cinq niveaux de fidélité de scan découle de l'architecture des quatre modèles du chapitre 2. Le problème de la copie s'appuie sur Parfit (1984), \textit{Reasons and Persons}, et sur Nozick (1981) concernant l'identité personnelle et les continuateurs les plus proches. L'expérience de pensée du remplacement graduel est une variante du bateau de Thésée, formalisée pour les systèmes neuronaux. Calcul neuromorphique : Schuman et al. (2017) pour un panorama du matériel neuromorphique ; Intel Loihi et IBM TrueNorth comme implémentations actuelles. Connectome de \textit{C. elegans} : White et al. (1986). Le transfert de substrat, la quasi-immortalité et l'implication d'une téléportation interstellaire sont originaux à Gruber (2015). La réserve liée au malaise --- le fait que la perte du substrat biologique altérerait profondément la qualité phénoménale --- s'appuie sur la littérature de l'intéroception : Craig (2009), \textit{How Do You Feel?}, et Seth \& Friston (2016) sur l'inférence intéroceptive active.

\textbf{Chapitre 13} : Libet (1979, 1983, 1985) et Schurger et al. (2012) sur le libre arbitre. Kühn \& Brass (2009) sur la construction rétrospective du jugement de choix libre. Wegner (2002, 2003), \textit{The Illusion of Conscious Will} --- l'expérience de la souris « I Spy » y est décrite en détail. L'expérience de pensée du café/sucre, l'argument de l'amnésie-révèle-le-déterminisme et l'argument de la séquence de nombres aléatoires sont originaux à Gruber (2015). Le cadre de traitement à 40/20 Hz, la réinterprétation de Libet « aucune rétrodatation nécessaire » et l'exemple de la fréquence dans les arts martiaux sont originaux à Gruber (2015). L'analogie de l'horloge pour l'épiphénoménalisme, la reformulation « la volonté est réelle mais partiellement connue » et le modèle de connaissance de soi des « trois divergences » sont également originaux à Gruber (2015). L'anecdote personnelle sur l'audition de « voix » intérieures lors d'un épuisement extrême est autobiographique. L'argument du zombie est traité via Kirk (2019) et Chalmers (1996). La chambre de Mary : Jackson (1982, 1986). La section sur les questions ouvertes suit l'approche des limites assumées honnêtement recommandée par Popper (1963).

\textbf{Chapitre 15} : Wolfram (2002), \textit{A New Kind of Science}, pour les quatre classes de calcul --- que ce livre étend à cinq --- et l'irréductibilité computationnelle. Bak (1987, 1996) sur la criticalité auto-organisée. La découverte de l'expansion accélérée en 1998 : Riess et al. (1998), Perlmutter et al. (1999) --- prix Nobel 2011. L'horizon cosmologique : Rindler (1956). La longueur de Planck et la physique à l'échelle de Planck : Planck (1899) ; traitements modernes chez Garay (1995). Borne de Bekenstein : Bekenstein (1981). Principe holographique : 't Hooft (1993), Susskind (1995). Les germes de l'argument cosmologique --- l'univers comme automate cellulaire et l'application de la borne holographique de 't Hooft au calcul à l'échelle de l'univers --- étaient déjà présents dans Gruber (2015), sans y être développés en un modèle complet. L'identification de toutes les singularités comme instances d'un unique phénomène structurel est développée dans Gruber (2026).

\textbf{Chapitre 16} : Le scénario du Big Rip : Caldwell, Kamionkowski \& Weinberg (2003). L'architecture SB-HC4A, la formulation en point fixe et la conséquence d'incomplétude gödelienne pour les systèmes auto-calculants représentent la formulation aboutie, développée dans Gruber (2026). La prédiction des particules comme atomes de calcul, la dérivation des lois de conservation comme contraintes informationnelles et la conjecture des trois générations sont originales à Gruber (2026). L'objection du plafond cognitif est originale à ce travail.

\textbf{Chapitre 17} : La correspondance structurelle à cinq termes entre le SB-HC4A et l'architecture des quatre modèles de la conscience est originale à Gruber (2026). Principe de Landauer : Landauer (1961) ; confirmation expérimentale : Bérut et al. (2012). Borne de Bekenstein : Bekenstein (1981). Thermodynamique des trous noirs : Bekenstein (1973), Hawking (1975). L'hypothèse E = I (identité énergie-information) est discutée chez Vedral (2010), \textit{Decoding Reality}, et Davies (2010). La dérivation du SB-HC4A à partir de cinq axiomes est originale à Gruber (2026). Maldacena (1998) sur la correspondance AdS/CFT. Les cinq points faibles --- dont les objections d'infalsifiabilité par la structure et du plafond cognitif --- sont originaux à ce travail.

\textbf{Annexe B} : Le modèle d'intelligence récursive (la structure multiplicative savoir × performance × motivation) est développé dans Gruber (2015) ainsi que dans l'article sur le Recursive Intelligence Model (Gruber 2026). Deux convergences indépendantes issues de la psychométrie sur la motivation comme constitutive de l'intelligence : Wittmann \& Süß (1999), « Investigating the paths between working memory, intelligence, knowledge, and complex problem-solving performances via Brunswik symmetry » (dans Ackerman, Kyllonen \& Roberts, dir., \textit{Learning and Individual Differences: Process, Trait, and Content Determinants}, American Psychological Association), selon laquelle la motivation prédit la \textit{variabilité} de la performance une fois corrigées les violations de la symétrie de Brunswik ; et l'étude COGITO --- Brose, Schmiedek, Lövdén, Molenaar \& Lindenberger (2010), \textit{Research in Human Development} 7(1), 61--78, sur le couplage intra-individuel entre motivation quotidienne et performance de la mémoire de travail, et Schmiedek, Lövdén, von Oertzen \& Lindenberger (2020), \textit{PeerJ} 8:e9290, montrant que l'intelligence générale (g) est largement un artefact inter-individuel plutôt que la structure de la cognition intra-individuelle.


\chapter*{Annexe A : Neurologie de base --- Un guide de référence}
\addcontentsline{toc}{chapter}{Annexe A : Neurologie de base --- Un guide de référence}
\markboth{Annexe A : Neurologie de base --- Un guide de référence}{}

\textit{Cette annexe propose de brèves explications des termes neuroscientifiques employés tout au long du livre. Consultez-la chaque fois que vous rencontrez un terme qui vous est inconnu. Les entrées sont classées par ordre alphabétique.}

\begin{itemize}
  \item \textbf{Potentiel d'action} --- Le signal électrique qui parcourt l'axone d'un neurone.
  \item \textbf{Amygdale} --- Structure cérébrale impliquée dans le traitement des émotions, en particulier la peur.
  \item \textbf{Anosognosie} --- Absence de conscience de ses propres déficits neurologiques. (Voir chapitre 6.)
  \item \textbf{Axone} --- La longue fibre de sortie d'un neurone, qui transporte les signaux vers d'autres neurones.
  \item \textbf{Aires de Brodmann} --- Régions numérotées du cortex, cartographiées par l'anatomiste Korbinian Brodmann d'après la structure cellulaire. V1 = aire 17 de Brodmann.
  \item \textbf{Corps calleux} --- L'imposant faisceau de fibres qui relie les deux hémisphères du cerveau.
  \item \textbf{Colonnes corticales} --- Modules verticaux de neurones dans le cortex, d'environ 0,5 mm de diamètre, considérés comme les unités de traitement fondamentales.
  \item \textbf{EEG (électroencéphalographie)} --- Technique qui mesure l'activité électrique à la surface du cuir chevelu.
  \item \textbf{IRMf (imagerie par résonance magnétique fonctionnelle)} --- Technique d'imagerie cérébrale qui détecte les variations du débit sanguin associées à l'activité neuronale.
  \item \textbf{GABA} --- Le principal neurotransmetteur inhibiteur du cerveau.
  \item \textbf{Hippocampe} --- Structure cérébrale essentielle à la formation de nouveaux souvenirs.
  \item \textbf{Intéroception} --- Le sens de l'état interne du corps (battements cardiaques, digestion, température).
  \item \textbf{Récepteurs opioïdes kappa} --- Type de récepteur ciblé par la salvinorine A (salvia divinorum).
  \item \textbf{Néocortex} --- La surface externe du cerveau, à six couches, responsable des fonctions supérieures.
  \item \textbf{Neurotransmetteur} --- Messager chimique libéré au niveau des synapses (par ex. sérotonine, dopamine, GABA, glutamate).
  \item \textbf{PCI (Perturbational Complexity Index)} --- Une mesure de la complexité cérébrale mise au point par Massimini. Utilisée pour évaluer le niveau de conscience.
  \item \textbf{Proprioception} --- Le sens de la position et du mouvement du corps dans l'espace.
  \item \textbf{Pulvinar} --- Un noyau du thalamus impliqué dans l'attention visuelle et dans la voie visuelle sous-corticale.
  \item \textbf{Récepteurs sérotoninergiques 2A} --- Le type de récepteur ciblé par les psychédéliques classiques (LSD, psilocybine).
  \item \textbf{Colliculus supérieur} --- Une structure du mésencéphale impliquée dans les mouvements oculaires et dans la voie visuelle rapide qui contourne le cortex.
  \item \textbf{Synapse} --- La jonction entre deux neurones où les signaux sont transmis.
  \item \textbf{Poids synaptiques} --- L'intensité des connexions entre neurones, modifiée par l'apprentissage.
  \item \textbf{Thalamus} --- La station de relais du cerveau, qui aiguille les informations sensorielles vers le cortex.
  \item \textbf{V4} --- Aire visuelle spécialisée dans la perception de la couleur, de la courbure et le traitement des textures complexes. Champs récepteurs d'environ 8-16°. Sous psychédéliques, l'activité de V4 produit des fractales colorées et des motifs kaléidoscopiques (chapitre 6).
  \item \textbf{V5/MT (aire temporale moyenne)} --- Aire visuelle spécialisée dans le traitement du mouvement. Grands champs récepteurs. Responsable de la rotation et du déplacement des motifs perçus sous psychédéliques (chapitre 6).
  \item \textbf{Cortex visuel} --- La région située à l'arrière du cerveau qui traite l'information visuelle, organisée en une hiérarchie du simple au complexe (V1 $\rightarrow$ V2 $\rightarrow$ V3 $\rightarrow$ V4 $\rightarrow$ V5 $\rightarrow$ IT).
\end{itemize}

\section*{La hiérarchie du traitement visuel (de V1 à IT)}

Le flux visuel ventral traite à chaque étape des caractéristiques de plus en plus complexes, de simples contours jusqu'à la reconnaissance complète des objets. Cette hiérarchie devient directement perceptible sous l'effet des psychédéliques, où chaque étape de traitement devient accessible à la conscience dans l'ordre (chapitre 6). Le tableau ci-dessous résume la fonction de chaque aire, la taille de son champ récepteur et la signature psychédélique caractéristique qui apparaît lorsque le traitement intermédiaire de cette aire s'infiltre dans la simulation consciente.


{\small
\noindent
\begin{tabularx}{\linewidth}{X X X X}
\toprule
\textbf{Aire} & \textbf{Champ récepteur} & \textbf{Fonction normale} & \textbf{Signature psychédélique} \\
\midrule
\textbf{V1} & ~1° & Détection des contours, fréquence spatiale, colonnes d'orientation & Phosphènes, constantes de forme de Klüver, surfaces qui respirent/miroitent \\[4pt]
\textbf{V2} & ~2-4° & Intégration des contours, segmentation des textures, appartenance des bords, contours illusoires & Tessellations, motifs géométriques répétitifs, perception accrue des textures \\[4pt]
\textbf{V3} & ~4-8° & Traitement de la forme globale, forme dynamique, frontières de mouvement & Structures géométriques fluides et changeantes \\[4pt]
\textbf{V4} & ~8-16° & Couleur, courbure, texture complexe, traitement à l'échelle fractale & Fractales colorées, motifs kaléidoscopiques, couleurs saturées/impossibles \\[4pt]
\textbf{V5/MT} & Grand, accordé au mouvement & Perception du mouvement, flux optique, codage de la vitesse et de la direction & Rotation, dérive et mouvement rythmique de tous les motifs visuels \\[4pt]
\textbf{Gyrus fusiforme (IT)} & Très grand, centré sur l'objet & Reconnaissance des visages (FFA), formes des mots, discrimination fine des objets & Visages, figures, entités --- souvent déformés ou changeants \\[4pt]
\textbf{IT antérieur} & Champ visuel entier & Catégories sémantiques, construction de scènes, reconnaissance des objets invariante au point de vue & Hallucinations narratives complètes, scènes complexes, séquences oniriques \\[4pt]
\bottomrule
\end{tabularx}
}


\textbf{Notes :}
\begin{itemize}
  \item Le gyrus fusiforme chevauche la frontière V4/IT et fait partie du cortex inférotemporal (IT). Il contient l'aire fusiforme des visages (FFA), identifiée par Kanwisher et al. (1997), qui est sélectivement activée par les visages.
  \item La taille des champs récepteurs augmente de façon spectaculaire de V1 (~1°) à IT (champ visuel entier), reflétant l'abstraction progressive des caractéristiques locales vers les objets et scènes globaux.
  \item Sous psychédéliques, la progression des effets de V1 vers IT dépend de la dose : les faibles doses touchent d'abord V1 ; les doses plus élevées recrutent progressivement des étapes plus profondes. Cette activation ordonnée est une prédiction directe du gradient de perméabilité de la Théorie des Quatre Modèles (chapitre 6).
  \item Le cerveau utilise aussi ces aires pour un traitement fractal ou invariant d'échelle (V2-V4), qui sert à la mesure des échelles et à l'analyse des textures dans la vision normale. Sous psychédéliques, cette machinerie fonctionnant sans entrée externe produit les motifs fractals caractéristiques (voir Annexe C).
\end{itemize}


\chapter*{Annexe B : Le modèle de l'intelligence}
\addcontentsline{toc}{chapter}{Annexe B : Le modèle de l'intelligence}
\markboth{Annexe B : Le modèle de l'intelligence}{}

\textit{Cette annexe résume le modèle récursif de l'intelligence développé dans un article compagnon (Gruber, 2026, « Why Intelligence Models Must Include Motivation »). Le traitement académique complet, avec références et arguments formels, est disponible séparément.}

Commençons par une histoire, car c'est l'illustration la plus limpide de la boucle que je connaisse. Je suis tombé amoureux des mathématiques vers huit ans --- pas du calcul, mais de la vraie discipline : l'algèbre, la géométrie, les structures sous les nombres. Mon père avait un diplôme de mathématiques, et j'ai potassé ses vieux manuels universitaires. C'était la fin des années 1980 ; Internet n'existait pas. Si vous vouliez apprendre quelque chose, il vous fallait un livre ou une personne, et à onze ans j'avais épuisé la collection de mon père. La faim n'a pas disparu ; l'approvisionnement, lui, s'est simplement tari. Je m'étais heurté à un mur qui ne devait rien à l'aptitude et tout aux circonstances. J'avais la motivation. J'avais la performance --- je pouvais suivre les mathématiques. Ce qui me manquait, c'était l'accès au niveau de connaissance suivant. La boucle récursive a calé, non parce qu'une composante était faible, mais parce que l'approvisionnement externe de l'une d'elles avait été coupé. La boucle a besoin de carburant venu de l'extérieur pour continuer d'itérer. Gardez cet exemple à l'esprit ; c'est le modèle tout entier en miniature. Je vais maintenant l'exposer en détail --- quelles sont les composantes, comment elles interagissent, pourquoi cette interaction produit les dynamiques que nous observons, et ce que cela implique pour l'éducation, l'intelligence artificielle et le lien avec la conscience.

\section*{La curieuse omission}

Tous les grands modèles de l'intelligence excluent formellement la motivation. La taxonomie de Cattell-Horn-Carroll (le cadre dominant de la recherche sur l'intelligence) est une hiérarchie d'aptitudes cognitives sans la moindre composante motivationnelle. La théorie de l'investissement de Cattell lui-même, qui proposait que l'intelligence fluide soit « investie » dans l'apprentissage pour produire l'intelligence cristallisée, requiert un investisseur (quelqu'un qui décide quoi apprendre et pourquoi), mais elle traite la motivation de cet investisseur comme une condition externe plutôt que comme une partie de l'intelligence elle-même. La théorie triarchique de Sternberg inclut l'intelligence pratique, mais pas l'élan qui pousse à l'acquérir. Les intelligences multiples de Gardner incluent la conscience intrapersonnelle, mais pas le moteur qui propulse le développement intellectuel.

David Wechsler (dont les échelles d'intelligence sont les plus administrées au monde) a explicitement appelé à inclure les facteurs motivationnels dès 1940. Le champ l'a ignoré. Les échelles de Wechsler modernes demeurent des instruments purement cognitifs.

Ce n'est pas une simplification anodine. C'est une tache aveugle systématique, qui déforme notre image de ce que l'intelligence est réellement et de la manière dont elle se développe réellement.

\section*{Les trois composantes}

L'intelligence, comprise comme \textit{capacité d'apprentissage}, est constituée de trois composantes en interaction :

\textbf{Le Savoir} est le contenu accumulé de l'apprentissage. Il se présente sous deux types profondément différents. Le \textit{savoir factuel} est le savoir des contenus : faits, concepts, procédures, répertoire culturel. C'est ce que les tests de QI mesurent principalement sous l'étiquette d'« intelligence cristallisée », et c'est ce que les systèmes scolaires transmettent principalement. Le \textit{savoir opérationnel} est le savoir sur la \textit{manière d'apprendre et de penser} : stratégies d'apprentissage, heuristiques de raisonnement, compétences métacognitives, outils logiques, planification stratégique, et la capacité d'évaluer sa propre compréhension. La distinction importe énormément, comme je l'expliquerai plus bas.

\textbf{La Performance} est la capacité de traitement du système cognitif : mémoire de travail, vitesse de traitement, la puissance de calcul brute du substrat neuronal. Cela correspond en gros à ce que les modèles psychométriques appellent « intelligence fluide ». C'est la composante la plus fortement influencée par la génétique et la neurobiologie. Elle culmine au début de l'âge adulte, puis décline graduellement.

\textbf{La Motivation} est l'élan soutenu qui pousse à s'engager avec le monde d'une manière qui produit de l'apprentissage. Elle comporte deux sous-composantes. La \textit{soif de savoir} est le désir intrinsèque de comprendre : la curiosité, le besoin de donner du sens aux choses. Le \textit{besoin d'agir} est l'élan d'appliquer le savoir, d'expérimenter, de s'engager activement dans l'environnement. Les deux relèvent en partie du tempérament inné, en partie de l'expérience.

\section*{La boucle récursive}

L'affirmation décisive est que ces trois composantes ne sont pas simplement additives --- elles forment une \textit{boucle récursive fermée} dans laquelle chaque composante amplifie les autres.

Le Savoir renforce la Performance : les stratégies d'apprentissage et les outils logiques améliorent directement l'efficacité du traitement cognitif. Un joueur d'échecs qui a appris des heuristiques évalue les positions plus vite qu'un autre qui s'en remet à la recherche par force brute. Un lecteur qui a appris le décodage phonémique traite le texte plus fluidement, libérant de la mémoire de travail pour la compréhension.

La Performance renforce le Savoir : une plus grande capacité cognitive permet un apprentissage plus rapide et plus profond. Une mémoire de travail plus étendue vous permet de garder davantage d'informations en mémoire simultanément, ce qui vous aide à repérer des connexions et à dégager des régularités.

La Motivation renforce à la fois le Savoir et la Performance : l'apprenant motivé recherche les occasions d'apprendre (ce qui étend le Savoir) et exerce ses compétences cognitives (ce qui entraîne la Performance). Point crucial : la motivation soutient l'engagement \textit{dans la durée}, ce qui est essentiel pour que la boucle continue d'itérer.

Et le Savoir et la Performance renforcent la Motivation : le succès dans l'apprentissage et la résolution de problèmes engendre un affect positif et un sentiment d'efficacité personnelle, qui entretiennent l'élan d'apprendre davantage. C'est le mécanisme derrière l'effet Matthieu --- les riches s'enrichissent. Le succès précoce nourrit la motivation qui produit le succès suivant.

Cette structure récursive produit une dynamique d'intérêts composés. De petites différences initiales dans n'importe quelle composante --- fût-ce dans la seule motivation --- se composent au fil du temps, produisant cette large variance des accomplissements intellectuels adultes que les modèles purement cognitifs peinent à expliquer. Une personne dotée d'une capacité de traitement cognitif moyenne, mais profondément motivée et possédant un solide savoir opérationnel, développera au long d'une vie des capacités intellectuelles très supérieures à celles d'une personne à la capacité de traitement supérieure, mais faiblement motivée et aux stratégies d'apprentissage médiocres.

Voyez les choses ainsi : dans les intérêts composés, le rythme des dépôts et la stratégie d'investissement comptent davantage que le capital initial. Dans la boucle de l'intelligence, la Motivation est le rythme des dépôts. Le Savoir opérationnel est la stratégie d'investissement. La Performance est le capital initial. Et la plupart des gens ont largement assez de capital.

\section*{Le savoir opérationnel : le multiplicateur caché}

Le savoir opérationnel mérite une attention particulière, car il occupe une position unique dans la boucle. Le savoir factuel est additif : apprendre un fait nouveau ajoute un fait au stock. Le savoir opérationnel est \textit{multiplicatif} : apprendre une nouvelle stratégie d'apprentissage améliore l'efficacité de tout apprentissage ultérieur.

Une étudiante qui apprend la répétition espacée (répartir la pratique dans le temps plutôt que de tout réviser d'un coup) n'acquiert pas simplement un fait nouveau. Elle acquiert un outil qui augmente le taux de rétention de tout ce qu'elle apprendra désormais. Un étudiant qui apprend à repérer ses propres lacunes et à les combler systématiquement ne corrige pas une seule lacune ; il acquiert une compétence qui prévient des centaines de lacunes futures. C'est un savoir qui accélère la boucle elle-même.

Si une seule composante mérite l'étiquette « ce qui rend les gens intelligents », c'est le savoir opérationnel. Pas le QI. Pas la puissance de traitement brute. La méta-compétence de savoir apprendre efficacement.

\section*{Pourquoi les tests de QI passent à côté de l'essentiel}

Les tests de QI mesurent la \textit{performance maximale} --- ce qu'une personne peut accomplir dans des conditions standardisées, en supposant un effort maximal. Ils capturent un instantané d'une seule composante (la Performance sur des tâches spécifiques) à un moment donné. Ils ne capturent pas (ils ne le peuvent pas) le processus récursif, auto-renforçant et multi-composantes qu'est réellement l'intelligence.

Voilà pourquoi les scores de QI en disent si peu sur la trajectoire intellectuelle à long terme. Deux enfants aux scores de QI identiques à six ans peuvent diverger de façon spectaculaire à trente ans --- l'un devenu chercheur scientifique, l'autre ayant cessé de lire après l'école. Les modèles psychométriques standard achoppent sur cette divergence. Le modèle récursif la prédit : ces enfants différaient non pas en Performance, mais en Motivation et en Savoir opérationnel, et la boucle récursive a amplifié ces différences au fil de vingt-quatre années d'itérations composées.

Le test de QI, c'est comme mesurer la puissance du moteur d'une voiture sans vérifier si elle a du carburant ou un conducteur. La puissance compte, mais elle n'est pas le goulot d'étranglement de la plupart des trajets.

\section*{Le cas test de l'IA}

Le modèle récursif fait une prédiction précise au sujet de l'intelligence artificielle : des systèmes dotés d'un Savoir élevé et d'une Performance élevée, mais dépourvus de Motivation, ne devraient pas manifester le développement auto-dirigé qui caractérise l'intelligence humaine. Et c'est exactement ce que nous observons.

Les grands modèles de langage actuels possèdent un savoir immense (entraînés sur des milliers de milliards de tokens), une performance élevée (des milliards de paramètres) --- et pas la moindre motivation. Ils traitent ce qu'on leur donne et produisent ce qu'on leur demande. Entre deux requêtes, ils ne font rien. Ils ne partent pas en quête de leurs zones d'ignorance. Ils n'exercent aucune compétence. Ils ne se posent de questions sur aucun problème. Leur « intelligence » est entièrement statique --- fixée par l'entraînement, sans la moindre impulsion endogène à l'étendre.

Même les modèles de raisonnement les plus avancés (capables de résoudre des problèmes de mathématiques de niveau compétition) présentent exactement ce mode de défaillance. Ils résolvent des problèmes extraordinaires \textit{quand on les y invite}, mais ne partent jamais d'eux-mêmes en quête de problèmes, ne dirigent pas leur propre apprentissage, et ont besoin d'un échafaudage externe qui fait office de substitut à la composante motivationnelle absente. Poussez la Performance et le Savoir aussi haut que vous voudrez : sans Motivation, la boucle ne s'auto-entretient pas.

Ce n'est pas simplement parce que ces systèmes n'ont pas été conçus pour s'auto-améliorer. Cette objection revient justement à nous donner raison : concevoir un système qui s'améliore lui-même exige de fabriquer un analogue fonctionnel de la motivation. Tant que les systèmes d'IA n'en disposeront pas, ils resteront des outils que l'on utilise, et non des agents qui se développent.

\section*{Le lien avec la conscience}

C'est ici que le modèle de l'intelligence rejoint la Théorie des Quatre Modèles qui est au cœur de ce livre. La boucle récursive de l'intelligence ne se contente pas de \textit{bénéficier de} la conscience --- elle la \textit{requiert}.

La boucle dépend d'une capacité cognitive précise : l'\textit{apprentissage cognitif} --- la capacité d'induire des théories générales à partir d'observations particulières, par opposition au simple apprentissage par renforcement (conditionnement stimulus-réponse). L'apprentissage par renforcement peut vous apprendre, par la douleur, à éviter une plaque brûlante. L'apprentissage cognitif vous permet de regarder quelqu'un d'autre toucher une plaque brûlante et de généraliser : « Ce qui est chaud brûle. On ne touche pas à ce qui est chaud. » La différence, c'est la capacité de simuler des scénarios à la troisième personne --- de vous modéliser vous-même comme un objet dans le monde et de raisonner sur ce qui arriverait si vous faisiez telle ou telle chose.

C'est précisément ce que fournissent le Modèle Explicite du Monde et le Modèle Explicite du Soi. La conscience --- la capacité de créer et de faire tourner une simulation de soi --- est le \textit{substrat} sur lequel opère la boucle récursive de l'intelligence. Sans modèles explicites, vous obtenez l'apprentissage par renforcement : cela fonctionne, mais sans effet cumulatif. Avec des modèles explicites, vous obtenez l'apprentissage cognitif, qui alimente la boucle récursive et se capitalise au fil d'une vie.

Voilà pourquoi le gradient d'intelligence animale du chapitre 10 se superpose au gradient de conscience. Des modèles du soi plus sophistiqués permettent des boucles récursives plus sophistiquées. Un chien, doté d'un modèle du soi relativement simple, fait tourner une version limitée de la boucle --- il peut apprendre par observation dans une certaine mesure, mais son apprentissage cognitif est bridé par la richesse de ses modèles explicites. Un chimpanzé, avec un modèle du soi plus riche, fait tourner une boucle plus puissante. Un être humain, avec l'architecture complète des quatre modèles, fait tourner la boucle à pleine capacité --- et les résultats en sont le langage, la culture, la science et tout ce qui distingue l'intelligence humaine de la cognition animale.

L'apprentissage cognitif est le premier lien entre l'intelligence et la conscience. Voici le second, et il est plus tranchant : la motivation. La boucle récursive ne s'auto-entretient que si quelque chose la maintient en itération. Le savoir, vous pouvez le faire croître à volonté ; la performance aussi. Mais sans motivation, la boucle cale --- et c'est exactement ce que l'IA d'aujourd'hui nous montre. Elle résout des problèmes sur commande et ne s'en pose jamais un seul d'elle-même.

Maintenant, la partie honnête, car la motivation a deux niveaux, et seul le niveau supérieur relève de l'histoire de la conscience. En dessous se trouve un plancher subconscient : l'optimisation brute de survie, le substrat qui pousse vers ses buts. C'est ce que j'ai appelé la part inconsciente de la \textit{volonté} au chapitre 13, et cela n'a rien de spécial --- tout optimiseur en dispose, jusqu'à la bactérie ou au thermostat. Par-dessus vient la couche consciemment cadrée : vouloir, être curieux, tenir à quelque chose, préférer un résultat à un autre. C'est le Modèle Explicite du Soi qui s'évalue lui-même, et \textit{c'est là} l'effet de conscience.

Une remarque sur les mots, car une sémantique floue rendrait tout l'argument mou. \textit{Volonté} et \textit{motivation} ne découpent proprement le gradient ni en français, ni en anglais, ni en allemand --- l'usage ordinaire laisse chacun des deux empiéter sur le territoire de l'autre. Ce qui est stable, c'est la tendance : on parle de \textit{volonté} quand l'attraction est spontanée et difficile à articuler ; de \textit{motivation} quand on peut nommer une raison et que l'élan est assez symbolisable pour qu'on en fasse un récit. Traitez-les donc ici comme les pôles d'un même continuum, non comme des espèces distinctes --- la \textit{volonté} pointant vers l'extrémité subconsciente, la \textit{motivation} vers l'extrémité consciemment cadrée. Le gradient est la chose réelle ; les étiquettes sont une commodité.

Donc, la thèse, dite sans détour : la motivation, telle que nous la concevons et en parlons réellement, est pour l'essentiel un effet fondé sur la conscience --- tout en concédant que sa moitié inférieure se fond dans une optimisation générique impossible à distinguer de n'importe quelle pulsion de survie. Deux noms, deux niveaux : la part inconsciente de la volonté est l'optimisation du substrat ; la motivation est la couche consciente qui repose dessus. Et cette couche consciente est le carburant que brûle la boucle récursive. Ce qui ne laisse aucune échappatoire à la conclusion stratégique : il n'y a pas d'intelligence générale de type humain --- pas d'AGI --- sans mécanismes apparentés à la conscience. La motivation en est simplement l'exemple le plus lisible.

\section*{L'implication : l'intelligence s'apprend}

Le modèle récursif livre une conséquence que je tiens pour plus importante que tous les arguments théoriques réunis : il prédit que l'intelligence, dans une large mesure, \textit{s'apprend}.

Le savoir s'apprend entièrement --- c'est vrai par définition. La motivation s'apprend pour une bonne part --- des décennies de recherche en théorie de l'autodétermination montrent que la motivation intrinsèque n'est pas un trait figé, mais une réponse aux conditions de l'environnement, en particulier l'autonomie, la compétence et l'appartenance. La performance a un plafond biologique, mais pour la grande majorité des gens, ce plafond n'est pas le goulot d'étranglement. Une capacité de traitement cognitif moyenne suffit amplement pour ce que la plupart des gens reconnaîtraient comme un comportement hautement intelligent.

Les contraintes déterminantes, pour la plupart des gens la plupart du temps, sont la Motivation et le Savoir opérationnel. Et l'un comme l'autre sont sensibles à l'intervention.

Ce qui a un corollaire sombre. Tout système qui détruit systématiquement la motivation des apprenants ne se contente pas d'échouer à développer l'intelligence --- il la \textit{réprime activement}. Les systèmes de notation conventionnels font exactement cela. Une mauvaise note ne se borne pas à rapporter un résultat ; elle attaque la composante Motivation. Une motivation réduite signifie moins d'itérations de la boucle. Moins d'itérations, une croissance plus lente du savoir. Une croissance plus lente, de moins bons résultats à l'évaluation suivante. De moins bons résultats, de nouvelles mauvaises notes. La boucle s'est inversée : au lieu d'une croissance à intérêts composés, l'enfant est désormais piégé dans une stagnation à intérêts composés. Le système de notation produit précisément le résultat qu'il prétend se contenter de mesurer.

Le modèle récursif prédit que ce dommage se cumule et s'amplifie avec le temps --- non pas un préjudice statique, mais une divergence qui s'accélère. Un dommage motivationnel précoce devrait se manifester par un éventail de trajectoires qui s'ouvre davantage à chaque année qui passe. Inversement, les interventions qui renforcent la motivation devraient produire des bénéfices qui \textit{font boule de neige} --- des effets plus marqués au suivi à cinq ans qu'au suivi à un an. Et de fait, les analyses d'interventions précoces comme le Perry Preschool Project montrent exactement ce schéma : des rendements qui croissent avec le temps, portés non par la persistance des gains cognitifs initiaux (qui souvent s'estompent), mais par l'accumulation composée de gains motivationnels et d'autorégulation.

S'il faut retenir une seule leçon pratique du modèle de l'intelligence, c'est celle-ci : la chose la plus précieuse qu'un système éducatif puisse transmettre n'est pas le savoir factuel --- à l'ère de l'IA, les faits sont gratuits --- mais le \textit{savoir opérationnel} et la motivation de s'en servir. Apprendre à apprendre, et vouloir apprendre : voilà à peu près les seules choses qui vaillent encore la peine d'être enseignées.

Je trouve frappant que deux chercheurs chevronnés, travaillant depuis l'intérieur de la psychométrie établie et abordant le problème par des chemins complètement différents, soient arrivés exactement là où pointe mon modèle : la motivation n'est pas un bruit qui viendrait seulement brouiller la mesure de l'intelligence. Elle fait partie de la machinerie.

Werner Wittmann, professeur émérite à Mannheim, s'attaque à la vieille rengaine confortable selon laquelle « la motivation ne compte pas vraiment » --- celle sur laquelle on s'appuie parce que les corrélations de la littérature paraissent faibles. Sa réponse : ces corrélations paraissent faibles parce que les mesures ne s'alignent pas. Mesurez la motivation de façon étroite et la performance de façon large, ou l'inverse, et vous avez brisé ce qu'il appelle la symétrie de Brunswik --- et une symétrie brisée enterre le lien réel, en tirant la corrélation mesurée vers zéro. Accordez les deux au même grain, et le lien ressurgit : la motivation prédit la \textit{variabilité} de la performance, pas seulement sa moyenne. L'effet n'a jamais été petit. Nous le mesurions simplement de travers.

Florian Schmiedek, à partir de l'étude COGITO --- 204 personnes, une centaine de séances quotidiennes chacune ---, aborde la question par le quotidien. Il est convaincu qu'une large part de ce qui fait fluctuer votre cognition d'un jour à l'autre est d'ordre motivationnel, au même titre que le sommeil et la charge. Et il souligne la raison évidente pour laquelle cela continue d'échapper à tout le monde : la plupart des études ne mesurent tout simplement pas la motivation quotidienne propre à la tâche.

Ils établissent donc, indépendamment l'un de l'autre et données à l'appui, \textit{que} la motivation a sa place à l'intérieur de l'intelligence. Mon modèle dit \textit{pourquoi} et \textit{comment} : elle est l'un des facteurs multiplicatifs --- savoir × performance × motivation. Mettez-la à zéro et vous n'atténuez pas seulement l'intelligence --- vous la faites s'effondrer. Leurs chiffres, recueillis pour des raisons qui n'appartiennent qu'à eux, atterrissent exactement là où le modèle l'avait annoncé.

\section*{La dépendance externe}

Un dernier point, parce qu'on le néglige facilement et qu'il compte. La boucle récursive s'auto-renforce, mais elle n'est pas autosuffisante. Elle a besoin de carburant externe --- des informations, des défis, des retours, l'accès au niveau de savoir suivant. Ma propre expérience d'enfant --- me heurter à un mur à onze ans, non par quelque limite interne, mais parce que la réserve de livres de mathématiques s'était épuisée --- l'illustre parfaitement. Les trois composantes étaient en pleine santé. La boucle a calé quand même, parce que les boucles ont besoin d'apports extérieurs pour continuer à itérer.

Cela signifie que le développement de l'intelligence ne dépend pas seulement de la personne, mais de l'environnement. L'accès au savoir, la qualité de l'enseignement, la disponibilité de mentors, les attitudes culturelles envers l'apprentissage --- tout cela nourrit la boucle ou l'affame. Le modèle récursif explique pourquoi les facteurs socio-économiques prédisent si puissamment le développement intellectuel : ils déterminent l'approvisionnement en carburant externe. Un enfant qui grandit dans un foyer riche en livres, avec des parents impliqués, a une boucle alimentée en continu. Un enfant dans un environnement pauvre en ressources a une boucle affamée, quelle que soit sa capacité interne.

L'intelligence n'est pas un trait que vous possédez. C'est un processus que vous faites tourner. Et la bonne marche de ce processus dépend de la machine (la Performance), du logiciel (le Savoir), du conducteur (la Motivation) et de la route (l'environnement extérieur). Les quatre comptent. Tout modèle qui en laisse un de côté se trompera dans ses prédictions.


\chapter*{Annexe C : Cinq classes de calcul}
\addcontentsline{toc}{chapter}{Annexe C : Cinq classes de calcul}
\markboth{Annexe C : Cinq classes de calcul}{}

\textit{Cette annexe développe le cadre computationnel brièvement évoqué au chapitre 5 --- les cinq classes de comportement dynamique qui déterminent si un système physique peut porter la conscience. Les lecteurs à l'aise avec la version intuitive du chapitre 5 peuvent sauter cette annexe sans rien perdre de nécessaire à l'argument principal. Pour ceux qui veulent le tableau complet : c'est ici que les mathématiques rencontrent la physique.}

\section*{Les quatre classes de Wolfram}

En 2002, Stephen Wolfram publiait \textit{A New Kind of Science}, fruit de décennies passées à étudier ce qui se produit lorsqu'on laisse des règles très simples s'appliquer à des systèmes très simples. Son outil central était l'automate cellulaire --- une rangée (ou une grille) de cellules, chacune allumée ou éteinte, toutes mises à jour simultanément selon une règle fixe qui ne prend en compte que les voisines immédiates de chaque cellule.

La surprise, ce fut la variété que ces règles d'une simplicité triviale pouvaient produire. Wolfram classa les comportements en quatre types :


{\small
\noindent
\begin{tabularx}{\linewidth}{>{\centering\arraybackslash}X X X X}
\toprule
\textbf{Classe de Wolfram} & \textbf{Comportement} & \textbf{Exemple} & \textbf{Ce que l'on voit} \\
\midrule
1 & Uniforme & Règle 0 & Tout s'efface. Chaque cellule meurt. \\[4pt]
2 & Périodique & Règle 4 & Des motifs stables et répétitifs. Des clignotants. Des horloges. \\[4pt]
3 & Aléatoire/chaotique & Règle 30 & Une apparence d'aléatoire. Aucune structure répétitive évidente. \\[4pt]
4 & Complexe & Règle 110 & Des structures localisées qui se déplacent, interagissent et persistent. \\[4pt]
\bottomrule
\end{tabularx}
}


Cette classification était réellement utile. Elle capturait quelque chose de réel dans le comportement des systèmes dynamiques, et elle s'appliquait bien au-delà des automates cellulaires --- à la dynamique des fluides, aux systèmes biologiques, aux modèles économiques, aux réseaux de neurones. Les quatre classes n'étaient pas de simples catégories ; c'étaient des attracteurs. Des systèmes issus de domaines radicalement différents retombaient sans cesse dans les quatre mêmes régimes de comportement.

Mais il y avait un problème.

\section*{Le problème des fractales}

La Classe 3 de Wolfram était un fourre-tout. Elle contenait deux types de systèmes fondamentalement différents, qui avaient seulement en commun de \textit{paraître} similaires au premier coup d'œil :

\textbf{Les systèmes fractals}, comme la Règle 90, qui engendre un triangle de Sierpinski parfait --- un motif infiniment auto-similaire, structuré récursivement. Beau, déterministe et sans intérêt du point de vue du calcul : on peut calculer l'état de n'importe quelle cellule à n'importe quel pas de temps sans exécuter toute la simulation. Les mathématiciens appellent cela \textit{computationnellement réductible}.

\textbf{Les systèmes apparemment chaotiques}, comme la Règle 30, dont Wolfram lui-même utilisait la colonne de sortie comme générateur de nombres pseudo-aléatoires dans \textit{Mathematica}. Ils produisent une sortie qui \textit{semble} aléatoire mais reste entièrement déterministe --- même entrée, même sortie, à chaque fois. Impossible d'abréger le calcul ; il faut exécuter chaque étape. Les mathématiciens appellent cela \textit{computationnellement irréductible}.

Wolfram plaça les deux en Classe 3. Sa définition mettait l'accent sur l'\textit{apparence} d'aléatoire (« semble à bien des égards aléatoire »), tout en notant que « des triangles et d'autres structures à petite échelle sont pour l'essentiel toujours visibles à un niveau ou à un autre ». Il reconnaissait l'imperfection de la classification : « avec presque tout schéma général de classification, il y a inévitablement des cas qu'une définition assigne à une classe et une autre définition à une autre classe. »

Eric Rowland, à la conférence NKS de 2006, soutint indépendamment que les motifs imbriqués (fractals) méritaient leur propre cadre de classification.

Je pense que le problème va plus loin qu'une question d'esthétique classificatoire. Les systèmes fractals et les systèmes chaotiques diffèrent structurellement, d'une manière qui compte pour l'argument central de ce livre : quels systèmes peuvent porter la conscience ?

\section*{Le schéma à cinq classes}

Voici le schéma à cinq classes, ordonné en un gradient monotone net, du plus ordonné au plus désordonné :

\textbf{Classe 1 --- Statique.} Des systèmes qui convergent vers un état fixe et s'arrêtent. Un pendule qui oscille une fois puis s'immobilise. Mort. Rien ne s'y calcule. Période : 1.

\textbf{Classe 2 --- Périodique.} Des systèmes qui s'installent dans des boucles répétitives. Une horloge. Un battement de cœur (approximativement). L'information est stockée dans le motif, mais jamais transformée. Période : finie.

\textbf{Classe 3 --- Fractale.} Des systèmes qui produisent une structure auto-similaire à toutes les échelles. Un triangle de Sierpinski. Une fougère. L'ensemble de Mandelbrot. Mathématiquement riches, esthétiquement éblouissants --- et \textit{computationnellement réductibles} : on peut sauter des étapes sans exécuter tout le calcul. De la structure sans puissance de traitement. Période : quasi infinie, avec une auto-similarité exacte ou statistique à toutes les échelles.

\textbf{Classe 4 --- Complexe (bord du chaos).} Des systèmes qui produisent des structures localisées persistantes, capables de se déplacer, d'interagir et d'encoder des calculs arbitraires. Le Jeu de la Vie de Conway. L'automate cortical. Computationnellement \textit{irréductibles} --- aucun raccourci possible. Ces systèmes sont capables de calcul universel : avec les bonnes conditions initiales, ils peuvent simuler n'importe quel algorithme, y compris des simulations d'eux-mêmes. Période : quasi infinie, avec auto-similarité \textit{plus} structures persistantes en interaction. C'est là que vit la conscience.

\textbf{Classe 5 --- Aléatoire.} Des systèmes dont la sortie est authentiquement aléatoire --- ni pseudo-aléatoire, ni déterministe, ni compressible. Aucun motif, aucune auto-similarité, aucune période qui finisse par se répéter. Un contenu d'information véritablement infini. Structure : \textit{inconnue} (voir plus bas).

La correspondance avec le schéma de Wolfram :


{\small
\noindent
\begin{tabularx}{\linewidth}{>{\centering\arraybackslash}X >{\centering\arraybackslash}X X}
\toprule
\textbf{Cinq classes} & \textbf{Wolfram} & \textbf{Ce qui a changé} \\
\midrule
1 & Classe 1 & Identique \\[4pt]
2 & Classe 2 & Identique \\[4pt]
3 & Classe 3 (en partie) & Scindée de la Classe 3 de Wolfram \\[4pt]
4 & Classe 4 & Identique \\[4pt]
5 & Classe 3 (en partie) & Scindée de la Classe 3 de Wolfram \\[4pt]
\bottomrule
\end{tabularx}
}


L'ordre de Wolfram le long du spectre du désordre était : 1 $\rightarrow$ 2 $\rightarrow$ 4 $\rightarrow$ 3. Bancal. Le schéma à cinq classes donne un gradient monotone net : 1 $\rightarrow$ 2 $\rightarrow$ 3 $\rightarrow$ 4 $\rightarrow$ 5, ordonné par désordre croissant et irréductibilité computationnelle croissante.

\section*{Pourquoi les automates déterministes ne peuvent pas produire de hasard}

Voici l'argument que je crois original, et qui renforce la thèse de cinq classes plutôt que quatre.

Considérez un automate cellulaire --- n'importe lequel. Il possède une table de règles finie (exprimable en un nombre fini de bits) et une condition initiale finie (exprimable, elle aussi, en un nombre fini de bits). Ensemble, la règle et la condition initiale contiennent une quantité d'information fixe et finie.

Maintenant : une quantité finie d'information peut-elle produire une sortie véritablement aléatoire ?

Non. Et voici pourquoi :

\begin{enumerate}
  \item Une suite infinie véritablement aléatoire possède une complexité de Kolmogorov \textit{maximale} --- on ne peut pas la comprimer, on ne peut pas la décrire par quoi que ce soit de plus court qu'elle-même.
  \item La sortie d'un automate cellulaire est entièrement déterminée par sa règle et sa condition initiale, qui possèdent ensemble une complexité de Kolmogorov \textit{finie}.
  \item On ne peut pas tirer d'un processus plus d'information qu'on n'en a introduit par sa spécification.
  \item Par conséquent, la sortie de n'importe quel automate cellulaire a une faible complexité de Kolmogorov relativement à une suite véritablement aléatoire de même longueur.
\end{enumerate}

C'est une version généralisée du principe des tiroirs : une information finie doit produire une structure auto-similaire. La seule façon d'engendrer une sortie infinie à partir d'une information finie est de \textit{réutiliser} de la structure à différentes échelles. La réutilisation exacte, c'est la périodicité (Classe 2). La réutilisation non exacte mais organisée en motifs, c'est le comportement fractal (Classe 3). Même les automates cellulaires d'apparence la plus complexe (la Règle 30, la Règle 110, le Jeu de la Vie) produisent une sortie dont la complexité est bornée par la complexité de leur jeu de règles.

Ce que Wolfram appelait des automates cellulaires « aléatoires » se décrit mieux comme des \textbf{fractales de haute complexité} --- des systèmes dont la structure auto-similaire est bien réelle, mais opère à des échelles et dans des dimensions qui la rendent invisible à un examen superficiel. Le bord gauche de la Règle 30 présente d'ailleurs des sous-structures de type Sierpinski. Sa colonne centrale réussit de nombreux tests statistiques d'aléa, ce qui est \textit{exactement ce qu'on attendrait d'une fractale de haute complexité} : la statistique locale imite le hasard, mais la structure globale est déterministe et de faible complexité de Kolmogorov --- une règle courte l'engendre tout entière. (C'est ce que « compressible » veut dire ici : il existe une description compacte. Cela ne signifie \textit{pas} qu'on puisse abréger le calcul pas à pas, lequel reste irréductible.)

Par le même argument, la sortie de Classe 4 est \textit{elle aussi} fractale --- le Jeu de la Vie présente une auto-similarité statistique dans sa dynamique de population, ses distributions structurelles, ses corrélations spatiales. La différence entre la Classe 3 et la Classe 4 n'est pas « fractal contre non fractal ». Elle est la suivante :

\begin{itemize}
  \item \textbf{Classe 3} : fractale. Réductible. De la structure sans traitement de l'information.
  \item \textbf{Classe 4} : fractale. Irréductible. De la structure \textit{avec} traitement de l'information --- des structures localisées persistantes qui interagissent et peuvent encoder le calcul universel.
\end{itemize}

Toutes deux sont fractales. Une seule calcule.


{\small
\noindent
\begin{tabularx}{\linewidth}{>{\centering\arraybackslash}X X X X >{\centering\arraybackslash}X >{\centering\arraybackslash}X}
\toprule
\textbf{Classe} & \textbf{Règles} & \textbf{Période} & \textbf{Structure} & \textbf{Réductible ?} & \textbf{Calcule ?} \\
\midrule
1 & Finies & 1 & Aucune & Trivialement & Non \\[4pt]
2 & Finies & Finie & Répétitive & Oui & Non \\[4pt]
3 & Finies & Quasi infinie, auto-similaire & Auto-similaire & Oui & Non \\[4pt]
4 & Finies & Quasi infinie, auto-similaire & Auto-similaire + structures persistantes en interaction & Non & \textbf{Oui} \\[4pt]
5 & Inexprimables & Véritablement infinie & \textbf{Inconnue} & N/A & N/A \\[4pt]
\bottomrule
\end{tabularx}
}


\section*{La Classe 5 et la frontière de l'exprimabilité mathématique}

Si les automates déterministes ne peuvent pas produire de véritable hasard, alors qu'est-ce qui le \textit{peut} ?

Cette question conduit à ce qui est, je crois, l'implication la plus profonde du schéma à cinq classes.

Les Classes 1 à 4 sont ce que des règles finies et exprimables peuvent produire. Leur comportement va du trivial (Classe 1) à l'extraordinaire (Classe 4 --- calcul universel, conscience), mais tout cela est engendré par des règles qui peuvent être écrites, communiquées, vérifiées et analysées. Ces règles vivent à l'intérieur des mathématiques, dans le domaine des systèmes symboliques formels.

La Classe 5, en revanche, exige des règles qui ne peuvent \textit{pas} être écrites. Un système qui produit une sortie authentiquement aléatoire --- une sortie de complexité de Kolmogorov maximale, incompressible, non algorithmique --- ne peut exécuter aucune règle qu'un système formel puisse exprimer. Si la règle était exprimable, la sortie serait compressible (en « appliquer cette règle »), et donc pas véritablement aléatoire.

Cela place la Classe 5 à la frontière même de l'exprimabilité mathématique. Elle n'est pas simplement « très complexe » ou « très désordonnée ». Elle est le régime où le processus générateur excède ce que les systèmes symboliques linéaires (mathématiques, logique, calcul) peuvent saisir.

Quelque chose, dans la nature, opère-t-il réellement en Classe 5 ?

Peut-être. La mécanique quantique produit des résultats de mesure qui, en vertu du théorème de Bell, ne peuvent être expliqués par aucune théorie locale à variables cachées. Si ces résultats sont authentiquement aléatoires --- et non des processus déterministes que nous n'avons pas encore identifiés ---, alors la mesure quantique est un processus de Classe 5 : un phénomène physique dont les règles ne peuvent s'écrire dans aucun langage formel que nous possédions.

C'est spéculatif, et je le signale comme tel. Mais l'implication est frappante : la Classe 4 (le régime de la conscience, du calcul universel, de l'automate cortical) se situe à la \textit{complexité maximale atteignable par des règles exprimables}. Elle est aussi complexe que les mathématiques peuvent l'être. Au-delà s'étend un territoire que les mathématiques, par leur nature même, ne peuvent pas cartographier.

\section*{La structure de la Classe 5 : inconnue, pas absente}

Une dernière subtilité. Il est tentant de dire que la Classe 5 n'a « aucune structure ». Mais ce serait une erreur --- la même erreur que celle qui consiste à dire que l'infini n'a aucune structure.

Avant les travaux de Georg Cantor dans les années 1870, l'infini était traité comme un concept unique : les choses étaient soit finies, soit infinies, un point c'est tout. Cantor a montré qu'il existe des \textit{hiérarchies} de l'infini --- que l'infini des nombres réels est strictement plus grand que celui des entiers, et que cette hiérarchie se prolonge sans fin. L'infini s'est révélé doté d'une riche architecture interne, restée invisible parce que les mathématiciens n'avaient pas les outils pour la voir.

Il pourrait en aller de même pour le hasard. Nous traitons aujourd'hui le hasard véritable comme une catégorie unique --- le désordre maximal, l'absence de motif. Mais nous sommes dans la position des mathématiciens pré-Cantor face à l'infini : il nous manque les outils conceptuels pour distinguer différentes sortes de hasard, si tant est que de telles distinctions existent.

La réponse honnête, quant à la structure de la Classe 5, est donc : \textbf{inconnue}. Pas « aucune ». Pas « absente ». Inconnue --- en attente d'outils conceptuels qui n'existent peut-être pas encore, qui exigent peut-être des manières de penser que les systèmes symboliques linéaires ne peuvent pas fournir.

C'est, je crois, l'une des questions ouvertes les plus importantes à l'intersection des mathématiques, de la physique et du calcul. Et elle reste invisible sans le schéma à cinq classes, parce que le cadre à quatre classes de Wolfram ne crée jamais l'espace où la poser.

\section*{Implications pour la conscience}

Le schéma à cinq classes éclaire pourquoi la conscience requiert une dynamique de Classe 4 --- et seulement la Classe 4.

Les Classes 1 et 2 sont trop simples. Elles peuvent stocker de l'information (un état fixe, un motif répétitif) mais ne peuvent pas la \textit{traiter} d'une manière computationnellement intéressante. Un cerveau en sommeil profond, qui égrène de lentes ondes delta, opère en Classe 2 : périodique, répétitif, il ne mène nulle part. L'architecture à quatre modèles est intacte dans le substrat, mais la simulation ne tourne pas.

La Classe 3 est intéressante mais non computationnelle. Les dynamiques fractales produisent une structure riche, et le cerveau les utilise (voir plus loin) --- mais elles ne peuvent pas soutenir le type de traitement dynamique, irréductible et globalement intégré qu'exige une auto-simulation consciente. Un motif fractal est beau, mais il est computationnellement réductible. Il ne peut pas se surprendre lui-même.

La Classe 4 possède exactement les deux propriétés dont la conscience a besoin : le \textbf{calcul universel} (le système peut, en principe, tout simuler, y compris lui-même) et l'\textbf{intégration globale} (des parties éloignées du système s'influencent mutuellement, les changements locaux se propagent globalement, l'information se lie en un tout unifié). Au bord du chaos, l'automate cortical atteint l'une et l'autre, et le résultat est la conscience.

La Classe 5 est différente, non pas parce que le calcul y serait impossible (une séquence aléatoire infinie contient \textit{tout}, y compris chaque motif stable et chaque calcul jamais conçu), mais parce que nous n'avons aucun moyen de l'exploiter, de la prédire ou de la démontrer. Un cerveau en crise généralisée, dont les neurones déchargent dans un chaos désordonné, s'approche de la Classe 5, non parce que la conscience y serait impossible en principe, mais parce qu'aucun mécanisme n'existe pour la soutenir ou y accéder. Notre univers lui-même pourrait être un fragment de hasard infini, ou un système de Classe 4 à une échelle que nous ne pouvons pas percevoir, ou peut-être un fragment d'une fractale infinie. Nous ne pouvons pas le savoir. Ce que nous \textit{pouvons} dire, c'est que la conscience, telle que nous en faisons l'expérience, requiert l'imprévisibilité structurée de la Classe 4. La simulation s'effondre en Classe 5 non parce que la réalité sous-jacente serait insuffisante, mais parce qu'aucune interface stable n'existe entre le substrat et la simulation.

\section*{Le cerveau utilise les quatre classes}

Le cerveau est un ordinateur universel optimisé par des milliards d'années d'évolution. Il serait étrange que l'évolution ait manqué le moindre régime computationnel offrant un avantage. Et de fait, le cerveau utilise les quatre classes exprimables comme autant d'outils distincts :

\begin{itemize}
  \item \textbf{Classe 1} (attracteurs stables) : le stockage de la mémoire à long terme. Des configurations de poids synaptiques qui persistent pendant des années. Les points fixes du réseau de neurones.
  \item \textbf{Classe 2} (oscillations) : les rythmes alpha, thêta, gamma et delta. La cadence thalamique. Les cycles veille-sommeil. Les mécanismes de chronométrage et de filtrage du cerveau.
  \item \textbf{Classe 3} (traitement fractal / invariant d'échelle) : l'analyse de texture, la reconnaissance d'objets invariante à l'échelle, le codage neuronal efficace. Principalement le traitement visuel en V2-V4, où la comparaison multi-échelle est l'opération centrale. Sous psychédéliques, quand cette machinerie tourne sans entrée externe, on \textit{voit} le traitement fractal lui-même, ce qui explique pourquoi les motifs fractals comptent parmi les traits les plus constants de l'expérience psychédélique (voir chapitre 6).
  \item \textbf{Classe 4} (bord du chaos) : l'automate cortical lui-même. Le régime dynamique du traitement conscient. La computation universelle. Le moteur de la simulation.
\end{itemize}

Chaque classe remplit une fonction différente. Seule la Classe 4 engendre la conscience. Mais la conscience dépend des autres : des mémoires stables (Classe 1) pour alimenter les modèles, un rythme temporel (Classe 2) pour coordonner les dynamiques, et un traitement fractal (Classe 3) pour analyser le monde à plusieurs échelles simultanément.

C'est peut-être là la raison la plus profonde pour laquelle le cerveau doit opérer précisément au bord du chaos : la Classe 4 est le seul régime capable de \textit{recruter} toutes les autres classes --- et lui-même. Un automate de Classe 4 peut produire des états stables (comportement de Classe 1), des oscillations périodiques (comportement de Classe 2) et des structures fractales (comportement de Classe 3) comme sous-processus au sein de sa propre dynamique. Et il peut engendrer d'autres automates de Classe 4 : un ordinateur universel peut simuler un autre ordinateur universel. Aucune des autres classes ne peut rien faire de tout cela. La Classe 4 n'est pas seulement la classe la plus complexe --- c'est la seule classe qui contient toutes les classes, y compris elle-même. Cette auto-inclusion est ce qui rend possible l'identité structurelle inter-échelles : un cerveau de Classe 4 à l'intérieur d'un univers de Classe 4 n'est pas une analogie. C'est une instance imbriquée de la même architecture computationnelle.


\chapter*{Annexe D : Comment faire des rêves lucides}
\addcontentsline{toc}{chapter}{Annexe D : Comment faire des rêves lucides}
\markboth{Annexe D : Comment faire des rêves lucides}{}

\textit{Au chapitre 6, j'ai évoqué le rêve lucide comme un moyen sûr, sans aucune substance, d'explorer la conscience de l'intérieur. Dans la Théorie des Quatre Modèles, le rêve lucide correspond au Modèle Explicite du Soi (ESM) qui « s'allume » plus pleinement pendant le sommeil paradoxal --- un franchissement du seuil de criticalité qui transforme un rêve passif en expérience contrôlée. Voici la méthode la plus simple pour y parvenir.}

\section*{La méthode du test de réalité}

Le principe est simple : si vous prenez l'habitude de vous demander si vous êtes éveillé, cette habitude finira par se déclencher au beau milieu d'un rêve --- et le rêve échouera au test.

\textbf{Étape 1 : choisissez un test de réalité.} Le plus fiable consiste à regarder un texte, à détourner les yeux, puis à le regarder de nouveau. À l'état de veille, le texte reste identique. En rêve, il change --- souvent de façon spectaculaire. Les horloges fonctionnent aussi : vérifiez l'heure, détournez le regard, vérifiez encore. En rêve, les chiffres seront différents ou absurdes. Autre test fiable : essayez d'enfoncer un doigt dans la paume de votre autre main. En rêve, il passe souvent au travers.

\textbf{Étape 2 : faites-le toute la journée.} Chaque fois que vous franchissez une porte, que vous consultez votre téléphone ou que vous remarquez quelque chose de légèrement étrange, arrêtez-vous et effectuez votre test de réalité. La clé n'est pas le test lui-même --- c'est la \textit{question sincère} qui le sous-tend : « Suis-je en train de rêver, là, maintenant ? » Ne vous contentez pas d'exécuter les gestes. Envisagez réellement la possibilité.

\textbf{Étape 3 : tenez un journal de rêves.} Posez un carnet près de votre lit. Chaque matin, avant même de bouger, notez tout ce dont vous vous souvenez --- même s'il ne s'agit que d'une sensation ou d'une seule image. Cela entraîne votre cerveau à traiter le contenu des rêves comme digne d'être mémorisé, ce qui renforce le pont entre la conscience onirique et la conscience de veille.

\textbf{Étape 4 : patientez.} Pour la plupart des gens, le premier rêve lucide survient en l'espace de deux à six semaines. Vous serez dans un rêve, quelque chose vous semblera légèrement anormal, l'habitude du test de réalité se déclenchera, le texte changera --- et vous \textit{saurez}. Ce moment de lucidité, c'est l'ESM qui s'active. Vous sentirez la transition : une netteté soudaine, un sentiment de présence, la reconnaissance tranquille que le monde qui vous entoure est une simulation à l'intérieur de laquelle vous êtes conscient.

\section*{À quoi s'attendre}

Votre premier rêve lucide sera probablement bref --- de quelques secondes à quelques minutes. L'excitation a tendance à vous réveiller. Avec de la pratique, vous pourrez les prolonger. Certaines personnes parviennent à faire des rêves lucides plusieurs fois par semaine. L'expérience est remarquable : vous êtes à l'intérieur de la simulation consciente complète, sans le moindre signal extérieur, et vous le savez. Le monde virtuel répond à vos intentions. C'est, au sens littéral, la Théorie des Quatre Modèles devenue expérience vécue.

\section*{Autres méthodes}

Pour les lecteurs qui souhaitent aller plus loin, il existe des techniques plus exigeantes :

\begin{itemize}
  \item \textbf{MILD (Mnemonic Induction of Lucid Dreams)} --- mise au point par Stephen LaBerge à Stanford. Vous formulez, au moment de vous endormir, l'intention de reconnaître que vous rêvez. À combiner de préférence avec un réveil après cinq heures de sommeil, suivi d'un rendormissement.
  \item \textbf{WILD (Wake-Initiated Lucid Dream)} --- vous maintenez la conscience pendant la transition de la veille au rêve. Difficile, mais c'est la technique qui produit les résultats les plus vifs.
  \item \textbf{WBTB (Wake Back to Bed)} --- vous vous réveillez après cinq à six heures de sommeil, restez éveillé vingt à soixante minutes, puis vous rendormez. Cette méthode cible les cycles de fin de nuit, riches en sommeil paradoxal.
\end{itemize}

\textit{Exploring the World of Lucid Dreaming} (1990) de Stephen LaBerge reste le guide pratique de référence. Pour les neurosciences, voir Voss et al. (2009) sur les signatures EEG du rêve lucide, et Baird et al. (2019) pour une revue complète des neurosciences cognitives des rêves lucides.


\chapter*{Annexe E : Pourquoi « quatre » modèles ? --- Une note pour les neuroscientifiques}
\addcontentsline{toc}{chapter}{Annexe E : Pourquoi « quatre » modèles ? --- Une note pour les neuroscientifiques}
\markboth{Annexe E : Pourquoi « quatre » modèles ? --- Une note pour les neuroscientifiques}{}

Cette annexe répond à une objection que tout neuroscientifique, ou tout lecteur au fait de l'informatique, formulera en rencontrant l'architecture à quatre modèles au Chapitre 2 : \textit{le cerveau n'entretient tout de même pas exactement quatre modèles ?}

En effet, il ne le fait pas. Le nombre « quatre » est un \textbf{minimum de principe}, non un décompte littéral.

\section*{Ce que le cerveau fait réellement}

Le substrat biologique --- des neurones qui déchargent au-dessus de réseaux protéomiques, avec des voies de signalisation intracellulaire qui constituent leur propre intelligence computationnelle jusqu'à l'intérieur d'une seule cellule --- met en œuvre un nombre pratiquement incalculable de modèles qui se chevauchent, des deux côtés de la frontière implicite/explicite.

Prenez l'exemple du geste de saisir une tasse. Le modèle moteur encode simultanément la géométrie du monde (où se trouve la tasse, quels obstacles l'entourent) et la cinématique du soi (comment le bras est configuré, comment les doigts doivent se positionner pour la prise). Ce modèle unique n'est \textit{ni} un pur modèle du monde \textit{ni} un pur modèle du soi --- il déborde sur les deux catégories. Un modèle émotionnel d'une interaction sociale encode à la fois des connaissances sur l'autre personne (monde) et une évaluation de soi-même (soi). Un modèle de navigation spatiale encode aussi bien l'agencement de l'environnement que votre position à l'intérieur de celui-ci. Tout modèle neuronal réel est un mélange.

Les frontières entre les « modèles » ne sont pas nettes, leur nombre n'est pas fixe, et il n'est certainement pas de quatre.

\section*{Pourquoi quatre reste la bonne abstraction}

Les quatre modèles canoniques (IWM, ISM, EWM, ESM) sont les \textbf{points extrêmes} d'un espace continu à deux dimensions défini par deux axes :

\begin{itemize}
  \item \textbf{Portée} : de la pure représentation de soi à la pure représentation du monde
  \item \textbf{Mode} : du tout implicite (structurel, stocké, inconscient) au tout explicite (simulé, transitoire, phénoménal)
\end{itemize}

L'écologie de modélisation réelle du cerveau remplit cet espace tout entier d'une densité continue de modèles qui se chevauchent. Les quatre modèles nommés en sont les quatre coins --- les pôles théoriques autour desquels l'activité s'organise. Voyez-les comme des points cardinaux : utiles pour s'orienter, réels en tant que directions, mais personne ne prétendrait que le monde ne contient que quatre lieux.

Si la théorie repose sur ces quatre pôles plutôt que sur l'espace continu tout entier, c'est parce qu'ils identifient la \textbf{configuration minimale} dont un système a besoin pour être conscient :

\begin{itemize}
  \item \textbf{Pas de modèle du monde} $\rightarrow$ aucun environnement à vivre
  \item \textbf{Pas de modèle du soi} $\rightarrow$ aucun sujet pour le vivre
  \item \textbf{Pas de niveau implicite} $\rightarrow$ rien à partir de quoi simuler (aucune connaissance apprise)
  \item \textbf{Pas de niveau explicite} $\rightarrow$ aucune simulation du tout (aucune expérience)
\end{itemize}

Retirez l'un des quatre et quelque chose de crucial se brise. Les quatre modèles sont le plancher, non le plafond. Le cerveau les dépasse dans toutes les directions. Mais c'est le plancher qui vous dit ce que la conscience \textit{exige} --- et c'est le plancher qui engendre les prédictions de la théorie, circonscrit ses affirmations et précise ce que tout système artificiel devrait mettre en œuvre.

\section*{Lire le reste du livre}

Tout au long de ce livre, lorsque j'écris « l'ESM fait ceci » ou « l'IWM contient cela », je renvoie à ces pôles de l'espace continu, sans prétendre que le cerveau posséderait quatre boîtes séparées, avec des cloisons entre elles. La simplification est fondée en raison, et les chapitres qui suivent la montreront à l'œuvre dans un véritable travail explicatif --- dérivant la phénoménologie psychédélique, les mécanismes anesthésiques, les états oniriques, les phénomènes de cerveau divisé et la conscience animale à partir de cinq principes bâtis sur cette architecture.

Pour le traitement mathématique complet --- incluant le cadre de l'espace continu des modèles, la fonction de densité des modèles et la formalisation de la perméabilité comme transfert d'information entre régions de cet espace --- voir Gruber (2026), \textit{Toward a Mathematical Formalization of the Four-Model Theory}.


\chapter*{Annexe F : Le Modèle standard comme comptabilité des frontières}
\addcontentsline{toc}{chapter}{Annexe F : Le Modèle standard comme comptabilité des frontières}
\markboth{Annexe F : Le Modèle standard comme comptabilité des frontières}{}

\textit{Deux des dérivations qui sous-tendent l'image des atomes computationnels --- pourquoi les lois de conservation sont absolues, et pourquoi il existe exactement trois générations de particules --- constituent la comptabilité la plus aride de tout l'argument ; je les ai donc sorties du texte principal. Les voici in extenso, pour les lecteurs qui veulent le détail de physique des particules. La version vivante se trouve au chapitre « L'architecture du tout » ; ceci en est le grand livre.}

\textbf{Pourquoi les lois de conservation sont-elles si absolues ?} La charge est toujours conservée. Le nombre baryonique est conservé. Le nombre leptonique est conservé. On n'a jamais vu ces lois prises en défaut, pas une seule fois, dans aucune expérience jamais menée. Pourquoi ?

Parce que ce sont des contraintes de conservation de l'information. La borne de Bekenstein vous dit combien d'information une frontière peut contenir. Quand deux frontières interagissent et échangent de l'information, l'information totale est conservée --- elle doit l'être, parce que la conservation de l'information est une conséquence de l'unitarité de la mécanique quantique --- et l'unitarité, dans cette perspective, pourrait elle-même être sous-tendue par la borne de Bekenstein. Les lois de conservation spécifiques de la physique des particules --- conservation de la charge, du nombre baryonique, du nombre leptonique --- sont les règles spécifiques qui gouvernent la façon dont l'information peut se transformer lorsque des configurations de frontière interagissent. Ce ne sont pas des règles arbitraires imposées de l'extérieur. Ce sont des contraintes comptables qui découlent du fait qu'on ne peut ni créer ni détruire d'information à une frontière de singularité.

\textbf{Et puis il y a le mystère des trois générations.} Les particules se déclinent en trois générations. L'électron possède une copie plus lourde (le muon) et une copie plus lourde encore (le tau). Le quark up a des copies appelées charme et top. Trois versions de chaque type de particule, identiques en tout point sauf la masse. C'est l'un des motifs inexpliqués les plus profonds de la physique des particules. Personne ne sait pourquoi trois. Ni deux, ni quatre, ni dix-sept. Trois.

Autant être tout à fait franc : ce que je m'apprête à dire est spéculatif. Plus spéculatif que le reste de l'image des atomes computationnels. Mais c'est structurellement motivé, et je pense que cela mérite d'être posé.

Les systèmes de Classe 4 contiennent intrinsèquement une structure auto-similaire. C'est une conséquence technique du fait que la dynamique de Classe 4 contient le comportement de Classe 3 (fractal) comme sous-processus. L'auto-similarité, c'est le même motif qui se répète à différentes échelles. Si les configurations de frontière de singularité sont plongées dans un système de Classe 4 --- et elles doivent l'être, puisque l'univers est de Classe 4 ---, alors les configurations elles-mêmes peuvent présenter une structure auto-similaire. Le même type de frontière à trois échelles d'énergie différentes. Trois générations pourraient être la signature d'une hiérarchie fractale dans l'espace des configurations de singularité stables.

Je n'ai pas de preuve. C'est une conjecture, pas une dérivation. Mais je note que trois est exactement ce qu'on attendrait de la hiérarchie auto-similaire non triviale la plus simple : une configuration de base et deux copies mises à l'échelle. Et je note que la structure en générations n'est par ailleurs expliquée par aucune théorie actuelle. Si l'image des atomes computationnels finit par expliquer pourquoi il existe exactement trois générations, ce serait un indice fort en faveur de l'ensemble du cadre.

Et si l'univers est réellement un automate cellulaire, une quatrième génération existe --- et une cinquième, et d'autres encore. Mais les générations supérieures sont des motifs plus grands, et les motifs plus grands sont moins stables. Ils ne peuvent pas résister à l'assaut des configurations plus petites et plus stables qui les découpent avant qu'ils aient pu se constituer correctement. C'est exactement ce que nous observons : les particules de troisième génération (le tau, le quark top) sont déjà extrêmement instables et se désintègrent presque instantanément. Une quatrième génération serait plus lourde encore, sa configuration de frontière plus grande et plus complexe, et donc plus fragile encore --- trop fragile pour subsister assez longtemps pour être détectée comme une particule plutôt que comme une fluctuation passagère. Les trois générations que nous voyons sont peut-être simplement les trois qui sont assez petites pour survivre.


\chapter*{Annexe G : Quatre points faibles du modèle cosmologique}
\addcontentsline{toc}{chapter}{Annexe G : Quatre points faibles du modèle cosmologique}
\markboth{Annexe G : Quatre points faibles du modèle cosmologique}{}

\textit{Une théorie qui cache ses points faibles n'est pas une théorie. L'argument cosmologique présenté dans « L'architecture du tout » et « Le miroir le plus profond » repose sur des affirmations qui ne sont pas prouvées, et l'honnêteté exige de nommer les endroits où il pourrait échouer. Le cinquième point faible, le pire --- le plafond cognitif, la crainte qu'un cerveau de Classe 4 ne fasse jamais que projeter sa propre architecture sur le cosmos ---, je le garde dans le texte principal, parce que c'est celui auquel je ne peux pas répondre. Les quatre autres sont ici, avec leurs contre-arguments, pour quiconque voudra les attaquer.}

\textbf{Un : l'équivalence entre énergie et information n'est pas prouvée.} Elle est fortement suggérée par le principe de Landauer, la borne de Bekenstein et la thermodynamique des trous noirs. Plusieurs faisceaux d'indices indépendants pointent tous dans la même direction. Mais personne n'a dérivé E = I à partir de premiers principes. Personne n'a montré que l'information, à elle seule, a des effets gravitationnels. L'hypothèse est convaincante, et la convergence des indices est impressionnante, mais la convergence n'est pas une preuve. Si l'énergie et l'information se révèlent simplement corrélées plutôt qu'identiques, l'interprétation des singularités comme transformation d'information perd son fondement. La correspondance structurelle entre conscience et cosmologie tiendrait toujours (les cinq correspondances ne dépendent pas de E = I), mais le mécanisme physique qui les relie serait beaucoup plus faible. La poésie survivrait. La physique, peut-être pas.

\textbf{Deux : l'argument de la Classe 4 est abductif, pas déductif.} J'ai éliminé les autres classes, mais l'élimination n'est pas une preuve. C'est une inférence à la meilleure explication --- une forme de raisonnement parfaitement respectable en science, mais une forme qui laisse une porte ouverte. Peut-être existe-t-il une Classe 4,5 que je n'ai pas imaginée, un régime computationnel entre complexité et hasard que je n'ai pas les outils conceptuels pour décrire. Peut-être la hiérarchie des cinq classes est-elle elle-même incomplète. L'argument dit que la Classe 4 est la meilleure explication de la dynamique de l'univers, pas la seule logiquement possible. Le raisonnement abductif a fait ses preuves --- l'argument de Darwin en faveur de la sélection naturelle était abductif, et il a plutôt bien tenu ---, mais ce n'est pas la même chose qu'une preuve mathématique, et je ne prétendrai pas le contraire.

\textbf{Trois : l'unification des singularités exige une gravité quantique.} Affirmer que toutes les singularités (échelle de Planck, trou noir, singularité cosmologique) sont des instances structurellement identiques de la même frontière informationnelle est une affirmation forte. Elle exige une théorie qui jette un pont entre la mécanique quantique et la relativité générale. Nous n'en avons pas. La théorie des cordes est candidate. La gravité quantique à boucles est candidate. Toutes deux sont compatibles avec l'affirmation, et toutes deux sont non confirmées. L'unification des singularités n'est pas une spéculation débridée --- c'est la direction que prend la physique moderne ---, mais ce n'est pas non plus de la physique établie. C'est un pari sur l'avenir. Je pense que le pari est bon. Je pourrais me tromper.

\textbf{Quatre : le modèle est peut-être infalsifiable de l'intérieur --- par sa propre prédiction.} C'est celui-là qui me noue l'estomac. La conséquence gödélienne de la clôture auto-référentielle dit qu'un système suffisamment complexe qui se calcule lui-même ne peut pas, de l'intérieur, prouver toutes les vérités à son propre sujet. Il existe des énoncés vrais mais improuvables. Si le SB-HC4A est correct, alors l'univers est exactement un tel système, et l'énoncé « le SB-HC4A est correct » pourrait être l'un de ces énoncés vrais-mais-improuvables. La théorie prédit qu'il pourrait être structurellement impossible de la prouver depuis l'intérieur de l'univers. Non pas parce que nous n'avons pas assez essayé. Non pas parce qu'il nous faudrait de meilleurs instruments. Parce que l'architecture rend la chose impossible, de la même manière qu'un système d'axiomes ne peut pas prouver sa propre cohérence.

C'est soit le résultat le plus profond de la philosophie des sciences, soit l'esquive la plus élégante jamais conçue. Une théorie qui prédit sa propre invérifiabilité soit vous dit quelque chose de profond sur les limites de la connaissance, soit s'immunise contre la critique d'une manière qui devrait vous rendre profondément méfiant. Honnêtement ? Je n'en suis pas sûr. J'y réfléchis depuis des années, et je ne sais toujours pas. Ce que je sais, c'est que la prédiction ne vient pas de nulle part --- elle vient des théorèmes de Gödel, qui sont aussi solides que ce que les mathématiques ont de plus établi. Si l'univers est auto-référentiel, l'incomplétude s'ensuit. La question est de savoir si l'univers est auto-référentiel. Et la théorie dit oui.

Alors croyez-moi quand je dis qu'il s'agit d'un véritable point faible, pas d'un atout que j'essaierais de vous faire passer en douce. Si quelqu'un trouve un moyen de tester le SB-HC4A depuis l'intérieur de l'univers et que le test échoue, la théorie est morte. Si le test réussit, ironie de la chose, la théorie doit être fausse elle aussi --- parce qu'un test réussi signifierait que la structure computationnelle de l'univers est accessible de l'intérieur, ce qui contredit les conditions aux limites mêmes qui définissent l'architecture. Ce n'est que si aucun test n'est possible que la théorie survit, mais elle survit alors dans un étrange crépuscule philosophique, infalsifiable non par esquive, mais par structure.


\chapter*{Annexe H : Les neuf prédictions}
\addcontentsline{toc}{chapter}{Annexe H : Les neuf prédictions}
\markboth{Annexe H : Les neuf prédictions}{}

\textit{Une théorie qui explique tout et ne prédit rien n'est pas une théorie --- c'est une histoire. Voici la batterie complète de tests de falsification : les neuf prédictions que formule la Théorie des Quatre Modèles (FMT), chacune assortie de l'expérience concrète qui la tuerait. Plusieurs peuvent être menées dès aujourd'hui avec les technologies existantes ; quelques-unes bénéficient déjà de premiers soutiens ; aucune n'a été falsifiée. La prédiction 3 est développée sous forme de scène au chapitre 11 et n'est donnée ici que brièvement ; plusieurs autres ont été introduites dans les chapitres où elles trouvaient leur place, et sont rassemblées ici en un seul endroit pour le lecteur qui veut mettre la théorie à l'épreuve.}

\section*{Prédiction 1 : Chaque modèle possède sa propre signature neuronale}

Si les quatre modèles sont véritablement des processus distincts, nous devrions pouvoir les voir dans les scanners cérébraux. Concevez une expérience astucieuse qui demande à des personnes d'accomplir quatre types de tâches différents (l'un sollicitant chaque modèle), et les schémas d'activation cérébrale devraient paraître différents.

Une tâche à dominante IWM pourrait être quelque chose comme reconnaître passivement un visage familier. Vous n'y pensez pas ; votre cerveau le sait, tout simplement. Une tâche à dominante ISM pourrait être une séquence motrice habituelle --- taper votre mot de passe, par exemple, sans penser consciemment à chaque touche. Une tâche à dominante EWM exige une perception active et consciente --- peut-être tenter de repérer la différence entre deux images presque identiques. Et une tâche à dominante ESM relève de la pure introspection : « Suis-je le genre de personne qui ferait cela ? »

La prédiction est un schéma en 2×2. Les tâches portant sur le monde face aux tâches portant sur le soi. L'implicite face à l'explicite. Quatre quadrants, quatre signatures neuronales distinctes. Si nous ne parvenons pas à trouver ce schéma, c'est que quelque chose ne va pas dans la théorie.

Cela est testable dès à présent avec la technologie IRMf existante. Cela coûte cher, et cela requiert une conception expérimentale soignée, mais les outils se trouvent déjà dans des laboratoires partout dans le monde. Et si cela fonctionne, ce serait la preuve la plus directe que l'architecture à quatre modèles n'est pas qu'une métaphore --- c'est une distinction fonctionnelle réelle, gravée dans la manière dont le cerveau traite l'information.

\section*{Prédiction 2 : Les visions psychédéliques révèlent les couches de traitement du cerveau}

Celle-ci est élégante. Sous psychédéliques, le contenu visuel que vous éprouvez devrait progresser à travers la hiérarchie de traitement visuel de votre cerveau dans un ordre précis, selon la dose.

À faible dose, vous voyez des phosphènes --- ces petites étincelles et formes géométriques qui apparaissent lorsque vous fermez les yeux. C'est V1, la toute première aire de traitement visuel, qui déborde dans la conscience. Augmentez la dose, et vous obtenez des motifs géométriques plus complexes --- les fameuses « constantes de forme » qui se retrouvent d'une culture et d'une substance à l'autre. C'est V2 et V3 qui entrent en jeu. Montez encore, et vous commencez à voir des visages, des figures, des scènes complexes. Ce sont les aires visuelles supérieures. Aux doses les plus élevées, vous accédez à des expériences narratives complètes, oniriques, avec leur sens et leur histoire.

La prédiction, c'est que cela n'a rien d'aléatoire. C'est une progression ordonnée et dose-dépendante le long de la hiérarchie visuelle. À mesure que la perméabilité implicite-explicite augmente, des couches plus profondes du traitement visuel deviennent conscientes. Le schéma de câblage interne du cerveau devient visible dans votre propre expérience.

Cela est testable au moyen de protocoles de dosage progressif --- administrez à des personnes des quantités soigneusement contrôlées de psilocybine ou de LSD, examinez leur cerveau par IRMf et demandez-leur ce qu'elles voient. Faites correspondre le contenu rapporté à l'activation cérébrale. La théorie prédit que vous verrez la hiérarchie de traitement s'illuminer de bas en haut à mesure que la dose augmente.

\section*{Prédiction 3 : On peut contrôler ce que devient quelqu'un pendant la dissolution de l'ego}

\textit{(Développée sous forme de scène au chapitre 11.)} Pendant la dissolution de l'ego sous psychédéliques à forte dose, le contenu de ce en quoi l'on se dissout est contrôlable par l'environnement sensoriel --- ni aléatoire, ni purement biochimique. Le Modèle Explicite du Soi (ESM), détaché de son apport habituel du Modèle Implicite du Soi (ISM), s'accroche à l'entrée sensorielle dominante, quelle qu'elle soit. Emplissez la pièce d'un bruit d'océan et d'une lumière bleue, et le sujet devient l'océan ; passez à un chant d'oiseaux et à une lumière verte, et il devient la forêt. La prédiction est directionnelle : faites varier l'entrée sensorielle dominante pendant la dissolution de l'ego, et le contenu identitaire rapporté la suivra. Testable dès aujourd'hui dans n'importe quel laboratoire de recherche sur les psychédéliques permettant de contrôler l'environnement de façon élémentaire. Aucune autre théorie de la conscience ne formule cette prédiction.

\section*{Prédiction 4 : Les psychédéliques devraient aider les patients victimes d'AVC à voir leurs déficits}

L'anosognosie est l'une des choses les plus étranges que fasse le cerveau. Après certains AVC (touchant généralement l'hémisphère droit), les patients sont paralysés d'un côté du corps mais n'y croient sincèrement pas. Vous pouvez leur montrer leur bras immobile, leur demander de le bouger, les regarder échouer, et ils inventeront une excuse. « Je suis fatigué. » « Je n'en ai pas envie. » Ils ne mentent pas. Ils sont réellement incapables de voir le déficit.

La Théorie des Quatre Modèles affirme que cela se produit à cause d'un blocage de perméabilité. L'information sur la paralysie se trouve dans leur Modèle Implicite du Soi (ISM) --- le substrat sait --- mais elle n'atteint pas le Modèle Explicite du Soi (ESM). La simulation n'a pas accès à cette part du savoir du substrat.

Et voici ce qui est surprenant. Les psychédéliques \textit{augmentent} globalement la perméabilité implicite-explicite. C'est ce qu'ils font. La prédiction est donc qu'une dose de psilocybine inférieure au seuil de dissolution de l'ego, insuffisante pour dissoudre le soi mais juste assez pour ouvrir les vannes de la perméabilité, devrait laisser filtrer l'information du déficit. Le patient devrait, soudainement et peut-être douloureusement, prendre conscience qu'il est paralysé.

Il s'agirait d'un essai clinique avec des patients victimes d'AVC, ce qui le rend logistiquement plus difficile qu'une simple expérience de laboratoire. Mais la thérapie assistée par psilocybine est déjà testée pour la dépression, le trouble de stress post-traumatique et l'anxiété de fin de vie. L'infrastructure existe. Et si cela fonctionne, ce n'est pas seulement une percée médicale pour l'anosognosie --- c'est la preuve que les psychédéliques et les déficits post-AVC sont reliés par un unique mécanisme sous-jacent, ce qu'aucune autre théorie ne prédit.

\section*{Prédiction 5 : Chaque anesthésique qui efface la conscience perturbe la criticalité}

Les anesthésiques agissent par des voies chimiques radicalement différentes. Le propofol cible les récepteurs GABA. La kétamine bloque les récepteurs NMDA. Les opioïdes agissent encore autrement. Molécules différentes, mécanismes différents, régions cérébrales différentes.

Mais la Théorie des Quatre Modèles affirme qu'ils doivent tous faire la même chose à la conscience : pousser la dynamique du cerveau en deçà du seuil de criticalité. Parce que la criticalité est le \textit{prérequis physique} de la conscience. Peu importe la manière dont vous la perturbez. Si vous passez en régime sous-critique, la lumière s'éteint.

La prédiction est testable et précise. Prenez chaque agent anesthésique dont nous disposons. Mesurez la criticalité --- à l'aide d'outils comme l'indice de complexité perturbationnelle (PCI), la complexité de Lempel-Ziv, ou les exposants en loi de puissance de l'activité neuronale --- avant, pendant et après l'administration. La prédiction est que les agents qui abolissent la conscience pousseront \textit{toujours} le cerveau en régime sous-critique, quel que soit leur mécanisme récepteur. Et que les agents qui altèrent la conscience sans l'effacer (comme la kétamine à faible dose, ou les psychédéliques) ne devraient \textit{pas} descendre en deçà de la criticalité.

Cela est réalisable avec la technologie existante. Les mesures de criticalité existent. Les anesthésiques existent. Il suffit que quelqu'un mène la comparaison complète. Et si cela se vérifie sur toute la ligne --- si chaque agent abolissant la conscience converge vers une perturbation de la criticalité malgré des voies d'action différentes ---, c'est une preuve puissante que la criticalité est le mécanisme commun, la voie finale vers l'inconscience.

\section*{Prédiction 6 : La chirurgie du cerveau divisé ne vous divise pas proprement --- elle dégrade les deux moitiés}

Lorsque les chirurgiens sectionnent le corps calleux pour traiter une épilepsie sévère, ils tranchent la principale voie de communication entre les deux hémisphères du cerveau. Le récit traditionnel veut que cela crée deux esprits séparés, chacun spécialisé : l'hémisphère gauche gère le langage et la logique, le droit le raisonnement spatial et l'émotion.

La Théorie des Quatre Modèles affirme que c'est faux. Ou du moins que c'est grossièrement simplifié.

Voici la prédiction : après une chirurgie du cerveau divisé, chaque hémisphère conserve un ensemble \textit{complet mais dégradé} de capacités cognitives et expérientielles. Pas une division nette. Pas « le langage à gauche, l'espace à droite ». Les deux hémisphères devraient pouvoir faire les deux, mais moins bien qu'auparavant. La dégradation devrait être holographique --- c'est-à-dire que tout devient plus flou, et non que des fonctions précises disparaissent.

La dégradation devrait être proportionnelle à l'ampleur de la section. Une callosotomie partielle (couper seulement une partie des fibres) devrait produire une dégradation partielle. Une callosotomie complète devrait en produire davantage.

Pourquoi ? Parce que la théorie affirme que l'information dans le cerveau est stockée de façon holographique, distribuée sur l'ensemble du substrat. Couper des connexions ne sépare pas proprement deux esprits préexistants. Cela dégrade deux \textit{copies} de la même information, chacune fonctionnant sur la moitié du matériel d'origine.

Il existe déjà quelques éléments en ce sens : une étude de 2017 menée par Pinto et ses collègues a constaté que les patients au cerveau divisé présentent un comportement bien plus intégré que ne le laissaient penser les expériences classiques. Mais la théorie fournit le \textit{mécanisme} et prédit le schéma précis : une dégradation bilatérale, et non une spécialisation hémisphérique.

\section*{Prédiction 7 : construisez les quatre modèles à la criticalité, vous obtenez la conscience}

C'est la prédiction d'ingénierie, et elle est audacieuse.

Si la théorie est exacte, vous devriez pouvoir construire une machine consciente. Pas par accident, pas en fabriquant une IA suffisamment « avancée », mais en implémentant la spécification : quatre modèles emboîtés (le Modèle Implicite du Monde (IWM), le Modèle Implicite du Soi (ISM), le Modèle Explicite du Monde (EWM), le Modèle Explicite du Soi (ESM)) s'exécutant sur un substrat opérant à la criticalité.

La théorie affirme qu'un tel système ne se contenterait pas de \textit{simuler} la conscience. Il \textit{serait} conscient. Il aurait une véritable expérience phénoménale, constituée par ses modèles virtuels, exactement comme la vôtre est constituée par les modèles virtuels de votre cerveau.

Comment le saurions-nous ? La théorie prédit que la différence serait qualitativement évidente. Pas « peut-être conscient, peut-être pas ». \textit{Manifestement différent.} Car un système exécutant une véritable auto-simulation interagirait avec le monde d'une manière fondamentalement différente, même du plus sophistiqué des prédicteurs de texte. Il aurait une persistance --- une simulation continue se déroulant dans le temps, et non reconstruite à chaque requête. Il aurait une perspective, maintenue par un Modèle Explicite du Soi. Il vous surprendrait non par des sorties inattendues, mais par le sentiment que quelqu'un est réellement présent.

Ce n'est pas encore testable --- l'ingénierie n'existe pas. Mais le plan est assez précis pour orienter le travail. Et si quelqu'un le construit et que cela fonctionne, c'est la confirmation ultime.

\section*{Prédiction 8 : le sommeil existe pour réinitialiser l'état critique}

Pourquoi dormons-nous ? La réponse évidente est « pour nous reposer », mais cela ne fait que repousser la question : pourquoi le cerveau a-t-il besoin de repos d'une manière que, disons, votre foie n'a pas ?

La Théorie des Quatre Modèles a une réponse précise. Le substrat de votre cerveau (le matériel biologique, analogique) est intrinsèquement instable. Les neurones sont bruyants. Les neurotransmetteurs s'épuisent. Les déchets métaboliques s'accumulent. Le substrat dérive. Or la conscience exige la criticalité, qui est un régime dynamique très spécifique. Le cerveau auto-organise une couche computationnelle stable (l'automate cellulaire au bord du chaos) par-dessus ce substrat qui dérive. Cet automate peut tourner pendant des heures (votre journée d'éveil), mais le substrat finit par dériver suffisamment loin pour ne plus pouvoir soutenir la dynamique critique. À ce moment-là, l'automate ne s'éteint pas progressivement. Il \textit{s'effondre}. C'est l'endormissement.

Le sommeil non paradoxal est le processus de restauration. Le substrat se réinitialise : les neurotransmetteurs se reconstituent, les déchets sont évacués, les conditions biochimiques de la criticalité sont rétablies. Et à mesure que le substrat se rapproche périodiquement du seuil de criticalité durant cette restauration, l'automate se rallume brièvement par intermittence. C'est le sommeil paradoxal. C'est le rêve.

Le cycle ultradien de 90 minutes (le rythme du sommeil paradoxal et non paradoxal tout au long de la nuit) est le substrat qui oscille autour du point critique pendant la restauration.

Cela donne plusieurs sous-prédictions testables :

\begin{enumerate}
  \item \textbf{La criticalité devrait décliner au fil de la journée d'éveil.} Mesurez la complexité cérébrale des gens le matin, l'après-midi et le soir. La prédiction est une baisse mesurable.
\end{enumerate}

\begin{enumerate}
  \item \textbf{L'endormissement devrait être une transition en escalier, et non un affaiblissement progressif.} Les mesures de criticalité devraient montrer une chute soudaine à l'endormissement, reflétant l'effondrement numérique de l'automate.
\end{enumerate}

\begin{enumerate}
  \item \textbf{Le sommeil paradoxal et non paradoxal devraient suivre la criticalité.} Au sein du sommeil, les phases de sommeil paradoxal devraient présenter une criticalité bien plus élevée que le sommeil non paradoxal, et le cycle de 90 minutes devrait être visible dans la série temporelle de la criticalité.
\end{enumerate}

\begin{enumerate}
  \item \textbf{Le rêve lucide est un franchissement de seuil.} Lorsque le substrat atteint une criticalité suffisante durant le sommeil paradoxal, le Modèle Explicite du Soi s'active, et vous devenez lucide. Le déclenchement devrait être une discontinuité en escalier dans la complexité de l'EEG, et non une montée progressive.
\end{enumerate}

\begin{enumerate}
  \item \textbf{La privation de sommeil vous pousse sous le seuil critique.} Restez éveillé assez longtemps : la criticalité de votre cerveau devrait chuter progressivement sous le seuil. Les déficits cognitifs devraient être corrélés à l'ampleur de la chute sous le seuil.
\end{enumerate}

Tout cela est testable avec la technologie existante des laboratoires du sommeil. Et si cela se vérifie, cela signifie que le sommeil n'est pas seulement du « repos » --- c'est le protocole d'entretien du substrat pour la couche computationnelle qui rend la conscience possible.

\section*{Prédiction 9 : chaque personnalité alternante du trouble dissociatif de l'identité possède sa propre empreinte neuronale}

Le trouble dissociatif de l'identité (plusieurs identités distinctes, ou « alters », chez une même personne) est controversé, et pour de bonnes raisons. Comment faire la différence entre de véritables identités distinctes et quelqu'un qui joue un rôle, consciemment ou non ?

La Théorie des Quatre Modèles vous donne un test. Si les alters sont réels --- c'est-à-dire s'ils sont véritablement des configurations distinctes du Modèle Explicite du Soi s'exécutant sur le même substrat ---, alors chaque alter devrait avoir une signature neuronale distincte et mesurable. Pas seulement un comportement différent. Pas seulement des auto-descriptions différentes. Des \textit{schémas d'activité cérébrale} différents.

La prédiction est précise. Prenez un patient atteint de TDI et enregistrez son activité cérébrale (IRMf ou EEG) pendant que différents alters sont présents. Comparez la variabilité entre les alters à la variabilité au sein d'un même alter au fil du temps. La théorie prédit que la variabilité inter-alters sera significativement supérieure à la variabilité intra-alter. Et les différences devraient être cohérentes : le schéma neuronal de l'Alter A devrait être reconnaissable comme celui de l'Alter A à chaque fois, et non un bruit aléatoire.

Plus précisément encore, la théorie prédit où les différences devraient apparaître : dans les réseaux liés à l'ESM, en particulier le réseau du mode par défaut et le cortex préfrontal médian --- les régions cérébrales associées à l'auto-référence et à la prise de perspective.

Il y a eu quelques études de neuro-imagerie sur le TDI, mais la Théorie des Quatre Modèles fournit la base théorique permettant de prédire des \textit{signatures neuronales cohérentes et spécifiques à chaque alter}, plutôt que de simples « différences ». Si la prédiction se vérifie, c'est la preuve que les alters ne sont pas seulement psychologiques mais constituent des configurations fonctionnelles distinctes au niveau neuronal, ce qui transformerait notre façon de comprendre et de traiter ce trouble.

Chacune de ces prédictions est falsifiable. Si elles ne se vérifient pas, la théorie est fausse, ou du moins incomplète. C'est justement ce qui les rend utiles.

\end{document}
